% !TeX program = pdflatex
% Compile this file twice with pdfLaTeX (TeX Live recommended).
% All mathematical text, bibliographies and printed code listings are inline.
% Volume labels are prefixed by manuscript; printed theorem numbering stays local.
\documentclass[11pt,letterpaper]{article}
\usepackage[T1]{fontenc}
\usepackage{lmodern,microtype}
\usepackage[letterpaper,left=0.96in,right=0.96in,top=0.86in,bottom=0.92in,headheight=14pt,headsep=15pt,footskip=29pt]{geometry}
\usepackage{amsmath,amssymb,amsthm,mathtools,array,booktabs,longtable}
\usepackage{xcolor,listings,fancyhdr}
\definecolor{ink}{HTML}{20363D}
\definecolor{linkink}{HTML}{205860}
\definecolor{ruleink}{HTML}{98A9A9}
\usepackage[unicode,hypertexnames=false,bookmarksopen=true,bookmarksopenlevel=0,bookmarksnumbered=true,colorlinks=true,linkcolor=linkink,urlcolor=linkink,citecolor=linkink]{hyperref}
\usepackage{bookmark}
\hypersetup{pdftitle={Swisher's Prime-Power Supercongruences: Collected Manuscripts},pdfsubject={Eight selected conjectures, reductions, proofs and working papers, 10--11 November 2026},pdfauthor={Keio research collection}}
\urlstyle{same}
\newtheorem{theorem}{Theorem}[section]
\newtheorem{lemma}[theorem]{Lemma}
\newtheorem{proposition}[theorem]{Proposition}
\newtheorem{corollary}[theorem]{Corollary}
\newtheorem{supplemma}{Lemma}
\numberwithin{equation}{section}
\newcommand{\vp}{\operatorname{ord}}
\newcommand{\Zp}{\mathbb Z_p}
\allowdisplaybreaks[2]
\emergencystretch=3em
\raggedbottom
\lstset{basicstyle=\ttfamily\scriptsize,breaklines=true,columns=fullflexible,keepspaces=true,showstringspaces=false}
\newcommand{\PaperShort}{Editorial guide}
\newcommand{\PaperID}{}
\pagestyle{fancy}\fancyhf{}
\fancyhead[L]{\small\sffamily\color{ink}\PaperShort}
\fancyhead[R]{\small\sffamily\color{ink}\PaperID}
\fancyfoot[C]{\small\thepage}
\renewcommand{\headrulewidth}{.3pt}
\newcommand{\RestoreCounters}{%
  \setcounter{section}{0}\setcounter{subsection}{0}\setcounter{equation}{0}%
  \counterwithout*{theorem}{section}\counterwithin*{theorem}{section}\counterwithout*{equation}{section}\counterwithin*{equation}{section}%
  \setcounter{theorem}{0}\setcounter{supplemma}{0}\setcounter{footnote}{0}%
  \renewcommand{\thesection}{\arabic{section}}%
}
\newcommand{\LocalAppendix}{\setcounter{section}{0}\setcounter{subsection}{0}\renewcommand{\thesection}{\Alph{section}}}
\newcommand{\SetPaper}[2]{%
  \renewcommand{\PaperID}{#1}\renewcommand{\PaperShort}{#2}%
  \pdfbookmark[0]{#1. #2}{volume-#1}\phantomsection\label{paper-#1}%
  \thispagestyle{plain}}
\newcommand{\PaperTitle}[3]{%
  \noindent{\sffamily\small\bfseries\color{linkink}\MakeUppercase{\if S\PaperID Supplement\else Manuscript\fi}\ \PaperID\par}
  \vspace{.65em}\noindent{\raggedright\color{ink}\LARGE\bfseries #1\par}
  \vspace{.65em}\noindent{\sffamily\small #2\hfill #3\par}
  \vspace{.4em}\noindent{\color{ruleink}\rule{\linewidth}{.45pt}}\par\smallskip}
\makeatletter
\newcommand{\LocalEntry}[3]{\@dottedtocline{1}{0pt}{2.3em}{\numberline{#2}\hyperref[#1]{#3}}{\pageref{#1}}}
\makeatother
\newenvironment{LocalContents}{\par\medskip\noindent{\sffamily\bfseries\color{ink}Contents of this manuscript}\par\smallskip\begingroup\small}{\endgroup\par}
\newcommand{\entry}[3]{%
  \noindent\makebox[.067\linewidth][l]{\sffamily\bfseries\color{linkink}#1}%
  \hyperref[paper-#1]{\textbf{#2}}\hfill\pageref{paper-#1}\par\nobreak
  \noindent\hspace*{.067\linewidth}{\small #3}\par\vspace{.55em}}
\newcommand{\vol}[1]{\hyperref[paper-#1]{#1}}
\begin{document}
\thispagestyle{plain}
\pdfbookmark[0]{Contents of the collection}{guide-collection}
\noindent{\sffamily\color{linkink}\small\bfseries COLLECTED MANUSCRIPTS\hfill 10--11 NOVEMBER 2026\par}
\vspace{1.1em}
\noindent{\fontsize{28}{32}\selectfont\bfseries\color{ink}Swisher's prime-power\\supercongruences\par}
\vspace{.5em}\noindent{\Large Eight conjectures from Table 2\par}
\vspace{.75em}\noindent{\color{ruleink}\rule{\linewidth}{.7pt}}\par
\medskip
\noindent This collection follows the eight targets A, C, D, E, F, G, J and K from the
first reductions to the final odd-level E.3 argument. The seven main manuscripts
appear in eight entries because the first circulated in two mathematically distinct
editions. A separately titled F.3 working proof follows as a supplement.
\par\medskip
\entry{I}{Targeted reductions for eight Swisher congruences}{First edition. Boundary sums, exact reductions and obstructions.}
\entry{II}{Cubic scaling and two Swisher supercongruences}{Revised first manuscript. All levels of C.3 and K.3 in the stated prime ranges.}
\entry{III}{All-level proofs for J.3 and the odd branch of F.3}{Half-cutoff scaling for J.3 and a short-cutoff proof for F.3.}
\entry{IV}{Supersingular advances for A.3, D.3 and G.3}{The second A.3 level, the D.3 residue branch and odd levels of G.3.}
\entry{V}{Exact interpolation and positive moments}{Completion of A.3 and the complementary D.3 and G.3 branches.}
\entry{VI}{Exact reductions for the remaining odd branch of E.3}{Three interpolation grids, a companion sum and a two-step residual.}
\entry{VII}{Finite polynomial separation and higher odd levels of E.3}{All odd levels from five onward and the third-level residual.}
\entry{VIII}{Reflection, rational rigidity and the odd branch of E.3}{The uniform third-level argument and the resulting all-odd-level statement.}
\entry{S}{A short-cutoff proof of the odd branch of F.3}{Working version with its own theorem and terminating identity.}
\clearpage
\thispagestyle{plain}
\pdfbookmark[0]{Guide to results and prime ranges}{guide-scope}
\noindent{\LARGE\bfseries\color{ink}Guide to results and prime ranges\par}
\medskip
\noindent The table locates the final theorem statements for the eight targets.
Each manuscript retains its own section and equation numbers. A statement
described as open in an early paper records its status at that stage.
\par\medskip
{\small\renewcommand{\arraystretch}{1.2}\noindent
\begin{tabular}{@{}p{.07\linewidth}p{.70\linewidth}p{.14\linewidth}@{}}\toprule
\textbf{Target}&\textbf{Range of the statements in the collected proofs}&\textbf{Paper}\\\midrule
A.3 & Both residue classes for all odd primes, including $3$ and $5$, at every prescribed level. & \vol{IV}, \vol{V}\\
C.3 & All $p>3$ and $r\ge1$, modulo $p^{4r}$. & \vol{II}\\
D.3 & $p\equiv1\pmod3$ at every level, modulo $p^{6r}$, and all prescribed levels for $p>3$, $p\equiv2\pmod3$. & \vol{IV}, \vol{V}\\
E.3 & $p>3$, $p\equiv2\pmod3$, every odd $r\ge3$, modulo $p^{3r-1}$. Paper VII treats $r\ge5$ and VIII completes $r=3$. & \vol{VI}--\vol{VIII}\\
F.3 & $p\equiv3\pmod4$, every odd $r\ge3$, including $p=3$. The proved exponent $3(r-1)/2$ attains the target $r$. & \vol{III}, \vol{S}\\
G.3 & All prescribed branches for $p>3$. The odd supersingular recurrence ($r\ge3$) also includes $p=3$. & \vol{IV}, \vol{V}\\
J.3 & All $p>3$ and $r\ge1$, with the half cutoff and modulus $p^{4r}$. & \vol{III}\\
K.3 & All $p>5$ and $r\ge1$, modulo $p^{4r}$, with the sign in the published table. & \vol{II}\\\bottomrule
\end{tabular}\par}
\medskip
\noindent\textbf{Small primes.} At $(p,r)=(3,1)$ the C.3 and J.3 differences are
$-27/16$ and $135/32$, respectively. Both have valuation $3$, below the required
$4$. For the even G.3 recurrence, $G(9)+27=3^5\cdot121019/2^{20}$ has
valuation $5$ where $6$ is required. See \vol{V}, \S9.
\par\medskip
\noindent\textbf{Relation to published work.} The directly proved E.3 and F.3
ranges above are the odd supersingular branches. For other branches see
Guo--Zudilin, \href{https://arxiv.org/abs/2001.02311}{\emph{Dwork-type supercongruences
through a creative $q$-microscope}}, \S3, and
Guo, \href{https://doi.org/10.1216/rmj.2025.55.147}{\emph{Rocky Mountain J. Math.}
55 (2025), 147--154}. The former also establishes B.3 and L.3. For I.3 see
Guo, \href{https://doi.org/10.1142/S1793042118501701}{\emph{Int. J. Number Theory}
15 (2019), 29--36}. The initial paper cites the separate H.3 manuscript [132]
as an input. The reference list for each paper remains with its text.
\par\medskip
\noindent\textbf{Editorial record.} The first edition was transcribed from its
13-page PDF and checked against the extracted original text. The seven later
entries through VIII use the deposited editable sources. Their proof and review
claims retain the status stated by the authors. This edition collates the texts
and checks the printed cross-references. Five certificate programs for E.3 were
rerun. Among the reproduced calculations are 158 higher-level and 56
fifth-level moment values, and nine direct evaluations at $r=3$.


%%%%%%%%%%%%%%%%%%%%%%%%%%%%%%%%%%%%%%%%%%%%%%%%%%%%%%%%%%%%%%%%%%%%%%%%%%
% I | Earlier edition, transcribed from the surviving 13-page PDF.
%%%%%%%%%%%%%%%%%%%%%%%%%%%%%%%%%%%%%%%%%%%%%%%%%%%%%%%%%%%%%%%%%%%%%%%%%%
\clearpage\begingroup
\SetPaper{I}{Targeted reductions for eight congruences}
\RestoreCounters
\counterwithout*{equation}{section}\renewcommand{\theequation}{\arabic{equation}}
\PaperTitle{Targeted Reductions for Eight Swisher Congruences\\[4pt]\large Transfer from Corpus [131] and [132], Exact Obstructions, and Remaining Boundary Sums}{Research working note}{10 November 2026}
\begin{abstract}
We investigate (A.3), (C.3), (D.3), (E.3), (F.3), (G.3), (J.3), and (K.3) using the excellent-Frobenius and jet-transport mechanisms in corpus manuscripts [131] and [132]. No new complete resolution of any of these eight conjectures is claimed. The main rigorous outputs are an all-level prefactor congruence and an equivalent residual problem for (K.3), an all-level diagonal congruence and boundary reduction for the odd branch of (F.3), and exact finite block reductions for (C.3) and (J.3). We also prove a parity-sensitive boundary-residual lemma, and give exact counterexamples to several tempting intermediate claims. The counterexamples concern proof transfers, not the original conjectures. Computation is reported separately from proof.
\end{abstract}
\noindent\textbf{Status.} This is an incomplete research note, with explicit proofs of the stated reductions. Historical novelty is not asserted. Classical identities and published input theorems retain their attribution. The corpus manuscripts are examined as proposed mathematical inputs; this note does not certify all of their geometric arguments.


\begin{LocalContents}
\LocalEntry{loc-I-1}{1}{Scope and outcome}
\LocalEntry{loc-I-2}{2}{What the corpus mechanisms do and do not transfer}
\LocalEntry{loc-I-3}{3}{An exact reduction of Swisher's \texorpdfstring{$(K.3)$}{(K.3)}}
\LocalEntry{loc-I-4}{4}{An exact boundary reduction of the odd branch of (F.3)}
\LocalEntry{loc-I-5}{5}{The C.3 and J.3 deformation reductions}
\LocalEntry{loc-I-6}{6}{An exact obstruction to a direct \texorpdfstring{$(A.3)$}{(A.3)} lift}
\LocalEntry{loc-I-7}{7}{Exact boundary residuals and an odd-level obstruction}
\LocalEntry{loc-I-8}{8}{Verification and the unresolved statements}
\end{LocalContents}
\clearpage

\section{Scope and outcome}\label{loc-I-1}
The excluded targets are (B.3), (H.3), (I.3), and (L.3).
We use Swisher's published Table~2~\cite{I-Swisher2015}, rather than
the initial arXiv version: the signs in (K.3) were corrected in the
published table. Throughout, $v_p(p)=1$ and $(a)_k$ denotes a rising
factorial. The natural non-exceptional ranges are stated with each result.
We do not infer a current exhaustive literature classification from a
failure to locate a later proof.

\begin{center}\small
\begin{tabular}{@{}lp{105mm}@{}}\toprule
Target & Outcome of this investigation\\\midrule
A.3 & The ordinary target is $p^{5r}$. A direct first-level Whipple
perturbation comparison fails already at $(p,r)=(5,2)$; the original
conjecture is not refuted.\\
C.3, J.3 & Exact terminating evaluations give a common finite three-term obstruction modulo $p^{4r}$. Its cancellation remains unproved.\\
D.3 & The ordinary target is $p^{6r}$. No all-level closure is obtained;
the boundary lemma also applies to its underlying sixth-power polynomial.\\
E.3 & No completion of the odd supersingular branch is obtained.
Its short cutoff does not satisfy the cubic residual precision used in [132].\\
F.3 & The odd-branch diagonal recurrence is proved for every level;
an explicit boundary recurrence remains.\\
G.3 & No complete closure is obtained. The fourth-power boundary
calculation identifies a transfer hypothesis that still needs repair.\\
K.3 & The entire prefactor recurrence is proved at modulus $p^{4r}$;
the remaining problem is a single $p^{3r}$ residual congruence.\\\bottomrule
\end{tabular}
\end{center}

The input corpus snapshot is commit
\texttt{91c0fa1f0872f070b001249d0ab7c47e3499a907} of
\url{https://github.com/QEEDDD/Verifiable-Mathematical-Reasoning-Corpus}.
We examined the TEX sources of [131], \emph{Excellent Frobenius Lifts
and Supercongruences Unified}, and [132], \emph{Swisher H3 Unified
Ordinary and Supersingular}. They are not modified here.


\section{What the corpus mechanisms do and do not transfer}\label{loc-I-2}
The modular cusp theorem in [131] (label \texttt{B:thm:main}) assumes
weight $w\in\{1,2\}$ and yields precision $p^{(w+1)r}$.
The higher-dimensional standard-family theorem yields $p^{2r}$.
An analytic evaluation at a unit argument requires its own integral
evaluation map, fixed-point condition, and multiplier identification.
These hypotheses cannot be replaced by a formal substitution.

There is an explicit obstruction already proved in [131]
(label \texttt{E:thm:sh-cube}). Let
\[
 F_d(u)=\sum_{k\ge0}\binom{2k}{k}^du^k,\quad K=p^r,\qquad
 E_r=F_d(u)F_{K/p}(\sigma(u))-F_d(\sigma(u))F_K(u),
\]
where $F_L$ is truncated at $k<L$ and $\sigma$ is the excellent lift
in that theorem. For $d\ge4$ and $p>d$,
\[
 [u^{K+1}]E_r\equiv2^{2d-1}d(d-3)p^{2r}\pmod{p^{2r+1}}.
\]
The right side is nonzero at this precision. Thus the full-series
congruence cannot simply be raised from $2r$ to $4r$ or $5r$.
For example, $d=4,p=7,r=1$ gives
\[
 [u^8]E_1=27435582641602576,\qquad v_7([u^8]E_1)=2.
\]
Even applying $1+4u\,d/du$ leaves an obstruction: its coefficient is
multiplied by $4K+5$, a $7$-adic unit. This excludes a naive global
derivative upgrade; it does not exclude additional cancellation at a
particular argument. This obstruction is from [131], not a new result
of the present note.

For clarity, there is an exact weighted transfer identity.
Write $D_r=F_r-\rho F_{r-1}\circ\sigma$,
$\theta=t\,d/dt$, and $\mu=\theta\sigma/\sigma$.
At a valid analytic fixed point $t_0$, if
\[
 c\theta\rho(t_0)=\rho(t_0)(\mu(t_0)-1),\qquad
 \lambda=\rho(t_0)\mu(t_0),
\]
then the product and chain rules give
\[
 (1+c\theta)F_r(t_0)-\lambda(1+c\theta)F_{r-1}(t_0)
 =(1+c\theta)D_r(t_0).
\]
Differentiation preserves an established integral modulus, but does
not itself increase it. The remaining weighted residual needs a
separate estimate. Likewise, the [132] Taylor-transport argument
requires the precise residual order of its differential equation;
Section~\ref{I-sec:parity-location} below calculates what happens at
the short odd cutoffs of other families.



\section{An exact reduction of Swisher's \texorpdfstring{$(K.3)$}{(K.3)}}\label{loc-I-3}
\label{I-sec:k-reduction}

For each odd positive integer $N$, put $h_N=(N-1)/2$ and define
\begin{align}
 K(N)&=\sum_{k=0}^{h_N}(42k+5)
               \frac{(1/2)_k^3}{k!^3 64^k},\label{I-eq:k-sum}\\
 B(N)&=\frac{12\displaystyle\prod_{j=0}^{N-1}(N+2j)}
                {2^{3N-2}(h_N!)^2},\label{I-eq:k-prefactor}\\
 U(N)&=\sum_{n=0}^{N}
                \frac{(-N/2)_n^2}{n!(-3N/2)_n}.\label{I-eq:k-residual}
\end{align}
All three quantities are rational numbers. In particular,
$K(1)=5$, $B(1)=6$, and $U(1)=5/6$. We use the positive sign in
Swisher's published Table~2: with $\chi_p=(-1)^{(p-1)/2}$, the target is
\begin{equation}\label{I-eq:k-target}
 K(p^r)\equiv\chi_p pK(p^{r-1})\pmod{p^{4r}},\qquad p>5.
\end{equation}
The first arXiv version has an erroneous additional minus sign in this
entry; the published version corrects it.

\begin{proposition}[Exact factorization]\label{I-prop:k-factorization}
For every odd positive integer $N$,
\begin{equation}\label{I-eq:k-factorization}
 K(N)=B(N)U(N).
\end{equation}
\end{proposition}
\begin{proof}
Use the Guillera WZ pair employed by Osburn and Zudilin in their proof
of $(K.2)$, namely
\begin{align*}
 F(n,k)&=(84n^2-56nk+4k^2+52n-12k+5)\\
 &\hspace{8mm}\cdot
 \frac{(-1)^k(1/2)_n(1/2)_{n+k}(1/2)_{n-k}^{2}}
 {2^{4n}(n!)^2(2n-k+1)!},\\
 G(n,k)&=
 \frac{64n^2(-1)^k(1/2)_n(1/2)_{n+k-1}(1/2)_{n-k}^{2}}
 {2^{4n}(n!)^2(2n-k)!}.
\end{align*}
Negative-index Pochhammer symbols are defined by reciprocals, and a
reciprocal factorial at a negative integer is zero. The pair satisfies
\[
 F(n,k-1)-F(n,k)=G(n+1,k)-G(n,k).
\]
Moreover $F(n,0)=(42n+5)(1/2)_n^3/(64^n n!^3)$.
Summing first over $0\le n\le h_N$ and then over $1\le k\le N$
gives, exactly as in Osburn--Zudilin's equation~(2.6),
\[
 K(N)=F(h_N,N)+\sum_{k=1}^{N}G(h_N+1,k).
\]
This telescoping identity requires only that $N$ be odd, not prime.
Direct simplification gives
\[
 F(h_N,N)=B(N),\qquad
 G(h_N+1,k)=B(N)
 \frac{(-N/2)_{N+1-k}^{2}}
 {(N+1-k)!(-3N/2)_{N+1-k}}.
\]
Changing the index to $n=N+1-k$ supplies the terms $n=1,\ldots,N$
in \eqref{I-eq:k-residual}; the boundary value supplies $n=0$.
\end{proof}

\begin{theorem}[All-level congruence for the prefactor]
\label{I-thm:k-prefactor}
For every prime $p\ge5$ and integer $r\ge1$,
\begin{equation}\label{I-eq:k-prefactor-congruence}
 B(p^r)\equiv\chi_p pB(p^{r-1})\pmod{p^{4r}},
 \qquad v_p(B(p^r))=r.
\end{equation}
\end{theorem}
\begin{proof}
Set $N=p^r$, $M=p^{r-1}$ and $t=(N-M)/2$, and write
\[
 V=\prod_{\substack{1\le j\le h_N\\p\nmid j}}j,
 \quad O=\prod_{\substack{1\le a<N\\a\ modd,\ p\nmid a}}a,
 \quad W=\prod_{\substack{0\le j<N\\p\nmid j}}(N+2j).
\]
Removing the factors divisible by $p$ from the numerator and the
factorials in \eqref{I-eq:k-prefactor} yields the exact equality
\begin{equation}\label{I-eq:k-prefactor-unit}
 \frac{B(N)}{pB(M)}=\frac{W}{2^{3(N-M)}V^2}.
\end{equation}
Pairing the factors in $W$ about $2N$ gives another exact equality,
\begin{equation}\label{I-eq:k-pairing}
 W=(-1)^t O^2
 \prod_{\substack{1\le a<N\\a\ modd,\ p\nmid a}}
       \left(1-\frac{4N^2}{a^2}\right).
\end{equation}
We claim that the last product is $1$ modulo $p^{3r}$.
Indeed, the inverse-square sum over all units modulo $N$ vanishes
modulo $N$. To see this, choose a unit $g$ with $g^2\not\equiv1\pmod p$;
multiplication by $g$ permutes those units and forces the sum to be
annihilated by the unit $g^{-2}-1$. Pairing $j$ with $N-j$ then shows
that the inverse-square sum over $1\le j\le h_N$, $p\nmid j$, also
vanishes modulo $N$. The sum over odd units $a<N$ is the full sum
minus one quarter of this half-interval sum, so it vanishes as well.
Consequently the linear term in the product in \eqref{I-eq:k-pairing}
is divisible by $p^{3r}$; every term involving two or more selected
factors is divisible by $p^{4r}$.

Let $C_N=\binom{N-1}{h_N}$. Its $p$-free quotient is
\begin{equation}\label{I-eq:k-morley-quotient}
 \frac{C_N}{C_M}=2^t\frac{O}{V}.
\end{equation}
Cai's prime-power extension of Morley's congruence
\cite[Corollary~3]{I-Cai2002} states that
\[
 \frac{C_N}{C_M}\equiv\chi_p4^{N-M}\pmod{p^{3r}}.
\]
Since $(-1)^t=\chi_p$, substituting \eqref{I-eq:k-pairing} and
\eqref{I-eq:k-morley-quotient} into \eqref{I-eq:k-prefactor-unit} gives
\[
 \frac{B(N)}{pB(M)}
 \equiv\chi_p\left(\frac{C_N}{C_M}\right)^2
                  2^{-4(N-M)}
 \equiv\chi_p\pmod{p^{3r}}.
\]
The exact quotient in \eqref{I-eq:k-prefactor-unit} is a $p$-adic unit.
Starting with $B(1)=6$, induction therefore gives $v_p(B(M))=r-1$.
Multiplication by $pB(M)$ raises the precision to $p^{4r}$ and proves
both assertions.
\end{proof}

\begin{proposition}[Integrality and equivalent remaining congruence]
\label{I-prop:k-equivalence}
For $p>5$ and $r\ge0$, $U(p^r)\in\mathbb Z_p^\times$ and
$U(p^r)\equiv5/6\pmod p$. For every $r\ge1$, \eqref{I-eq:k-target}
holds if and only if
\begin{equation}\label{I-eq:k-residual-congruence}
 \boxed{\quad U(p^r)\equiv U(p^{r-1})\pmod{p^{3r}}.\quad}
\end{equation}
\end{proposition}
\begin{proof}
For $N=p^r$ and $1\le n\le N$, write the summand in $U(N)$ as
\begin{equation}\label{I-eq:k-summand-product}
 f(N,n)=-\frac{N}{6n}
 \prod_{j=1}^{n-1}
 \frac{(1-N/(2j))^2}{1-3N/(2j)}.
\end{equation}
For $j<N$, $v_p(N/j)\ge1$, so every factor in this product is a
$p$-adic unit congruent to $1$ modulo $p$. It follows that
$v_p(f(N,n))=r-v_p(n)\ge0$. Only $n=N$ survives modulo $p$ among
$n\ge1$, and that term is $-1/6$ modulo $p$. Thus
$U(N)\equiv1-1/6=5/6\pmod p$.

Finally the exact identity \eqref{I-eq:k-factorization} gives
\begin{align*}
 K(N)-\chi_p pK(M)
 &=B(N)\bigl(U(N)-U(M)\bigr)\\
 &\quad+\bigl(B(N)-\chi_p pB(M)\bigr)U(M).
\end{align*}
The second term belongs to $p^{4r}\mathbb Z_p$ by
Theorem~\ref{I-thm:k-prefactor} and the integrality just proved.
The valuation of $B(N)$ is exactly $r$, proving the equivalence.
\end{proof}

\noindent\textbf{Unresolved step.}
Equation~\eqref{I-eq:k-residual-congruence} is not proved here. The reduction removes a prefactor whose required precision is now established for all levels; it does not establish the remaining sum congruence or claim novelty for the reformulation. Splitting \eqref{I-eq:k-summand-product} into $p$-divisible and $p$-free indices produces correction products, but termwise estimates do not attain the required exponent $3r$. An additional cancellation in their total sum remains necessary.
\subsection{A further harmonic reduction of the K residual}
This subsection uses the $U(N)$ of the preceding K reduction. Put
\[
 P_j(N)=\prod_{k=1}^j\frac{(1-N/(2k))^2}{1-3N/(2k)},\qquad P_0(N)=1.
\]
For every odd positive integer $N$,
\begin{equation}\label{I-eq:Uharmonic}
 U(N)=\frac43-\frac12P_{N-1}(N)+\frac13T(N),\qquad
 T(N)=N^2\sum_{n=1}^{N-1}\frac{P_{n-1}(N)}{n^2(1-3N/(2n))}.
\end{equation}
Indeed the $n$th summand of $U$ is $-NP_{n-1}(N)/(6n)$, and
\[
 f(x)=\frac{(1-x/2)^2}{1-3x/2},\qquad
 f(x)-1=\frac{x}{2}+\frac{x^2}{1-3x/2}.
\]
Summing the resulting difference $P_n-P_{n-1}$ for $1\le n<N$
and then adding the $n=N$ term proves \eqref{I-eq:Uharmonic}.

\begin{proposition}
For $p>5$, $N=p^r$, and $M=p^{r-1}$,
\[
 P_{N-1}(N)\equiv P_{M-1}(M)\pmod{p^{3r}}.
\]
Consequently (K.3) is equivalent to
$T(p^r)\equiv T(p^{r-1})\pmod{p^{3r}}$.
The last congruence remains unproved.
\end{proposition}
\begin{proof}
Remove the factors indexed by multiples of $p$ from the quotient of
the two products. What remains is $\prod_{1\le k<N,\ p\nmid k}f(N/k)$.
On the full unit residue set modulo $N$,
\[
 H_2=\sum_{\substack{1\le k<N\\p\nmid k}}k^{-2}\in p^r\mathbb Z_p,
 \qquad
 H_1=\sum_{\substack{1\le k<N\\p\nmid k}}k^{-1}\in p^{2r}\mathbb Z_p.
\]
For the first claim multiply the residues by a unit whose square is
not $1$ modulo $p$. For the second, pair $k$ with $N-k$ and expand
modulo $N^2$, obtaining $2H_1\equiv-NH_2\pmod{N^2}$.
Now
\[
 \log f(x)=\sum_{\ell\ge1}\frac{3^\ell-2}{\ell2^\ell}x^\ell.
\]
The first two terms summed over the units lie in $p^{3r}\mathbb Z_p$
by these harmonic bounds. Every later term has valuation at least
$\ell r-v_p(\ell)\ge3r$. The convergent logarithm and exponential
therefore prove the product congruence. Equation~\eqref{I-eq:Uharmonic}
and the K equivalence above prove the last assertion, since $3$ is a unit.
\end{proof}

\section{An exact boundary reduction of the odd branch of (F.3)}\label{loc-I-4}
Let $p\equiv3\pmod4$ be a prime. For odd $r\ge1$ put
\[
 N_r=(p^r-3)/4,\qquad
 S_r=\sum_{n=0}^{N_r}(-1)^n(8n+1)\frac{(1/4)_n^3}{(n!)^3}.
\]
Define
\[
 D_r=(8N_r+1)\frac{(1/4)_{2N_r}}{(N_r!)^2},\qquad
 B_r=\sum_{k=1}^{N_r}G(N_r+1,k),
\]
where an empty sum is zero and
\[
 G(n,k)=(-1)^{n+k}4
 \frac{(1/4)_n^2(1/4)_{n+k-1}}
 {(n-1)!^2(n-k)!(1/4)_k^2}.
\]
\begin{proposition}
For every such $p,r$, $S_r=D_r+B_r$. For odd $r\ge3$,
\[
 v_p(D_r)=r-1,\qquad D_r\equiv p^2D_{r-2}\pmod{p^r}.
\]
Consequently the odd branch of Swisher's original (F.3) is equivalent to
\[
 \boxed{B_r\equiv p^2B_{r-2}\pmod{p^r}}\qquad(r\ge3\text{ odd}).
\]
This boxed boundary recurrence is not proved here.
\end{proposition}
\begin{proof}
Set $a=1/4$ and
\[
 F(n,k)=(-1)^{n+k}(8n+1)
 \frac{(a)_n^2(a)_{n+k}}{(n!)^2(n-k)!(a)_k^2}.
\]
Use $1/m!=0$ for negative integers $m$, and set $G(0,k)=0$.
The WZ identity
\[
 F(n,k-1)-F(n,k)=G(n+1,k)-G(n,k)
\]
is an exact rational identity. Indeed its four quotients after division
by $F(n,k)$ are
\[
 \frac{F(n,k-1)}{F(n,k)}=-\frac{(a+k-1)^2}{(a+n+k-1)(n-k+1)},
\]
\[
 \frac{G(n+1,k)}{F(n,k)}=-\frac{(a+n)^2}{(2n+a)(n-k+1)},\qquad
 \frac{G(n,k)}{F(n,k)}=\frac{n^2}{(a+n+k-1)(2n+a)}.
\]
Substitution proves the identity, including the boundary cases by the
zero convention. Summing first over $0\le n\le N$, then over
$1\le k\le N$, gives
\[
 \sum_{n=0}^NF(n,0)=F(N,N)+\sum_{k=1}^NG(N+1,k).
\]
Since $F(N,N)=(8N+1)(a)_{2N}/(N!)^2$, this proves the decomposition.
The WZ pair appears as Lemma 3.1 in Jana--Kalita,
\emph{Proof of a supercongruence conjecture of (F.3) of Swisher using
the WZ-method}, arXiv:2011.02762v1. That article was withdrawn for
errors; only the identity, independently verified above, is used here.

Write $R_r=(1/4)_{2N_r}/(N_r!)^2$. Let $\Gamma_p$ be Morita's gamma
function. Extracting the factors divisible by $p$ twice from the rising
factorial and the factorial gives the exact formula
\[
 \frac{R_r}{p^2R_{r-2}}=
 \left(\frac{p^{r-2}}2-\frac54\right)
 \left(\frac{p^{r-2}}2-\frac14\right)
 \frac{\Gamma_p(p^r/2-5/4)\Gamma_p(p^{r-1}/2+1/4)}
 {\Gamma_p(1/4)\Gamma_p(3/4)
  \Gamma_p((p^r+1)/4)^2\Gamma_p((p^{r-1}+3)/4)^2}.
\tag{*}
\]
For completeness, in $(1/4)_{(p^r-3)/2}$ the first extraction leaves
$p^{(p^{r-1}-1)/2}(3/4)_{(p^{r-1}-1)/2}$. The second leaves
$p^{(p^{r-2}+1)/2}(1/4)_{(p^{r-2}+1)/2}$. Removing the final two
factors of this last rising factorial gives the two prefactors in (*).
In $N_r!$, the two factorials left are indexed by
$(p^{r-1}-1)/4$ and $N_{r-2}$. Comparing the powers of $p$ gives
exactly $p^2$. The gamma signs in the numerator multiply to $+1$
because $(p^r-3)/2$ and $(p^{r-1}-1)/2$ are even; denominator signs
are squared.

All the factors on the right of (*) are $p$-adic units: $p\ne5$.
Modulo $p$, its value is
\[
 \frac5{16}\frac{\Gamma_p(-5/4)}
 {\Gamma_p(1/4)^2\Gamma_p(3/4)^3}=1.
\]
The equality follows from
$\Gamma_p(-5/4)=\frac{16}5\Gamma_p(3/4)$ and
$\Gamma_p(1/4)^2\Gamma_p(3/4)^2=1$.
These are the gamma functional equation and reflection formula.
The base $R_1$ is a unit: $(p-3)/2<(3p-1)/4$ is below the first
$p$-divisible factor in $(1/4)_j$, and $N_1<p$.
Induction using (*) proves $v_p(R_r)=r-1$ and
$R_r\equiv p^2R_{r-2}\pmod{p^r}$. Finally,
$8N_r+1=2p^r-5$ is a unit and is congruent to
$2p^{r-2}-5$ modulo $p$. Multiplying proves both claims for $D_r$.
Subtracting the diagonal recurrence from the exact decomposition proves
the stated equivalence.
\end{proof}

\paragraph{A precise obstruction to deleting the boundary.}
The withdrawn article's Lemma 3.2 asserts
$v_p(G((p^r+1)/4,k))\ge3(r-1)/2$ throughout the relevant range.
For $p=7,r=3,k=13$, this is false:
\[
 v_7(G(86,13))=1,
 \qquad \frac{G(86,13)}7\equiv34\pmod{49}.
\]
For an entirely integer valuation check, let
$A(n)=\sum_{j=0}^{n-1}v_7(4j+1)$ and
$C(n)=v_7(n!)$. Then
\[
 A(86)=14,\quad A(98)=16,\quad A(13)=3,\quad
 C(85)=13,\quad C(73)=11,
\]
so $v_7G=2(14)+16-2(13)-11-2(3)=1$.
Direct rational arithmetic gives
\[
 D_3\equiv245,\qquad B_3\equiv49,\qquad S_3\equiv294\pmod{343}.
\]
Thus the boundary is essential, even though the original conjecture is
satisfied in this test. These computations diagnose the invalid termwise
bound, and do not constitute a proof of the new boundary recurrence.

\section{The C.3 and J.3 deformation reductions}\label{loc-I-5}
For $p>3$, put $N=p^r$, $n_r=(N-1)/2$, and
\[
 c_k=(4k+1)\frac{\binom{2k}{k}^{4}}{256^k},\qquad
 j_k=(6k+1)\frac{\binom{2k}{k}^{3}}{256^k},\quad
 C_r=\sum_{k=0}^{n_r}c_k,\quad J_r=\sum_{k=0}^{n_r}j_k.
\]
Here $C_0=J_0=1$, and $\varepsilon=(-1)^{(p-1)/2}$.
The respective targets are
$C_r-pC_{r-1}\in p^{4r}\mathbb Z_p$ and
$J_r-\varepsilon pJ_{r-1}\in p^{4r}\mathbb Z_p$.
Guo's published all-level result has modulus $p^{3r}$ for both sums
(\url{https://arxiv.org/abs/1910.07551}).

For $d=2,4$, define the finite rational product
\[
 R_d(k,x)=\prod_{a=0}^{k-1}
 \frac{1-x^2/(2a+1)^2}{1-x^2/[d(a+1)]^2}.
\]
Dougall's terminating very-well-poised summation and the terminating
identity used in Long's Lemma 4.2, respectively, give for every positive
odd integer $N$ the exact identities
\begin{align}
 \sum_{k=0}^{(N-1)/2}c_k R_2(k,N)&=N,\\
 \sum_{k=0}^{(N-1)/2}j_k R_4(k,N)&=(-1)^{(N-1)/2}N.
\end{align}
The second identity is recorded at
\url{https://msp.org/pjm/2011/249-2/pjm-v249-n2-p09-p.pdf},
pp.~414--415; replacing the letter $p$ there by an odd integer does not
change the terminating evaluation. The first is the specialization
$a=b=1/2$, $c=(1+N)/2$, and terminating parameter $-(N-1)/2$ of
Dougall's formula. Its right-hand side is
\[
 \frac{(3/2)_{(N-1)/2}((1-N)/2)_{(N-1)/2}}
 {(1)_{(N-1)/2}(1-N/2)_{(N-1)/2}}=N.
\]
These exact classical evaluations are not proofs of the desired
undeformed recurrences.

\subsection{An exact finite separation of the factors divisible by \texorpdfstring{$p$}{p}}
The following reduction avoids any unjustified integral Taylor expansion.
Write $k=pt+e$, $0\le e<p$, put $h=(p-1)/2$, $M=p^{r-1}$,
and $\delta_e=1$ if $e>h$, and $0$ otherwise. Splitting the products
according to divisibility by $p$ proves exactly
\[
 R_d(k,N)=R_d(t,M)
 \left(1-\frac{M^2}{(2t+1)^2}\right)^{\delta_e}U_d(k,N),
\]
where
\[
 U_d(k,N)=
 \frac{\prod_{0\le a<k,\ p\nmid 2a+1}(1-N^2/(2a+1)^2)}
 {\prod_{1\le a\le k,\ p\nmid a}(1-N^2/(da)^2)}.
\]
Every denominator in this last product is a $p$-adic unit. Consequently,
with
\[
 W_d(k)=\sum_{\substack{0\le a<k\\p\nmid 2a+1}}\frac1{(2a+1)^2}
 -\frac1{d^2}\sum_{\substack{1\le a\le k\\p\nmid a}}\frac1{a^2},
\]
one has the rigorous, uniform expansion
\[
 U_d(k,N)\equiv1-N^2W_d(k)\pmod{p^{4r}}.
\]
The other factors above are integral for $k\le n_r$. In the last block,
$t=(M-1)/2$, only $e\le h$ occurs, so the potentially zero extra factor
is never included. This proves that this expansion can safely be
multiplied by all the remaining factors and summed.

Here is an entirely finite expression for the remaining obstruction.
Use $(w,d,\eta)=(c,2,1)$ or $(j,4,\varepsilon)$ and set
\[
 L_t=R_d(t,M),\quad z_t=\frac{M^2}{(2t+1)^2},\quad
 E_t=\{e:0\le e<p,\ pt+e\le n_r\},
\]
\[
 V_t=\sum_{e\in E_t}w_{pt+e},\qquad
 B_t=\sum_{\substack{e\in E_t\\e>h}}w_{pt+e},\qquad
 H_t=\sum_{e\in E_t}w_{pt+e}(1-z_t)^{\delta_e}W_d(pt+e).
\]
Then the exact terminating identities at levels $r$ and $r-1$, followed
by the preceding expansion, prove
\begin{align}
 \sum_{k=0}^{n_r}w_k-\eta p\sum_{k=0}^{n_{r-1}}w_k
 \equiv{}&\sum_{t=0}^{n_{r-1}}(1-L_t)(V_t-\eta p w_t)
 +\sum_{t=0}^{n_{r-1}}L_t z_t B_t\notag\\
 &+N^2\sum_{t=0}^{n_{r-1}}L_tH_t\pmod{p^{4r}}.
\end{align}
Indeed the evaluated deformation at level $r$ is $\eta^rN$, the
lower-level evaluation multiplied by $\eta p$ is the same, and the
$k$-block of the higher-level evaluation is
$L_t(V_t-z_tB_t-N^2H_t)$ modulo $p^{4r}$. Rearrangement gives the display.
All quantities in this congruence are $p$-integral. Thus either original
conjecture is equivalent to the vanishing modulo $p^{4r}$ of its displayed
finite obstruction. We have not proved this vanishing. In particular,
none of its three summands may be suppressed without a separate bound.
For example, at $(p,r)=(5,2)$ their valuations are $(7,7,7)$ for C
and $(6,6,6)$ for J, while the valuation of their total is $8$ in both
cases. At $(5,3)$ the three valuations are $(9,9,10)$ for C and
$(10,10,10)$ for J, while both totals have valuation $12$. These exact
rational checks exhibit the cross-term cancellation that a proof must
control; a termwise $4r$ bound is already false in these examples.

\subsection{Taylor formulation and its precision limitation}
For comparison put
\[
 a^{X}_{r,j}=[x^{2j}]\sum_{k=0}^{n_r}w_kR_d(k,x),
 \qquad X=C,J.
\]
For each fixed $r$, this expansion converges at $x=p^r$: every denominator
has $p$-adic valuation at most $r-1$, so $x^2/[d(a+1)]^2\in p^2\mathbb Z_p$.
One obtains the exact convergent identity
\[
 \Delta_X(r)=-\sum_{j\ge1}p^{2rj}
 \left(a^X_{r,j}-\eta p^{1-2j}a^X_{r-1,j}\right).
\]
For $r=0$, take $a_{0,0}=1$ and $a_{0,j}=0$ for $j>0$.
The quadratic coefficient is explicitly
\[
 a^X_{r,1}=\sum_{k=0}^{n_r}w_k
 \left(\frac1{d^2}H_k^{(2)}-
 \sum_{a=0}^{k-1}\frac1{(2a+1)^2}\right).
\]
The coefficients need not be integral uniformly in $r$. Exact rational
calculation for $p=5$ gives $v_5(a^C_{r,1})=v_5(a^J_{r,1})=2,1,0,-1$
for $r=1,2,3,4$, respectively. Thus simply replacing $p$ by $p^r$ in the
first-level integral Taylor argument is invalid. The finite factor
separation above is one way to retain correct valuation bookkeeping.

\subsection{Finite calculations, not an all-level proof}
The program \texttt{probe.py} checks the original recurrences for
$p\in\{5,7,11,13,17,19,23,29,31\}$ and $1\le r\le5$ with
$p^r\le200000$, at precision $p^{4r+4}$: 36 pairs $(p,r)$ for each
conjecture. All passed. Except for $p=7$, the tested defects had exact
valuation $4r$; at $p=7$ the C defect had valuation $6$ at $r=1$ and
$4r+1$ thereafter, while the J defect had valuation $4r+1$ throughout.
The additional observed relation
\[
 \frac{\Delta_C(r)}{p^{4r}}\equiv\frac76 B_{p-3},\qquad
 \frac{\Delta_J(r)}{p^{4r}}\equiv\varepsilon^r\frac7{24}B_{p-3}
 \pmod p
\]
is an unproved pattern in this finite sample, not a theorem or an assertion of novelty. The directly computed leading digits are stored
in \texttt{probe.jsonl}; \texttt{validate.py} independently checks all 72 Bernoulli leading-digit comparisons, eight exact rational instances of the finite block identity, and the two negative-valuation examples.

\section{An exact obstruction to a direct \texorpdfstring{$(A.3)$}{(A.3)} lift}\label{loc-I-6}
\label{I-sec:a-obstruction}

Define, for odd $N$,
\[
 A(N)=\sum_{k=0}^{(N-1)/2}(4k+1)(-1)^k\frac{(1/2)_k^5}{k!^5}.
\]
The four-root perturbation in Swisher's first-level argument, with $p$
replaced by $N$, is the following rational number:
\begin{equation}\label{I-eq:a-perturbation}
 \widetilde A(N)=\sum_{k=0}^{(N-1)/2}
 (4k+1)(-1)^k\frac{(1/2)_k^5}{k!^5}
 \frac{\displaystyle\prod_{j=0}^{k-1}
       \left(1-\frac{N^4}{(2j+1)^4}\right)}
      {\displaystyle\prod_{j=1}^{k}
       \left(1-\frac{N^4}{16j^4}\right)}.
\end{equation}
It is the terminating very-well-poised ${}_6F_5(-1)$ with upper
parameters $5/4,1/2,(1-iN)/2,(1+iN)/2,(1-N)/2,(1+N)/2$
and lower parameters $1/4,1+iN/2,1-iN/2,1+N/2,1-N/2$.
Whipple's transformation and the classical gamma reflection identity
give the exact identity
\begin{equation}\label{I-eq:a-whipple}
 \widetilde A(N)=(-1)^{(N-1)/2}N\,
 {}_3F_2\!\left[
 \begin{matrix}1/2,(1-N)/2,(1+N)/2\\
                 1+iN/2,1-iN/2\end{matrix};1\right].
\end{equation}
The series on the right terminates at $(N-1)/2$. In particular it is a
\emph{perturbed} cubic hypergeometric sum; it is not the unperturbed
sum in $(H.3)$.

A direct extension of the first-level comparison would require
\[
 \widetilde A(p^r)\equiv A(p^r)\pmod{p^{5r}}.
\]
This proposed intermediate statement is false. Exact rational
arithmetic at $p=5$, $r=2$ gives
\begin{equation}\label{I-eq:a-bridge-failure}
 \widetilde A(25)-A(25)\equiv7062500\pmod{5^{10}},
 \qquad
 v_5\bigl(\widetilde A(25)-A(25)\bigr)=6.
\end{equation}
Equivalently, its quotient by $5^6$ is $2$ modulo $5$.
This is a counterexample to the proposed intermediate comparison,
\emph{not} a counterexample to $(A.3)$.

The calculation can be reproduced without complex arithmetic: if
$t_k$ is the $k$th term of \eqref{I-eq:a-perturbation} after deleting
$4k+1$, then $t_0=1$ and
\[
 \frac{t_k}{t_{k-1}}=
 -\frac{(2k-1)\bigl((2k-1)^4-N^4\bigr)}
        {2k\bigl(16k^4-N^4\bigr)}.
\]
For the eight tested pairs $p=5,13$, $r=1,2,3$, and $p=17$, $r=1,2$,
the exact valuation of $\widetilde A(p^r)-A(p^r)$ is $r+4$.
These finite observations are not asserted as a general valuation theorem.

There is a corresponding obstacle to replacing the right side of
\eqref{I-eq:a-whipple} by the unperturbed $(H.3)$ sum. Writing
\[
 H(N)=\sum_{k=0}^{(N-1)/2}\frac{(1/2)_k^3}{k!^3},
\]
the six exact tests $p=5,13$, $r=1,2,3$ give
\[
 v_p\bigl(A(p^r)-p^rH(p^r)\bigr)=r+3,
 \qquad
 v_p\bigl(\widetilde A(p^r)-p^rH(p^r)\bigr)=r+3.
\]
Thus an all-level theorem for $H$ does not by itself control the
parameter corrections required for $A$. Terms with indices divisible
by $p$ introduce lower-level denominator valuations, and their
corrections must be transported between levels before the exponent
$5r$ can be justified.


\section{Exact boundary residuals and an odd-level obstruction}\label{loc-I-7}\label{I-sec:parity-location}

The following elementary statement isolates a hypothesis needed to transfer
the jet argument of corpus [132]. It is a result about finite hypergeometric
polynomials, not a proof of any remaining Swisher conjecture.

\begin{proposition}[Boundary residual and parity of its valuation]
Let $d\geq3$, let $p>d$ be a prime with $p\equiv-1\pmod d$, and let
$h\geq1$ be an integer. Set
\[
 c_k=\left(\frac{(1/d)_k}{k!}\right)^h,\qquad
 F_N(z)=\sum_{k=0}^{N}c_kz^k,\qquad
 \theta=z\frac{d}{dz},\qquad
 \mathcal L_{d,h}=\theta^h-z(\theta+1/d)^h.
\]
For every $N\geq0$ there is the exact identity
\begin{equation}\label{I-eq:boundary-general}
 \mathcal L_{d,h}F_N(z)
   =-(N+1/d)^h c_Nz^{N+1}.
\end{equation}
If $s\geq1$ is odd and
\[
 N_- = \frac{p^s-d+1}{d},\qquad
 N_+ = \frac{(d-1)p^s-1}{d},
\]
then
\begin{align}
 v_p\!\left(\frac{(1/d)_{N_-}}{N_-!}\right)&=\frac{s-1}{2},
 &v_p\!\left((N_-+1/d)^h c_{N_-}\right)&=\frac{h(s-1)}2,
 \label{I-eq:short-valuation}\\
 v_p\!\left(\frac{(1/d)_{N_+}}{N_+!}\right)&=0,
 &v_p\!\left((N_++1/d)^h c_{N_+}\right)&=hs.
 \label{I-eq:long-valuation}
\end{align}
If $s\geq2$ is even and $N_0=(p^s-1)/d$, then
\begin{equation}\label{I-eq:even-valuation}
 v_p\!\left(\frac{(1/d)_{N_0}}{N_0!}\right)=0,
 \qquad
 v_p\!\left((N_0+1/d)^h c_{N_0}\right)=hs.
\end{equation}
\end{proposition}

\begin{proof}
The coefficient recurrence
$k^hc_k=(k-1+1/d)^hc_{k-1}$ cancels every interior coefficient
in $\mathcal L_{d,h}F_N$, giving \eqref{I-eq:boundary-general}.

For the valuation of $(1/d)_N/N!$, count the factors divisible by
$t=p^e$ in $\prod_{j=0}^{N-1}(dj+1)$ and subtract $\lfloor N/t\rfloor$.
The least nonnegative solution to $dj+1\equiv0\pmod t$ is
\[
 j_e=\begin{cases}
 (t-1)/d,&e\text{ even},\\
 ((d-1)t-1)/d,&e\text{ odd}.
 \end{cases}
\]
Writing $N=qt+R$ with $0\leq R<t$, the difference of the two counts
is $1$ if $j_e<R$, and $0$ otherwise.

For odd $s$ and $N=N_-$, the remainder for $1\leq e<s$ is
\[
 R=\begin{cases}
 (d-1)(t-1)/d,&e\text{ even},\\
 (t-d+1)/d,&e\text{ odd}.
 \end{cases}
\]
Thus the count difference is $1$ exactly when $e$ is even.
There is no contribution for $e\geq s$, since both $N_-<p^s$ and
$d(N_--1)+1=p^s-2d+2<p^s$. This proves the first identity
of \eqref{I-eq:short-valuation}. Moreover
$N_-+1/d=(p^s-d+2)/d$ is a $p$-adic unit because $p>d$.
The second identity follows.

For $N=N_+$, its remainder modulo $p^e$ equals $j_e$ for
$1\leq e\leq s$, so every count difference is zero. No larger
power contributes: $N_+<p^s$, and the numerator factors are
less than $(d-1)p^s<p^{s+1}$. Also
$N_++1/d=(d-1)p^s/d$, of valuation $s$.
This proves \eqref{I-eq:long-valuation}.

For even $s$ and $N=N_0$, the remainder again equals $j_e$ for
$1\leq e\leq s$; there are no larger contributions. Since
$N_0+1/d=p^s/d$, \eqref{I-eq:even-valuation} follows.
\end{proof}

\paragraph{Two exact examples.}
For $d=3$, $p=5$, and $s=3$, the two cutoffs are $N_-=41$ and
$N_+=83$. Their top-coefficient ratios have valuations $1$ and $0$,
respectively. The boundary residual coefficients therefore have
valuations $h$ and $3h$; for cubic sums these are $3$ and $9$.
For $d=4$, $p=7$, and $s=3$, the cutoffs are $N_-=85$ and
$N_+=257$, with the same respective ratio valuations $1$ and $0$
and boundary residual valuations $h$ and $3h$.

\paragraph{Consequence for transferring the method.}
For the unweighted hypergeometric polynomial underlying the short odd
cutoffs, the residual has exact valuation $h(s-1)/2$, not $hs$.
Consequently one cannot import the $p^{3s}$ residual hypothesis of
corpus [132] unchanged into the cubic cases $d=3,4$.
The statement also records the boundary valuation for exponents $h=4$
or $h=6$ if those polynomials occur in a proposed argument. It does
not identify their differential modules or establish a Frobenius
congruence for them.

The longer odd cutoff restores the boundary precision $hs$. A possible
repair must therefore establish both the appropriate Frobenius theorem
for that polynomial and a sufficiently precise relation between its
long and short truncations. Neither follows from the valuation
calculation alone. The linear weight and specialization in each
Swisher sum also have to be accounted for separately.


\section{Verification and the unresolved statements}\label{loc-I-8}
All finite checks use integer modular arithmetic with explicit $p$-factor extraction or exact rational arithmetic. No floating-point tolerance is used. The supplied \texttt{verify\_reductions.py} independently checks 26 K factorizations, eight K prefactor cases, four F diagonal recurrences, five F decompositions, the explicit false F termwise bound, the false A perturbation bridge, and five parameter sets in the parity lemma. Every check passes.

The larger C/J program checks 36 pairs $(p,r)$ for each conjecture, and all 72 original congruences pass. Its separate validator checks the finite block reduction in eight cases and confirms the displayed nonintegral Taylor coefficient. These tests support the calculations; they do not prove the remaining universal statements.

The exact new work target is now especially explicit for K and F:
\[
 U(p^r)-U(p^{r-1})\in p^{3r}\mathbb Z_p,\quad
 \text{equivalently}\quad T(p^r)-T(p^{r-1})\in p^{3r}\mathbb Z_p,
\]
and
\[
 B_r-p^2 B_{r-2}\in p^r\mathbb Z_p\quad \text{for the F boundary sum and odd } r\ge3.
\]
Neither is established by the preceding product or diagonal congruences. For C and J the displayed finite three-term obstruction must vanish as a sum; its components need not separately have the target valuation. For A, a parameter-deformation correction must be transported between levels. For D, E, and G this investigation has not supplied the required all-level cancellation. These precise limitations are part of the conclusions, not omitted proof steps.

\begin{thebibliography}{9}
\bibitem{I-Swisher2015} H. Swisher, \emph{On the supercongruence conjectures
of van Hamme}, Research in the Mathematical Sciences 2 (2015), 18.
Published Table 2. \url{https://doi.org/10.1186/s40687-015-0037-6}.
\bibitem{I-OZ} R. Osburn and W. Zudilin, \emph{On the (K.2)
supercongruence of Van Hamme}, J. Math. Anal. Appl. 433 (2016), 706--711.
\url{https://arxiv.org/abs/1504.01976}.
\bibitem{I-Cai2002} T. Cai, \emph{A congruence involving the quotients of
Euler and its applications (I)}, Acta Arithmetica 103.4 (2002), 313--320,
Corollary 3. \url{https://www.impan.pl/shop/publication/transaction/download/product/83800}.
\bibitem{I-Long2011} L. Long, \emph{Hypergeometric evaluation identities
and supercongruences}, Pacific J. Math. 249 (2011), 405--418.
\url{https://msp.org/pjm/2011/249-2/pjm-v249-n2-p09-p.pdf}.
\bibitem{I-Guo} V. J. W. Guo, \emph{$q$-Analogues of Dwork-type
supercongruences}, J. Math. Anal. Appl. 487 (2020), 124022.
\url{https://arxiv.org/abs/1910.07551}.
\bibitem{I-JKwithdrawn} A. Jana and G. Kalita, \emph{Proof of a
supercongruence conjecture of (F.3) of Swisher using the WZ-method},
arXiv:2011.02762v1 (2020), withdrawn in 2024 with the comment
``There are mistakes in the article.'' Only its independently checked
WZ identity is used here. \url{https://arxiv.org/abs/2011.02762}.
\bibitem{I-Corpus} VMRC, manuscripts [131] and [132], snapshot recorded
in Section 1. \url{https://github.com/QEEDDD/Verifiable-Mathematical-Reasoning-Corpus}.
\end{thebibliography}

\clearpage\endgroup

%%%%%%%%%%%%%%%%%%%%%%%%%%%%%%%%%%%%%%%%%%%%%%%%%%%%%%%%%%%%%%%%%%%%%%%%%%
% II | Swisher_Eight_Targets_Reductions.tex
%%%%%%%%%%%%%%%%%%%%%%%%%%%%%%%%%%%%%%%%%%%%%%%%%%%%%%%%%%%%%%%%%%%%%%%%%%
\clearpage\begingroup
\SetPaper{II}{Cubic scaling and the proofs of C.3 and K.3}
\RestoreCounters
\counterwithout*{equation}{section}\renewcommand{\theequation}{\arabic{equation}}
\PaperTitle{Cubic Scaling and Two Swisher Supercongruences\\[5pt]\large Complete Proofs of (C.3) and (K.3), with Six Remaining Targets}{Research working note}{10 November 2026}

\begin{abstract}
We prove a cubic scaling theorem for finite products of an arbitrary
power series $f\in1+X\mathbb Z_p[[X]]$, for every prime $p>3$.
As a consequence, a truncated generalized hypergeometric sum whose
parameters are all proportional to $p^r$ is stable modulo $p^{3r}$
when its scale is divided by $p$. An exact WZ factorization and Cai's
prime-power Morley congruence then prove Swisher's published (K.3),
with its full modulus $p^{4r}$, for every $p>5$ and $r\ge1$.
A second exact WZ identity converts (C.3) to a weighted convolution;
alternating unit-prefix sums and summation by parts prove its full
$p^{4r}$ recurrence for every $p>3$ and $r\ge1$.
We retain the rigorous reductions for (A.3), (D.3), (E.3),
(F.3), (G.3), and (J.3), and sharpen the remaining J target to a
normalized parity-product stability statement. Those six targets
are not claimed solved. The scaling proofs are elementary and
do not depend on the geometric claims of corpus [131] or [132].
\end{abstract}
\noindent\textbf{Status.} This revision contains complete proofs of
(C.3) and (K.3), including every all-level estimate used in their reductions.
The argument has undergone independent internal checks; it has not
undergone external peer review. Historical priority is not asserted.
Classical identities and published input theorems retain their attribution.
The corpus manuscripts were consulted as proposed mathematical inputs;
this note does not certify all of their geometric arguments.

\begin{LocalContents}
\LocalEntry{loc-II-1}{1}{Scope and outcome}
\LocalEntry{loc-II-2}{2}{Cubic scaling for cutoff product sums}
\LocalEntry{loc-II-3}{3}{A complete proof of Swisher's \texorpdfstring{$(K.3)$}{(K.3)}}
\LocalEntry{loc-II-4}{4}{A complete proof of the C.3 recurrence}
\LocalEntry{loc-II-5}{5}{What the corpus mechanisms do and do not transfer}
\LocalEntry{loc-II-6}{6}{An exact boundary reduction of the odd branch of (F.3)}
\LocalEntry{loc-II-7}{7}{The C.3 and J.3 deformation reductions}
\LocalEntry{loc-II-8}{8}{An exact obstruction to a direct \texorpdfstring{$(A.3)$}{(A.3)} lift}
\LocalEntry{loc-II-9}{9}{Exact boundary residuals and an odd-level obstruction}
\LocalEntry{loc-II-10}{10}{Verification and the six remaining targets}
\end{LocalContents}
\clearpage

\section{Scope and outcome}\label{loc-II-1}
The excluded targets are (B.3), (H.3), (I.3), and (L.3).
We use Swisher's published Table~2~\cite{II-Swisher2015}, rather than
the initial arXiv version: the signs in (K.3) were corrected in the
published table. Throughout, $v_p(p)=1$ and $(a)_k$ denotes a rising
factorial. The natural non-exceptional ranges are stated with each result.
We do not infer a current exhaustive literature classification from a
failure to locate a later proof.

\begin{center}\small
\begin{tabular}{@{}lp{105mm}@{}}\toprule
Target & Outcome of this investigation\\\midrule
A.3 & The ordinary target is $p^{5r}$. A direct first-level Whipple
perturbation comparison fails already at $(p,r)=(5,2)$; the original
conjecture is not refuted.\\
C.3 & Complete proof at $p^{4r}$, for $p>3$ and every $r\ge1$.\\
J.3 & The endpoint multiplier is eliminated, leaving a single normalized parity-product recurrence. This remains unproved.\\
D.3 & The ordinary target is $p^{6r}$. No all-level closure is obtained;
the boundary lemma also applies to its underlying sixth-power polynomial.\\
E.3 & No completion of the odd supersingular branch is obtained.
Its short cutoff does not satisfy the cubic residual precision used in [132].\\
F.3 & The odd-branch diagonal recurrence is proved for every level;
an explicit boundary recurrence remains.\\
G.3 & No complete closure is obtained. The fourth-power boundary
calculation identifies a transfer hypothesis that still needs repair.\\
K.3 & The entire prefactor recurrence is proved at modulus $p^{4r}$;
the residual recurrence is now also proved at $p^{3r}$, giving a complete proof of (K.3).\\\bottomrule
\end{tabular}
\end{center}

The input corpus snapshot is commit
\texttt{91c0fa1f0872f070b001249d0ab7c47e3499a907} of
\url{https://github.com/QEEDDD/Verifiable-Mathematical-Reasoning-Corpus}.
We examined the TEX sources of [131], \emph{Excellent Frobenius Lifts
and Supercongruences Unified}, and [132], \emph{Swisher H3 Unified
Ordinary and Supersingular}. They are not modified here.

\section{Cubic scaling for cutoff product sums}\label{loc-II-2}
\begin{theorem}[Cubic cutoff-product scaling]\label{II-thm:product-scaling}
Let $p>3$ be prime, $\kappa\in\mathbb Z_p$, and
$f(X)\in1+X\mathbb Z_p[[X]]$. For $L$ a positive power of $p$,
including $L=1$, define
\[
\mathcal S(L)=1+\kappa\sum_{n=1}^{L}\frac{L}{n}
 \prod_{j=1}^{n-1}f(L/j).
\]
Every evaluation of $f$ is at a point of $p\mathbb Z_p$, so all
products converge and all summands are $p$-integral. For every
$r\geq1$,
\[
\mathcal S(p^r)\equiv\mathcal S(p^{r-1})\pmod{p^{3r}}.
\]
The theorem requires no value $f(1)$.
\end{theorem}
\begin{proof}
Set $N=p^r$, $M=p^{r-1}$ and write
\[
f(X)=1+\delta X+\varepsilon X^2+O(X^3),\qquad
\eta=2\varepsilon-\delta^2,
\qquad g(X)=\frac{f(X)-1-\delta X}{X^2}\in\mathbb Z_p[[X]].
\]
If $j<L$, then $v_p(L/j)\geq1$, so $f(L/j)$ is a well-defined
unit and $L/n$ is integral for $1\leq n\leq L$. This proves the
integrality assertion.

Define only the products
\[
P_m=\prod_{j=1}^{m}f(M/j)\quad(0\leq m<M),\qquad P_0=1,
\]
and set $t_m=\kappa(M/m)P_{m-1}$ for $1\leq m\leq M$.
There is no need to define $P_M$; the series $f$ need not converge
at $1$. We have
\[
v_p(t_m)\geq r-1-v_p(m),
\qquad
v_p(P_m-P_{m-1})\geq r-1-v_p(m)\quad(1\leq m<M).
\]

For $0\leq m\leq M$, introduce the complete unit-prefix sums
\[
H_m=\sum_{\substack{1\leq j<pm\\p\nmid j}}j^{-1},\qquad
J_m=\sum_{\substack{1\leq j<pm\\p\nmid j}}j^{-2},
\qquad H_0=J_0=0.
\]
For $m\geq1$, let $e_m=1+v_p(m)$. Complete unit residues modulo
$p^{e_m}$ give $J_m\in p^{e_m}\mathbb Z_p$: multiplication by a
unit whose square is not $1$ modulo $p$ proves the vanishing of the
sum of inverse squares. Pairing $j$ with $pm-j$ gives
$2H_m\equiv-pmJ_m\pmod{p^{2e_m}}$. Hence
\begin{equation}\label{II-eq:gauss-prefix}
v_p(H_m)\geq2e_m,\qquad v_p(J_m)\geq e_m.
\end{equation}

Let $R_m$ be the product of $f(N/j)$ over $1\leq j<pm$, $p\nmid j$.
Let $R_{m,b}$ be the corresponding product with $j<pm+b$, for
$0\leq m<M$ and $1\leq b<p$. The expansion
\[
f(x)=1+\delta x+\varepsilon x^2+O(x^3),
\qquad 2\varepsilon-\delta^2=\eta,
\]
and finite multiplication give
\begin{align*}
R_m&\equiv1+\delta NH_m+
 \frac{N^2}{2}(\delta^2H_m^2+\eta J_m)\pmod{N^3},\\
R_{m,b}&\equiv1+\delta N\left(H_m+
 \sum_{a=1}^{b-1}\frac1{pm+a}\right)\pmod{N^2}.
\end{align*}
All omitted coefficients are integral, since the indices in these
unit products are prime to $p$.

Separate the terms of $\mathcal S(N)$ according to divisibility of
$n$ by $p$. The defining product decomposition gives exactly
\[
\mathcal S(N)-\mathcal S(M)
=\sum_{m=1}^{M}t_m(R_m-1)
+\kappa N\sum_{m=0}^{M-1}P_m
 \sum_{b=1}^{p-1}\frac{R_{m,b}}{pm+b}.
\]
Writing $C_m=H_m^2-J_m$, the inner sum satisfies
\[
\sum_{b=1}^{p-1}\frac{R_{m,b}}{pm+b}
\equiv H_{m+1}-H_m+\frac{\delta N}{2}(C_{m+1}-C_m)
\pmod{N^2}.
\]
It follows that
\begin{align}\label{II-eq:gauss-block}
\mathcal S(N)-\mathcal S(M)\equiv{}&
\delta N\sum_{m=1}^{M}t_mH_m
+\kappa N\sum_{m=0}^{M-1}P_m(H_{m+1}-H_m)\notag\\
&+\frac{N^2}{2}\sum_{m=1}^{M}t_m(\delta^2H_m^2+\eta J_m)
+\frac{\kappa\delta N^2}{2}
 \sum_{m=0}^{M-1}P_m(C_{m+1}-C_m)
\pmod{N^3}.
\end{align}
We now estimate these contributions without dividing by $\delta$.

The exact identity $f(x)-1=\delta x+x^2g(x)$ implies,
for $1\leq m<M$,
\[
\delta t_m
=\kappa(P_m-P_{m-1})-\kappa W_m,
\qquad
W_m=(M/m)^2g(M/m)P_{m-1}.
\]
The argument $M/m$ belongs to $p\mathbb Z_p$, and
$v_p(W_m)\geq2(r-1-v_p(m))$. Summation by parts, with the
$m=M$ term kept separate, turns the first line of
\eqref{II-eq:gauss-block} into
\[
\kappa N(1+\delta)P_{M-1}H_M
-\kappa N\sum_{m=1}^{M-1}W_mH_m.
\]
The boundary term belongs to $p^{3r}\mathbb Z_p$, since
$v_p(H_M)\geq2r$. Every remaining product satisfies
\[
v_p(W_mH_m)\geq2(r-1-v_p(m))+2(1+v_p(m))=2r.
\]
Thus the entire first line vanishes modulo $p^{3r}$.

For the third contribution, \eqref{II-eq:gauss-prefix} gives
\[
v_p\bigl(t_m(\delta^2H_m^2+\eta J_m)\bigr)
\geq(r-1-v_p(m))+(1+v_p(m))=r
\]
for all $1\leq m\leq M$. It therefore vanishes after multiplication
by $N^2/2$.

Finally, $v_p(C_m)\geq1+v_p(m)$ for $m\geq1$, and
\[
\sum_{m=0}^{M-1}P_m(C_{m+1}-C_m)
=P_{M-1}C_M-
 \sum_{m=1}^{M-1}(P_m-P_{m-1})C_m.
\]
The boundary term is divisible by $p^r$, as is every summand on the
right. The fourth contribution is therefore also divisible by
$p^{3r}$. This proves the theorem. The argument includes $r=1$:
then $M=1$, and the sums indexed by $1\leq m<M$ are empty.
\end{proof}


\begin{corollary}[Hypergeometric parameter scaling]
\label{II-thm:gauss-scaling}
Let $p>3$ and $d\geq1$. Let $\alpha_1,\ldots,\alpha_d\in\mathbb Z_p$
and $\gamma_1,\ldots,\gamma_{d-1}\in\mathbb Z_p^\times$. Put
\[
\mathcal U(L)=\sum_{n=0}^{L}
\frac{\prod_{i=1}^{d}(-\alpha_i L)_n}
 {n!\prod_{j=1}^{d-1}(-\gamma_j L)_n}.
\]
For every $r\geq1$,
\[
\mathcal U(p^r)\equiv\mathcal U(p^{r-1})\pmod{p^{3r}}.
\]
The products indexed by $j$ are empty when $d=1$.
\end{corollary}
\begin{proof}
Take
\[
\kappa=-\frac{\prod_i\alpha_i}{\prod_j\gamma_j},\qquad
f(X)=\frac{\prod_i(1-\alpha_iX)}{\prod_j(1-\gamma_jX)}.
\]
For $1\leq n\leq L$, cancelling the nonzero integer factors gives
exactly
\[
\frac{\prod_i(-\alpha_iL)_n}
 {n!\prod_j(-\gamma_jL)_n}
=\kappa\frac{L}{n}\prod_{k=1}^{n-1}f(L/k).
\]
Every denominator is nonzero, since its factors are $-\gamma_jL$
and $k-\gamma_jL$ with $0<k<L$; their valuations are respectively
$v_p(L)$ and $v_p(k)<v_p(L)$. The formula also holds if a numerator
parameter is zero. Apply Theorem~\ref{II-thm:product-scaling}.
\end{proof}

In particular, the corollary applies to
$\sum_{n=0}^{L}(-\alpha L)_n(-\beta L)_n/(n!(-\gamma L)_n)$
for arbitrary $\alpha,\beta\in\mathbb Z_p$ and
$\gamma\in\mathbb Z_p^\times$. No restriction on $\gamma-1$ or
$\gamma-\alpha-\beta$ occurs.

\begin{corollary}[The K residual]\label{II-cor:K-residual}
For $p>3$ and $r\geq1$, put
\[
U(L)=\sum_{n=0}^{L}\frac{(-L/2)_n^2}{n!(-3L/2)_n}.
\]
Then $U(p^r)\equiv U(p^{r-1})\pmod{p^{3r}}$.
\end{corollary}
\begin{proof}
Apply Corollary~\ref{II-thm:gauss-scaling} with
$\alpha_1=\alpha_2=1/2$ and $\gamma_1=3/2$.
\end{proof}

\paragraph{A classical specialization.}
For $f(X)=1+X$ and $\kappa=1$, the sum is
$\mathcal S(N)=\binom{2N}{N}$, by the hockey-stick binomial identity.
Theorem~\ref{II-thm:product-scaling} therefore contains the prime-power
central-binomial instance of Jacobsthal's congruence. The proof uses
the classical separation of indices divisible by $p$ and unit-index
harmonic cancellation; comparable tools occur in
Osburn--Sahu--Straub~\cite{II-OSS}. No historical novelty of these tools
is claimed.
\section{A complete proof of Swisher's \texorpdfstring{$(K.3)$}{(K.3)}}\label{loc-II-3}
\label{II-sec:k-reduction}

For each odd positive integer $N$, put $h_N=(N-1)/2$ and define
\begin{align}
 K(N)&=\sum_{k=0}^{h_N}(42k+5)
               \frac{(1/2)_k^3}{k!^3 64^k},\label{II-eq:k-sum}\\
 B(N)&=\frac{12\displaystyle\prod_{j=0}^{N-1}(N+2j)}
                {2^{3N-2}(h_N!)^2},\label{II-eq:k-prefactor}\\
 U(N)&=\sum_{n=0}^{N}
                \frac{(-N/2)_n^2}{n!(-3N/2)_n}.\label{II-eq:k-residual}
\end{align}
All three quantities are rational numbers. In particular,
$K(1)=5$, $B(1)=6$, and $U(1)=5/6$. We use the positive sign in
Swisher's published Table~2: with $\chi_p=(-1)^{(p-1)/2}$, the target is
\begin{equation}\label{II-eq:k-target}
 K(p^r)\equiv\chi_p pK(p^{r-1})\pmod{p^{4r}},\qquad p>5.
\end{equation}
The first arXiv version has an erroneous additional minus sign in this
entry; the published version corrects it.

\begin{proposition}[Exact factorization]\label{II-prop:k-factorization}
For every odd positive integer $N$,
\begin{equation}\label{II-eq:k-factorization}
 K(N)=B(N)U(N).
\end{equation}
\end{proposition}
\begin{proof}
Use the Guillera WZ pair employed by Osburn and Zudilin in their proof
of $(K.2)$, namely
\begin{align*}
 F(n,k)&=(84n^2-56nk+4k^2+52n-12k+5)\\
 &\hspace{8mm}\cdot
 \frac{(-1)^k(1/2)_n(1/2)_{n+k}(1/2)_{n-k}^{2}}
 {2^{4n}(n!)^2(2n-k+1)!},\\
 G(n,k)&=
 \frac{64n^2(-1)^k(1/2)_n(1/2)_{n+k-1}(1/2)_{n-k}^{2}}
 {2^{4n}(n!)^2(2n-k)!}.
\end{align*}
Negative-index Pochhammer symbols are defined by reciprocals, and a
reciprocal factorial at a negative integer is zero. The pair satisfies
\[
 F(n,k-1)-F(n,k)=G(n+1,k)-G(n,k).
\]
Moreover $F(n,0)=(42n+5)(1/2)_n^3/(64^n n!^3)$.
Summing first over $0\le n\le h_N$ and then over $1\le k\le N$
gives, exactly as in Osburn--Zudilin's equation~(2.6),
\[
 K(N)=F(h_N,N)+\sum_{k=1}^{N}G(h_N+1,k).
\]
This telescoping identity requires only that $N$ be odd, not prime.
Direct simplification gives
\[
 F(h_N,N)=B(N),\qquad
 G(h_N+1,k)=B(N)
 \frac{(-N/2)_{N+1-k}^{2}}
 {(N+1-k)!(-3N/2)_{N+1-k}}.
\]
Changing the index to $n=N+1-k$ supplies the terms $n=1,\ldots,N$
in \eqref{II-eq:k-residual}; the boundary value supplies $n=0$.
\end{proof}

\begin{theorem}[All-level congruence for the prefactor]
\label{II-thm:k-prefactor}
For every prime $p\ge5$ and integer $r\ge1$,
\begin{equation}\label{II-eq:k-prefactor-congruence}
 B(p^r)\equiv\chi_p pB(p^{r-1})\pmod{p^{4r}},
 \qquad v_p(B(p^r))=r.
\end{equation}
\end{theorem}
\begin{proof}
Set $N=p^r$, $M=p^{r-1}$ and $t=(N-M)/2$, and write
\[
 V=\prod_{\substack{1\le j\le h_N\\p\nmid j}}j,
 \quad O=\prod_{\substack{1\le a<N\\a\ \mathrm{odd},\ p\nmid a}}a,
 \quad W=\prod_{\substack{0\le j<N\\p\nmid j}}(N+2j).
\]
Removing the factors divisible by $p$ from the numerator and the
factorials in \eqref{II-eq:k-prefactor} yields the exact equality
\begin{equation}\label{II-eq:k-prefactor-unit}
 \frac{B(N)}{pB(M)}=\frac{W}{2^{3(N-M)}V^2}.
\end{equation}
Pairing the factors in $W$ about $2N$ gives another exact equality,
\begin{equation}\label{II-eq:k-pairing}
 W=(-1)^t O^2
 \prod_{\substack{1\le a<N\\a\ \mathrm{odd},\ p\nmid a}}
       \left(1-\frac{4N^2}{a^2}\right).
\end{equation}
We claim that the last product is $1$ modulo $p^{3r}$.
Indeed, the inverse-square sum over all units modulo $N$ vanishes
modulo $N$. To see this, choose a unit $g$ with $g^2\not\equiv1\pmod p$;
multiplication by $g$ permutes those units and forces the sum to be
annihilated by the unit $g^{-2}-1$. Pairing $j$ with $N-j$ then shows
that the inverse-square sum over $1\le j\le h_N$, $p\nmid j$, also
vanishes modulo $N$. The sum over odd units $a<N$ is the full sum
minus one quarter of this half-interval sum, so it vanishes as well.
Consequently the linear term in the product in \eqref{II-eq:k-pairing}
is divisible by $p^{3r}$; every term involving two or more selected
factors is divisible by $p^{4r}$.

Let $C_N=\binom{N-1}{h_N}$. Its $p$-free quotient is
\begin{equation}\label{II-eq:k-morley-quotient}
 \frac{C_N}{C_M}=2^t\frac{O}{V}.
\end{equation}
Cai's prime-power extension of Morley's congruence
\cite[Corollary~3]{II-Cai2002} states that
\[
 \frac{C_N}{C_M}\equiv\chi_p4^{N-M}\pmod{p^{3r}}.
\]
Since $(-1)^t=\chi_p$, substituting \eqref{II-eq:k-pairing} and
\eqref{II-eq:k-morley-quotient} into \eqref{II-eq:k-prefactor-unit} gives
\[
 \frac{B(N)}{pB(M)}
 \equiv\chi_p\left(\frac{C_N}{C_M}\right)^2
                  2^{-4(N-M)}
 \equiv\chi_p\pmod{p^{3r}}.
\]
The exact quotient in \eqref{II-eq:k-prefactor-unit} is a $p$-adic unit.
Starting with $B(1)=6$, induction therefore gives $v_p(B(M))=r-1$.
Multiplication by $pB(M)$ raises the precision to $p^{4r}$ and proves
both assertions.
\end{proof}

\begin{proposition}[Integrality and the residual criterion]
\label{II-prop:k-equivalence}
For $p>5$ and $r\ge0$, $U(p^r)\in\mathbb Z_p^\times$ and
$U(p^r)\equiv5/6\pmod p$. For every $r\ge1$, \eqref{II-eq:k-target}
holds if and only if
\begin{equation}\label{II-eq:k-residual-congruence}
 \boxed{\quad U(p^r)\equiv U(p^{r-1})\pmod{p^{3r}}.\quad}
\end{equation}
\end{proposition}
\begin{proof}
For $N=p^r$ and $1\le n\le N$, write the summand in $U(N)$ as
\begin{equation}\label{II-eq:k-summand-product}
 f(N,n)=-\frac{N}{6n}
 \prod_{j=1}^{n-1}
 \frac{(1-N/(2j))^2}{1-3N/(2j)}.
\end{equation}
For $j<N$, $v_p(N/j)\ge1$, so every factor in this product is a
$p$-adic unit congruent to $1$ modulo $p$. It follows that
$v_p(f(N,n))=r-v_p(n)\ge0$. Only $n=N$ survives modulo $p$ among
$n\ge1$, and that term is $-1/6$ modulo $p$. Thus
$U(N)\equiv1-1/6=5/6\pmod p$.

Finally the exact identity \eqref{II-eq:k-factorization} gives
\begin{align*}
 K(N)-\chi_p pK(M)
 &=B(N)\bigl(U(N)-U(M)\bigr)\\
 &\quad+\bigl(B(N)-\chi_p pB(M)\bigr)U(M).
\end{align*}
The second term belongs to $p^{4r}\mathbb Z_p$ by
Theorem~\ref{II-thm:k-prefactor} and the integrality just proved.
The valuation of $B(N)$ is exactly $r$, proving the equivalence.
\end{proof}

\begin{theorem}[Swisher's (K.3)]\label{II-thm:K-complete}
For every prime $p>5$ and every integer $r\ge1$,
\[
\boxed{\displaystyle
\sum_{k=0}^{(p^r-1)/2}(42k+5)\frac{(1/2)_k^3}{64^k(k!)^3}
\equiv (-1)^{(p-1)/2}p
\sum_{k=0}^{(p^{r-1}-1)/2}(42k+5)\frac{(1/2)_k^3}{64^k(k!)^3}
\pmod{p^{4r}}.}
\]
\end{theorem}
\begin{proof}
Corollary~\ref{II-cor:K-residual} establishes exactly the residual
congruence in Proposition~\ref{II-prop:k-equivalence}, for every level.
More explicitly, set $N=p^r$ and $M=N/p$. In the exact identity
\[
K(N)-\chi_p pK(M)
=B(N)(U(N)-U(M))+(B(N)-\chi_p pB(M))U(M),
\]
the first term belongs to $p^{4r}\mathbb Z_p$ because
$v_p(B(N))=r$ and $U(N)-U(M)\in p^{3r}\mathbb Z_p$.
The second belongs to the same ideal by
Theorem~\ref{II-thm:k-prefactor} and $U(M)\in\mathbb Z_p$.
This includes $r=1$, where $K(1)=5$ and $U(1)=5/6$.
\end{proof}

\subsection{A further harmonic reduction of the K residual}
This subsection uses the $U(N)$ of the preceding K reduction. Put
\[
 P_j(N)=\prod_{k=1}^j\frac{(1-N/(2k))^2}{1-3N/(2k)},\qquad P_0(N)=1.
\]
For every odd positive integer $N$,
\begin{equation}\label{II-eq:Uharmonic}
 U(N)=\frac43-\frac12P_{N-1}(N)+\frac13T(N),\qquad
 T(N)=N^2\sum_{n=1}^{N-1}\frac{P_{n-1}(N)}{n^2(1-3N/(2n))}.
\end{equation}
Indeed the $n$th summand of $U$ is $-NP_{n-1}(N)/(6n)$, and
\[
 f(x)=\frac{(1-x/2)^2}{1-3x/2},\qquad
 f(x)-1=\frac{x}{2}+\frac{x^2}{1-3x/2}.
\]
Summing the resulting difference $P_n-P_{n-1}$ for $1\le n<N$
and then adding the $n=N$ term proves \eqref{II-eq:Uharmonic}.

\begin{proposition}
For $p>5$, $N=p^r$, and $M=p^{r-1}$,
\[
 P_{N-1}(N)\equiv P_{M-1}(M)\pmod{p^{3r}}.
\]
Consequently (K.3) is equivalent to
$T(p^r)\equiv T(p^{r-1})\pmod{p^{3r}}$.
The last congruence is now a consequence of Theorem~\ref{II-thm:K-complete}.
\end{proposition}
\begin{proof}
Remove the factors indexed by multiples of $p$ from the quotient of
the two products. What remains is $\prod_{1\le k<N,\ p\nmid k}f(N/k)$.
On the full unit residue set modulo $N$,
\[
 H_2=\sum_{\substack{1\le k<N\\p\nmid k}}k^{-2}\in p^r\mathbb Z_p,
 \qquad
 H_1=\sum_{\substack{1\le k<N\\p\nmid k}}k^{-1}\in p^{2r}\mathbb Z_p.
\]
For the first claim multiply the residues by a unit whose square is
not $1$ modulo $p$. For the second, pair $k$ with $N-k$ and expand
modulo $N^2$, obtaining $2H_1\equiv-NH_2\pmod{N^2}$.
Now
\[
 \log f(x)=\sum_{\ell\ge1}\frac{3^\ell-2}{\ell2^\ell}x^\ell.
\]
The first two terms summed over the units lie in $p^{3r}\mathbb Z_p$
by these harmonic bounds. Every later term has valuation at least
$\ell r-v_p(\ell)\ge3r$. The convergent logarithm and exponential
therefore prove the product congruence. Equation~\eqref{II-eq:Uharmonic}
and the K equivalence above prove the last assertion, since $3$ is a unit.
\end{proof}


% Standalone proof fragment. This is a reduction, not a proof of (F.3).
\section{A complete proof of the C.3 recurrence}\label{loc-II-4}\label{II-sec:C-complete}
For an odd positive integer $N$, write $n=(N-1)/2$ and
\[
 b_k=\frac{(1/2)_k}{k!}=\frac{\binom{2k}{k}}{4^k},\qquad
 C(N)=\sum_{k=0}^{n}(4k+1)b_k^4.
\]
The following finite identity is exact over the rational numbers:
\begin{equation}\label{II-Cnew:conv}
 C(N)=2N^2b_n^2\sum_{l=0}^{n}\frac{b_l b_{N-l}}{N-2l}.
\end{equation}
To prove it, put $B(s,k)=b_s^2b_{s+k}b_{s-k}$ for $s\ge k\ge0$,
and set $B(s,k)=0$ for $k>s$. Define
\[
 F(s,k)=(4s+1)B(s,k),\qquad
 G(s,k)=\frac{16s^2(s+k)}{(2k-1)(2s+2k-1)}B(s,k),\quad k\ge1.
\]
Direct cancellation of successive factorials verifies
\[
 F(s,k-1)-F(s,k)=G(s+1,k)-G(s,k).
\]
The boundary $s=k-1$ also follows by direct substitution, and $G(0,k)=0$.
Sum first over $0\le s<h$, then over $1\le k\le h$, where
$h=(N+1)/2$. The left side is $C(N)$, because $F(s,h)=0$.
On reversing $k=h-l$, the right side becomes
\[
 \sum_{k=1}^hG(h,k)
 =8h^2b_h^2\sum_{l=0}^{h-1}\frac{b_l b_{N-l}}{N-2l}
 =2N^2b_n^2\sum_{l=0}^{n}\frac{b_l b_{N-l}}{N-2l},
\]
since $b_h=N b_n/(N+1)$. This proves \eqref{II-Cnew:conv}.

Define
\[
 Q_k(N)=\prod_{j=1}^{k}(1-2N/j)^{(-1)^j},\qquad
 R_l(N)=\frac{Q_{2l-1}(N)}{1-N/(2l)}\quad(l\ge1),
\]
\begin{equation}\label{II-Cnew:U}
 U_C(N)=1+\frac{N^2}{2}\sum_{l=1}^{(N-1)/2}\frac{R_l(N)}{l^2}.
\end{equation}
Using
\[
 \frac{b_l b_{N-l}}{b_N}=
 \frac{\binom Nl^2}{\binom{2N}{2l}},
\]
and expanding the two binomial products, \eqref{II-Cnew:conv} becomes
\begin{equation}\label{II-Cnew:factor}
 C(N)=A_C(N)U_C(N),\qquad A_C(N)=2N b_{(N-1)/2}^{2}b_N.
\end{equation}
For $N=p^r$, all factors in $Q_k(N)$ and $R_l(N)$ are $p$-adic units,
because their indices are strictly below $N$.

\subsection{Two harmonic estimates}
Let $p>3$. If $L$ is positive and $a=v_p(L)$, the unit square sum
\[
 \sum_{\substack{1\le j\le L\\p\nmid j}}j^{-2}
\]
is divisible by $p^a$. This follows by grouping into complete residue
systems modulo $p^a$; the sum of the inverse squares of the units modulo
$p^a$ is zero modulo $p^a$. In particular, for
\[
 H_l=\sum_{\substack{1\le j\le l\\p\nmid j}}j^{-2},
 \qquad e_l=\max\{v_p(l),v_p(l+1),v_p(2l+1)\},
\]
one has $v_p(H_l)\ge e_l$. The cases $p^{e_l}\mid l$ and
$p^{e_l}\mid l+1$ use the complete-residue estimate (the endpoint
$l+1$ is omitted when divisible by $p$). In the third case, reflection
$j\mapsto2l+1-j$ gives
$2H_l\equiv\sum_{1\le j<2l+1,p\nmid j}j^{-2}\equiv0\pmod{p^{e_l}}$.
In particular $H_{(p^r-1)/2}\in p^r\mathbb Z_p$.

A second estimate concerns an even integer $L$ with $a=v_p(L)\ge1$:
\[
 S_i(L)=\sum_{\substack{1\le j<L\\p\nmid j}}\frac{(-1)^j}{j^i}.
\]
Then $S_2(L)\in p^a\mathbb Z_p$ and
$S_1(L)\in p^{2a}\mathbb Z_p$. For the first assertion write
$j=u+tp^a$, with $1\le u<p^a$ a unit. The number $L/p^a$ is even,
so the alternating signs cancel modulo $p^a$ for every residue $u$.
For the second, reflection about $L/2$ preserves the sign and gives
\[
 2S_1(L)\equiv-LS_2(L)\pmod{p^{2a}}.
\]

\subsection{The residual factor is stable to three times the level}
We prove
\begin{equation}\label{II-Cnew:scaling}
 U_C(p^r)\equiv U_C(p^{r-1})\pmod{p^{3r}}.
\end{equation}
Put $N=p^r$, $M=N/p$, and $n=(N-1)/2$.
First consider the part of \eqref{II-Cnew:U} with $p\nmid l$.
A direct quotient gives
\[
 \frac{R_{l+1}(N)}{R_l(N)}=
 \frac{1-N/l}{1-2N/(2l+1)}
 \frac{1-N/(2l)}{1-N/(2l+2)}.
\]
All four factors are units and hence
\[
 R_{l+1}(N)-R_l(N)\in p^{r-e_l}\mathbb Z_p.
\]
Summation by parts therefore proves
\[
 \sum_{\substack{1\le l\le n\\p\nmid l}}
 \frac{R_l(N)}{l^2}
 =R_n(N)H_n-\sum_{l=1}^{n-1}H_l(R_{l+1}(N)-R_l(N))
 \in p^r\mathbb Z_p.
\]
Multiplication by $N^2/2$ shows that all these indices together
contribute zero modulo $p^{3r}$.

For the remaining indices write $l=pb$, where
$1\le b\le(M-1)/2$. Oddness of $p$ preserves the parity of the
factors whose index is divisible by $p$. Thus, exactly,
\[
 R_{pb}(N)=R_b(M)\mathcal A_b,\qquad
 \mathcal A_b=\prod_{\substack{1\le j<2pb\\p\nmid j}}
 (1-2N/j)^{(-1)^j}.
\]
Set $a=v_p(2pb)=1+v_p(b)\le r-1$. Expansion to second order,
with an integral remainder because every denominator $j$ is a unit,
gives
\[
 \mathcal A_b\equiv
 1-2NS_1(2pb)+2N^2\bigl(S_1(2pb)^2-S_2(2pb)\bigr)
 \pmod{p^{3r}}.
\]
The preceding harmonic estimates show
$\mathcal A_b-1\in p^{r+2a}\mathbb Z_p$.
Meanwhile $M^2/(2b^2)$ has valuation $2(r-a)$. Therefore each
individual difference
\[
 \frac{N^2}{2(pb)^2}R_{pb}(N)
 -\frac{M^2}{2b^2}R_b(M)
\]
is divisible by $p^{3r}$. Together with the unit-index part this proves
\eqref{II-Cnew:scaling}, also for $r=1$ when the divisible-index part is empty.

\subsection{The prefactor and the original C.3 statement}
Let $B_N=\binom{N-1}{(N-1)/2}$ and $D_N=\binom{2N}{N}$.
Cai's prime-power Morley congruence and Jacobsthal's congruence give,
for $N=p^r$, $M=N/p$,
\[
 \frac{B_N}{B_M}\equiv\varepsilon4^{N-M}\pmod{p^{3r}},
 \qquad \frac{D_N}{D_M}\equiv1\pmod{p^{3r}},
 \qquad\varepsilon=(-1)^{(p-1)/2}.
\]
The first is Cai\textquotesingle s Corollary 3~\cite{II-Cai2002}.
The second also follows from Theorem 3 of that paper with $u=2,v=1$.
Since $b_{(N-1)/2}=B_N/2^{N-1}$ and $b_N=D_N/4^N$, these imply
\[
 \frac{A_C(N)}{pA_C(M)}\equiv1\pmod{p^{3r}}.
\]
Both $A_C(N)/N$ and $U_C(N)$ are $p$-integral units. Combining this
prefactor congruence with \eqref{II-Cnew:scaling} and $v_p(A_C(N))=r$ gives
\[
 \boxed{C(p^r)\equiv pC(p^{r-1})\pmod{p^{4r}}\qquad(p>3,\ r\ge1).}
\]
This is exactly the original half-cutoff C.3 recurrence. For clarity,
the final implication is
\[
C(N)-pC(M)=A_C(N)(U_C(N)-U_C(M))
 +(A_C(N)-pA_C(M))U_C(M).
\]
Both terms belong to $p^{4r}\mathbb Z_p$, by the established valuations.
The argument includes $r=1$, with $A_C(1)=U_C(1)=C(1)=1$.
The restriction $p>3$ is essential: $C(3)-3=-27/16$ is not divisible
by $3^4$.
\section{What the corpus mechanisms do and do not transfer}\label{loc-II-5}
The modular cusp theorem in [131] (label \texttt{B:thm:main}) assumes
weight $w\in\{1,2\}$ and yields precision $p^{(w+1)r}$.
The higher-dimensional standard-family theorem yields $p^{2r}$.
An analytic evaluation at a unit argument requires its own integral
evaluation map, fixed-point condition, and multiplier identification.
These hypotheses cannot be replaced by a formal substitution.

There is an explicit obstruction already proved in [131]
(label \texttt{E:thm:sh-cube}). Let
\[
 F_d(u)=\sum_{k\ge0}\binom{2k}{k}^du^k,\quad K=p^r,\qquad
 E_r=F_d(u)F_{K/p}(\sigma(u))-F_d(\sigma(u))F_K(u),
\]
where $F_L$ is truncated at $k<L$ and $\sigma$ is the excellent lift
in that theorem. For $d\ge4$ and $p>d$,
\[
 [u^{K+1}]E_r\equiv2^{2d-1}d(d-3)p^{2r}\pmod{p^{2r+1}}.
\]
The right side is nonzero at this precision. Thus the full-series
congruence cannot simply be raised from $2r$ to $4r$ or $5r$.
For example, $d=4,p=7,r=1$ gives
\[
 [u^8]E_1=27435582641602576,\qquad v_7([u^8]E_1)=2.
\]
Even applying $1+4u\,d/du$ leaves an obstruction: its coefficient is
multiplied by $4K+5$, a $7$-adic unit. This excludes a naive global
derivative upgrade; it does not exclude additional cancellation at a
particular argument. This obstruction is from [131], not a new result
of the present note.

For clarity, there is an exact weighted transfer identity.
Write $D_r=F_r-\rho F_{r-1}\circ\sigma$,
$\theta=t\,d/dt$, and $\mu=\theta\sigma/\sigma$.
At a valid analytic fixed point $t_0$, if
\[
 c\theta\rho(t_0)=\rho(t_0)(\mu(t_0)-1),\qquad
 \lambda=\rho(t_0)\mu(t_0),
\]
then the product and chain rules give
\[
 (1+c\theta)F_r(t_0)-\lambda(1+c\theta)F_{r-1}(t_0)
 =(1+c\theta)D_r(t_0).
\]
Differentiation preserves an established integral modulus, but does
not itself increase it. The remaining weighted residual needs a
separate estimate. Likewise, the [132] Taylor-transport argument
requires the precise residual order of its differential equation;
Section~\ref{II-sec:parity-location} below calculates what happens at
the short odd cutoffs of other families.


\section{An exact boundary reduction of the odd branch of (F.3)}\label{loc-II-6}
Let $p\equiv3\pmod4$ be a prime. For odd $r\ge1$ put
\[
 N_r=(p^r-3)/4,\qquad
 S_r=\sum_{n=0}^{N_r}(-1)^n(8n+1)\frac{(1/4)_n^3}{(n!)^3}.
\]
Define
\[
 D_r=(8N_r+1)\frac{(1/4)_{2N_r}}{(N_r!)^2},\qquad
 B_r=\sum_{k=1}^{N_r}G(N_r+1,k),
\]
where an empty sum is zero and
\[
 G(n,k)=(-1)^{n+k}4
 \frac{(1/4)_n^2(1/4)_{n+k-1}}
 {(n-1)!^2(n-k)!(1/4)_k^2}.
\]
\begin{proposition}
For every such $p,r$, $S_r=D_r+B_r$. For odd $r\ge3$,
\[
 v_p(D_r)=r-1,\qquad D_r\equiv p^2D_{r-2}\pmod{p^r}.
\]
Consequently the odd branch of Swisher's original (F.3) is equivalent to
\[
 \boxed{B_r\equiv p^2B_{r-2}\pmod{p^r}}\qquad(r\ge3\text{ odd}).
\]
This boxed boundary recurrence is not proved here.
\end{proposition}
\begin{proof}
Set $a=1/4$ and
\[
 F(n,k)=(-1)^{n+k}(8n+1)
 \frac{(a)_n^2(a)_{n+k}}{(n!)^2(n-k)!(a)_k^2}.
\]
Use $1/m!=0$ for negative integers $m$, and set $G(0,k)=0$.
The WZ identity
\[
 F(n,k-1)-F(n,k)=G(n+1,k)-G(n,k)
\]
is an exact rational identity. Indeed its four quotients after division
by $F(n,k)$ are
\[
 \frac{F(n,k-1)}{F(n,k)}=-\frac{(a+k-1)^2}{(a+n+k-1)(n-k+1)},
\]
\[
 \frac{G(n+1,k)}{F(n,k)}=-\frac{(a+n)^2}{(2n+a)(n-k+1)},\qquad
 \frac{G(n,k)}{F(n,k)}=\frac{n^2}{(a+n+k-1)(2n+a)}.
\]
Substitution proves the identity, including the boundary cases by the
zero convention. Summing first over $0\le n\le N$, then over
$1\le k\le N$, gives
\[
 \sum_{n=0}^NF(n,0)=F(N,N)+\sum_{k=1}^NG(N+1,k).
\]
Since $F(N,N)=(8N+1)(a)_{2N}/(N!)^2$, this proves the decomposition.
The WZ pair appears as Lemma 3.1 in Jana--Kalita,
\emph{Proof of a supercongruence conjecture of (F.3) of Swisher using
the WZ-method}, arXiv:2011.02762v1. That article was withdrawn for
errors; only the identity, independently verified above, is used here.

Write $R_r=(1/4)_{2N_r}/(N_r!)^2$. Let $\Gamma_p$ be Morita's gamma
function. Extracting the factors divisible by $p$ twice from the rising
factorial and the factorial gives the exact formula
\[
 \frac{R_r}{p^2R_{r-2}}=
 \left(\frac{p^{r-2}}2-\frac54\right)
 \left(\frac{p^{r-2}}2-\frac14\right)
 \frac{\Gamma_p(p^r/2-5/4)\Gamma_p(p^{r-1}/2+1/4)}
 {\Gamma_p(1/4)\Gamma_p(3/4)
  \Gamma_p((p^r+1)/4)^2\Gamma_p((p^{r-1}+3)/4)^2}.
\tag{*}
\]
For completeness, in $(1/4)_{(p^r-3)/2}$ the first extraction leaves
$p^{(p^{r-1}-1)/2}(3/4)_{(p^{r-1}-1)/2}$. The second leaves
$p^{(p^{r-2}+1)/2}(1/4)_{(p^{r-2}+1)/2}$. Removing the final two
factors of this last rising factorial gives the two prefactors in (*).
In $N_r!$, the two factorials left are indexed by
$(p^{r-1}-1)/4$ and $N_{r-2}$. Comparing the powers of $p$ gives
exactly $p^2$. The gamma signs in the numerator multiply to $+1$
because $(p^r-3)/2$ and $(p^{r-1}-1)/2$ are even; denominator signs
are squared.

All the factors on the right of (*) are $p$-adic units: $p\ne5$.
Modulo $p$, its value is
\[
 \frac5{16}\frac{\Gamma_p(-5/4)}
 {\Gamma_p(1/4)^2\Gamma_p(3/4)^3}=1.
\]
The equality follows from
$\Gamma_p(-5/4)=\frac{16}5\Gamma_p(3/4)$ and
$\Gamma_p(1/4)^2\Gamma_p(3/4)^2=1$.
These are the gamma functional equation and reflection formula.
The base $R_1$ is a unit: $(p-3)/2<(3p-1)/4$ is below the first
$p$-divisible factor in $(1/4)_j$, and $N_1<p$.
Induction using (*) proves $v_p(R_r)=r-1$ and
$R_r\equiv p^2R_{r-2}\pmod{p^r}$. Finally,
$8N_r+1=2p^r-5$ is a unit and is congruent to
$2p^{r-2}-5$ modulo $p$. Multiplying proves both claims for $D_r$.
Subtracting the diagonal recurrence from the exact decomposition proves
the stated equivalence.
\end{proof}

\paragraph{A precise obstruction to deleting the boundary.}
The withdrawn article's Lemma 3.2 asserts
$v_p(G((p^r+1)/4,k))\ge3(r-1)/2$ throughout the relevant range.
For $p=7,r=3,k=13$, this is false:
\[
 v_7(G(86,13))=1,
 \qquad \frac{G(86,13)}7\equiv34\pmod{49}.
\]
For an entirely integer valuation check, let
$A(n)=\sum_{j=0}^{n-1}v_7(4j+1)$ and
$C(n)=v_7(n!)$. Then
\[
 A(86)=14,\quad A(98)=16,\quad A(13)=3,\quad
 C(85)=13,\quad C(73)=11,
\]
so $v_7G=2(14)+16-2(13)-11-2(3)=1$.
Direct rational arithmetic gives
\[
 D_3\equiv245,\qquad B_3\equiv49,\qquad S_3\equiv294\pmod{343}.
\]
Thus the boundary is essential, even though the original conjecture is
satisfied in this test. These computations diagnose the invalid termwise
bound, and do not constitute a proof of the new boundary recurrence.

% Exact finite identity and the remaining F.3 target. No closure is claimed.
\subsection{A normalized finite identity for the remaining odd (F.3) branch}
For $N\ge0$, define
\[
 S_N=\sum_{k=0}^N(-1)^k(8k+1)\frac{(1/4)_k^3}{(k!)^3},\qquad
 D_N=(8N+1)\frac{(1/4)_{2N}}{(N!)^2}.
\]
Then the following identity is exact:
\begin{equation}\label{II-eq:F-normal-form}
 S_N=D_NV_N,\qquad
 V_N=\sum_{n=0}^N\frac{(-N-1/4)_n^2}{n!(-2N-1/4)_n}.
\end{equation}
To prove it, use the verified WZ identity of the preceding reduction,
including its last boundary term, to write
$S_N=\sum_{k=1}^{N+1}G(N+1,k)$. Its final term is
$G(N+1,N+1)=D_N$. Reversing the index gives, for $0\le n\le N$,
\[
 \frac{G(N+1,N+1-n)}{D_N}
 =\frac{(-N-1/4)_n^2}{n!(-2N-1/4)_n}.
\]
This follows directly from
$(a)_{m-n}/(a)_m=(-1)^n/(1-a-m)_n$, so summing proves
\eqref{II-eq:F-normal-form} without any limiting argument.

Let $p\equiv3\pmod4$, $r\ge3$ odd, $L=p^r$, $M=p^{r-2}$,
and $N_X=(X-3)/4$. The already proved diagonal statement says
\[
 v_p(D_{N_L})=r-1,\qquad
 \frac{D_{N_L}}{p^2D_{N_M}}\in1+p\mathbb Z_p.
\]
Consequently a sufficient remaining theorem, stated with all its
integrality obligations, is
\begin{equation}\label{II-eq:F-remaining}
 V_{N_{p^s}}\in\mathbb Z_p\quad(s\ge1\text{ odd}),\qquad
 V_{N_L}-V_{N_M}\in p\mathbb Z_p\quad(r\ge3\text{ odd}).
\end{equation}
It is also necessary for the full odd branch of (F.3): its first-level
$S_{N_p}$ and $D_{N_p}$ are integral and $D_{N_p}$ is a unit, so the
original recurrence inductively gives the displayed integrality and
normalized congruence. Thus \eqref{II-eq:F-remaining}, jointly over all odd
levels, is equivalent to that branch. It remains unproved here.

\paragraph{The precise parameter orbit.}
In the normalized sum the upper parameter is
$a_L=1/2-L/4$ and the lower parameter is $c_L=5/4-L/2$.
For Dwork's dash $x'=(x+t)/p$, $0\le t<p$, $x+t\in p\mathbb Z_p$,
two dash steps give
\[
 a_L\longmapsto\frac12-\frac{L}{4p}
     \longmapsto\frac12-\frac{M}{4}=a_M,
\]
\[
 c_L\longmapsto\frac34-\frac{L}{2p}
     \longmapsto\frac14-\frac{M}{2}=c_M-1.
\]
The representatives used are $(p-1)/2$ at both upper steps,
$(3p-5)/4$ at the first lower step, and $(p-3)/4$ at the second.
The integer shift in the lower parameter cannot be discarded.

There is an exact contiguous relation that accounts for this shift.
For $a=-N-1/4$, $c=-2N-1/4$ and
$t_n=(a)_n^2/[n!(c)_n]$, the coefficient recurrence implies
\[
 \frac34\sum_{n=0}^N n t_n+a^2V_N=\frac1{16}t_N.
\]
Indeed $n(c+n-1)t_n=(a+n-1)^2t_{n-1}$; summing and subtracting
the shifted sums leaves the terminal contribution.
Therefore
\begin{equation}\label{II-eq:F-contiguous}
 \sum_{n=0}^{N_M}\frac{(a_M)_n^2}{n!(c_M-1)_n}
 =-\frac{(M+1)^2}{3(1-2M)}V_{N_M}
   +\frac{t_{N_M}}{3(1-2M)}.
\end{equation}
This finite relation is valid as a rational identity. It does not by
itself prove a Frobenius congruence for the sums.

\paragraph{Why direct Lucas reduction does not close the problem.}
Although the desired normalized sums are experimentally integral,
their summands need not be integral. For instance, at $p=7,r=3$,
$N_L=85$, the minimum valuation of a summand in $V_{85}$ is $-1$.
A termwise reduction modulo $p$ is therefore invalid. A proof must
establish cancellation of those negative-valuation terms before using
any finite-field reduction or Dwork comparison. The pure scaling
lemma for parameters all proportional to $p^r$ does not apply directly
to $a_L,c_L$, whose nonzero constant parts and lower-parameter shift
are displayed above.
\section{The C.3 and J.3 deformation reductions}\label{loc-II-7}
For $p>3$, put $N=p^r$, $n_r=(N-1)/2$, and
\[
 c_k=(4k+1)\frac{\binom{2k}{k}^{4}}{256^k},\qquad
 j_k=(6k+1)\frac{\binom{2k}{k}^{3}}{256^k},\quad
 C_r=\sum_{k=0}^{n_r}c_k,\quad J_r=\sum_{k=0}^{n_r}j_k.
\]
Here $C_0=J_0=1$, and $\varepsilon=(-1)^{(p-1)/2}$.
The respective targets are
$C_r-pC_{r-1}\in p^{4r}\mathbb Z_p$ and
$J_r-\varepsilon pJ_{r-1}\in p^{4r}\mathbb Z_p$.
Guo's published all-level result has modulus $p^{3r}$ for both sums
(\url{https://arxiv.org/abs/1910.07551}).

For $d=2,4$, define the finite rational product
\[
 R_d(k,x)=\prod_{a=0}^{k-1}
 \frac{1-x^2/(2a+1)^2}{1-x^2/[d(a+1)]^2}.
\]
Dougall's terminating very-well-poised summation and the terminating
identity used in Long's Lemma 4.2, respectively, give for every positive
odd integer $N$ the exact identities
\begin{align}
 \sum_{k=0}^{(N-1)/2}c_k R_2(k,N)&=N,\\
 \sum_{k=0}^{(N-1)/2}j_k R_4(k,N)&=(-1)^{(N-1)/2}N.
\end{align}
The second identity is recorded at
\url{https://msp.org/pjm/2011/249-2/pjm-v249-n2-p09-p.pdf},
pp.~414--415; replacing the letter $p$ there by an odd integer does not
change the terminating evaluation. The first is the specialization
$a=b=1/2$, $c=(1+N)/2$, and terminating parameter $-(N-1)/2$ of
Dougall's formula. Its right-hand side is
\[
 \frac{(3/2)_{(N-1)/2}((1-N)/2)_{(N-1)/2}}
 {(1)_{(N-1)/2}(1-N/2)_{(N-1)/2}}=N.
\]
These exact classical evaluations are not proofs of the desired
undeformed recurrences.

\subsection{An exact finite separation of the factors divisible by \texorpdfstring{$p$}{p}}
The following reduction avoids any unjustified integral Taylor expansion.
Write $k=pt+e$, $0\le e<p$, put $h=(p-1)/2$, $M=p^{r-1}$,
and $\delta_e=1$ if $e>h$, and $0$ otherwise. Splitting the products
according to divisibility by $p$ proves exactly
\[
 R_d(k,N)=R_d(t,M)
 \left(1-\frac{M^2}{(2t+1)^2}\right)^{\delta_e}U_d(k,N),
\]
where
\[
 U_d(k,N)=
 \frac{\prod_{0\le a<k,\ p\nmid 2a+1}(1-N^2/(2a+1)^2)}
 {\prod_{1\le a\le k,\ p\nmid a}(1-N^2/(da)^2)}.
\]
Every denominator in this last product is a $p$-adic unit. Consequently,
with
\[
 W_d(k)=\sum_{\substack{0\le a<k\\p\nmid 2a+1}}\frac1{(2a+1)^2}
 -\frac1{d^2}\sum_{\substack{1\le a\le k\\p\nmid a}}\frac1{a^2},
\]
one has the rigorous, uniform expansion
\[
 U_d(k,N)\equiv1-N^2W_d(k)\pmod{p^{4r}}.
\]
The other factors above are integral for $k\le n_r$. In the last block,
$t=(M-1)/2$, only $e\le h$ occurs, so the potentially zero extra factor
is never included. This proves that this expansion can safely be
multiplied by all the remaining factors and summed.

Here is an entirely finite expression for the defect. For C its vanishing at the target precision follows from Section~\ref{II-sec:C-complete}; for J it remains open in this note.
Use $(w,d,\eta)=(c,2,1)$ or $(j,4,\varepsilon)$ and set
\[
 L_t=R_d(t,M),\quad z_t=\frac{M^2}{(2t+1)^2},\quad
 E_t=\{e:0\le e<p,\ pt+e\le n_r\},
\]
\[
 V_t=\sum_{e\in E_t}w_{pt+e},\qquad
 B_t=\sum_{\substack{e\in E_t\\e>h}}w_{pt+e},\qquad
 H_t=\sum_{e\in E_t}w_{pt+e}(1-z_t)^{\delta_e}W_d(pt+e).
\]
Then the exact terminating identities at levels $r$ and $r-1$, followed
by the preceding expansion, prove
\begin{align}
 \sum_{k=0}^{n_r}w_k-\eta p\sum_{k=0}^{n_{r-1}}w_k
 \equiv{}&\sum_{t=0}^{n_{r-1}}(1-L_t)(V_t-\eta p w_t)
 +\sum_{t=0}^{n_{r-1}}L_t z_t B_t\notag\\
 &+N^2\sum_{t=0}^{n_{r-1}}L_tH_t\pmod{p^{4r}}.
\end{align}
Indeed the evaluated deformation at level $r$ is $\eta^rN$, the
lower-level evaluation multiplied by $\eta p$ is the same, and the
$k$-block of the higher-level evaluation is
$L_t(V_t-z_tB_t-N^2H_t)$ modulo $p^{4r}$. Rearrangement gives the display.
All quantities in this congruence are $p$-integral. Thus either original
conjecture is equivalent to the vanishing modulo $p^{4r}$ of its displayed
finite obstruction. Its vanishing for C now follows from the convolution proof above; its vanishing for J remains unproved. In particular,
none of its three summands may be suppressed without a separate bound.
For example, at $(p,r)=(5,2)$ their valuations are $(7,7,7)$ for C
and $(6,6,6)$ for J, while the valuation of their total is $8$ in both
cases. At $(5,3)$ the three valuations are $(9,9,10)$ for C and
$(10,10,10)$ for J, while both totals have valuation $12$. These exact
rational checks exhibit the cross-term cancellation that a proof must
control; a termwise $4r$ bound is already false in these examples.

\subsection{Taylor formulation and its precision limitation}
For comparison put
\[
 a^{X}_{r,j}=[x^{2j}]\sum_{k=0}^{n_r}w_kR_d(k,x),
 \qquad X=C,J.
\]
For each fixed $r$, this expansion converges at $x=p^r$: every denominator
has $p$-adic valuation at most $r-1$, so $x^2/[d(a+1)]^2\in p^2\mathbb Z_p$.
One obtains the exact convergent identity
\[
 \Delta_X(r)=-\sum_{j\ge1}p^{2rj}
 \left(a^X_{r,j}-\eta p^{1-2j}a^X_{r-1,j}\right).
\]
For $r=0$, take $a_{0,0}=1$ and $a_{0,j}=0$ for $j>0$.
The quadratic coefficient is explicitly
\[
 a^X_{r,1}=\sum_{k=0}^{n_r}w_k
 \left(\frac1{d^2}H_k^{(2)}-
 \sum_{a=0}^{k-1}\frac1{(2a+1)^2}\right).
\]
The coefficients need not be integral uniformly in $r$. Exact rational
calculation for $p=5$ gives $v_5(a^C_{r,1})=v_5(a^J_{r,1})=2,1,0,-1$
for $r=1,2,3,4$, respectively. Thus simply replacing $p$ by $p^r$ in the
first-level integral Taylor argument is invalid. The finite factor
separation above is one way to retain correct valuation bookkeeping.

\subsection{Finite calculations, not an all-level proof}
The program \texttt{probe.py} checks the original recurrences for
$p\in\{5,7,11,13,17,19,23,29,31\}$ and $1\le r\le5$ with
$p^r\le200000$, at precision $p^{4r+4}$: 36 pairs $(p,r)$ for each
conjecture. All passed. Except for $p=7$, the tested defects had exact
valuation $4r$; at $p=7$ the C defect had valuation $6$ at $r=1$ and
$4r+1$ thereafter, while the J defect had valuation $4r+1$ throughout.
The additional observed relation
\[
 \frac{\Delta_C(r)}{p^{4r}}\equiv\frac76 B_{p-3},\qquad
 \frac{\Delta_J(r)}{p^{4r}}\equiv\varepsilon^r\frac7{24}B_{p-3}
 \pmod p
\]
is not proved here. The J relation is a specialization of Sun\textquotesingle s 2019 conjecture, quoted by Mao--Sun~\cite{II-MaoSun}; see the attribution below. The directly computed leading digits are stored
in \texttt{probe.jsonl}; \texttt{validate.py} independently checks all 72 Bernoulli leading-digit comparisons, eight exact rational instances of the finite block identity, and the two negative-valuation examples.

\subsection{A WZ boundary form for the J sum}
For a positive odd integer $N$, write $n=(N-1)/2$ and
\[
 J(N)=\sum_{a=0}^{n}(6a+1)\frac{\binom{2a}{a}^3}{256^a}.
\]
The classical WZ pair used by Mao--Sun~\cite[Section 3]{II-MaoSun} is
$F(a,k)=(6a-2k+1)B(a,k)$ and
$G(a,k)=16a^2B(a,k)/(2a+2k-1)$, where
\[
 B(a,k)=\frac{(1/2)_a(1/2)_{a+k}(1/2)_{a-k}}
 {4^a(a!)^2(a-k)!(1/2)_k}.
\]
The convention is $F=G=0$ for $a<k$. Direct cancellation of adjacent
factors verifies
$F(a,k-1)-F(a,k)=G(a+1,k)-G(a,k)$.
Summing first over $0\le a\le n$ and then over $1\le k\le n+1$ gives
\[
 J(N)=\sum_{k=1}^{n+1}G(n+1,k).
\]
Put
\[
 E_N=\frac{N}{2^{3N-3}}\binom{2N-1}{N-1}\binom{N-1}{(N-1)/2},
 \qquad
 T_N=\sum_{l=0}^{(N-1)/2}
 \frac{(1/2)_l(-N/2)_l}{l!(1/2-N)_l}.
\]
Then the exact identity
\[
 \boxed{J(N)=E_NT_N}
\]
follows by reversing $k=n+1-l$. Indeed $G(n+1,n+1)=E_N$ and the ratio
of successive reversed terms is
\[
 \frac{(2l+1)(N-2l)}{(2N-2l-1)(2l+2)}.
\]
This is the ratio of the displayed ${}_2F_1$ summands.

For $N=p^r$, $M=N/p$, and $p>3$, both binomial factors in $E_N/N$
are $p$-adic units by Lucas's theorem. Let
\[
 U_{p,r}=\frac{E_N}{pE_M}\in\mathbb Z_p^\times.
\]
Consequently the original J.3 statement is exactly equivalent to
\[
 T_N-\varepsilon U_{p,r}^{-1}T_M\in p^{3r}\mathbb Z_p,
 \qquad \varepsilon=(-1)^{(p-1)/2}.
\]
The known $p^{3r}$ result for $J$ gives only $p^{2r}$ precision in this
last display. No $p^{3r}$ proof for this transformed sum is established
here. The boundary identity combines the endpoint and WZ correction
before estimating them, and lowers the hypergeometric order, but does
not by itself supply the missing precision.

\subsection{Prior source of the refined J conjecture}
The leading Bernoulli digit observed computationally is a special case
of an existing stronger conjecture. Mao--Sun, arXiv:2109.09877, p.~2,
quote Z.-W. Sun's 2019 Conjecture 22: for every positive odd $m$,
\[
 \frac{16^{m-1}}{(pm)^4\binom{m-1}{(m-1)/2}^{3}}
 \bigl(J(pm)-\varepsilon pJ(m)\bigr)
 \equiv\varepsilon\frac7{24}B_{p-3}\pmod p.
\]
Their paper proves the case $m=1$. With $m=p^{r-1}$, Lucas's theorem
and Fermat's theorem turn this conjecture into
\[
 \frac{J(p^r)-\varepsilon pJ(p^{r-1})}{p^{4r}}
 \equiv\varepsilon^r\frac7{24}B_{p-3}\pmod p.
\]
Thus the observed all-level Bernoulli pattern should be attributed to
that conjecture, and should not be described as a new conjecture.

\subsection{An exact parity-product embedding}
The half-length sum can be embedded in a full-length product sum, but
the resulting factors depend on parity. Define
\[
 f_{\rm odd}(x)=\frac1{1-2x},\qquad f_{\rm even}(x)=1-x,
 \qquad P_k(N)=\prod_{j=1}^{k}f_{j\bmod2}(N/j).
\]
For $0\le k<N=p^r$, every argument $N/j$ belongs to $p\mathbb Z_p$,
so these products are integral units. Duplication of the factors gives
\[
 \frac{(1/2)_l(-N/2)_l}{l!(1/2-N)_l}
 =-\frac{N}{2l}P_{2l-1}(N)\qquad(l\ge1).
\]
Since $P_{2l}-P_{2l-1}=-(N/(2l))P_{2l-1}$, one obtains
\[
 \boxed{T_N=\sum_{k=0}^{N-1}(-1)^kP_k(N).}
\]
This has exactly $N$ terms. Under a split $j=pm+b$ with odd $p$, the
factors indexed by multiples of $p$ reproduce the same parity pattern
at $M=N/p$. The unit factors still have alternating linear
coefficients $-1$ and $2$, however, so the constant-function
cutoff-product scaling theorem cannot be applied without proving a
separate parity-dependent version. No such extension is asserted here.
\subsection{The J endpoint multiplier can be eliminated}
Put $B_N=\binom{N-1}{(N-1)/2}$ and $D_N=\binom{2N-1}{N-1}$.
For $N=p^r$, $M=N/p$, and $p>3$, the prime-power form of Morley's
congruence gives
\[
 B_N/B_M\equiv\varepsilon\,4^{N-M}\pmod{p^{3r}},
 \qquad \varepsilon=(-1)^{(p-1)/2}.
\]
This is Cai's Corollary 3~\cite{II-Cai2002}.
Cai's Theorem 3 with $n=N$, $u=2$, $v=1$ gives $D_N/D_M\equiv1\pmod{p^{3r}}$.
Consequently the exact WZ prefactor
\[
 E_N=ND_NB_N/2^{3N-3}
\]
satisfies
\[
 \frac{E_N}{pE_M}\equiv\varepsilon 2^{M-N}\pmod{p^{3r}}.
\]
Both $E_N/N$ and $T_N$ are $p$-adic units. For the former this follows
from Lucas's theorem. For the latter, each term after the initial $1$
in the exact product formula below is divisible by $p$.
Thus Swisher's full J.3 recurrence is equivalent to
\[
 T_{p^r}\equiv2^{p^r-p^{r-1}}T_{p^{r-1}}\pmod{p^{3r}}.
\]
Equivalently, $2^{-p^r}T_{p^r}$ must be stable modulo $p^{3r}$.
This eliminates all gamma or factorial prefactors; the remaining
period-two product stability is still to be proved.
\subsection{A terminating Euler coefficient transform for the J residual}
For odd $N$, set $h=(N-1)/2$ and $c=1/2-N$. Define
\[
 T_N=\sum_{l=0}^{h}\frac{(1/2)_l(-N/2)_l}{l!(c)_l}.
\]
Euler's hypergeometric transformation gives the formal identity
\[
 {}_2F_1(1/2,-N/2;c;z)
 =(1-z)^{-N/2}{}_2F_1(-N,-h;c;z).
\]
The second hypergeometric series is a polynomial of degree $h$.
Since a partial sum of Taylor coefficients through degree $h$ is the
coefficient of $z^h$ after division by $1-z$, it follows that
\begin{align*}
 T_N
 &=[z^h](1-z)^{-N/2-1}{}_2F_1(-N,-h;c;z)\\
 &=\frac{(1+N/2)_h}{h!}
 {}_3F_2\!\left[\begin{matrix}-h,-h,-N\\c,c\end{matrix};1\right].
\end{align*}
Indeed, the coefficient ratio needed in the convolution is
\[
 \frac{(1+N/2)_{h-j}/(h-j)!}{(1+N/2)_h/h!}
 =\frac{(-h)_j}{(c)_j}.
\]
The resulting ${}_3F_2$ is zero-balanced: the sum of its lower
parameters equals the sum of its upper parameters. It does not meet
the usual Saalsch\"utz balance condition, which requires their
difference to be $1$. Its parameters $-h=(1-N)/2$ and $c=1/2-N$
also retain constant half shifts. Consequently this exact transform
does not by itself bring the J residual within the pure parameter
scaling theorem. No all-level congruence for this transformed sum is
claimed here.
\section{An exact obstruction to a direct \texorpdfstring{$(A.3)$}{(A.3)} lift}\label{loc-II-8}
\label{II-sec:a-obstruction}

Define, for odd $N$,
\[
 A(N)=\sum_{k=0}^{(N-1)/2}(4k+1)(-1)^k\frac{(1/2)_k^5}{k!^5}.
\]
The four-root perturbation in Swisher's first-level argument, with $p$
replaced by $N$, is the following rational number:
\begin{equation}\label{II-eq:a-perturbation}
 \widetilde A(N)=\sum_{k=0}^{(N-1)/2}
 (4k+1)(-1)^k\frac{(1/2)_k^5}{k!^5}
 \frac{\displaystyle\prod_{j=0}^{k-1}
       \left(1-\frac{N^4}{(2j+1)^4}\right)}
      {\displaystyle\prod_{j=1}^{k}
       \left(1-\frac{N^4}{16j^4}\right)}.
\end{equation}
It is the terminating very-well-poised ${}_6F_5(-1)$ with upper
parameters $5/4,1/2,(1-iN)/2,(1+iN)/2,(1-N)/2,(1+N)/2$
and lower parameters $1/4,1+iN/2,1-iN/2,1+N/2,1-N/2$.
Whipple's transformation and the classical gamma reflection identity
give the exact identity
\begin{equation}\label{II-eq:a-whipple}
 \widetilde A(N)=(-1)^{(N-1)/2}N\,
 {}_3F_2\!\left[
 \begin{matrix}1/2,(1-N)/2,(1+N)/2\\
                 1+iN/2,1-iN/2\end{matrix};1\right].
\end{equation}
The series on the right terminates at $(N-1)/2$. In particular it is a
\emph{perturbed} cubic hypergeometric sum; it is not the unperturbed
sum in $(H.3)$.

A direct extension of the first-level comparison would require
\[
 \widetilde A(p^r)\equiv A(p^r)\pmod{p^{5r}}.
\]
This proposed intermediate statement is false. Exact rational
arithmetic at $p=5$, $r=2$ gives
\begin{equation}\label{II-eq:a-bridge-failure}
 \widetilde A(25)-A(25)\equiv7062500\pmod{5^{10}},
 \qquad
 v_5\bigl(\widetilde A(25)-A(25)\bigr)=6.
\end{equation}
Equivalently, its quotient by $5^6$ is $2$ modulo $5$.
This is a counterexample to the proposed intermediate comparison,
\emph{not} a counterexample to $(A.3)$.

The calculation can be reproduced without complex arithmetic: if
$t_k$ is the $k$th term of \eqref{II-eq:a-perturbation} after deleting
$4k+1$, then $t_0=1$ and
\[
 \frac{t_k}{t_{k-1}}=
 -\frac{(2k-1)\bigl((2k-1)^4-N^4\bigr)}
        {2k\bigl(16k^4-N^4\bigr)}.
\]
For the eight tested pairs $p=5,13$, $r=1,2,3$, and $p=17$, $r=1,2$,
the exact valuation of $\widetilde A(p^r)-A(p^r)$ is $r+4$.
These finite observations are not asserted as a general valuation theorem.

There is a corresponding obstacle to replacing the right side of
\eqref{II-eq:a-whipple} by the unperturbed $(H.3)$ sum. Writing
\[
 H(N)=\sum_{k=0}^{(N-1)/2}\frac{(1/2)_k^3}{k!^3},
\]
the six exact tests $p=5,13$, $r=1,2,3$ give
\[
 v_p\bigl(A(p^r)-p^rH(p^r)\bigr)=r+3,
 \qquad
 v_p\bigl(\widetilde A(p^r)-p^rH(p^r)\bigr)=r+3.
\]
Thus an all-level theorem for $H$ does not by itself control the
parameter corrections required for $A$. Terms with indices divisible
by $p$ introduce lower-level denominator valuations, and their
corrections must be transported between levels before the exponent
$5r$ can be justified.

\section{Exact boundary residuals and an odd-level obstruction}\label{loc-II-9}\label{II-sec:parity-location}

The following elementary statement isolates a hypothesis needed to transfer
the jet argument of corpus [132]. It is a result about finite hypergeometric
polynomials, not a proof of any remaining Swisher conjecture.

\begin{proposition}[Boundary residual and parity of its valuation]
Let $d\geq3$, let $p>d$ be a prime with $p\equiv-1\pmod d$, and let
$h\geq1$ be an integer. Set
\[
 c_k=\left(\frac{(1/d)_k}{k!}\right)^h,\qquad
 F_N(z)=\sum_{k=0}^{N}c_kz^k,\qquad
 \theta=z\frac{d}{dz},\qquad
 \mathcal L_{d,h}=\theta^h-z(\theta+1/d)^h.
\]
For every $N\geq0$ there is the exact identity
\begin{equation}\label{II-eq:boundary-general}
 \mathcal L_{d,h}F_N(z)
   =-(N+1/d)^h c_Nz^{N+1}.
\end{equation}
If $s\geq1$ is odd and
\[
 N_- = \frac{p^s-d+1}{d},\qquad
 N_+ = \frac{(d-1)p^s-1}{d},
\]
then
\begin{align}
 v_p\!\left(\frac{(1/d)_{N_-}}{N_-!}\right)&=\frac{s-1}{2},
 &v_p\!\left((N_-+1/d)^h c_{N_-}\right)&=\frac{h(s-1)}2,
 \label{II-eq:short-valuation}\\
 v_p\!\left(\frac{(1/d)_{N_+}}{N_+!}\right)&=0,
 &v_p\!\left((N_++1/d)^h c_{N_+}\right)&=hs.
 \label{II-eq:long-valuation}
\end{align}
If $s\geq2$ is even and $N_0=(p^s-1)/d$, then
\begin{equation}\label{II-eq:even-valuation}
 v_p\!\left(\frac{(1/d)_{N_0}}{N_0!}\right)=0,
 \qquad
 v_p\!\left((N_0+1/d)^h c_{N_0}\right)=hs.
\end{equation}
\end{proposition}

\begin{proof}
The coefficient recurrence
$k^hc_k=(k-1+1/d)^hc_{k-1}$ cancels every interior coefficient
in $\mathcal L_{d,h}F_N$, giving \eqref{II-eq:boundary-general}.

For the valuation of $(1/d)_N/N!$, count the factors divisible by
$t=p^e$ in $\prod_{j=0}^{N-1}(dj+1)$ and subtract $\lfloor N/t\rfloor$.
The least nonnegative solution to $dj+1\equiv0\pmod t$ is
\[
 j_e=\begin{cases}
 (t-1)/d,&e\text{ even},\\
 ((d-1)t-1)/d,&e\text{ odd}.
 \end{cases}
\]
Writing $N=qt+R$ with $0\leq R<t$, the difference of the two counts
is $1$ if $j_e<R$, and $0$ otherwise.

For odd $s$ and $N=N_-$, the remainder for $1\leq e<s$ is
\[
 R=\begin{cases}
 (d-1)(t-1)/d,&e\text{ even},\\
 (t-d+1)/d,&e\text{ odd}.
 \end{cases}
\]
Thus the count difference is $1$ exactly when $e$ is even.
There is no contribution for $e\geq s$, since both $N_-<p^s$ and
$d(N_--1)+1=p^s-2d+2<p^s$. This proves the first identity
of \eqref{II-eq:short-valuation}. Moreover
$N_-+1/d=(p^s-d+2)/d$ is a $p$-adic unit because $p>d$.
The second identity follows.

For $N=N_+$, its remainder modulo $p^e$ equals $j_e$ for
$1\leq e\leq s$, so every count difference is zero. No larger
power contributes: $N_+<p^s$, and the numerator factors are
less than $(d-1)p^s<p^{s+1}$. Also
$N_++1/d=(d-1)p^s/d$, of valuation $s$.
This proves \eqref{II-eq:long-valuation}.

For even $s$ and $N=N_0$, the remainder again equals $j_e$ for
$1\leq e\leq s$; there are no larger contributions. Since
$N_0+1/d=p^s/d$, \eqref{II-eq:even-valuation} follows.
\end{proof}

\paragraph{Two exact examples.}
For $d=3$, $p=5$, and $s=3$, the two cutoffs are $N_-=41$ and
$N_+=83$. Their top-coefficient ratios have valuations $1$ and $0$,
respectively. The boundary residual coefficients therefore have
valuations $h$ and $3h$; for cubic sums these are $3$ and $9$.
For $d=4$, $p=7$, and $s=3$, the cutoffs are $N_-=85$ and
$N_+=257$, with the same respective ratio valuations $1$ and $0$
and boundary residual valuations $h$ and $3h$.

\paragraph{Consequence for transferring the method.}
For the unweighted hypergeometric polynomial underlying the short odd
cutoffs, the residual has exact valuation $h(s-1)/2$, not $hs$.
Consequently one cannot import the $p^{3s}$ residual hypothesis of
corpus [132] unchanged into the cubic cases $d=3,4$.
The statement also records the boundary valuation for exponents $h=4$
or $h=6$ if those polynomials occur in a proposed argument. It does
not identify their differential modules or establish a Frobenius
congruence for them.

The longer odd cutoff restores the boundary precision $hs$. A possible
repair must therefore establish both the appropriate Frobenius theorem
for that polynomial and a sufficiently precise relation between its
long and short truncations. Neither follows from the valuation
calculation alone. The linear weight and specialization in each
Swisher sum also have to be accounted for separately.

\section{Verification and the six remaining targets}\label{loc-II-10}
The C and K proofs are symbolic and hold for
all permitted primes and levels. Finite computation is an additional
check, not a substitute for any step of the proof.
The supplied \texttt{audit\_blocks.py} checks all complete-prefix
harmonic bounds, both exact summation-by-parts identities, and the
valuations of the linear, quadratic, and expansion-remainder terms.
All 13 cases pass, including $(p,r)=(5,4)$, $(7,4)$, and $(11,3)$;
the largest cutoff is $7^4=2401$.
The separate \texttt{verify\_scaling.py} performs 126 exact rational
tests of the general theorem, including a zero linear coefficient,
a function singular at $1$, and $\kappa=0$. It also tests the original
K.3 and its residual in 12 cases. All pass.
An independent \texttt{audit\_c.py} checks the C factorization, every
full/half-prefix harmonic bound, every alternating-prefix bound,
each divisible-index product identity, the exact summation by parts,
and the prefactor congruence in the same 13 prime/level cases.
All pass, with maximum cutoff $7^4=2401$. The C WZ identity was also
verified symbolically, and its finite convolution was checked for
every odd $N$ from $1$ to $101$.

The original \texttt{verify\_reductions.py} checks 26 exact K
factorizations, eight K prefactor cases, four F diagonal recurrences,
five F decompositions, the false F termwise bound, the false A
perturbation bridge, and five parameter sets in the parity lemma.
Every check passes. The C/J validator retains its 72 original
congruence and Bernoulli-digit comparisons and eight finite block
identity checks. These computations remain supporting evidence only
for the refined Bernoulli-digit statements and the unproved J recurrence.

For J.3 the remaining statement is now precisely
\[
2^{-p^r}T_{p^r}-2^{-p^{r-1}}T_{p^{r-1}}
\in p^{3r}\mathbb Z_p.
\]
Its factors depend on parity. Theorem~\ref{II-thm:product-scaling}
assumes one common function $f$ and does not establish this variant.
For the odd F branch, the remaining boundary target is
\[
B_r-p^2B_{r-2}\in p^r\mathbb Z_p\qquad(r\ge3\text{ odd}).
\]
For C, the earlier finite three-term obstruction now vanishes as a
consequence of Section~\ref{II-sec:C-complete}; its components need not
individually have the target valuation.
For A, parameter-deformation corrections must be transported between
levels. For D, E, and G, the required all-level cancellation has not
been supplied. Thus C.3 and K.3 are the two targets completed in this
revision; the other six remain outside the proved conclusions of this note.

\begin{thebibliography}{99}
\bibitem{II-Swisher2015} H. Swisher, \emph{On the supercongruence conjectures
of van Hamme}, Research in the Mathematical Sciences 2 (2015), 18.
Published Table 2. \url{https://doi.org/10.1186/s40687-015-0037-6}.
\bibitem{II-OZ} R. Osburn and W. Zudilin, \emph{On the (K.2)
supercongruence of Van Hamme}, J. Math. Anal. Appl. 433 (2016), 706--711.
\url{https://arxiv.org/abs/1504.01976}.
\bibitem{II-Cai2002} T. Cai, \emph{A congruence involving the quotients of
Euler and its applications (I)}, Acta Arithmetica 103.4 (2002), 313--320,
Corollary 3. \url{https://www.impan.pl/shop/publication/transaction/download/product/83800}.
\bibitem{II-Long2011} L. Long, \emph{Hypergeometric evaluation identities
and supercongruences}, Pacific J. Math. 249 (2011), 405--418.
\url{https://msp.org/pjm/2011/249-2/pjm-v249-n2-p09-p.pdf}.
\bibitem{II-Guo} V. J. W. Guo, \emph{$q$-Analogues of Dwork-type
supercongruences}, J. Math. Anal. Appl. 487 (2020), 124022.
\url{https://arxiv.org/abs/1910.07551}.
\bibitem{II-JKwithdrawn} A. Jana and G. Kalita, \emph{Proof of a
supercongruence conjecture of (F.3) of Swisher using the WZ-method},
arXiv:2011.02762v1 (2020), withdrawn in 2024 with the comment
``There are mistakes in the article.'' Only its independently checked
WZ identity is used here. \url{https://arxiv.org/abs/2011.02762}.
\bibitem{II-MaoSun} G.-S. Mao and Z.-W. Sun,
\emph{Proof of a conjectural supercongruence modulo $p^5$},
arXiv:2109.09877 (2021). \url{https://arxiv.org/abs/2109.09877}.
\bibitem{II-OSS} R. Osburn, B. Sahu, and A. Straub,
\emph{Supercongruences for sporadic sequences},
Proc. Edinburgh Math. Soc. 59 (2016), 503--518.
\url{https://doi.org/10.1017/S0013091515000255}.
\bibitem{II-Corpus} VMRC, manuscripts [131] and [132], snapshot recorded
in Section 1. \url{https://github.com/QEEDDD/Verifiable-Mathematical-Reasoning-Corpus}.
\end{thebibliography}

\clearpage\endgroup

%%%%%%%%%%%%%%%%%%%%%%%%%%%%%%%%%%%%%%%%%%%%%%%%%%%%%%%%%%%%%%%%%%%%%%%%%%
% III | Swisher_J3_and_F3_Odd_All_Levels.tex
%%%%%%%%%%%%%%%%%%%%%%%%%%%%%%%%%%%%%%%%%%%%%%%%%%%%%%%%%%%%%%%%%%%%%%%%%%
\clearpage\begingroup
\SetPaper{III}{J.3 and the odd branch of F.3}
\RestoreCounters
\counterwithout*{equation}{section}\renewcommand{\theequation}{\arabic{equation}}
\counterwithout*{theorem}{section}\renewcommand{\thetheorem}{\arabic{theorem}}
\PaperTitle{All-Level Proofs for Swisher J.3\\and the Remaining Odd Branch of F.3}{Research working manuscript}{10 November 2026}

\begin{abstract}
We prove two all-level supercongruences from Swisher's prime-power
table. For J.3, an elementary telescoping identity, a general
half-cutoff product-scaling lemma, and Cai's prime-power Morley
congruence give the original modulus $p^{4r}$ for every prime $p>3$.
For the remaining odd branch of F.3, creative microscoping with an
adapted short cutoff gives the stronger modulus $p^{3(r-1)/2}$ for
every prime $p\equiv3\pmod4$ and odd $r\ge3$, hence the required
$p^r$. We explicitly account for the parameter denominator,
cyclotomic multiplicities, and the integral specialization at $q=1$.
The E.3 odd branch and the A.3, D.3, and G.3 families remain outside
the conclusions. Exact computations accompany the proofs. This is a
research working manuscript; no external review or historical
priority is claimed.
\end{abstract}

\section{Statement and scope}\label{loc-III-1}
For a positive odd integer $N$, define
\[
 J(N)=\sum_{k=0}^{(N-1)/2}(6k+1)\frac{\binom{2k}{k}^3}{256^k}
 =\sum_{k=0}^{(N-1)/2}(6k+1)\frac{(1/2)_k^3}{(k!)^3\,4^k}.
\]
Here $(a)_0=1$, $(a)_k=a(a+1)\cdots(a+k-1)$, and $J(1)=1$.
Congruences between rational numbers are interpreted in $\Zp$.

\begin{theorem}\label{III-thm:J}
For every prime $p>3$ and every integer $r\ge1$,
\[
 \boxed{J(p^r)\equiv(-1)^{(p-1)/2}pJ(p^{r-1})
             \pmod{p^{4r}}.}
\]
\end{theorem}
This is the half-cutoff J.3 recurrence in the published table of
Swisher~\cite{III-Swisher}. The proof covers both residue classes modulo
$4$ and includes the first level $r=1$. It does not replace the
half-cutoff by a full cutoff. For this J.3 proof, the substantive imported arithmetic
result is Cai's prime-power Morley congruence~\cite[Corollary 3]{III-Cai}.
The finite identity and the residual scaling estimate are proved below.

The second main conclusion concerns the exact remaining odd branch
of F.3. Write
\[
 F_r(p)=\sum_{k=0}^{(p^r-3)/4}(-1)^k(8k+1)
                         \frac{(1/4)_k^3}{(k!)^3}
 \qquad(p\equiv3\pmod4,\ r\text{ odd}).
\]
We prove in Section~\ref{III-sec:F} that
\[
 \boxed{F_r(p)\equiv p^2F_{r-2}(p)\pmod{p^{3(r-1)/2}}
        \qquad(r\ge3\text{ odd}).}
\]
Since $3(r-1)/2\ge r$, this attains the original odd-branch target.
This includes $p=3$. Ordinary primes and even levels of F.3 are not
reproved here; the theorem addresses the remaining range specified
for this research round.

\section{The exact finite factorization}\label{loc-III-2}
For $n\ge0$ put $b_n=(1/2)_n/n!$ and
\[
 q(n,k)=\frac{(-n-1/2)_k^2}{(k!)^2},\qquad
 V_n=\sum_{k=0}^{n}q(n,k),\qquad
 B_n=\frac{(2n+1)b_n}{4^n}.
\]
We prove $B_nV_n=\sum_{k=0}^{n}(6k+1)b_k^3/4^k$.
For $n\ge1$ define
\[
 R(n,k)=\frac{4k^2(4k-6n-3)}{(2n+1)^3}.
\]
The two term ratios
\[
 \frac{q(n,k+1)}{q(n,k)}=
 \frac{(2k-2n-1)^2}{4(k+1)^2},\qquad
 \frac{q(n-1,k)}{q(n,k)}=
 \frac{(2n+1-2k)^2}{(2n+1)^2}
\]
show, by direct polynomial cancellation, that
\[
 q(n,k+1)R(n,k+1)-q(n,k)R(n,k)
 =q(n,k)-\frac{8n}{2n+1}q(n-1,k).
\]
Summing over $0\le k\le n$, and using
\[
 R(n,0)=0,\qquad q(n-1,n)=b_n^2,\qquad
 q(n,n+1)R(n,n+1)=\frac{1-2n}{2n+1}b_n^2,
\]
gives
\[
 V_n-\frac{8n}{2n+1}V_{n-1}
 =\frac{6n+1}{2n+1}b_n^2.
\]
Since $B_n/B_{n-1}=(2n+1)/(8n)$, multiplication by $B_n$
yields
\[
 B_nV_n-B_{n-1}V_{n-1}=(6n+1)\frac{b_n^3}{4^n}.
\]
The initial value is $B_0V_0=1$, proving the asserted identity by
induction. Taking $N=2n+1$ gives
\[
 J(N)=\frac{N}{4^{N-1}}\binom{N-1}{(N-1)/2}
       \sum_{k=0}^{(N-1)/2}\frac{(-N/2)_k^2}{(k!)^2}.
\]
This certificate proves the identity over the rational numbers, without
an appeal to convergence or cancellation of terminating parameters.
\section{A half-cutoff square-weighted scaling lemma}\label{loc-III-3}
Let $p>3$, $f(X)\in1+X\mathbb Z_p[[X]]$, and $\kappa\in\mathbb Z_p$.
For an odd prime power $L=p^s$ put
\[
 R_k(L)=\prod_{j=1}^{k-1}f(L/j),\qquad
 V_f(L)=1+\kappa L^2\sum_{k=1}^{(L-1)/2}\frac{R_k(L)}{k^2}.
\]
The empty sum at $L=1$ is zero. All product arguments lie in
$p\mathbb Z_p$, and all terms of $V_f(L)$ are integral. The following congruence is the general scaling theorem:
\begin{equation}\label{III-Jnew:halfscale}
 V_f(p^r)\equiv V_f(p^{r-1})\pmod{p^{3r}}\qquad(r\ge1).
\end{equation}

Write $N=p^r$, $M=N/p$, and $n=(N-1)/2$. Define
\[
 H_k=\sum_{\substack{1\le j\le k\\p\nmid j}}j^{-2}.
\]
Complete unit residue sums show $v_p(H_k)\ge v_p(k)$.
Also $H_n\in p^r\mathbb Z_p$: reflection $j\mapsto N-j$
identifies the two half sums modulo $p^r$, while their complete sum
is zero modulo $p^r$. Here the inverse-square sum on the units modulo
$p^a$ is zero modulo $p^a$ for every $a\ge1$, since $p>3$.
Since $f(X)-1\in X\mathbb Z_p[[X]]$,
\[
 v_p(R_{k+1}(N)-R_k(N))\ge r-v_p(k).
\]
Summation by parts therefore gives
\[
 \sum_{\substack{1\le k\le n\\p\nmid k}}\frac{R_k(N)}{k^2}
 =R_n(N)H_n-\sum_{k=1}^{n-1}H_k(R_{k+1}(N)-R_k(N))
 \in p^r\mathbb Z_p.
\]
After multiplication by $\kappa N^2$, the unit-index contribution
is zero modulo $p^{3r}$.

For a remaining index $k=pb$, $1\le b\le(M-1)/2$, splitting the
product into unit and divisible indices gives the exact identity
\[
 R_{pb}(N)=R_b(M)A_b,\qquad
 A_b=\prod_{\substack{1\le j<pb\\p\nmid j}}f(N/j).
\]
Set $a=v_p(pb)=1+v_p(b)\le r-1$. For $L=pb$ define
\[
 h_i(L)=\sum_{\substack{1\le j<L\\p\nmid j}}j^{-i}.
\]
Then $h_2(L)\in p^a\mathbb Z_p$ by complete residue sums, and
reflection $j\mapsto L-j$ gives
\[
 2h_1(L)\equiv-Lh_2(L)\pmod{p^{2a}},
 \qquad h_1(L)\in p^{2a}\mathbb Z_p.
\]
Write $f(X)=1+dX+eX^2+O(X^3)$. A finite multiplication, with an
integral remainder, gives
\[
 A_b\equiv1+dNh_1(L)+\frac{N^2}{2}
 \bigl(d^2h_1(L)^2+(2e-d^2)h_2(L)\bigr)
 \pmod{p^{3r}}.
\]
Consequently $A_b-1\in p^{r+2a}\mathbb Z_p$: the displayed
linear, quadratic, and remainder terms have valuations at least
$r+2a$, $2r+a$, and $3r$, respectively. Since
$v_p(M^2/b^2)=2(r-a)$, each term
\[
 \kappa\frac{N^2}{(pb)^2}R_{pb}(N)
 -\kappa\frac{M^2}{b^2}R_b(M)
\]
belongs to $p^{3r}\mathbb Z_p$. This proves
\eqref{III-Jnew:halfscale}, including $r=1$, for which the divisible-index
sum is empty.

\subsection{Application to the transformed J sum}
Take $f(X)=(1-X/2)^2$ and $\kappa=1/4$. Then
\[
 V_f(N)=\sum_{k=0}^{(N-1)/2}\frac{(-N/2)_k^2}{(k!)^2},
\]
because for $k\ge1$
\[
 \frac{(-N/2)_k^2}{(k!)^2}
 =\frac{N^2}{4k^2}\prod_{j=1}^{k-1}(1-N/(2j))^2.
\]
The exact finite transformation proved separately gives
\begin{equation}\label{III-Jnew:factor}
 J(N)=B(N)V_f(N),\qquad
 B(N)=\frac{N}{4^{N-1}}\binom{N-1}{(N-1)/2}.
\end{equation}
Cai's prime-power Morley congruence~\cite[Corollary 3]{III-Cai} implies
\[
 B(p^r)\equiv\varepsilon p B(p^{r-1})\pmod{p^{4r}},
 \qquad \varepsilon=(-1)^{(p-1)/2}.
\]
More explicitly, for $N=p^r$ and $M=p^{r-1}$, Cai's result is
\[
 \frac{\binom{N-1}{(N-1)/2}}{\binom{M-1}{(M-1)/2}}
 \equiv\varepsilon4^{N-M}\pmod{p^{3r}}.
\]
The central binomial coefficients here are units: for example,
$\binom{N-1}{h}=\prod_{j=1}^h(N-j)/j$ is a product of units,
since $v_p(N-j)=v_p(j)$ for $1\le j<N$. Thus $v_p(B(N))=r$,
and the displayed congruence is equivalent to
$B(N)/(pB(M))\equiv\varepsilon\pmod{p^{3r}}$. Combining this with \eqref{III-Jnew:halfscale} yields
\[
 \boxed{J(p^r)\equiv\varepsilon pJ(p^{r-1})
        \pmod{p^{4r}}\qquad(p>3,\ r\ge1).}
\]
The finite transformation and the scaling proof are independent of
any conjectural deformation estimate.

For clarity, the final error splits as
\[
 \begin{aligned}
 J(N)-\varepsilon pJ(M)
 ={}&B(N)\bigl(V_f(N)-V_f(M)\bigr)\\
 &+\bigl(B(N)-\varepsilon pB(M)\bigr)V_f(M).
 \end{aligned}
\]
The two terms lie in $p^{4r}\Zp$ separately. This completes the proof
of Theorem~\ref{III-thm:J}; no inference from a finite numerical range is used.

\section{What the general lemma adds}\label{loc-III-4}
The product-scaling estimate is not restricted to the particular
quadratic polynomial needed for J.3. For any fixed $\lambda\in\Zp$,
\[
 Q_\lambda(p^r):=\sum_{k=0}^{(p^r-1)/2}
                    \frac{(-\lambda p^r)_k^2}{(k!)^2}
 \equiv Q_\lambda(p^{r-1})\pmod{p^{3r}}.
\]
Indeed take $f(X)=(1-\lambda X)^2$ and $\kappa=\lambda^2$;
the case $\lambda=0$ is immediate. More generally, $f$ may be any
power series in $1+X\Zp[[X]]$, with the same $p^{3r}$ modulus.
Convergence of each factor is automatic because $p^r/j\in p\Zp$
throughout its finite product. The proof uses only inverse-square
unit sums, reflection, and summation by parts.

\section{The remaining odd branch of F.3}\label{loc-III-5}\label{III-sec:F}
Let $(x;q)_k=\prod_{i=0}^{k-1}(1-xq^i)$ and
$[m]_q=(1-q^m)/(1-q)$. Products with several first arguments mean the
product of the respective shifted factorials. Define
\[
 T_L(a,q)=\sum_{k=0}^L c_k(a,q),\qquad
 c_k(a,q)=(-1)^k[8k+1]_q q^k
 \frac{(aq,q/a,q,-q^2;q^4)_k}{(aq^4,q^4/a,q^4,-q^3;q^4)_k}.
\]
The parameter $a$ is an indeterminate throughout every root-of-unity
calculation. In particular, parameter denominators do not vanish there.

\begin{theorem}\label{III-Fodd:main}
Let $p\equiv3\pmod4$ be prime and $r=2s+1\ge3$. Put
$N=(p^r-3)/4$ and $M=(p^{r-2}-3)/4$. Then
\[
 \sum_{k=0}^{N}(-1)^k(8k+1)\frac{(1/4)_k^3}{k!^3}
 \equiv p^2\sum_{k=0}^{M}(-1)^k(8k+1)\frac{(1/4)_k^3}{k!^3}
 \pmod {p^{3s}}.
\]
Since $3s\ge2s+1=r$, this implies the odd branch of Swisher F.3.
\end{theorem}

\subsection{The terminating evaluation}
We use the following terminating very-well-poised basic hypergeometric
evaluation, with $v=4t+1$:
\begin{equation}\label{III-Fodd:jackson}
 T_t(q^{v},q)=T_t(q^{-v},q)=(-1)^t q^{-2t}[v]_q.
\end{equation}
Here is an exact general identity from which it follows:
\begin{equation}\label{III-Fodd:jackson-general}
 \sum_{k=0}^{t}[2mk+1]_q
 \frac{(q^{1-v},q^{1+v},q/c,q;q^m)_k}
 {(q^{m-v},q^{m+v},cq^m,q^m;q^m)_k}
 c^k q^{(m-2)k}
 = (c/q)^t\frac{(q^2/c;q^m)_t}{(cq^m;q^m)_t}[v]_q,
 \quad v=mt+1.
\end{equation}
To specify the input identity completely, the terminating
very-well-poised ${}_6\phi_5$ summation is
\[
 \sum_{k=0}^t\frac{1-AQ^{2k}}{1-A}
 \frac{(A,B,C,Q^{-t};Q)_k}
 {(Q,AQ/B,AQ/C,AQ^{t+1};Q)_k}
 \left(\frac{AQ^{t+1}}{BC}\right)^k
 =\frac{(AQ,AQ/(BC);Q)_t}{(AQ/B,AQ/C;Q)_t}.
\]
Take $Q=q^m$, $A=q$, $B=q^{1+v}$, $C=q/c$, and $v=mt+1$.
The argument becomes $cq^{m-2}$. On the right, reversal of the
two finite products with negative powers gives
\[
 \frac{(q^{m+1};q^m)_t}{(q^{m-v};q^m)_t}
 =(-1)^tq^{mt(t-1)/2+t}[v]_q,
\]
\[
 (cq^{m-v-1};q^m)_t
 =(-c)^tq^{-mt(t-1)/2-2t}(q^2/c;q^m)_t,
\]
which proves \eqref{III-Fodd:jackson-general} from that standard summation.
This is also the specialization used in Guo--Zudilin,
\emph{Dwork-type supercongruences
through a creative $q$-microscope}, arXiv:2001.02311~\cite{III-GZ}, formula (3.23)
and the proof of Theorem 3.6. Its substitution $m=4,c=-q^{-1}$ gives
\eqref{III-Fodd:jackson}. Equation \eqref{III-Fodd:jackson-general} is an exact identity
at the two parameter specializations, rather than a congruence.

\begin{lemma}[Root blocks]\label{III-Fodd:blocks}
Suppose $h>1$, $h\equiv1\pmod4$, and $\zeta$ is a primitive $h$th root
of unity. Write $t=(h-1)/4$. For $b\ge0$ and $0\le j<h$,
\begin{equation}\label{III-Fodd:block-formula}
 \lim_{q\to\zeta}c_{bh+j}(a,q)=
 \begin{cases}
 (-1)^b\dfrac{(1/4)_b}{b!}c_j(a,\zeta),&0\le j\le t,\\
 0,&t<j<h.
 \end{cases}
\end{equation}
Furthermore $c_{t-j}(a,\zeta)=-c_j(a,\zeta)$ for $0\le j\le t$.
Consequently, if $L\bmod h\ge t$, then $T_L(a,q)$ vanishes modulo
$\Phi_h(q)$, in the local ring with coefficient field $\mathbb Q(a)$.
\end{lemma}
\begin{proof}
Set
\[
 U_k=\frac{(aq,q/a;q^4)_k}{(aq^4,q^4/a;q^4)_k},\qquad
 Q_k=\frac{(q;q^4)_k}{(q^4;q^4)_k},\qquad
 V_k=\frac{(-q^2;q^4)_k}{(-q^3;q^4)_k}.
\]
In a complete block of length $h$, the numerator and denominator of
$U$ each run over the same root residues, so their ratio is one.
The analogous full-block ratio for $V$ is one because both products
equal $2$. In $Q$, the numerator zeros occur at indices $t+uh$ and
the denominator zeros at indices $(u+1)h-1$. Their limiting ratio in
the $u$th block is $(4u+1)/(4u+4)$; the products of nonzero factors
cancel by permutation of the nontrivial roots. Thus $b$ full blocks
contribute $(1/4)_b/b!$. The external $(-1)^kq^k$ contributes
$(-1)^b$. A final partial block with $j>t$ has one extra numerator
zero. This proves \eqref{III-Fodd:block-formula}.

For completeness, reversing the products in the first short block
gives, with all the following expressions evaluated at $q=\zeta$,
\[
 \frac{U_{t-j}}{U_j}=q^{6j-4t(t+1)},\qquad
 \frac{Q_{t-j}}{Q_j}=(-1)^tq^{3j-2t(t+1)},\qquad
 \frac{V_{t-j}}{V_j}=q^{j-2t^2-t}.
\]
These follow by replacing the reversed factors $q^{h-u}$ by $q^{-u}$;
all denominators here are nonzero for generic $a$.
Also
$[8(t-j)+1]_q=-q^{-8j-1}[8j+1]_q$.
Combining these identities and the factor $(-1)^{t-2j}q^{t-2j}$
from the external sign and power yields
\[
 c_{t-j}(a,q)=-q^{-(2t+1)h}c_j(a,q)=-c_j(a,q).
\]
The short sum is zero, and \eqref{III-Fodd:block-formula} proves the assertion
for every cutoff whose last block includes this short block.
\end{proof}

\subsection{Parameter interpolation and the common denominator}
Set $n=p^2$, so $n\equiv1\pmod8$, and write
\[
 N=nM+\frac{3(n-1)}4,\qquad
 \Delta(a,q)=T_N(a,q)-q^{(1-n)/2}[n]_qT_M(a,q^n).
\]
For $0\le j\le M$, let $A_j=(4j+1)n$.
At $a=q^{\pm A_j}$, the left series terminates at
$nj+(n-1)/4\le N$, while the right series terminates at $j$.
By \eqref{III-Fodd:jackson}, the left value is
$(-1)^{nj+(n-1)/4}q^{(1-A_j)/2}[A_j]_q$.
The right value is
$q^{(1-n)/2}[n]_q(-1)^j q^{-2nj}[4j+1]_{q^n}$.
They agree since $(n-1)/4$ is even and
$[n]_q[4j+1]_{q^n}=[A_j]_q$. Therefore
\[
 P(a,q)=\prod_{j=0}^{M}(1-aq^{A_j})(a-q^{A_j})
\]
divides the parameter numerator of $\Delta$.
Explicitly a common parameter denominator is
\[
 D(a,q)=(aq^4,q^4/a;q^4)_N.
\]
It also clears the right denominator: for $1\le i\le M$, the
exponent $4ni$ is among the exponents $4j$ with $1\le j\le N$.
Hence $a^ND\Delta$ is a polynomial in $a$ with coefficients in
$\mathbb Q(q)$ and is divisible by $P$ in $\mathbb Q(q)[a]$.
The leading coefficient of $P$ is a signed power of $q$.

\subsection{Preserving an extra cyclotomic factor}
Fix an even $\ell\in\{2,4,\ldots,2s\}$, and put $h=p^\ell$.
Let $R_h=\mathbb Q[q,q^{-1}]_{(\Phi_h)}$.
For any odd integer $d$, the ratio $Q_k=(q;q^4)_k/(q^4;q^4)_k$
has nonnegative $\Phi_d$ valuation: in each period $d$ its numerator
zero occurs at some index $0\le j_0<d$, whereas its denominator zero
occurs at index $d-1$. The factors $1+q^{4i+3}$ have no zero at an
odd-order root. The same facts apply to $q^n$ at a primitive
$p^\ell$ root (at $\ell=2$ the limit is $q^n=1$ and $Q_k$ has the
finite value $(1/4)_k/k!$). Thus
$a^ND\Delta\in R_h[a]$.

At a primitive $h$th root, the left sum vanishes by Lemma
\ref{III-Fodd:blocks}, since
\[
 N\bmod h=\frac{3(h-1)}4\ge\frac{h-1}4.
\]
For $\ell=2$, the right side vanishes due to $[n]_q$ and the remaining
sum is finite for generic $a$. For $\ell>2$, its base $q^n$ has order
$h/n$, and
$M\bmod(h/n)=3(h/n-1)/4$, so Lemma \ref{III-Fodd:blocks} again applies.
Consequently $a^ND\Delta\in\Phi_hR_h[a]$.

Since $P$ has unit leading coefficient in $R_h$, polynomial division
shows that $H=a^ND\Delta/P$ belongs to $R_h[a]$. Reducing the equality
$PH=a^ND\Delta$ modulo $\Phi_h$, the polynomial $P$ is nonzero
(indeed its leading coefficient remains a unit). The residue ring
is a polynomial ring over a field and has no zero divisors; hence
$H\in\Phi_hR_h[a]$. Specializing $a=1$ gives
\begin{equation}\label{III-Fodd:multiplicity}
 \vp_{\Phi_h}\Delta(1,q)\ge1+
 \vp_{\Phi_h}P(1,q)-\vp_{\Phi_h}D(1,q).
\end{equation}

\subsection{Exact root counts}
Put $u=h/n=p^{\ell-2}$. Then $u\equiv1\pmod4$ (including $u=1$),
and
\[
 \#\{0\le j\le M:u\mid4j+1\}
 =\lfloor M/u\rfloor+1=\lfloor N/h\rfloor+1.
\]
Indeed the first solution in a period is $(u-1)/4$, while
$M\bmod u=3(u-1)/4$; for $u=1$ both residues are zero.
The last floor equality follows from
$N=nM+3(n-1)/4$ and
$n\cdot3(u-1)/4+3(n-1)/4=3(h-1)/4<h$.
Therefore
\[
 \vp_{\Phi_h}P(1,q)=2(\lfloor N/h\rfloor+1),\qquad
 \vp_{\Phi_h}D(1,q)=2\lfloor N/h\rfloor.
\]
Equation \eqref{III-Fodd:multiplicity} gives at least three factors for every
even $\ell$, i.e.
\begin{equation}\label{III-Fodd:qresult}
 \vp_{\Phi_{p^{2j}}}\Delta(1,q)\ge3\quad(1\le j\le s).
\end{equation}

\subsection{The specialization at \texorpdfstring{$q=1$}{q=1}}
This point is needed in addition to the local multiplicities.
At $a=1$ the Pochhammer portion of each summand is
$Q_k^3V_k$. The preceding zero-count argument shows that after
cancellation $Q_k$ has no odd-order cyclotomic denominator.
Its remaining denominator factors come from $1-q^{4i}$, and so are
even-order cyclotomic polynomials; their values at $1$ are $1$ or $2$.
Likewise every factor of $1+q^{4i+3}$ has even order and value at $1$
equal to $1$ or $2$. Thus there is a Laurent polynomial denominator
$B(q)\in\mathbb Z[q,q^{-1}]$ such that $B(1)$ is a power of $2$,
$A(q)=B(q)\Delta(1,q)\in\mathbb Z[q,q^{-1}]$, and $B$ has no
odd-order cyclotomic factor. This remains true for the right side
after substituting $q\mapsto q^n$ because $n$ is odd: an odd power
cannot turn an even-order root into an odd-order root, and the value
at $q=1$ of the substituted denominator is unchanged.
By \eqref{III-Fodd:qresult}, $A$ is divisible over $\mathbb Q[q,q^{-1}]$ by
$\prod_{j=1}^s\Phi_{p^{2j}}(q)^3$. Since these factors are monic,
Gauss's lemma gives the same divisibility over
$\mathbb Z[q,q^{-1}]$. Their values at $q=1$ are all $p$.
Finally,
\[
 T_L(1,1)=\sum_{k=0}^L(-1)^k(8k+1)\frac{(1/4)_k^3}{k!^3},
 \qquad q^{(1-n)/2}[n]_q\big|_{q=1}=p^2.
\]
Since $B(1)$ is a $p$-adic unit, \eqref{III-Fodd:qresult} proves
Theorem \ref{III-Fodd:main}.

\LocalAppendix
\section{An exact representation for E.3 and F.3}\label{loc-III-A}\label{III-app:EF}
The following identity provides an auxiliary integral representation.
It is not used to replace the cyclotomic proof of F.3 above, and it
does not close the remaining E.3 branch.

\begin{proposition}
For $m\ge0$ and a parameter $\alpha$ at which the expressions are
defined, put
\[
 S_\alpha(m)=\sum_{k=0}^m(-1)^k\frac{2k+\alpha}{\alpha}
                         \frac{(\alpha)_k^3}{(k!)^3}.
\]
Then the following identity of rational functions holds:
\[
 \boxed{S_\alpha(m)=\frac{(1+\alpha)_m}{m!}
      \sum_{k=0}^m\frac{(-m-\alpha)_k(\alpha)_k}{(k!)^2}.}
\]
In particular it holds for $\alpha=1/3$ and $\alpha=1/4$.
\end{proposition}
\begin{proof}
Set $q(m,k)=(-m-\alpha)_k(\alpha)_k/(k!)^2$ and
$R(m,k)=-k^2/(m+\alpha)^2$. Direct cancellation gives
\[
 q(m,k+1)R(m,k+1)-q(m,k)R(m,k)
 =q(m,k)-\frac{m}{m+\alpha}q(m-1,k).
\]
This is an identity of rational functions, so it suffices to check
it at generic parameter values. Write $b_m=(\alpha)_m/m!$ and
$W_m=\sum_{k=0}^m q(m,k)$. The two endpoint identities are
\[
 q(m-1,m)=(-1)^m b_m^2,
 \qquad q(m,m+1)R(m,m+1)=(-1)^m b_m^2.
\]
Summing the certificate for $0\le k\le m$ therefore yields
\[
 W_m-\frac{m}{m+\alpha}W_{m-1}
 =(-1)^m\frac{2m+\alpha}{m+\alpha}b_m^2.
\]
Now $P_m=(1+\alpha)_m/m!$ satisfies
$P_m/P_{m-1}=(m+\alpha)/m$ and
$P_m=(m+\alpha)b_m/\alpha$. Multiplication by $P_m$
gives precisely the $m$th summand of $S_\alpha(m)$ as an
increment. The initial value is $P_0W_0=1$.
\end{proof}

If $\alpha\in\Zp$, every term of $W_m$ is $p$-integral: each
factor $(\beta)_k/k!=(-1)^k\binom{-\beta}{k}$ lies in $\Zp$
for $\beta\in\Zp$. This follows, for example, by approximating
$-\beta$ by nonnegative integers in $\Zp$ and using continuity
of the fixed binomial polynomial.

For the short odd-level branches the parameters become
\[
 \begin{array}{c|c|c|c}
 &\alpha&m&-m-\alpha\\\hline
 E.3&1/3&(p^r-2)/3&(1-p^r)/3\\
 F.3&1/4&(p^r-3)/4&1/2-p^r/4
 \end{array}
\]
Here E.3 has $p\equiv2\pmod3$, F.3 has $p\equiv3\pmod4$,
and $r\ge3$ is odd. In contrast to the J residual, these upper
parameters are not both divisible by $p^r$. Consequently the
square-product scaling theorem proved above cannot simply be applied.
The original targets for these two branches are respectively
\[
 \begin{aligned}
 S_{1/3}((p^r-2)/3)&\equiv
 p^2S_{1/3}((p^{r-2}-2)/3)\pmod{p^{3r-1}},\\
 S_{1/4}((p^r-3)/4)&\equiv
 p^2S_{1/4}((p^{r-2}-3)/4)\pmod{p^r}.
 \end{aligned}
\]
The F.3 target has been proved in Section~\ref{III-sec:F}. The E.3 target
remains unproved here: its precision $3r-1$ is not supplied by the
integral transformation or by the F.3 argument. No closure of A.3,
D.3, E.3, or G.3 is claimed in this manuscript.

\section{Verification and reproducibility}\label{loc-III-B}
The accompanying archive contains the single source file, this PDF,
and exact-arithmetic verification scripts with their recorded output.
The checks include the finite J identity, the telescoping certificates,
all harmonic-prefix and product-error bounds used in the proof, the
Morley prefactor precision, and the final J recurrence. The principal
scaling audit tests $p=5,7$ through $r=4$, $p=11$ through $r=3$,
and $p=13,17,19$ through $r=2$, for 17 prime-power cases. The
general E/F identity is checked at $\alpha=1/2,1/3,1/4$ and
$0\le m\le30$. For F.3, the archive additionally records exact rational checks of
the stronger integer congruence, an exact calculation modulo
$\Phi_9^3$ (with nonzero remainder modulo $\Phi_9^4$), and a separate
calculation modulo $\Phi_{49}^3$ over $\mathbb F_{1009}$ that exercises
cyclotomic factor cancellation in a denominator. The finite-field
check is only a diagnostic, not a characteristic-zero proof.
These tests can detect implementation and transcription errors; the all-level conclusions rest on the preceding proofs.

The manuscript is independent of any unverified geometric lifting
assertion in earlier research notes or corpus entries. The present
proof of J.3 uses only the explicitly stated elementary lemmas and
the published Cai congruence. The F.3 proof uses the standard
terminating very-well-poised summation and the explicit local-ring
argument above. A separate assessment of historical
priority and external mathematical review remains appropriate.

\begin{thebibliography}{9}
\bibitem{III-Swisher}
H. Swisher, \emph{On the supercongruence conjectures of van Hamme},
Research in the Mathematical Sciences \textbf{2}, article 18 (2015).
\href{https://doi.org/10.1186/s40687-015-0037-6}{doi:10.1186/s40687-015-0037-6}.
See the published prime-power conjecture table, entry J.3.
\bibitem{III-Cai}
T. Cai, \emph{A congruence involving the quotients of Euler and its
applications (I)}, Acta Arithmetica \textbf{103}(4) (2002), 313--320.
\href{https://www.impan.pl/shop/publication/transaction/download/product/83800}{Publisher's full text}.
See Corollary 3 for the prime-power Morley congruence.
\bibitem{III-GZ}
V. J. W. Guo and W. Zudilin,
\emph{Dwork-type supercongruences through a creative $q$-microscope},
arXiv:2001.02311v4 (2020).
\href{https://arxiv.org/abs/2001.02311}{arXiv record}.
See formula (3.23) and the proof of Theorem 3.6 for the terminating
evaluation and the creative-microscoping framework. The short odd
cutoff and its denominator bookkeeping are proved here.
\end{thebibliography}

\clearpage\endgroup

%%%%%%%%%%%%%%%%%%%%%%%%%%%%%%%%%%%%%%%%%%%%%%%%%%%%%%%%%%%%%%%%%%%%%%%%%%
% IV | Swisher_ADG_Supersingular_Advances.tex
%%%%%%%%%%%%%%%%%%%%%%%%%%%%%%%%%%%%%%%%%%%%%%%%%%%%%%%%%%%%%%%%%%%%%%%%%%
\clearpage\begingroup
\SetPaper{IV}{Supersingular advances for A.3, D.3 and G.3}
\RestoreCounters
\PaperTitle{Supersingular Advances for Swisher's A.3, D.3 and G.3\\
\large A second-level congruence, a complete D.3 residue branch,\\
and every odd level of a G.3 residue branch}{Research manuscript}{10 November 2026}

\begin{abstract}
We establish three precisely delimited results for Swisher's prime-power
supercongruences. For A.3 we prove the second-level congruence modulo
$p^8$ for every prime $p\equiv3\pmod4$. For D.3 we prove the complete
$p\equiv2\pmod3$, $p>3$, branch, including both parities of the level.
For G.3 we prove the complete odd-level branch for primes
$p\equiv3\pmod4$. The all-level estimates behind the latter results are
$v_p(D_r)\ge5\lfloor r/2\rfloor$ and
$v_p(G_r)\ge3\lfloor r/2\rfloor$.
The proofs combine exact terminating hypergeometric evaluations,
cyclotomic interpolation at colliding nodes, and a finite interval
partition respecting prime-power digits. The first odd level requires
separate local block identities. No ordinary-prime all-level result is
claimed here, and the odd supersingular branch of E.3 remains unresolved.
\end{abstract}

\noindent\textbf{Status.} This is a research manuscript supported by internal
cross-checks and reproducible exact arithmetic. It has not undergone
external peer review. No claim of historical priority is made.
\setcounter{tocdepth}{1}

\begin{LocalContents}
\LocalEntry{loc-IV-1}{1}{Statements and exact scope}
\LocalEntry{loc-IV-2}{2}{An even-level divisibility theorem for (D.3)}
\LocalEntry{loc-IV-3}{3}{All supersingular levels: a digit-interval argument}
\LocalEntry{loc-IV-4}{4}{The odd third level of D.3 for all supersingular primes}
\LocalEntry{loc-IV-5}{5}{The first odd supersingular instance of G.3}
\LocalEntry{loc-IV-6}{6}{Every odd supersingular level of G.3}
\LocalEntry{loc-IV-7}{7}{The second supersingular level of A.3}
\LocalEntry{loc-IV-8}{8}{Verification and remaining targets}
\end{LocalContents}
\clearpage


\section{Statements and exact scope}\label{loc-IV-1}
We use $(a)_k=\prod_{j=0}^{k-1}(a+j)$, $(a)_0=1$, and the valuation
$v_p$ normalized by $v_p(p)=1$, with $v_p(0)=+\infty$.
All sums in this paper are rational numbers; congruences refer to their
images in $\mathbb Z_p$. The labels A.3, D.3 and G.3 refer to
Section~7 of Swisher~\cite{IV-Swisher}. We call the indicated nontrivial
residue classes \emph{supersingular} only as convenient terminology
for the branches of those conjectures.

For odd $N$, put
\[
 A(N)=\sum_{k=0}^{(N-1)/2}(-1)^k(4k+1)\frac{(1/2)_k^5}{k!^5}.
\]
For primes $p>3$ with $p\equiv2\pmod3$, put
\[
 D_r=\sum_{k=0}^{d_r}(6k+1)\frac{(1/3)_k^6}{k!^6},\qquad
 d_r=\begin{cases}(p^r-1)/3&r\text{ even},\\
                  (p^r-2)/3&r\text{ odd}.
       \end{cases}
\]
For primes $p\equiv3\pmod4$, put
\[
 G_r=\sum_{k=0}^{g_r}(8k+1)\frac{(1/4)_k^4}{k!^4},\qquad
 g_r=\begin{cases}(p^r-1)/4&r\text{ even},\\
                  (p^r-3)/4&r\text{ odd}.
       \end{cases}
\]
The displayed definitions include $D_0=G_0=1$.

\begin{theorem}[Second supersingular level of A.3]\label{IV-main:A}
For every prime $p\equiv3\pmod4$,
\[
 A(p^2)\equiv p^4\pmod{p^8}.
\]
\end{theorem}

\begin{theorem}[Complete supersingular branch of D.3]\label{IV-main:D}
Let $p>3$ be prime with $p\equiv2\pmod3$. For every $r\ge1$,
\[
 v_p(D_r)\ge5\lfloor r/2\rfloor.
\]
Moreover, all the supersingular assertions of D.3 hold:
\begin{align*}
 D_2&\equiv0\pmod{p^4},\\
 D_r&\equiv p^4D_{r-2}\pmod{p^{2r+1}}
       &&(r\ge4\text{ even}),\\
 D_r&\equiv-p^5D_{r-2}\pmod{p^{2r}}
       &&(r\ge3\text{ odd}).
\end{align*}
\end{theorem}

\begin{theorem}[Every odd supersingular level of G.3]\label{IV-main:G}
Let $p\equiv3\pmod4$ be prime. For every $r\ge1$,
\[
 v_p(G_r)\ge3\lfloor r/2\rfloor.
\]
For every odd $r\ge3$,
\[
 G_r\equiv-p^3G_{r-2}\pmod{p^{r+1}}.
\]
The prime $p=3$ is included.
\end{theorem}

The valuation estimates do not assert exact valuations. They suffice
for the original high-level targets because both terms in the proposed
recurrence then vanish at the required precision. At $r=3$ the leading
term and sign require extra work. This distinction is essential: the
same valuation argument does not give the stronger even-level target
in G.3.

\begin{center}
\small
\begin{tabular}{@{}p{.08\linewidth}p{.42\linewidth}p{.43\linewidth}@{}}
\toprule
Label & Established in this manuscript & Not established here \\
\midrule
A.3 & $p\equiv3\pmod4$, $r=2$, modulo $p^8$
& Higher supersingular levels and ordinary higher-level targets \\[3pt]
D.3 & Every required level for $p>3$, $p\equiv2\pmod3$
& Ordinary-prime higher-level target modulo $p^{6r}$ \\[3pt]
G.3 & $p\equiv3\pmod4$, every odd $r\ge3$
& Supersingular even-level recurrence in general; ordinary higher levels \\[3pt]
E.3 & An exact short-cutoff deformation is recorded in the companion bundle
& The remaining odd-level target modulo $p^{3r-1}$ \\
\bottomrule
\end{tabular}
\end{center}

The table describes the scope of this manuscript, not a claim that
all omitted base cases are unknown in the literature. Earlier work in
the surrounding project and the excluded labels are not re-proved here.
The proofs below are independent of the all-level Frobenius
and period-map results in the motivating corpus~\cite{IV-Corpus}.

\paragraph{Proof organization.}
Sections~\ref{IV-sec:D-even}--\ref{IV-sec:D-cubic} prove Theorem~\ref{IV-main:D}.
The even-level argument uses the specific published parametric input
of Ni--Wang~\cite{IV-NiWang}, reproduced by Wei~\cite{IV-Wei}; the subsequent
interval and cubic-block arguments are given explicitly.
Sections~\ref{IV-sec:G-r3}--\ref{IV-sec:G-all} prove Theorem~\ref{IV-main:G}
from terminating summation formulas and gamma reflection.
Section~\ref{IV-sec:a-second-level} proves Theorem~\ref{IV-main:A}, using a
shifted WZ identity and removal of a four-root deformation.
The standard gamma estimates used there are specified from
Long--Ramakrishna~\cite{IV-LongRamakrishna}.

% D.3, supersingular even branch.  Proof fragment; 10 November 2026.
\section{An even-level divisibility theorem for (D.3)}\label{loc-IV-2}
\label{IV-sec:D-even}
Write $(z;Q)_k=\prod_{j=0}^{k-1}(1-zQ^j)$, and use a list of
arguments for a product of such symbols.  Put $[m]_q=(1-q^m)/(1-q)$ and
\[
 c_k(a,b;q)=[6k+1]_q q^{3k}
 \frac{(aq,q/a,bq,q/b;q^3)_k(q;q^3)_k^2}
 {(aq^3,q^3/a,bq^3,q^3/b;q^3)_k(q^3;q^3)_k^2},
 \qquad S_L(a,b;q)=\sum_{k=0}^L c_k(a,b;q).
\]
For a positive integer $h>1$ coprime to $3$, let $v_h$ denote the
valuation at the irreducible polynomial $\Phi_h(q)$ on $\mathbb Q(q)$.
The same symbol denotes its Gauss extension when independent parameters
are present.  All congruences in this section are in this local sense.

\paragraph{The short-block input.}
We use the following published parametric congruence, with its exact
cutoff.  If $h\equiv2\pmod3$ and $t=(2h-1)/3$, then
\begin{equation}\label{IV-D4:base2}
 S_t(a,b;q)\equiv0\pmod{[h]_q\Phi_h(q)}.
\end{equation}
This is Ni--Wang, \emph{Two $q$-supercongruences from Watson's
transformation}, Rev. R. Acad. Cienc. Exactas F\'is. Nat., Ser. A Mat.
\textbf{116} (2022), Article 30, Theorem 2.3, with $d=3,r=1$.
The exact specialization is reproduced as equation (2.12) in
Chuanan Wei, \emph{$q$-Supercongruences from Jackson's ${}_8\phi_7$
summation and Watson's ${}_8\phi_7$ transformation}, J. Combin.
Theory Ser. A \textbf{204} (2024), 105853,
\url{https://arxiv.org/abs/2104.07025}.
Since $t<h$, all the displayed denominator factors remain
$\Phi_h$-units when $b=1$ (with $a$ generic).  Consequently
\eqref{IV-D4:base2} implies $v_h(S_t(a,1;q))\ge2$ without a loss of
precision upon that specialization.
For $h\equiv1\pmod3$ and $t=(h-1)/3$, we also use the simpler
short-block congruence $S_t(1,1;q)\equiv0\pmod{\Phi_h(q)}$.
It follows either by term reversal, or from Ni--Wang Lemma 2.2 as
reproduced in Wei equations (2.7)--(2.9); here again $t<h$ makes
all the specialized denominators units.

\paragraph{A uniform parametric block lemma.}
If $h\equiv2\pmod3$, $t=(2h-1)/3$ and $B\ge0$, then
\begin{equation}\label{IV-D4:block2}
 v_h\bigl(S_{Bh+t}(a,1;q)\bigr)\ge2.
\end{equation}
Here $a$ is an independent parameter.  We give the shift argument in
detail, since a short-block result alone would not imply this claim.
Set $Q=q^3$ and define
\[
 A_j(x,a,b;q)=\frac{1-x^2q^{6j+1}}{1-q}\,q^{3j}
 \frac{(x^2q,axq,xq/a,bxq,xq/b,q;Q)_j}
 {(x^2q^3,axq^3,xq^3/a,bxq^3,xq^3/b,q^3;Q)_j},
\]
\[
 \mathcal A_t(x,a,b;q)=\sum_{j=0}^t A_j(x,a,b;q).
\]
These expressions are regular near $x=b=1$ in the $\Phi_h$-local
ring with $a$ generic.  In particular, their denominators at
$x=b=1$ are units, since $j\le t<h$.
We claim that
\begin{equation}\label{IV-D4:central}
 \mathcal A_t(x,a,b;q)\equiv0\pmod{\Phi_h(q)}
\end{equation}
as a rational function of the independent parameters $x,a,b$,
and with denominators regular at $x=b=1$.
To verify this claim, apply the terminating balanced Jackson
${}_8\phi_7$ summation with base $Q$, central parameter $A=x^2q$,
and its five additional upper parameters
\[
 b_1=axq,\quad c=xq/a,\quad d=bxq,\quad
 e=xq^{1+2h}/b,\quad f=q^{1-2h}=Q^{-t}.
\]
The balancing condition is exact:
$b_1cdef=x^4q^5=A^2Q$.  Modulo $\Phi_h(q)$ the resulting
finite sum, multiplied by $(1-A)/(1-q)$, is precisely
$\mathcal A_t(x,a,b;q)$.  Jackson's product side is
\[
 \frac{1-A}{1-q}
 \frac{(AQ,AQ/(b_1c),AQ/(b_1d),AQ/(cd);Q)_t}
 {(AQ/b_1,AQ/c,AQ/d,AQ/(b_1cd);Q)_t}.
\]
It contains the numerator factor $(q^2;Q)_t$, which vanishes
modulo $\Phi_h$: its factor with index $(h-2)/3$ is $1-q^h$.
The four denominator arguments are $xq^3/a,axq^3,xq^3/b,q/(bx)$.
They stay units near $x=b=1$, because the first occurrence of
$1-q^{mh}$ in $(q;Q)_j$ is at index $t$, outside the product
for $j\le t$.  This proves \eqref{IV-D4:central}, including the
regularity assertion.  In particular,
\begin{equation}\label{IV-D4:derivative}
 v_h\bigl(\partial_x \mathcal A_t(1,a,1;q)\bigr)\ge1.
\end{equation}

For a block index $u\ge0$, put $x_u=q^{3uh}$ and
\[
 E_u=q^{3uh}
 \frac{(aq,q/a,q,q;Q)_{uh}(q;Q)_{uh}^{2}}
 {(aq^3,q^3/a,q^3,q^3;Q)_{uh}(q^3;Q)_{uh}^{2}}.
\]
Each complete block of $h$ factors in $(q;Q)/(q^3;Q)$ has one
numerator zero and one denominator zero; hence $E_u$ is a
$\Phi_h$-unit for generic $a$.  For $0\le j\le t$, direct
cancellation of shifted factorials gives
\[
 c_{uh+j}(a,1;q)=E_u A_j(x_u,a,1;q)R_j(x_u;q),
\]
where
\[
 R_j(x;q)=\frac{(xq;Q)_j^2}{(q;Q)_j(x^2q;Q)_j}
          \frac{(q^3;Q)_j(x^2q^3;Q)_j}{(xq^3;Q)_j^2}.
\]
All these factors are regular near $x=1$; direct differentiation
gives $R_j(1;q)=1$ and $\partial_xR_j(1;q)=0$.
Since $v_h(x_u-1)\ge1$ when $u>0$, it follows that
$R_j(x_u;q)\equiv1\pmod{\Phi_h^2}$; the assertion is exact
for $u=0$.  Equations \eqref{IV-D4:base2}, \eqref{IV-D4:derivative}
and Taylor expansion now give
\[
 \mathcal A_t(x_u,a,1;q)\equiv
 \mathcal A_t(1,a,1;q)+(x_u-1)\partial_x \mathcal A_t(1,a,1;q)
 \equiv0\pmod{\Phi_h^2}.
\]
Thus every block prefix from $uh$ through $uh+t$ has valuation
at least two.  In the rest of a block, $t<j<h$, the quotient
$(q;Q)_{uh+j}/(q^3;Q)_{uh+j}$ has valuation one; its fourth
power occurs in $c_{uh+j}(a,1;q)$.  Consequently those remaining
terms even have valuation at least four.  Adding the complete
blocks and the final prefix proves \eqref{IV-D4:block2}.

\paragraph{Interpolation at colliding nodes.}
Under the same assumptions on $h,t,B$, we prove
\begin{equation}\label{IV-D4:four}
 v_h\bigl(S_{Bh+t}(1,1;q)\bigr)\ge4.
\end{equation}
Put $L=Bh+t$ and
\[
 D(a)=(aq^3,q^3/a;Q)_L,\qquad
 F(a)=a^L D(a)S_L(a,1;q).
\]
Then $F(a)$ is a polynomial in $a$ with coefficients in the
$\Phi_h$-local ring.  Indeed, $D$ clears all the parameter
denominators, the factor $a^L$ clears all negative powers of $a$,
and the remaining $(q;Q)/(q^3;Q)$ quotients are regular at
$\Phi_h$.  By \eqref{IV-D4:block2} every coefficient of $F$ has
valuation at least two (equivalently, its Gauss valuation is
at least two).

For $0\le i\le B$ define
\[
 A_i=(3i+2)h,\qquad t_i=(A_i-1)/3=ih+t,
 \qquad \alpha_{i,+}=q^{A_i},\quad\alpha_{i,-}=q^{-A_i}.
\]
The nodes are distinct in $\mathbb Q(q)$, and $A_i\equiv1\pmod3$.
At either node the sum terminates at $t_i\le L$, and the exact
Jackson evaluation is
\begin{equation}\label{IV-D4:values}
 S_L(q^{\pm A_i},1;q)
 =[A_i]_q\left\{\frac{(q^2;Q)_{t_i}}{(q^3;Q)_{t_i}}\right\}^{3}.
\end{equation}
For completeness, this is Jackson's identity with central parameter
$q$, other upper parameters $q,q,q,q^{1+A_i},q^{1-A_i}$ and
base $Q$; the factor $(q^4;Q)_{t_i}/(q;Q)_{t_i}$ telescopes
to $[A_i]_q$.  No parameter denominator vanishes identically
because $3\nmid A_i$.  The same evaluation is the specialization
$b=1$ of Wei equation (2.3), with the termination exponent
allowed to be any positive integer $A_i\equiv1\pmod3$.
The quotient in braces has valuation one: it has $i+1$
numerator zeros and $i$ denominator zeros.  Therefore the
right side of \eqref{IV-D4:values} has valuation four.
At each node, $v_h(D(\alpha_{i,\pm}))=2B$, since among
$1\le j\le L$ precisely $B$ integers satisfy $h\mid3j$.
It follows that
\begin{equation}\label{IV-D4:Fvalues}
 v_h(F(\alpha_{i,\pm}))\ge2B+4.
\end{equation}

Let $P(a)=\prod_{i=0}^B(a-\alpha_{i,+})(a-\alpha_{i,-})$.
Division by this monic polynomial gives $F=PH+R$, with
$\deg R<2B+2$.  Since the coefficients of $P$ are integral
in the local ring, those of $H$ and $R$ still have valuation
at least two.  Thus
\[
 v_h(P(1)H(1))\ge(2B+2)+2=2B+4.
\]
The nodes all approach $1$ modulo $\Phi_h$, so ordinary
unit-separated interpolation is unavailable.  Nevertheless,
the precise Lagrange weights at the point $1$ are units:
\[
 R(1)=\sum_{\alpha}F(\alpha)
       \prod_{\beta\ne\alpha}\frac{1-\beta}{\alpha-\beta}.
\]
For every node $\beta$, $v_h(1-\beta)=1$, and for any two
distinct nodes $\alpha,\beta$ we have
$v_h(\alpha-\beta)=1$, since their nonzero exponents and
exponent differences are multiples of $h$.
Every displayed Lagrange weight consequently has valuation
zero.  By \eqref{IV-D4:Fvalues}, $v_h(R(1))\ge2B+4$, and hence
$v_h(F(1))\ge2B+4$.  Finally $v_h(D(1))=2B$; dividing proves
\eqref{IV-D4:four}.

\paragraph{The complementary roots.}
If $h\equiv1\pmod3$, $t=(h-1)/3$, and $B\ge0$, then
\begin{equation}\label{IV-D4:one}
 v_h(S_{Bh+t}(1,1;q))\ge1.
\end{equation}
Indeed, the short block has valuation at least one by the
short-block input above.  In a block $uh+j$, $0\le j\le t$,
the shifted quotient factors reduce modulo $\Phi_h$ to the
unshifted ones, multiplied by a common regular block factor.
The weight and $q^{3k}$ reduce to their unshifted values too.
For $t<j<h$, each summand has six numerator zeros in excess
of its denominator zeros, and so is zero modulo $\Phi_h$.
This proves \eqref{IV-D4:one} block by block.

\paragraph{Specialization and the original conjecture.}
Let $p>3$ be prime with $p\equiv2\pmod3$, let $s\ge1$,
and put $L=(p^{2s}-1)/3$.  For $h=p^\ell$, $1\le\ell\le2s$,
the residue of $L$ modulo $h$ is $(2h-1)/3$ if $\ell$ is odd,
and $(h-1)/3$ if $\ell$ is even.  Equations \eqref{IV-D4:four}
and \eqref{IV-D4:one} therefore give the rational-function
congruence
\[
 S_L(1,1;q)\equiv0\pmod{
 \prod_{\substack{1\le\ell\le2s\\\ell\text{ odd}}}\Phi_{p^\ell}(q)^4
 \prod_{\substack{1\le\ell\le2s\\\ell\text{ even}}}\Phi_{p^\ell}(q)}.
\]
There is no loss of $p$-adic precision at $q=1$.
To see this explicitly, the reduced denominator of each
$(q;q^3)_k/(q^3;q^3)_k$ has only cyclotomic factors whose
orders are divisible by $3$.  For any order $d$ coprime to
$3$, the numerator's single zero in a period occurs before
the denominator's zero, so its valuation is nonnegative.
For $3\mid d$, a remaining factor $\Phi_d(1)$ is either $1$
or $3$; thus a common reduced denominator for the finitely
many summands has value a power of $3$ at $1$.
After clearing this denominator, the numerator lies in
$\mathbb Z[q]$.  Division by the monic integral cyclotomic
product above preserves integral coefficients.  Since
$\Phi_{p^\ell}(1)=p$, taking $q\to1$ proves
\begin{equation}\label{IV-D4:valuation}
 \boxed{\quad
 v_p\left(\sum_{k=0}^{(p^{2s}-1)/3}
       (6k+1)\frac{(1/3)_k^6}{k!^6}\right)\ge5s.
 \quad}
\end{equation}
Denote this sum by $D_{2s}$.  In particular $D_2\equiv0\pmod{p^5}$,
which includes the original base target modulo $p^4$.
For $s\ge2$, both $D_{2s}$ and $p^4D_{2s-2}$ are divisible
by $p^{4s+1}$, because $5s\ge4s+1$ and
$4+5(s-1)=5s-1\ge4s+1$.  Hence
\[
 \boxed{D_{2s}\equiv p^4D_{2s-2}\pmod{p^{4s+1}}\qquad(s\ge2),}
\]
which proves the full even supersingular branch of (D.3).
The next two sections supply the odd levels. The ordinary-prime
branch is outside the present theorem.

% Companion to even_q_proof.tex.  The odd valuation bound uses a finite
% interval partition, not an unproved transfer of the even recurrence.
\section{All supersingular levels: a digit-interval argument}\label{loc-IV-3}
\label{IV-sec:D-all}
Retain $c_k(1,1;q)$ and the notation of the preceding proof.  Fix a
prime $p>3$ with $p\equiv2\pmod3$, and put
\[
 n=p^2,\qquad d=(n-1)/3,\qquad t_j=(n^j-1)/3
       =d n^{j-1}+t_{j-1}\quad(j\ge1),\qquad t_0=0.
\]
Let
\[
 P_j(q)=\prod_{\substack{1\le\ell\le2j\\\ell\text{ odd}}}
                 \Phi_{p^\ell}(q)^4
          \prod_{\substack{1\le\ell\le2j\\\ell\text{ even}}}
                 \Phi_{p^\ell}(q),\qquad P_0=1.
\]
For every integer $u\ge0$, define the interval
\[
 I(u,j)=\{un^j,un^j+1,\ldots,un^j+t_j\}.
\]
The preceding proof implies the following shifted-prefix statement:
\begin{equation}\label{IV-D4:shiftedprefix}
 P_j(q)\ \Bigm|\ \sum_{k\in I(u,j)}c_k(1,1;q).
\end{equation}
Divisibility here and below is after clearing a denominator in
$\mathbb Z[q]$ whose value at $1$ is a power of $3$.
Indeed, for any $h=p^\ell$, $\ell\le2j$, the right endpoint
$un^j+t_j$ has the same residue modulo $h$ as $t_j$, namely the
appropriate short-block cutoff in the preceding proof.  The endpoint
$un^j-1$ has residue $h-1$.  The congruences in that proof also
hold for this latter residue: the extra terms beyond the short
cutoff have at least six numerator zeros in excess of denominator
zeros.  Subtracting the two initial sums gives
\eqref{IV-D4:shiftedprefix}; for $u=0$ the subtracted sum is empty.

For an interval of length $n^j$, also introduce the sets
\[
 \begin{split}
 B(u,j)&=\{un^j,\ldots,(u+1)n^j-1\},\\
 T(u,j)&=\{un^j+t_j+1,\ldots,(u+1)n^j-1\},\\
 J(u,j)&=\{un^j,\ldots,un^j+2t_j\}.
 \end{split}
\]
At level zero, $B(u,0)=J(u,0)=I(u,0)=\{u\}$ and $T(u,0)$
is empty.  The following are exact disjoint interval partitions
for $j\ge1$:
\begin{align}
 B(u,j)&=I(u,j)\,\dot\cup\,T(u,j),\label{IV-D4:partitionB}\\
 T(u,j)&=T(nu+d,j-1)
       \,\dot\cup\!\!\!\bigcup_{c=d+1}^{n-1}B(nu+c,j-1),
                                                    \label{IV-D4:partitionT}\\
 J(u,j)&=I(u,j)\,\dot\cup\,T(nu+d,j-1)
       \,\dot\cup\!\!\!\bigcup_{c=d+1}^{2d-1}B(nu+c,j-1)
       \,\dot\cup\,J(nu+2d,j-1).\label{IV-D4:partitionJ}
\end{align}
These identities follow immediately from
$t_j=dn^{j-1}+t_{j-1}$ and $2t_j=2dn^{j-1}+2t_{j-1}$.
In \eqref{IV-D4:partitionJ}, every piece other than $I(u,j)$
lies in the tail $T(u,j)$.

Apply these rules recursively, stopping at every set $I(v,k)$.
They give finite disjoint partitions of both $B(u,j)$ and
$J(u,j)$ into such sets.  A leaf $I(v,k)$ has undergone exactly
$j-k$ descents, through tails at levels $k+1,k+2,\ldots,j$.
For a level $\ell$ on that path, every index $m$ in the leaf
has remainder greater than $t_\ell$ modulo $n^\ell$.
Since $n^\ell\equiv1\pmod3$, it follows that
\[
 v_{\Phi_{n^\ell}}\left(\frac{(q;q^3)_m}{(q^3;q^3)_m}\right)=1.
\]
Consequently each summand belonging to that leaf, and hence its
sum, has a factor $\Phi_{n^\ell}(q)^6$ for every such ancestor
level.  Combining these distinct high-order factors with
\eqref{IV-D4:shiftedprefix} gives
\begin{equation}\label{IV-D4:leafproduct}
 P_k(q)\prod_{\ell=k+1}^{j}\Phi_{p^{2\ell}}(q)^6
       \ \Bigm|\ \sum_{m\in I(v,k)}c_m(1,1;q).
\end{equation}
There is no assumption of integral comaximality between different
cyclotomic polynomials here.  Clear a common denominator whose
value at $1$ is a power of $3$.  The integral numerator is
divisible over $\mathbb Q[q]$ by every indicated irreducible
factor; hence it is divisible by their product.  Since that
product is monic and integral, its quotient is again integral.

Taking $q\to1$ in \eqref{IV-D4:leafproduct} shows that every leaf
contributes a rational number of $p$-adic valuation at least
\[
 5k+6(j-k)=5j+(j-k)\ge5j.
\]
Therefore the original hypergeometric sums over each of
$B(u,j)$ and $J(u,j)$ have valuation at least $5j$, uniformly
in $u\ge0$.

Finally write $p=3m+2$.  For an odd level $r=2s+1$,
\[
 \frac{p^{2s+1}-2}{3}=mn^s+2t_s.
\]
The full summation interval is the disjoint union of
$B(0,s),\ldots,B(m-1,s)$ and $J(m,s)$.  Thus, together with
the even-level theorem, we have proved
\begin{equation}\label{IV-D4:allvaluation}
 \boxed{\quad v_p(D_r)\ge5\lfloor r/2\rfloor
       \qquad(p>3,\ p\equiv2\pmod3,\ r\ge1).\quad}
\end{equation}
For odd $r=2s+1\ge5$, both $D_r$ and $p^5D_{r-2}$ have
valuation at least $5s\ge4s+2=2r$.  Hence
\[
 \boxed{D_r\equiv-p^5D_{r-2}\pmod{p^{2r}}
       \qquad(r\ge5\text{ odd},\ p\equiv2\pmod3).}
\]
At this stage the only supersingular recurrence still to prove
is the odd level $r=3$, treated in the next section.
At that level \eqref{IV-D4:allvaluation} proves only $D_3\in p^5
\mathbb Z_p$; the remaining required leading congruence is
$D_3\equiv-p^5D_1\pmod{p^6}$.

\section{The odd third level of D.3 for all supersingular primes}\label{loc-IV-4}
\label{IV-sec:D-cubic}
Let $p>3$ be prime, $p\equiv2\pmod3$. Write
\[
 m=(p-2)/3,\quad \tau=(2p-1)/3,\quad q=p^2,\quad n=(q-1)/3=pm+\tau,
 \qquad b_k=(1/3)_k/k!.
\]
We prove
\begin{equation}\label{IV-dcubic:main}
 \sum_{k=0}^{(p^3-2)/3}(6k+1)b_k^6
 \equiv-p^5\sum_{k=0}^{(p-2)/3}(6k+1)b_k^6\pmod{p^6}.
\end{equation}
This is precisely the odd $r=3$ instance of Swisher D.3.

\subsection{An exact cubic Dougall model}
For $N\equiv1\pmod3$, put $t=(N-1)/3$ and define
\[
 T_3(N)=\sum_{k=0}^t(6k+1)b_k^6
 \prod_{i=0}^{k-1}\frac{1-N^3/(3i+1)^3}{1+N^3/(3i+3)^3}.
\]
Let $\omega^2+\omega+1=0$. Dougall's terminating very-well-poised
${}_7F_6$ summation, with central parameter $1/3$ and the five other
upper parameters
\[
 (1-N)/3,\quad(1-N\omega)/3,\quad(1-N\omega^2)/3,\quad1/3,\quad1/3,
\]
gives
\begin{equation}\label{IV-dcubic:product}
 T_3(N)=
 \frac{(4/3)_t((2-N)/3)_t((2+N\omega)/3)_t((2+N\omega^2)/3)_t}
 {(1+N\omega/3)_t(1+N\omega^2/3)_t(1)_t((1-N)/3)_t}.
\end{equation}
Indeed the five parameters have sum $1+2/3$, as required by the
terminating balanced Dougall identity, and $(1-N)/3=-t$.
The factors with the three cubic roots combine to the displayed
cubic deformation. Equivalently, reversal of the two real shifted
factorials yields the rational identity
\begin{equation}\label{IV-dcubic:rational}
 T_3(N)=N\frac{(1/3)_t(2/3)_t}{(t!)^2}
 \prod_{j=0}^{t-1}
 \frac{(3j+2)^2-N(3j+2)+N^2}
 {(3j+3)^2-N(3j+3)+N^2}.
\end{equation}
These are exact finite identities.

\subsection{A uniform first-digit cancellation}
Define, for $z\in p\mathbb Z_p$,
\[
 L(z)=\sum_{b=0}^{\tau}(6(z+b)+1)
 \left(\frac{(1/3+z)_b}{(1+z)_b}\right)^6.
\]
Then
\begin{equation}\label{IV-dcubic:lowblock}
 L(z)\in p^3\mathbb Z_p\qquad(z\in p\mathbb Z_p).
\end{equation}
First $L(0)\in p^3\mathbb Z_p$. To see this apply the cubic model
with $N=2p$ and $t=\tau<p$. Every factor $3i+1$ and $3i+3$ with
$0\le i<\tau$ is a $p$-adic unit: a possible multiple $p$ has the
wrong residue modulo $3$, and the numerator endpoint $2p$ is
excluded. Hence $L(0)-T_3(2p)\in p^3\mathbb Z_p$.
In \eqref{IV-dcubic:rational}, $(1/3)_\tau/\tau!$ is a unit,
$(2/3)_\tau/\tau!$ has valuation one, and the quadratic-product
ratio has valuation two. The latter assertion follows because its
only numerator index divisible by $p$ is $j=m$, where the quadratic
is $3p^2$, whereas every denominator is a unit. Therefore
$v_p(T_3(2p))=4$, proving the required bound for $L(0)$.

Put $u(z)=(1/3+z)_\tau/(1+z)_\tau$. Reversing the finite products
gives the exact rational identity
\[
 L(z)=-u(z)^6L(-z-2p/3).
\]
Thus $L(-2p/3)\in p^3\mathbb Z_p$. At the midpoint $z=-p/3$,
product reversal gives $u(-p/3)=(-1)^\tau$, so $L(-p/3)=0$ exactly.
All denominators in $L(z)$ are units for $z\in p\mathbb Z_p$.
Consequently the Taylor expansion of $L(px)$ modulo $p^3$ is a
polynomial of degree at most two in $x$, with coefficients in
$\mathbb Z_p$. Its values at the three points
$0,-1/3,-2/3$ vanish modulo $p^3$. The Vandermonde determinant of
these points is a $p$-adic unit, proving \eqref{IV-dcubic:lowblock}.

\subsection{The unit two-digit blocks}
For $0\le c\le m$, define
\[
 K_c=\sum_{\substack{0\le j\le n\\b_{qc+j}\in\mathbb Z_p^\times}}
 (6(qc+j)+1)\left(\frac{b_{qc+j}}{b_c}\right)^6.
\]
For $0\le C\le m$ put $h=3C+1$ and
\[
 E_c(h)=\prod_{i=0}^{c-1}
 \frac{1-h^3/(3i+1)^3}{1+h^3/(3i+3)^3}.
\]
Cubing permutes $\mathbb F_p^\times$. Thus each $E_c(h)$ for
$c\le C$ is a unit: a numerator zero would require $h=3i+1\pmod p$,
which is impossible for $i<C$; a denominator zero would require
$h+3i+3=p$, impossible modulo $3$. The positive sum in the latter
case is less than $2p$.

We claim
\begin{equation}\label{IV-dcubic:block}
 T_3(qh)\equiv\sum_{c=0}^{C}b_c^6E_c(h)K_c\pmod{p^6}.
\end{equation}
The unit indices in the large cutoff are exactly
\[
 k=qc+pa+b,\qquad0\le c\le C,\quad0\le a\le m,\quad0\le b\le\tau.
\]
This follows by two applications of the generalized-binomial Lucas
criterion; the dash parameters are $1/3\mapsto2/3\mapsto1/3$.
All other summands have $b_k^6\in p^6\mathbb Z_p$ and can be omitted:
their deformation factors are units by the same cubic-injectivity
argument for the factors divisible by $p^2$ and are $1$ modulo $p^3$
for all remaining factors.

For a unit index as displayed, the deformation factors divisible
by $p^2$ before cancellation give exactly $E_c(h)$. The factors
with valuation one give a product $P_{c,a}(h)$ which is $1$ modulo
$p^3$ and is independent of $b$. Indeed numerator multiples of
$p$ have indices $\tau+pt$ and denominator multiples have indices
$p-1+pt$; both lists have exactly $pc+a$ entries below $k$, before
removing those divisible by $p^2$. Finally the factors with
valuation zero are $1$ modulo $p^6$.
For fixed $a,c$, the sum of the undeformed summands over the low
digit $b$ equals
\[
 b_{qc+pa}^6L(qc+pa)\in p^3\mathbb Z_p
\]
by \eqref{IV-dcubic:lowblock}. Multiplication by
$P_{c,a}(h)-1\in p^3\mathbb Z_p$ therefore changes it by an element
of $p^6\mathbb Z_p$. This proves \eqref{IV-dcubic:block}.

\subsection{Two-digit scaling of the exact model}
For every $0\le C\le m$, $h=3C+1$, one has
\begin{equation}\label{IV-dcubic:scale}
 T_3(p^2h)\equiv-p^5T_3(h)\pmod{p^6}.
\end{equation}
Here are the product counts and sign in full. Work in the unramified
quadratic extension containing $\omega$ and use
\eqref{IV-dcubic:product}, with $(4/3)_t=N(1/3)_t$.
The large cutoff is $t=p^2C+n$. For a fixed integral $\beta$, extraction
of the $p$-factors from $(\alpha+p^2h\beta)_t$ twice gives
\[
 (\alpha+p^2h\beta)_t
 =p^{L_\alpha+C}(\alpha+h\beta)_C U_{\alpha,\beta},
 \quad L_{1/3}=L_1=pC+m,\quad L_{2/3}=pC+m+1,
\]
where $U_{\alpha,\beta}$ is a unit and its residue modulo $p$ is
independent of $\beta$. For $\alpha=1/3,2/3,1$, the first zero
indices are respectively $\tau,m,p-1$; the new parameters after
one extraction are $2/3+ph\beta,1/3+ph\beta,1+ph\beta$.
The second extraction has exactly $C$ terms in each case.
Substitution in \eqref{IV-dcubic:product} leaves $T_3(h)$, $p^2$ from
the leading $N$, and $p^3$ from the three numerator parameters
$2/3$ versus the three denominator parameters $1$. Hence
$T_3(p^2h)=p^5T_3(h)U$.

The two $1/3$ unit products cancel modulo $p$. The other three
unit-product ratios are identical; complete residue blocks cancel,
and their remaining portions are
\[
 U\equiv(AB)^3\pmod p,\qquad
 A=\frac{\prod_{0\le j\le2m,\ j\ne m}(j+2/3)}{(2m+1)!},
 \qquad B=\frac{(1/3)_{m+1}}{m!}.
\]
The numerator residues in $AB$, after clearing $3^{p-1}$, are
exactly $1,\ldots,p-1$. Wilson's and Fermat's theorems give
\[
 AB\equiv-\frac1{m!(2m+1)!}=(-1)^m=-1\pmod p,
\]
since $m$ is odd. Thus $U\equiv-1\pmod p$. The coarse sum $T_3(h)$
is integral because all its cubic deformation denominators are
units, as already proved. This establishes \eqref{IV-dcubic:scale}.

\subsection{Triangular recovery and the original sum}
The exact small-scale identity
\[
 T_3(h)=\sum_{c=0}^{C}(6c+1)b_c^6E_c(h)
\]
together with \eqref{IV-dcubic:block} and \eqref{IV-dcubic:scale} yields
\[
 \sum_{c=0}^{C}b_c^6E_c(h)\{K_c+p^5(6c+1)\}\equiv0\pmod{p^6}.
\]
The coefficient $b_C^6E_C(h)$ is a unit. Induction on $C$ therefore
proves $K_c\equiv-p^5(6c+1)\pmod{p^6}$ for every $c\le m$.
Finally $(p^3-2)/3=q m+2n$. In each $p^2$ block the unit
low-digit indices are exactly $j=pa+b$ with $a\le m$ and
$b\le\tau$, hence satisfy $j\le n$. The additional final tail
$n<j\le2n$ therefore contains no unit $b_{qm+j}$, and all its
summands lie in $p^6\mathbb Z_p$. Thus the original sum partitions
into the stated unit blocks plus summands divisible by $p^6$. It is
congruent to $\sum_{c=0}^m b_c^6K_c$, which proves
\eqref{IV-dcubic:main}.

\section{The first odd supersingular instance of G.3}\label{loc-IV-5}
\label{IV-sec:G-r3}
Let $p\equiv3\pmod4$ be prime and put
\[
S(n)=\sum_{k=0}^{n}(8k+1)\left(\frac{(1/4)_k}{k!}\right)^4.
\]
Then
\begin{equation}\label{IV-G:r3}
S((p^3-3)/4)\equiv-p^3S((p-3)/4)\pmod{p^4}.
\end{equation}

For $p\ge7$, write $m=(p-3)/4$, $t=(3p-1)/4$,
$b_j=(1/4)_j/j!$, $c_j=(3/4)_j/j!$.
We use Morita's gamma function, the continuous extension to
$\mathbb Z_p$ of
\[
\Gamma_p(n)=(-1)^n\prod_{\substack{1\le j<n\\p\nmid j}}j
\qquad(n\ge1).
\]
It is unit-valued and satisfies
\[
\frac{\Gamma_p(x+1)}{\Gamma_p(x)}=
\begin{cases}-x&x\in\mathbb Z_p^\times,\\-1&x\in p\mathbb Z_p,\end{cases}
\qquad
\Gamma_p(x)\Gamma_p(1-x)=(-1)^{a_0(x)},
\]
where $a_0(x)$ is the representative of $x$ modulo $p$ in
$\{1,\ldots,p\}$.

We require the following elementary local polynomial property.
For every $a\in\mathbb Z_p$, the function $\Gamma_p(a+pX)$
has, modulo $p^4$, a polynomial representative of degree at most
three whose coefficient of $X^j$ belongs to $p^j\mathbb Z_p$.
The same is true of finite products and unit-valued quotients of
these functions. Here is a direct proof, avoiding any assertion
about the coefficients of an infinite Taylor series. The functional
equation gives
\[
\frac{\Gamma_p(a+p(X+1))}{\Gamma_p(a+pX)}
=-\prod_{\substack{0\le i<p\\a+i\in\mathbb Z_p^\times}}(a+i+pX).
\]
The right side lies in $1+p\mathbb Z_p$. Put
$C=-\prod_{a+i\in\mathbb Z_p^\times}(a+i)$ and
$A_j=\sum_{a+i\in\mathbb Z_p^\times}(a+i)^{-j}$.
For $j=1,2,3$ these complete unit power sums are zero modulo $p$.
The logarithm of the displayed ratio consequently satisfies
\[
\log C+pA_1X-\frac{p^2A_2}{2}X^2\pmod{p^4}.
\]
The cubic term is in $p^4\mathbb Z_p$, and all higher terms
are also in $p^4\mathbb Z_p$ for $p\ge7$.
Summing this difference equation for nonnegative integer $X$ gives
\[
\log\frac{\Gamma_p(a+pX)}{\Gamma_p(a)}
\equiv X\log C+pA_1\frac{X(X-1)}2
-\frac{p^2A_2}{2}\frac{X(X-1)(2X-1)}6\pmod{p^4}.
\]
The degree-$j$ coefficients on the right belong to $p^j\mathbb Z_p$.
Exponentiation preserves this coefficient grading modulo $p^4$;
continuity extends the result to all $X\in\mathbb Z_p$.
Products and unit-valued inverses preserve it as well.
Throughout the proof, congruences between the resulting polynomials
may be read coefficientwise: since their degrees are at most three
and $p\ge7$, evaluation at $X=0,1,2,3$ has a unit Vandermonde
determinant. Thus congruence at all $X\in\mathbb Z_p$ implies the
corresponding coefficient congruence.

Define analytic functions on $\mathbb Z_p$ by
\[
F(X)=-\frac{\Gamma_p(1/4+pX)}
 {\Gamma_p(1/4)\Gamma_p(1+pX)},\qquad
J_0(X)=-\frac{\Gamma_p(3/4+pX)}
 {\Gamma_p(3/4)\Gamma_p(1+pX)},
\]
\[
L(X)=\sum_{a=0}^{t}(8a+1+8pX)
 \left(\frac{(1/4+pX)_a}{(1+pX)_a}\right)^4,
\quad H(X)=F(X)^4L(X),
\]
\[
V(X)=\sum_{b=0}^{m}(8b+3+8pX)
 \left(\frac{(3/4+pX)_b}{(1+pX)_b}\right)^4,
\quad J(X)=J_0(X)^4V(X).
\]
All finite-product denominators are units. Also $F(0)=J_0(0)=1$.
Reversing the finite sums, then applying
$\Gamma_p(x)\Gamma_p(1-x)=(-1)^{a_0(x)}$, gives the exact reflections
\begin{equation}\label{IV-G:reflect}
H(X)=-H(-X-3/4),\qquad J(X)=-J(-X-1/4).
\end{equation}
Here the gamma reflection signs disappear after taking fourth powers.
For example the reversed finite-product ratio in $L$ is
\[
\frac{(-t-pX)_a}{(3/4-t-pX)_a}
=\frac{(1/4-p(X+3/4))_a}{(1-p(X+3/4))_a}.
\]
Moreover
$F(X)((1/4+pX)_t/(1+pX)_t)/F(-X-3/4)$ has fourth power one.
The proof for $J$ is identical, with $t$ replaced by $m$.

Set $L_0=L(0)$ and $V_0=V(0)$. We first establish
\begin{equation}\label{IV-G:two-base-values}
L_0\equiv3p\frac{(1/2)_t}{t!}\pmod{p^3},\qquad
V_0\equiv p\frac{(3/2)_m}{m!}\pmod{p^2}.
\end{equation}
Use the terminating very-well-poised Dougall identity
\[
\sum_{k=0}^{M}\frac{2k+A}{A}
\frac{(A)_k(B)_k(C)_k(-M)_k}
{k!(1+A-B)_k(1+A-C)_k(1+A+M)_k}
=\frac{(1+A)_M(1+A-B-C)_M}
{(1+A-B)_M(1+A-C)_M}.
\]
With $(A,B,C,M)=(1/4,1/4,(1+3p)/4,t)$, its left side is
\[
\sum_{a=0}^{t}(8a+1)
\frac{(1/4)_a^2((1-3p)/4)_a((1+3p)/4)_a}
{(a!)^2(1-3p/4)_a(1+3p/4)_a},
\]
and its right side simplifies exactly to $3p(1/2)_t/t!$.
The displayed sum minus $L_0$, modulo $p^3$, is
\[
\frac{9p^2}{16}\sum_{a=0}^{t}(8a+1)b_a^4 R_a,\qquad
R_a=\sum_{j=0}^{a-1}\left((j+1)^{-2}-(j+1/4)^{-2}\right).
\]
This latter sum vanishes modulo $p$: reversal $a\mapsto t-a$ gives
$b_{t-a}^4\equiv b_a^4$, $8(t-a)+1\equiv-(8a+1)$,
and $R_{t-a}\equiv R_a$. The last identity follows by reversing the
summands in $R_t$, which also gives $R_t\equiv0$. The analogous
Dougall specialization with
$(A,B,C,M)=(3/4,3/4,(3+p)/4,m)$, multiplied by $3$, differs
from $V_0$ by a multiple of $p^2$ and evaluates exactly as
$p(3/2)_m/m!$.
This proves \eqref{IV-G:two-base-values}. In particular,
$L_0\in p^2\mathbb Z_p$ and $V_0\in p\mathbb Z_p$.

The Taylor property and the first reflection in \eqref{IV-G:reflect}
now imply that there is $E\in\mathbb Z_p$ such that
\begin{equation}\label{IV-G:H-cubic}
H(X)\equiv\frac{L_0}{3}(8X+3)
 +p^3E\big((8X+3)^3-9(8X+3)\big)\pmod{p^4}.
\end{equation}
Indeed $H$ modulo $p^2$ is affine with constant term zero, so its
reflection forces its linear coefficient to vanish too. Thus $H$ is
coefficientwise divisible by $p^2$. Modulo $p^4$ it has degree at most
three, its cubic coefficient is divisible by $p^3$, and reflection
makes it an odd polynomial in $8X+3$. Evaluating at $X=0$ fixes the
remaining linear coefficient. Similarly,
\begin{equation}\label{IV-G:J-linear}
J(X)\equiv(8X+1)V_0\pmod{p^2}.
\end{equation}

We next prove the uniform block identity, for each integer $c\ge0$,
\begin{equation}\label{IV-G:uniform-block}
\sum_{j=0}^{p^2-1}(8p^2c+8j+1)
 \left(\frac{b_{p^2c+j}}{b_c}\right)^4
 \equiv -p^3(8c+1)\pmod{p^4}.
\end{equation}
Write $j=pb+a$, $0\le a,b<p$. Removing numerator and denominator
factors divisible by $p$ twice shows that only $a\le t$, $b\le m$
can contribute modulo $p^4$; all other ratios $b_{p^2c+j}/b_c$ are
in $p\mathbb Z_p$. On the contributing rectangle one has exactly
\[
b_{pB+a}=F(B)c_B\frac{(1/4+pB)_a}{(1+pB)_a},\quad B=pc+b,
\]
\[
c_{pc+b}=J_0(c)b_c\frac{(3/4+pc)_b}{(1+pc)_b}.
\]
Consequently the block in \eqref{IV-G:uniform-block} is congruent to
\[
J_0(c)^4\sum_{b=0}^{m}
 \left(\frac{(3/4+pc)_b}{(1+pc)_b}\right)^4 H(pc+b).
\]
The $p^3E$ term from \eqref{IV-G:H-cubic} vanishes in this expression
modulo $p^4$: reducing its coefficient modulo $p$, the reversal
$b\mapsto m-b$ fixes $c_b^4$ and negates $8b+3$, hence annihilates
every odd polynomial in $8b+3$. The remaining expression is
\[
\frac{L_0}{3}J(c)
\equiv\frac{L_0V_0}{3}(8c+1)\pmod{p^4}.
\]
Finally \eqref{IV-G:two-base-values} and Wilson's theorem yield
\[
\frac{L_0V_0}{3}\equiv
p^2\frac{(1/2)_t}{t!}\frac{(3/2)_m}{m!}
\equiv-p^3\pmod{p^4}.
\]
For completeness, after removing the unique factor $p$ in
$\binom{2t}{t}$, the last unit quotient is
\[
\frac1p\frac{(1/2)_t}{t!}\frac{(3/2)_m}{m!}
\equiv-\frac{((2m+1)!)^2}{(t!)^2(m!)^2\,4^{t+m}}
\equiv-1\pmod p.
\]
Here $t+m=p-1$, $(t!m!)^2\equiv1$, and
$((2m+1)!)^2\equiv1$ since $2m+1=(p-1)/2$ is odd.
This proves \eqref{IV-G:uniform-block}.

The cutoff $(p^3-3)/4=p^2m+(3p^2-3)/4$ contains all contributing
low-digit pairs in each block $c=0,\ldots,m$, because their largest
index is $pm+t=(p^2-1)/4$. All omitted or added terms have fourth
powers divisible by $p^4$, and $b_c$ is a unit for $0\le c\le m$.
Multiplying \eqref{IV-G:uniform-block} by $b_c^4$ and summing proves
\eqref{IV-G:r3}. For $p=3$, direct evaluation gives
\[
S(6)=\frac{19474365987782658225}{18446744073709551616}
\equiv54=-27\pmod{81}.
\]
Since $S((3-3)/4)=S(0)=1$, this proves \eqref{IV-G:r3} for $p=3$
as well.

\section{Every odd supersingular level of G.3}\label{loc-IV-6}
\label{IV-sec:G-all}
For an odd prime $p\equiv3\pmod4$, set
\[
G_{2s}=\sum_{k=0}^{(p^{2s}-1)/4}(8k+1)\frac{(1/4)_k^4}{k!^4},\qquad
G_{2s+1}=\sum_{k=0}^{(p^{2s+1}-3)/4}(8k+1)\frac{(1/4)_k^4}{k!^4}.
\]
We prove
\begin{equation}\label{IV-G:all-valuation}
v_p(G_r)\ge3\lfloor r/2\rfloor.
\end{equation}
Together with \eqref{IV-G:r3}, namely
$G_3\equiv-p^3G_1\pmod{p^4}$, this proves the full odd supersingular
branch of G.3:
\begin{equation}\label{IV-G:all-odd}
G_r\equiv-p^3G_{r-2}\pmod{p^{r+1}}
\qquad(r\ge3\text{ odd},\ p\equiv3\pmod4).
\end{equation}
Write $(z;Q)_k=\prod_{i=0}^{k-1}(1-zQ^i)$, with a list of
arguments denoting the product of the corresponding symbols.
Set $Q=q^4$, $[j]=(1-q^j)/(1-q)$ and
\[
g_k(a;q)=[8k+1]q^{2k}
 \frac{(aq,q/a,q,q;Q)_k}{(aq^4,q^4/a,q^4,q^4;Q)_k},
\qquad T_L(a;q)=\sum_{k=0}^L g_k(a;q).
\]
The terminating very-well-poised ${}_6\phi_5$ identity is
\[
\sum_{k=0}^{t}\frac{1-AQ^{2k}}{1-A}
\frac{(A,B,C,Q^{-t};Q)_k}
{(Q,AQ/B,AQ/C,AQ^{t+1};Q)_k}
\left(\frac{AQ^{t+1}}{BC}\right)^k
=\frac{(AQ,AQ/(BC);Q)_t}{(AQ/B,AQ/C;Q)_t}.
\]
For $v=4t+1$, substitute $A=q$, $B=q^{1+v}$, $C=q$ and
$Q=q^4$. The argument of the sum is $q^2$, and reversing the
negative powers on the product side gives
\begin{equation}\label{IV-G:q-Jackson}
T_L(q^{\pm v};q)=q^{-t}[v]\frac{(q^2;Q)_t}{(q^4;Q)_t}
\qquad(L\ge t).
\end{equation}
In detail, the product side before reversal is
$(q^5,q^{3-v};Q)_t/(q^{4-v},q^4;Q)_t$;
$(q^{3-v};Q)_t/(q^{4-v};Q)_t
=q^{-t}(q^2;Q)_t/(q;Q)_t$ and
$(q^5;Q)_t/(q;Q)_t=[v]$.
The numerator parameter $q^{1-v}=Q^{-t}$ terminates the series;
no denominator vanishes identically since $v\equiv1\pmod4$.

\paragraph{Cyclotomic prefix congruences.}
Let $h>1$ be odd, let $\epsilon\in\{1,3\}$ satisfy
$h\equiv\epsilon\pmod4$, and set
\[
t=(\epsilon h-1)/4,\qquad L=Bh+t\quad(B\ge0),\qquad
w_h=\begin{cases}1&h\equiv1\pmod4,\\2&h\equiv3\pmod4.\end{cases}
\]
We claim
\begin{equation}\label{IV-G:prefix}
v_{\Phi_h}(T_L(1;q))\ge w_h.
\end{equation}
Work in the local ring of $\mathbb Q(q)$ at $\Phi_h(q)$.
Set
\[
D(a)=(aq^4,q^4/a;Q)_L,\qquad F(a)=a^LD(a)T_L(a;q).
\]
Then $F$ is a polynomial in $a$ over this local ring: $D$ clears
the parameter denominators, $a^L$ clears negative powers, and
$(q;Q)_k/(q^4;Q)_k$ has nonnegative $\Phi_h$ valuation for every
$k$. Indeed in each period the numerator zero precedes the
denominator zero.

For $0\le i\le B$ put
\[
v_i=(4i+\epsilon)h,\quad t_i=(v_i-1)/4=ih+t,
\quad\alpha_{i,\pm}=q^{\pm v_i}.
\]
These $2B+2$ nodes are distinct. At each node
\eqref{IV-G:q-Jackson} has valuation $w_h$. To see this, the
quotient $(q^2;Q)_{t_i}/(q^4;Q)_{t_i}$ has valuation zero
when $h\equiv1\pmod4$, and one when $h\equiv3\pmod4$:
the numerator zero occurs at index $(h-1)/2$, while the
denominator zero occurs at index $h-1$. The factor $[v_i]$
has valuation one in both cases. Also
\[
v_{\Phi_h}(D(\alpha_{i,\pm}))=2B,
\]
since among $1\le j\le L$ exactly $B$ integers satisfy
$h\mid4j$. Therefore $v_{\Phi_h}(F(\alpha_{i,\pm}))\ge2B+w_h$.

Let $P(a)=\prod_{i=0}^{B}(a-\alpha_{i,+})(a-\alpha_{i,-})$.
Monic division gives $F=PH+R$ with integral coefficients and
$\deg R<2B+2$. Hence $v_{\Phi_h}(P(1)H(1))\ge2B+2$.
The exact Lagrange formula is
\[
R(1)=\sum_{\alpha}F(\alpha)
 \prod_{\beta\ne\alpha}\frac{1-\beta}{\alpha-\beta}.
\]
Every factor in every displayed weight has valuation zero:
$1-\beta$ and $\alpha-\beta$ both have valuation exactly one.
Thus $v_{\Phi_h}(R(1))\ge2B+w_h$. Since $w_h\le2$ and
$v_{\Phi_h}(D(1))=2B$, division proves \eqref{IV-G:prefix}.
The same congruence holds at cutoff $Bh+h-1$, because every
extra term past $Bh+t$ has valuation at least four after $a=1$.

\paragraph{Integral specialization.}
After reduction, the denominator of each
$(q;q^4)_k/(q^4;q^4)_k$ has only cyclotomic factors of even
order. All such factors evaluate at $q=1$ to either $1$ or $2$.
Consequently a common denominator for all the specialized terms
$g_k(1;q)$ has value a power of two at $1$. After clearing it,
all cyclotomic divisibilities below are divisibilities of integral
polynomials. Products of distinct monic cyclotomic factors can
therefore be combined before setting $q=1$, without any assertion
of integral comaximality.

\paragraph{The digit-interval partition.}
Put $n=p^2$, $d=(n-1)/4$, $t_j=(n^j-1)/4$ and $t_0=0$. Define
\[
P_j(q)=\prod_{\substack{1\le\ell\le2j\\\ell\text{ odd}}}
\Phi_{p^\ell}(q)^2
\prod_{\substack{1\le\ell\le2j\\\ell\text{ even}}}
\Phi_{p^\ell}(q),\qquad P_0=1.
\]
For every $u\ge0$ let
\[
I(u,j)=[un^j,un^j+t_j]\cap\mathbb Z.
\]
Subtracting two initial prefix congruences from \eqref{IV-G:prefix}
gives
\begin{equation}\label{IV-G:shifted-prefix}
P_j(q)\ \Bigm|\ \sum_{k\in I(u,j)}g_k(1;q).
\end{equation}
Indeed for $h=p^\ell$, $\ell\le2j$, the endpoint $un^j+t_j$
has residue $(3h-1)/4$ if $\ell$ is odd and $(h-1)/4$ if
$\ell$ is even. The left preceding endpoint $un^j-1$ has
residue $h-1$, covered by the last sentence after
\eqref{IV-G:prefix}; for $u=0$ the subtracted sum is empty.

Also define the full block, tail, and three-quarter prefix
\[
\begin{split}
B(u,j)&=[un^j,(u+1)n^j-1]\cap\mathbb Z,\\
T(u,j)&=[un^j+t_j+1,(u+1)n^j-1]\cap\mathbb Z,\\
J(u,j)&=[un^j,un^j+3t_j]\cap\mathbb Z.
\end{split}
\]
At level zero $I=B=J=\{u\}$ and $T$ is empty. For $j\ge1$,
$t_j=dn^{j-1}+t_{j-1}$ gives the disjoint partitions
\[
B(u,j)=I(u,j)\,\dot\cup\,T(u,j),
\]
\[
T(u,j)=T(nu+d,j-1)\,\dot\cup\!
\bigcup_{c=d+1}^{n-1} B(nu+c,j-1),
\]
\[
J(u,j)=I(u,j)\,\dot\cup\,T(nu+d,j-1)
\,\dot\cup\!\bigcup_{c=d+1}^{3d-1} B(nu+c,j-1)
\,\dot\cup\,J(nu+3d,j-1).
\]
Every non-$I$ piece lies in its parent tail. Recursing and
stopping at each $I(v,k)$ partitions both $B(u,j)$ and $J(u,j)$
into finitely many such leaves. A leaf at depth $j-k$ lies in
a tail at every ancestor scale $\ell=k+1,\ldots,j$.
Thus each index in it has remainder greater than $t_\ell$
modulo $p^{2\ell}$, and each summand carries the additional
factor $\Phi_{p^{2\ell}}(q)^4$. Combining these ancestor factors
with \eqref{IV-G:shifted-prefix} yields
\[
P_k(q)\prod_{\ell=k+1}^{j}\Phi_{p^{2\ell}}(q)^4
\ \Bigm|\ \sum_{a\in I(v,k)}g_a(1;q).
\]
All factors here have distinct orders. Clearing the power-of-two
specialization denominator and dividing by the monic integral
product justifies setting $q=1$. Each leaf then has $p$-adic
valuation at least
\[
3k+4(j-k)=3j+(j-k)\ge3j.
\]
The same is true for the sums over $B(u,j)$ and $J(u,j)$.

Finally, writing $p=4m+3$, the odd cutoff satisfies
\[
(p^{2s+1}-3)/4=mn^s+3t_s.
\]
Its interval is the union of $B(0,s),\ldots,B(m-1,s)$ and
$J(m,s)$, giving $v_p(G_{2s+1})\ge3s$. The even level follows
directly from \eqref{IV-G:prefix}, since $P_s(1)=p^{3s}$.
This proves \eqref{IV-G:all-valuation}. For odd $r=2s+1\ge5$, both
$G_r$ and $p^3G_{r-2}$ have valuation at least $3s\ge2s+2=r+1$,
so their signed recurrence follows. Equation \eqref{IV-G:r3} completes \eqref{IV-G:all-odd}, including $p=3$.
For that prime the interval decomposition above remains valid with
$m=0$, so the final interval consists only of $J(0,s)$.

\section{The second supersingular level of A.3}\label{loc-IV-7}
\label{IV-sec:a-second-level}
Write $b_k=(1/2)_k/k!$ and
\[
 A(N)=\sum_{k=0}^{(N-1)/2}(-1)^k(4k+1)b_k^5
 \qquad(N\text{ odd}).
\]
\begin{theorem}\label{IV-thm:a-r2}
For every prime $p\equiv3\pmod4$,
\[
 \boxed{A(p^2)\equiv p^4\pmod{p^8}.}
\]
Thus the entire $r=2$ supersingular branch of Swisher's A.3 holds.
\end{theorem}

\subsection{A shifted short-block divisibility lemma}
For $p\ge7$, $p\equiv3\pmod4$, put $h=(p-1)/2$ and define
\[
 F_p(x)=\sum_{j=0}^{h}(-1)^j(4j+1+4x)
       \left(\frac{(x+1/2)_j}{(x+1)_j}\right)^5.
\]
This rational function has a power series in $\mathbb Z_p[[x]]$.
\begin{lemma}\label{IV-lem:a-shifted-block}
For $x\in p\mathbb Z_p$, $F_p(x)\in p^3\mathbb Z_p$ and
$F'_p(x)\in p^2\mathbb Z_p$. Moreover, for $0\le a\le h$,
\[
 F_p((h-a)p)\equiv-F_p(ap)\pmod{p^4}.
\]
\end{lemma}
\begin{proof}
For any integer $t$, set
\[
 R_t(x)=\frac{(x+1/2)_t}{(x+1)_t},\qquad
 c_k=\frac{(3/4)_k}{(1/4)_k}.
\]
Negative indices mean reciprocals; more explicitly, for $d\ge1$,
\[
 R_{-d}(x)=\frac{\prod_{i=0}^{d-1}(x-i)}
                       {\prod_{i=1}^{d}(x+1/2-i)}.
\]
The shifted WZ pair is
\begin{align*}
 \mathcal F(j,k;x)&=(-1)^{j+k}(4x+4j+1)c_k
               R_j(x)R_{j+k}(x)^2R_{j-k}(x)^2,\\
 \mathcal G(j,k;x)&=8(x+j)(-1)^{j+k}
       \frac{(2k-1)c_k}{4k-1}
       R_j(x)R_{j+k-1}(x)^2R_{j-k}(x)^2.
\end{align*}
Adjacent-factor cancellation verifies the rational-function identity
\[
 \mathcal F(j,k-1;x)-\mathcal F(j,k;x)
 =\mathcal G(j+1,k;x)-\mathcal G(j,k;x).
\]
The displayed formulas give a direct verification: divide the
identity by $\mathcal F(j,k;x)$ and substitute the adjacent
Pochhammer ratios, then clear denominators. Since this is an identity
of rational functions, it remains valid when a numerator factor
vanishes. Summing over $0\le j\le h$ and $1\le k\le h+1$ gives
\begin{equation}\label{IV-eq:a-shift-telescope}
 F_p(x)=\sum_{j=0}^{h}\mathcal F(j,h+1;x)
       +\sum_{k=1}^{h+1}\bigl(\mathcal G(h+1,k;x)-
                              \mathcal G(0,k;x)\bigr).
\end{equation}

First take $x=ap$, $0\le a\le h$. Every expression in
\eqref{IV-eq:a-shift-telescope} is well-defined. The elementary factor
counts give
\[
\begin{array}{c|c}
\text{factor}&\text{lower bound on its }p\text{-valuation}\\ \hline
R_t(x),\ 0\le t\le h&0\quad\text{(in fact a unit)}\\
R_t(x),\ h+1\le t\le p-1&1\\
R_p(x)&0\\
R_{-d}(x),\ 1\le d\le h&1\\
R_{-(h+1)}(x)&0
\end{array}
\]
Indeed, the positive factors acquire the numerator $x+p/2$ at
$t=h+1$ and the denominator $x+p$ at $t=p$. The negative factors
contain $x$ in their numerator and first acquire a nonunit denominator
$x-p/2$ at $d=h+1$. The factors $a+1$ and $a-1/2$ are units for
$0\le a\le h$; the possible extra divisibility of $a+1/2$ only
improves the bounds. If $a=0$, some factors vanish exactly, which
also improves the estimates.

There is exactly one factor of $p$ in $c_{h+1}$, so
$v_p(c_{h+1})=1$. For
\[
 d_k=\frac{(2k-1)c_k}{4k-1}\quad(1\le k\le h+1),
\]
we have $v_p(d_k)\ge0$ and $v_p(d_{h+1})=2$.
The only possible nonunit $4k-1=p$ cancels the final numerator
factor in $c_k$; no denominator factor $4i+1$ in $c_k$ is divisible
by $p$, since $p\equiv3\pmod4$ and $4i+1<2p$.

These counts show that every term on the right of
\eqref{IV-eq:a-shift-telescope} belongs to $p^3\mathbb Z_p$.
For $\mathcal F(j,h+1;x)$, the interior indices $0<j<h$ even have
valuation at least $5$, while each endpoint has valuation at least
$3$. For the upper $\mathcal G$ boundary, $k\le h$ supplies one
factor from $R_{h+1}$ and two from $R_{h+k}$; at $k=h+1$ the two
missing factors are supplied by $d_{h+1}$. For the lower boundary,
$k\le h$ supplies $x$ and $R_{-k}^2$; at $k=h+1$ use $x$ and
$d_{h+1}$. Hence $F_p(ap)\in p^3\mathbb Z_p$ for $0\le a\le h$.

Write $F_p(x)=\sum_{j\ge0}f_jx^j$ with $f_j\in\mathbb Z_p$.
Modulo $p^3$, the three evaluations $a=0,1,2$ say that
$f_0+pf_1a+p^2f_2a^2$ vanishes at three distinct residues. The
Vandermonde determinant is $2$, a unit, so
\[
 f_0\in p^3\mathbb Z_p,\qquad f_1\in p^2\mathbb Z_p,
 \qquad f_2\in p\mathbb Z_p.
\]
All higher terms then imply the asserted divisibility of $F_p(x)$
and $F'_p(x)$ throughout $p\mathbb Z_p$.

Reverse the defining sum with $j\mapsto h-j$. The exact identity is
\[
 F_p(x)=R_h(x)^5 F_p(-p/2-x),
\]
because $h$ is odd. Also $R_h(x)\equiv b_h\equiv-1\pmod p$.
Since both short sums are divisible by $p^3$, this gives
$F_p(-p/2-x)\equiv-F_p(x)\pmod{p^4}$.
For $x=ap$ the difference between $(h-a)p$ and $-p/2-ap$ is
$p^2/2$. Taylor expansion, using $F'_p\in p^2\mathbb Z_p$ and
integrality of all the other Taylor coefficients, changes $F_p$ by
a multiple of $p^4$. This proves the last assertion.
\end{proof}

\subsection{Removing the four-root deformation at the second level}
Let $N=p^2$ and put
\[
 D_k(N)=
 \frac{\prod_{j=0}^{k-1}(1-N^4/(2j+1)^4)}
      {\prod_{j=1}^{k}(1-N^4/(16j^4))},\qquad
 T(N)=\sum_{k=0}^{(N-1)/2}(-1)^k(4k+1)b_k^5D_k(N).
\]
All the factors in $D_k(N)$ are integral units, and
$D_k(N)-1\in p^4\mathbb Z_p$.
Write $k=ap+b$ with $0\le b<p$. If $b>h$, the permitted cutoff
forces $a<h$, and the central-binomial factor $b_k$ has valuation
$1$. Hence these terms contribute to $A(N)-T(N)$ only modulo $p^9$.
For the remaining indices, $0\le a,b\le h$, splitting the
$p$-divisible factors from the others gives
\[
 D_{ap+b}(p^2)\equiv1+p^4L_a\pmod{p^8},\qquad
 L_a=\sum_{j=1}^{a}\frac1{(2j)^4}
              -\sum_{j=0}^{a-1}\frac1{(2j+1)^4}.
\]
The singular numerator indices are exactly $j=up+h$, $0\le u<a$;
the singular denominator indices are $j=up$, $1\le u\le a$.
All other factors are $1$ modulo $p^8$, and products of two
singular corrections are already multiples of $p^8$.

Define the weighted block
\[
 S_a=\sum_{b=0}^{h}(-1)^{ap+b}(4ap+4b+1)b_{ap+b}^5
     =(-1)^a b_{ap}^5F_p(ap).
\]
Lemma~\ref{IV-lem:a-shifted-block} gives $S_a\in p^3\mathbb Z_p$.
Lucas's congruence gives $b_{ap}\equiv b_a\pmod p$, and reversal
of the half-binomial products gives $b_{h-a}\equiv-b_a\pmod p$.
Together with the last assertion of the lemma and the oddness of
$h$, these imply
\[
 S_{h-a}\equiv-S_a\pmod{p^4}.
\]
On the other hand,
\[
 L_{h-a}\equiv L_a\pmod p.
\]
To see this, write $L_a=\sum_{j=1}^{2a}(-1)^j/j^4$.
Reflection $j\mapsto p-j$ reverses parity and preserves inverse
fourth powers modulo $p$, so $L_h\equiv0$ and
$L_h-L_{h-a}\equiv-L_a$. Since $h$ is odd, pairing $a$ with
$h-a$ has no fixed point. Consequently
\[
 \sum_{a=0}^{h}L_aS_a\equiv0\pmod{p^4}.
\]
It follows that
\begin{equation}\label{IV-eq:a-r2-delete}
 A(p^2)-T(p^2)\equiv-p^4\sum_{a=0}^{h}L_aS_a
                  \equiv0\pmod{p^8}.
\end{equation}

\subsection{The explicit second-level model}
For every odd positive integer $L$, define $T(L)$ by the same
finite deformation sum, replacing $N$ by $L$. Its upper parameters
are $5/4,1/2,(1-L)/2,(1+L)/2,(1-iL)/2,(1+iL)/2$ and its lower
parameters are $1/4,1+L/2,1-L/2,1+iL/2,1-iL/2$.
Whipple's terminating ${}_6F_5(-1)$ transformation gives
\[
 T(L)=\frac{\Gamma(1+L/2)\Gamma(1-L/2)}
                {\Gamma(3/2)\Gamma(1/2)}
 {}_3F_2\!\left[
 \begin{matrix}(1-L)/2,(1+L)/2,1/2\\1+iL/2,1-iL/2\end{matrix};1
 \right].
\]
The series terminates at $(L-1)/2$. The classical Whipple
${}_3F_2$ evaluation, whose upper pair sums to $1$ and whose lower
pair sums to twice the third upper parameter plus $1$, gives
\begin{equation}\label{IV-eq:a-r2-classical-gamma}
 T(L)=
 \frac{2\Gamma(1+L/2)\Gamma(1-L/2)
          \Gamma(1+iL/2)\Gamma(1-iL/2)}
 {\displaystyle\prod_{\varepsilon,\eta\in\{\pm1\}}
          \Gamma((3+\varepsilon L+\eta iL)/4)}.
\end{equation}
These are the classical identities applied in H.~Swisher,
\emph{On the supercongruence conjectures of van Hamme},
Research in the Mathematical Sciences \textbf{2} (2015), article~18,
Section~6, equations~(38)--(40),
\url{https://doi.org/10.1186/s40687-015-0037-6}.
They are exact terminating identities before any prime-power
specialization. Gamma shifts and reflection now yield
\[
 T(L)=
 \begin{cases}
 L\displaystyle\prod_{j=1}^{(L-1)/4}
       \frac{(4j-2)^2+L^2}{(4j)^2+L^2},&L\equiv1\pmod4,\\[2mm]
 -L^3\displaystyle\frac{\prod_{j=1}^{(L-3)/4}((4j)^2+L^2)}
       {\prod_{j=0}^{(L-3)/4}((4j+2)^2+L^2)},&L\equiv3\pmod4.
 \end{cases}
\]
For the second formula, put $L=4t+3$ and $z=iL/4$ in
\eqref{IV-eq:a-r2-classical-gamma}. The shift, reflection and duplication formulas reduce it to
$Lz^2(1-z)_t(1+z)_t/((1/2-z)_{t+1}(1/2+z)_{t+1})$, which is exactly
the displayed product.

For $p\equiv3\pmod4$ and $L=p^s$, separate the $p$-divisible
factors in either product and let $V(L)$ denote the remaining unit
quotient. Direct counting gives
\[
 T(L)T(L/p)=-p^3(L/p)^4V(L).
\]
For $L\equiv1\pmod4$ the odd-factor numerator has one more divisible
factor than the even-factor denominator; for $L\equiv3\pmod4$
the two counts are equal. These are precisely the two cases of the
preceding equality. Thus
\[
 T(p^2)=p^4\frac{V(p^2)}{V(p)}.
\]
The exact unit-product gamma formula in both cases is
\[
 V(L)=\frac1{16}\left(
 \prod_{\varepsilon,\eta\in\{\pm1\}}
 \Gamma_p\left(\frac34+\frac{\varepsilon+\eta i}{4}L\right)
 \right)^{-1}.
\]
Here $i^2=-1$ lies in the unramified quadratic extension
$K=\mathbb Q_p(i)$, with valuation ring $\mathcal O_K$. The required
extension of $\Gamma_p$ is justified explicitly by
L.~Long and R.~Ramakrishna, \emph{Some supercongruences occurring in
truncated hypergeometric series}, arXiv:1403.5232v2, Proposition~13
and equations~(2.2)--(2.3),
\url{https://arxiv.org/abs/1403.5232}.
For rational $a\in\mathbb Z_p$, their coefficient bounds are
\[
 c_k(a)=\frac{\Gamma_p^{(k)}(a)}{k!\Gamma_p(a)},\qquad
 c_k(a)\in\mathbb Z_p\ (k<p),\qquad
 v_p(c_k(a))\ge-k\left(\frac1p+\frac1{p-1}\right).
\]
The series $\Gamma_p(a)\sum_{k\ge0}c_k(a)z^k$ converges for
$z\in p\mathcal O_K$ and defines the gamma values used here.
The step and reflection identities on these extended disks are
also equations~(2.2)--(2.3) of that source. In particular the
half-argument pair has value $(-1)^{(p+1)/2}=1$, and
$\Gamma_p(1+iL/4)\Gamma_p(1-iL/4)=1$. The factor $1/16$ in $V(L)$
records the single excess quadratic denominator among its unit
factors in either residue case; it cancels in $V(p^2)/V(p)$.

For $p\ge7$ and $z\in p\mathcal O_K$, those bounds imply
\[
 \frac{\Gamma_p(a+z)}{\Gamma_p(a)}
 \equiv1+c_1(a)z+c_2(a)z^2+c_3(a)z^3
 \pmod{p^4\mathcal O_K}.
\]
Indeed terms $4\le k<p$ are integral multiples of $p^4$;
for $k\ge p$ their valuations are at least
$k(1-1/p-1/(p-1))\ge p-2-1/(p-1)>4$.
This argument includes $p=7$ and does not require the stronger
$p\ge11$ version of Long--Ramakrishna's Theorem~14.
The four shifts form one orbit under multiplication by $i$, so
\[
 V(p)\equiv V(p^2)\equiv\frac1{16}\Gamma_p(3/4)^{-4}\pmod{p^4}.
\]
Both are units. Hence $T(p^2)\equiv p^4\pmod{p^8}$, and
\eqref{IV-eq:a-r2-delete} proves the theorem for $p\ge7$.
Finally, the sole remaining prime $p=3$ is checked exactly:
\[
 A(9)-3^4=-\frac{419492893}{34359738368}\,3^8.
\]
The denominator is a power of $2$, completing the proof.

\section{Verification and remaining targets}\label{loc-IV-8}
The preceding arguments are mathematical proofs, not inferences from
finite numerical tests. The companion code supplies independent checks
of the delicate identities, signs, cutoffs, and intermediate valuation
bounds. It uses integer arithmetic, exact rational arithmetic, and
modular arithmetic with explicit removal of prime factors.

\begin{center}
\small
\begin{tabular}{@{}p{.15\linewidth}p{.77\linewidth}@{}}
\toprule
Check & Recorded coverage \\
\midrule
A.3 & 17,752 shifted WZ identities; 101 exact model identities;
all intermediate bounds for $p=7,11,19,23,31,43$;
the separate exact $p=3$ calculation. \\[4pt]
D.3 & 20 exact cubic-model identities; all coarse cutoffs
$0\le C\le(p-2)/3$ for $p=5,11,17,23,29,41,47$,
including low blocks, model scaling and triangular targets. \\[4pt]
G.3 & 292 normalized two-digit blocks, comprising 275,476 ratio terms,
for $p=7,11,19,23,31,43$; 224 cyclotomic prefix and derivative
evaluations in exact finite fields; 36 interval partitions with
49,176 leaves. \\
\bottomrule
\end{tabular}
\end{center}

\paragraph{What is still missing.}
The A.3 argument removes a fourth-order deformation at the second
level. It does not bound its accumulated error at the required
$p^{5r-2}$ precision for arbitrary $r\ge3$. The ordinary D.3 target
has precision $p^{6r}$; the supersingular vanishing method proved here
does not address that unit-root recurrence. The even supersingular
G.3 target is $p^{4r-2}$, which is much stronger than the divisibility
$p^{3\lfloor r/2\rfloor}$ proved here. The ordinary G.3 branch also
needs a separate residual-sum scaling argument.

For E.3, an exact short-cutoff deformation was derived, but removal
of that deformation at the required odd-level precision remains
unproved. The companion bundle retains that formula and an ordinary
G.3 transformation as explicitly unfinished research checkpoints.
Neither checkpoint is used to establish any theorem in this paper.

\paragraph{Reproducibility.}
The companion archive contains this complete single-file \LaTeX\ source,
the PDF, the proof fragments, the exact-check scripts and their recorded
JSON results. Its README gives working-directory and dependency
instructions. The original statements are taken from Swisher's
published Section~7; published inputs are identified at their point
of use. Internal checking does not replace external mathematical review.

\begin{thebibliography}{9}
\bibitem{IV-Swisher}
H. Swisher, \emph{On the supercongruence conjectures of van Hamme},
Research in the Mathematical Sciences \textbf{2} (2015), Article 18.
\url{https://doi.org/10.1186/s40687-015-0037-6}.

\bibitem{IV-NiWang}
H.-X. Ni and L.-Y. Wang, \emph{Two $q$-supercongruences from Watson's
transformation}, Revista de la Real Academia de Ciencias Exactas,
Fisicas y Naturales. Serie A. Matematicas \textbf{116} (2022), Article 30.
Theorem 2.3 is also reproduced in \cite{IV-Wei}, equation (2.12).

\bibitem{IV-Wei}
C. Wei, \emph{$q$-Supercongruences from Jackson's ${}_8\phi_7$
summation and Watson's ${}_8\phi_7$ transformation}, Journal of
Combinatorial Theory, Series A \textbf{204} (2024), 105853.
\url{https://arxiv.org/abs/2104.07025} (version 4).

\bibitem{IV-LongRamakrishna}
L. Long and R. Ramakrishna, \emph{Some supercongruences occurring in
truncated hypergeometric series}, Advances in Mathematics
\textbf{290} (2016), 773--808.
\url{https://arxiv.org/abs/1403.5232} (version 2).

\bibitem{IV-Corpus}
Verifiable Mathematical Reasoning Corpus, entries [131] and [132],
motivating research material; the present proofs are independent of
the results in these entries.
\url{https://github.com/QEEDDD/Verifiable-Mathematical-Reasoning-Corpus}.
\end{thebibliography}

\clearpage\endgroup

%%%%%%%%%%%%%%%%%%%%%%%%%%%%%%%%%%%%%%%%%%%%%%%%%%%%%%%%%%%%%%%%%%%%%%%%%%
% V | Swisher_Remaining_Branches_Further_Advances.tex
%%%%%%%%%%%%%%%%%%%%%%%%%%%%%%%%%%%%%%%%%%%%%%%%%%%%%%%%%%%%%%%%%%%%%%%%%%
\clearpage\begingroup
\SetPaper{V}{Exact interpolation and positive moments}
\RestoreCounters
\PaperTitle{Exact Interpolation and Positive Moments\\
for Swisher Supercongruences\\
\large Complete A.3 and the remaining D.3 and G.3 branches}{Research manuscript}{10 November 2026}

\begin{abstract}
We prove all residue branches and all levels of Swisher's A.3,
including the primes $3$ and $5$. We also prove every ordinary-prime
level of D.3 and G.3, with their full respective moduli $p^{6r}$
and $p^{4r}$, and every even supersingular G.3 level for $p>3$,
with modulus $p^{4r-2}$. The central tool is an exact interpolation formula for
root-of-unity deformations of terminating hypergeometric sums.
Every positive interpolation moment admits a polynomial formula
whose coefficients have only powers of two in their denominators.
These identities remove all unit nodes together and control the
entire analytic correction when the remaining nodes are divided
by $p$. A separate termwise interpolation argument proves the
supersingular A.3 branch, using an elementary four-factor gamma
product that also covers $p=3$. A companion family with parameter
$3/4$ supplies the intermediate step for even supersingular G.3.
The odd supersingular E.3 branch is not claimed.
\end{abstract}
\noindent\textbf{Status.} This is a research manuscript with complete
proofs of the explicitly stated results, independent internal
mathematical checks, and reproducible exact arithmetic. It has not
undergone external peer review. No historical-priority assertion is made.
\setcounter{tocdepth}{1}

\begin{LocalContents}
\LocalEntry{loc-V-1}{1}{Statements and exact scope}
\LocalEntry{loc-V-2}{2}{A.3 in every supersingular level by exact interpolation}
\LocalEntry{loc-V-3}{3}{Positive interpolation moments and the ordinary branch of A.3}
\LocalEntry{loc-V-4}{4}{Exact interpolation and the ordinary D model}
\LocalEntry{loc-V-5}{5}{Every positive fifth moment of the D interpolating measure}
\LocalEntry{loc-V-6}{6}{The complete ordinary D.3 recursion}
\LocalEntry{loc-V-7}{7}{The entire ordinary branch of G.3 by positive cubic moments}
\LocalEntry{loc-V-8}{8}{The even supersingular branch of G.3 through an alternate family}
\LocalEntry{loc-V-9}{9}{What remains and how the results were checked}
\end{LocalContents}
\clearpage


\section{Statements and exact scope}\label{loc-V-1}
Write $(a)_k=\prod_{j=0}^{k-1}(a+j)$ and $(a)_0=1$. We normalize
$v_p(p)=1$ and put $v_p(0)=+\infty$. Congruences are interpreted
in $\mathbb Q_p$, using the indicated valuation of the difference.
Throughout, $p$ denotes an odd prime.
The labels refer to Swisher's Section~7, Table~2~\cite{V-Swisher}.
The words ordinary and supersingular name the indicated residue
branches; no geometric assertion is implicit in this terminology.
For the indicated positive integers define
\begin{align*}
 A(N)&=\sum_{k=0}^{(N-1)/2}(-1)^k(4k+1)\frac{(1/2)_k^5}{(k!)^5}
 &&(N\text{ odd}),\\
 D(N)&=\sum_{k=0}^{(N-1)/3}(6k+1)\frac{(1/3)_k^6}{(k!)^6}
 &&(N\equiv1\pmod3),\\
 G(N)&=\sum_{k=0}^{(N-1)/4}(8k+1)\frac{(1/4)_k^4}{(k!)^4}
 &&(N\equiv1\pmod4).
\end{align*}
\begin{theorem}[All of A.3]
For every prime $p\equiv1\pmod4$ and every $r\ge1$,
\[
 A(p^r)\equiv-p\Gamma_p(1/4)^4 A(p^{r-1})\pmod{p^{5r}}.
\]
For every prime $p\equiv3\pmod4$,
\[
 A(p)\equiv0\pmod{p^3},\qquad
 A(p^r)\equiv p^4A(p^{r-2})\pmod{p^{5r-2}}\quad(r\ge2).
\]
These statements include $p=5$ and $p=3$, respectively.
\end{theorem}
\begin{theorem}[All ordinary levels of D.3]
For every prime $p\equiv1\pmod3$ and every $r\ge1$,
\[
 D(p^r)\equiv-p\Gamma_p(1/3)^9D(p^{r-1})\pmod{p^{6r}}.
\]
\end{theorem}
\begin{theorem}[All ordinary levels of G.3]
For every prime $p\equiv1\pmod4$ and every $r\ge1$,
\[
 G(p^r)\equiv p\lambda_pG(p^{r-1})\pmod{p^{4r}},\qquad
 \lambda_p=-(-1)^{(p-1)/4}\Gamma_p(1/2)\Gamma_p(1/4)^2.
\]
\end{theorem}
\begin{theorem}[All even supersingular levels of G.3]
For every prime $p>3$ with $p\equiv3\pmod4$ and every $s\ge1$,
\[
 G(p^{2s})\equiv-p^3G(p^{2s-2})\pmod{p^{8s-2}}.
\]
\end{theorem}
\begin{theorem}[Exact supersingular valuations]
For $p\equiv3\pmod4$ and $r\ge0$,
$v_p(A(p^r))=2r+(r\bmod2)$. For $p>3$, $p\equiv3\pmod4$,
and $s\ge1$,
\[
 v_p(G(p^{2s}))=3s.
\]
\end{theorem}

\begin{center}\small
\begin{tabular}{@{}p{.12\textwidth}p{.39\textwidth}p{.40\textwidth}@{}}
\toprule
Label & Branch & Scope proved here\\
\midrule
A.3 & Both residue classes modulo $4$ & Every target level, including $p=3,5$.\\
D.3 & $p\equiv1\pmod3$ & Every level, modulus $p^{6r}$.\\
G.3 & $p\equiv1\pmod4$ & Every level, modulus $p^{4r}$.\\
G.3 & $p>3$, $p\equiv3\pmod4$, even levels & Every target level, modulus $p^{4r-2}$.\\
E.3 & $p\equiv2\pmod3$, odd $r\ge3$ & Remains unresolved.\\
\bottomrule
\end{tabular}
\end{center}
The earlier companion manuscript~\cite{V-Companion} established the $p\equiv2\pmod3$
branch of D.3 and the odd supersingular branch of G.3. Those results
are not reused as hypotheses in the new proofs. Combining the
ordinary D.3 theorem here with that earlier result completes D.3
within this research series. Similarly, the two new G.3 branches
and the earlier odd-level result complete G.3 for $p>3$.
The present document does not repeat the earlier supersingular
D and odd-level G proofs. B.3, H.3, I.3 and L.3 are outside
the scope of this continuation.

\paragraph{Why moments are needed.}
Some individual interpolation terms have negative valuations at
split primes. The proofs therefore do not infer integrality of a
sum from an unjustified termwise bound. Instead every positive
moment is evaluated exactly; an Euler-exponent limit removes unit
nodes together, and the remaining analytic correction is expanded
against the already-controlled moments. The polynomial operator
used for these moment identities has unit triangular diagonal over
$\mathbb Z[1/2]$, so division by a varying odd prime never occurs.
The standard terminating evaluations used below are recorded in
\cite{V-DLMF}, equations 16.4.7, 16.4.9 and 16.4.10.

\section{A.3 in every supersingular level by exact interpolation}\label{loc-V-2}
\label{V-sec:a-supersingular-all}
Set $b_k=(1/2)_k/k!$ and
\[
 A(N)=\sum_{k=0}^{(N-1)/2}(-1)^k(4k+1)b_k^5\qquad(N\text{ positive and odd}).
\]
\begin{theorem}\label{V-thm:a-supersingular-all}
For every prime $p\equiv3\pmod4$ and every $r\ge2$,
\[
 A(p^r)\equiv p^4A(p^{r-2})\pmod{p^{5r-2}}.
\]
The same argument also gives $A(p)\equiv0\pmod{p^3}$.
\end{theorem}

\subsection{An exact finite interpolation formula}
For odd $N$, write $h=(N-1)/2$ and define
\[
 S_N(z)=\sum_{k=0}^h(-1)^k(4k+1)b_k^5
 \frac{\prod_{j=0}^{k-1}(1-z^4/(2j+1)^4)}
 {\prod_{j=1}^k(1-z^4/(2j)^4)}.
\]
Put $T(H)=S_H(H)$ for every positive odd $H$.  The rational function
$S_N(z)$ depends only on $y=z^4$, and multiplication by
$D_N(y)=\prod_{j=1}^h(1-y/(2j)^4)$ gives a polynomial of degree at most
$h$. At each of the $h+1$ distinct nodes $y=H^4$, $H=1,3,\ldots,N$,
the sum terminates at $(H-1)/2$, so $S_N(H)=T(H)$. Lagrange
interpolation at zero consequently gives the exact rational identity
\begin{equation}\label{V-eq:a-lagrange-exact}
 A(N)=\sum_{\substack{1\le H\le N\\H\ \mathrm{odd}}}W_N(H)T(H),
 \qquad
 W_N(H)=\frac{\prod_{\substack{2\le E<N\\E\ \mathrm{even}}}(1-H^4/E^4)}
 {\prod_{\substack{1\le K\le N\\K\ \mathrm{odd},\ K\ne H}}(1-H^4/K^4)}.
\end{equation}
All denominators here are nonzero rational numbers. In particular,
there is no limiting convention at $H=N$.

\begin{lemma}\label{V-lem:a-weight-scaling}
Let $N=p^r$, where $p$ is an odd prime, and let $H\le N$ be odd with
$s=v_p(H)\ge1$. Put $M=N/p$. Then
\[
 \frac{W_N(H)}{W_M(H/p)}\equiv1\pmod{p^{r+4s}}.
\]
\end{lemma}
\begin{proof}
In the exact product \eqref{V-eq:a-lagrange-exact}, separate indices
divisible by $p$. Because $p$ is odd, division by $p$ preserves parity;
these factors are exactly $W_M(H/p)$. The other factors give
\[
 \frac{W_N(H)}{W_M(H/p)}=
 \frac{\prod_{\substack{2\le E<N\\E\ \mathrm{even},\ p\nmid E}}(1-H^4/E^4)}
 {\prod_{\substack{1\le K<N\\K\ \mathrm{odd},\ p\nmid K}}(1-H^4/K^4)}.
\]
The map $E\mapsto K=N-E$ bijects these two index sets. Each paired
factor differs from $1$ by
\[
 \frac{1-H^4/E^4}{1-H^4/K^4}-1
 =\frac{H^4(E^4-K^4)}{E^4(K^4-H^4)}.
\]
Its denominator is a unit. Since $K=N-E$, the numerator has valuation
at least $4s+r$. Multiplying proves the assertion, including $H=N$.
\end{proof}

\begin{lemma}\label{V-lem:a-weight-valuation}
Suppose $p\equiv3\pmod4$, $N=p^r$, and $H\le N$ is odd. If
$s=v_p(H)$, then
\[
 v_p(W_N(H))=5(r-s).
\]
\end{lemma}
\begin{proof}
An elementary factorization of differences of fourth powers gives
\begin{equation}\label{V-eq:a-weight-factorization}
 W_N(H)=2(N/H)^4b_{(N-1)/2}^{\,4}
 b_{(N-H)/2}b_{(N+H)/2}
 \frac{\prod_{\substack{2\le E<N\\E\ \mathrm{even}}}(E^2+H^2)}
 {\prod_{\substack{1\le K\le N\\K\ \mathrm{odd},K\ne H}}(K^2+H^2)}.
\end{equation}
For clarity, the difference-factor identity used here is
\[
 \frac{\prod_{E\ \mathrm{even},\,2\le E<N}(E^2-H^2)}
 {\prod_{K\ \mathrm{odd},\,1\le K\le N,\,K\ne H}(K^2-H^2)}
 =2b_{(N-H)/2}b_{(N+H)/2};
\]
the remaining normalization is $(Nb_{(N-1)/2}/H)^4$.
Since $-1$ is a nonsquare modulo $p$, a quadratic factor $X^2+H^2$
has valuation $2\min(v_p(X),s)$. At every division by $p$, the numbers
of divisible numerator and denominator factors in the final quotient
in \eqref{V-eq:a-weight-factorization} are equal: both are
$(p^{r-1}-1)/2$, because the denominator omits $H$. Repeating until
$H/p^s$ is a unit proves that this quotient is a unit.
Also $b_{(N-1)/2}$ and $b_N$ are units, by Lucas's theorem. If $H<N$,
put $k=(N-H)/2$, so $v_p(k)=v_p(N-k)=s<r$. The identity
\[
 \frac{b_kb_{N-k}}{b_N}
 =\frac{\binom Nk^2}{\binom{2N}{2k}}
\]
shows that $v_p(b_kb_{N-k})=r-s$. Indeed
$v_p\binom Nk=r-s$, and
$v_p\binom{2N}{2k}=r-s$: in the latter equality use
$\binom{2N}{2k}=(N/k)\binom{2N-1}{2k-1}$, whose last factor is a unit
by Lucas because $2k-1<N$ and all the lower $r$ base-$p$ digits of
$2N-1$ equal $p-1$. Formula \eqref{V-eq:a-weight-factorization} now gives
$4(r-s)+(r-s)$. At $H=N$ all factors except $b_N$ are units, giving
valuation zero as required.
\end{proof}

\subsection{The exact Whipple model at arbitrary odd nodes}
The classical terminating Whipple transformations give
\[
 T(L)=\frac{2\Gamma(1+L/2)\Gamma(1-L/2)
 \Gamma(1+iL/2)\Gamma(1-iL/2)}
 {\prod_{\varepsilon,\eta\in\{\pm1\}}
 \Gamma((3+\varepsilon L+\eta iL)/4)}.
\]
These are exactly the finite identities in Swisher, Section~6,
equations~(38)--(40), before specializing $L$ to a prime. Gamma shifts
and reflection yield, for every positive odd $L$,
\begin{equation}\label{V-eq:a-node-model}
 T(L)=\begin{cases}
 L\displaystyle\prod_{j=1}^{(L-1)/4}
       \frac{(4j-2)^2+L^2}{(4j)^2+L^2},&L\equiv1\pmod4,\\[2mm]
 -L^3\displaystyle\frac{\prod_{j=1}^{(L-3)/4}((4j)^2+L^2)}
 {\prod_{j=0}^{(L-3)/4}((4j+2)^2+L^2)},&L\equiv3\pmod4.
 \end{cases}
\end{equation}
For $p\equiv3\pmod4$ and any odd $L$ divisible by $p$, extraction of
the divisible quadratic factors in these two cases gives
\begin{equation}\label{V-eq:a-node-product}
 T(L)T(L/p)=-p^3(L/p)^4V(L),\qquad
 V(L)=\frac1{16}\left(\prod_{\varepsilon,\eta\in\{\pm1\}}
 \Gamma_p\left(\frac34+\frac{\varepsilon+\eta i}{4}L\right)\right)^{-1}.
\end{equation}
Here $i^2=-1$ is in the unramified quadratic extension of $\mathbb Q_p$.
Here is an explicit verification for arbitrary $L=pU$, rather than
only for prime powers. Write $p=4q+3$ and $z=iL/4$. If $L=4t+1$,
then $U=4v+3$ and $t=pv+3q+2$. The divisible numerator indices
in the first line of \eqref{V-eq:a-node-model} are
$j=(p+1)/2+pw$, $0\le w\le v$; the divisible denominator indices
are $j=pw$, $1\le w\le v$. Their quotient is
$-p^2U^3/T(U)$. If $L=4t+3$, then $U=4v+1$ and $t=pv+q$.
The divisible numerator indices are $j=pw$, $1\le w\le v$;
the divisible denominator indices, written as $(4j+2)^2+L^2$,
are $j=(p-1)/2+pw$, $0\le w<v$. Their quotient is $U/T(U)$.
Retaining the unit quotient $V(L)$ and multiplying by the external
factor $L$, respectively $-L^3$, proves the first identity in \eqref{V-eq:a-node-product}.

For the unit quotient, the corresponding Pochhammer pairs are
$(1/2\pm z)_t/(1\pm z)_t$ in the first case, and
$(1\pm z)_t/(1/2\pm z)_{t+1}$ in the second. Converting each raw
quadratic to these normalized factors contributes $16^{-1}$:
in the first case the unit numerator count is one smaller, and
in the second the total denominator count is one larger while
the divisible counts agree. The exact unit-factor formula
\[
 \prod_{\substack{0\le j<m\\a+j\text{ a unit}}}(a+j)
 =(-1)^m\frac{\Gamma_p(a+m)}{\Gamma_p(a)}
\]
therefore expresses $V(L)$, in either case, as $1/16$ times
\[
 \frac{\Gamma_p((L+1)/4+z)\Gamma_p((L+1)/4-z)}
 {\Gamma_p((L+3)/4+z)\Gamma_p((L+3)/4-z)}.
\]
The constant gamma pairs at $1/2\pm z$ and at $1\pm z$ both equal
$1$, because $p\equiv3\pmod4$. Reflect the two numerator arguments
to $(3-L)/4\mp z$. The reflection signs multiply to $1$, as the
two arguments belong to the same residue disk. This proves the
second identity in \eqref{V-eq:a-node-product}.
\subsection{An integral four-orbit gamma estimate, including $p=3$}
We use the local analytic extension of Morita's gamma function to
$a+p\mathcal O_{\mathbb Q_p(i)}$. It is available for every odd prime:
Dwork, \emph{A note on the $p$-adic gamma function}, Groupe d'\'etude
d'analyse ultram\'etrique \textbf{9} (1981--1982), no.~3, expos\'e J5, pp.~J1--J10,
\url{https://numdam.org/item/GAU_1981-1982__9_3_A6_0.pdf},
explicitly assumes $p\ne2$ and gives the analytic radius
$p^{-1/p-1/(p-1)}>p^{-1}$, which contains each of these disks.
The step and reflection formulas extend to these disks by the
analytic identity theorem: on each disk they agree at the rational
$p$-adic points, which have an accumulation point in that disk.
For $p\ge5$, the same extension and identities are also recorded
in Long--Ramakrishna, arXiv:1403.5232v2, Proposition~13 and equations
(2.2)--(2.3). We do not use their $p\ge5$ assumption for the case $3$.

\begin{lemma}\label{V-lem:a-four-orbit-gamma}
For any odd prime $p$, $a\in\mathbb Z_p$, and $z\in p^s\mathcal O$,
$s\ge1$, where $\mathcal O$ is the ring of integers of a finite
extension containing $i$, one has
\[
 \frac{\prod_{\varepsilon,\eta\in\{\pm1\}}
 \Gamma_p(a+(\varepsilon+\eta i)z/4)}{\Gamma_p(a)^4}
 \equiv1\pmod{p^{4s}\mathcal O}.
\]
\end{lemma}
\begin{proof}
Let $\mathcal C=\{(\varepsilon+\eta i)/4:\varepsilon,\eta=\pm1\}$.
The step and reflection identities give
$\Gamma_p(x)\Gamma_p(-x)=1$ for $x\in p\mathcal O$.
Because $\mathcal C$ consists of two opposite pairs,
$\prod_{c\in\mathcal C}\Gamma_p(cz)=1$.
For every positive integer $m$, the step formula therefore gives
\[
 \frac{\prod_{c\in\mathcal C}\Gamma_p(m+cz)}{\Gamma_p(m)^4}
 =\prod_{\substack{1\le t<m\\p\nmid t}}
      \prod_{c\in\mathcal C}(1+cz/t)
 =\prod_{\substack{1\le t<m\\p\nmid t}}
      \left(1+\frac{z^4}{64t^4}\right)
 \equiv1\pmod{p^{4s}\mathcal O}.
\]
Choose positive integers $m_j\to a$ $p$-adically. Eventually these
centers lie in the same residue disk as $a$, and $m_j+cz\to a+cz$
inside that one analytic disk. Local analytic continuity consequently
allows passage to the limit in the last congruence; the ideal
$p^{4s}\mathcal O$ is closed. This proves the lemma, including $p=3$.
\end{proof}
Taking $a=3/4$ and $z=L$ in this lemma and then using
\eqref{V-eq:a-node-product} proves, whenever $v_p(L)=s\ge1$,
\[
 V(L)\equiv\frac1{16}\Gamma_p(3/4)^{-4}\pmod{p^{4s}}.
\]
If $p\nmid u$, every quadratic factor in $T(u)$ is a unit. Taking
valuations in \eqref{V-eq:a-node-product} and then iterating gives
\begin{equation}\label{V-eq:a-node-valuation}
 v_p(T(p^su))=2s+(s\bmod2).
\end{equation}
Dividing two consecutive instances of \eqref{V-eq:a-node-product} also
gives, for $s\ge2$,
\begin{equation}\label{V-eq:a-node-two-step}
 \frac{T(p^su)}{p^4T(p^{s-2}u)}
 =\frac{V(p^su)}{V(p^{s-1}u)}
 \equiv1\pmod{p^{4s-4}}.
\end{equation}

\subsection{Transporting the interpolating sum}
Apply \eqref{V-eq:a-lagrange-exact} at $N=p^r$. A node $H$ of valuation
$s$ contributes a term of valuation
\[
 v_p(W_N(H)T(H))=5(r-s)+2s+(s\bmod2)
 =5r-3s+(s\bmod2).
\]
The terms with $s=0,1$ therefore vanish modulo $p^{5r-2}$. Every
remaining node has $H=p^2J$ with $1\le J\le p^{r-2}$ odd. Applying
Lemma~\ref{V-lem:a-weight-scaling} twice gives
\[
 \frac{W_{p^r}(H)}{W_{p^{r-2}}(H/p^2)}
 \equiv1\pmod{p^{r+4s-5}}.
\]
This precision is at least $4s-4$, since the difference is $r-1$. Combining it with
\eqref{V-eq:a-node-two-step} shows that
\[
 W_{p^r}(H)T(H)-p^4W_{p^{r-2}}(H/p^2)T(H/p^2)
\]
has valuation at least
\[
 (5r-3s+(s\bmod2))+(4s-4)
 =5r+s+(s\bmod2)-4\ge5r-2,
\]
because $s\ge2$. Summing over the surviving nodes proves the theorem.
At $r=1$, the same valuation formula gives at least $3$ for every
node, proving the initial vanishing.

\begin{corollary}
For every prime $p\equiv3\pmod4$ and every $r\ge0$,
\[
 v_p(A(p^r))=2r+(r\bmod2).
\]
\end{corollary}
\begin{proof}
The endpoint $H=p^r$ in \eqref{V-eq:a-lagrange-exact} has precisely
that valuation. Each other node has $s<r$ and valuation
$5r-3s+(s\bmod2)$, which is strictly larger: decreasing $s$ by
one increases this expression by either $2$ or $4$. Therefore the
endpoint cannot cancel with the other terms.
\end{proof}

\paragraph{Scope of the result.}
This proves the complete $p\equiv3\pmod4$ branch of A.3, including
$p=3$ and every level. The $p\equiv1\pmod4$ branch is proved in the
following section by a different argument: the valuation of the
interpolation weights used here relies on $-1$ being a nonsquare.

\section{Positive interpolation moments and the ordinary branch of A.3}\label{loc-V-3}
Retain the exact interpolation formula
\[
A(N)=\sum_{H=1,3,\ldots,N}W_N(H)T(H),\qquad
b_k=(1/2)_k/k!,\quad h=(N-1)/2.
\]
All identities below are exact over $\mathbb Q$ before prime-power
specialization.

\subsection{The positive moments}
There are polynomials $P_j(X)\in\mathbb Z[1/2][X]$ for $j\ge0$ such that
\begin{equation}\label{V-A6:moment}
\sum_{H=1,3,\ldots,N}W_N(H)T(H)H^{4j}
=(-1)^hN^5b_h^5P_{j-1}(N)\qquad(j\ge1).
\end{equation}
Here $P_0=1$, and every $P_j$ is invariant under $X\mapsto-1-X$.
In particular, for every odd prime $p$, $N=p^r$ and $j\ge1$, the
left side of \eqref{V-A6:moment} belongs to $p^{5r}\mathbb Z_p$.

To prove this, write $y=z^4$ and put
\[
D_N(y)=\prod_{a=1}^h(1-y/(2a)^4),\qquad
Q_N(y)=\prod_{a=0}^h(y-(2a+1)^4),\qquad
F_N(y)=D_N(y)S_N(z).
\]
The established interpolation says $F_N(H^4)=D_N(H^4)T(H)$.
The partial-fraction identity, or Lagrange interpolation of the
remainder of $y^jF_N(y)$ modulo $Q_N(y)$, gives
\[
\sum_H W_N(H)T(H)H^{4j}
=-Q_N(0)[y^{-1}]\frac{y^{j-1}F_N(y)}{Q_N(y)}.
\]
Cancellation of the first $k$ odd factors in the summand defining
$S_N$ gives
\[
\frac{F_N(y)}{Q_N(y)}
=\frac{(-1)^h}{\prod_{a=1}^h(2a)^4}
\sum_{k=0}^h(-1)^k(4k+1)b_k
\frac{\prod_{a=k+1}^h(y-(2a)^4)}
{\prod_{a=k}^h(y-(2a+1)^4)}.
\]
Define the formal power series
\[
\mathcal R_h(x)=\sum_{k=0}^h(-1)^k(4k+1)b_k
\frac{\prod_{a=k+1}^h(1-(2a)^4x)}
{\prod_{a=k}^h(1-(2a+1)^4x)}.
\]
It follows that the moment is $N^4b_h^4[x^{j-1}]\mathcal R_h(x)$.
The exact recurrence is
\[
(1-(2h+1)^4x)\mathcal R_h(x)
=(1-(2h)^4x)\mathcal R_{h-1}(x)+(-1)^h(4h+1)b_h.
\]
The initial series is $\mathcal R_0(x)=(1-x)^{-1}$.
We now prove
\[
\mathcal R_h(x)=(-1)^h(2h+1)b_h
\sum_{j\ge0}P_j(2h+1)x^j.
\]
The coefficient recurrence required for this assertion is
\begin{equation}\label{V-A6:P-rec}
X P_j(X)+(X-1)P_j(X-2)
=X^5P_{j-1}(X)+(X-1)^5P_{j-1}(X-2)\qquad(j\ge1).
\end{equation}
For $j=0$, $P_0=1$ gives the inhomogeneous constant coefficient
$2X-1$.

Here is the polynomial existence and integrality argument.
Over $R=\mathbb Z[1/2]$, the polynomials invariant under
$X\mapsto-1-X$ are exactly $R[X(X+1)]$. The polynomials
anti-invariant under $X\mapsto1-X$ are exactly
$(2X-1)R[X(X-1)]$. The operator
\[
\mathcal L P=XP(X)+(X-1)P(X-2)
\]
maps the first space into the second. In the bases
$[X(X+1)]^d$ and $(2X-1)[X(X-1)]^d$, it is triangular with
diagonal entries one: both leading terms are $2X^{2d+1}$.
It is therefore an isomorphism on each finite degree filtration,
over $R$. If $P_{j-1}$ is invariant, the right side of
\eqref{V-A6:P-rec} is anti-invariant because the exponent $5$ is odd.
Thus \eqref{V-A6:P-rec} uniquely defines an invariant $P_j\in R[X]$.
Moreover $\deg P_j\le4j$. At $X=1$ the same recurrence gives
$P_j(1)=P_{j-1}(1)=1$, which supplies the initial series at $h=0$.
The recurrence for $\mathcal R_h$ now proves the asserted coefficient
formula by induction on $h$, and hence \eqref{V-A6:moment}.

\subsection{Removing the unit nodes}
For $N=p^r$, let $\mu_N(H)=W_N(H)T(H)$. For positive integers
$e_v=(p-1)p^v$, one has
\[
H^{4e_v}\longrightarrow
\begin{cases}1,&p\nmid H,\\0,&p\mid H\end{cases}
\quad\text{in }\mathbb Q_p.
\]
Since the sum has finitely many nodes, \eqref{V-A6:moment} and closedness
of $p^{5r}\mathbb Z_p$ imply
\begin{equation}\label{V-A6:units}
\sum_{\substack{1\le H\le N\\H\text{ odd},\ p\nmid H}}
\mu_N(H)\in p^{5r}\mathbb Z_p.
\end{equation}
Individual unit nodes need not satisfy this bound; their exact
moment cancellation is essential.

\subsection{Ordinary nodes divisible by $p$}
Let $p\equiv1\pmod4$, so $p\ge5$, and fix $i\in\mathbb Z_p$ with
$i^2=-1$. For every positive odd $L$ divisible by $p$, the two exact
Whipple product forms imply
\begin{equation}\label{V-A6:ordinary-node}
T(L)=-p\Gamma_p(1/4)^4\,T(L/p)\,U(L),
\qquad
U(L)=\prod_{\varepsilon,\eta\in\{\pm1\}}
\frac{\Gamma_p(1/4+(\varepsilon+\eta i)L/4)}{\Gamma_p(1/4)}.
\end{equation}
For completeness, write $p=4q+1$ and $L=pU$.
When $L\equiv1\pmod4$, write $U=4v+1$; the product length is
$(L-1)/4=pv+q$. The divisible numerator indices are
$j=(p+1)/2+pw$, $0\le w<v$, and the divisible denominator
indices are $j=pw$, $1\le w\le v$. After removing the powers of
$p$, their quotient is exactly $T(U)/U$, and the external factor
$L$ contributes $p$. When $L\equiv3\pmod4$, write $U=4v+3$;
the product length is $(L-3)/4=pv+3q$. The divisible numerator
indices are $j=pw$, $1\le w\le v$, while the denominator indices
are $j=(p-1)/2+pw$, $0\le w\le v$. Their normalized quotient
is $T(U)/(-U^3)$, with one extra denominator power $p^2$; hence
the external factor $-L^3$ again leaves $pT(U)$. In both cases the unit
numerator and denominator counts agree, so no factor $1/16$ occurs.
The unit quotient is
\[
-\left\{\prod_{\varepsilon,\eta}
\Gamma_p(3/4+(\varepsilon+\eta i)L/4)\right\}^{-1}.
\]
Indeed the pair at $1/2\pm iL/4$ has reflection value $-1$, whereas
the pair at $1\pm iL/4$ has value $1$. Reflection of the four
arguments changes the displayed inverse product into the product at
$1/4$ in \eqref{V-A6:ordinary-node}; the four reflection signs multiply
to one. These identities hold even when $T(L/p)$ is nonintegral.

The scaled gamma correction has a restricted power series
\begin{equation}\label{V-A6:gamma-coeff}
U(pX)=1+\sum_{j\ge1}u_jX^{4j},\qquad
u_j\in p^4\mathbb Z_p,\quad v_p(u_j)\longrightarrow+\infty.
\end{equation}
To justify the entire series, not only its first coefficient, use
Long--Ramakrishna's bounds
\[
c_k(a)=\frac{\Gamma_p^{(k)}(a)}{k!\Gamma_p(a)},\qquad
c_k(a)\in\mathbb Z_p\ (k<p),\qquad
v_p(c_k(a))\ge-k\left(\frac1p+\frac1{p-1}\right).
\]
The four shifts form an orbit under multiplication by $i$, so only
degrees divisible by four survive. Degree four has valuation at
least four because $4<p$. A term of degree $4j\ge8$ in the product
of four scaled series has valuation at least
\[
4j\left(1-\frac1p-\frac1{p-1}\right)
\ge\frac{11}{20}\,4j>4.
\]
These lower bounds tend to infinity, proving \eqref{V-A6:gamma-coeff},
including $p=5$.

Now let $N=p^r$, $M=p^{r-1}$ and $H=pJ$. The exact weight quotient is
\[
R_N(J)=\frac{W_N(pJ)}{W_M(J)}
=\frac{\prod_{\substack{2\le E<N\\E\text{ even},p\nmid E}}
 (1-p^4J^4/E^4)}
{\prod_{\substack{1\le K<N\\K\text{ odd},p\nmid K}}
 (1-p^4J^4/K^4)}.
\]
As a series in $J^4$, it has constant term one and its coefficient
of degree $j\ge1$ belongs to $p^{4j}\mathbb Z_p$. Thus the product
\[
R_N(J)U(pJ)=1+\sum_{j\ge1}e_{N,j}J^{4j}
\]
has $e_{N,j}\in p^4\mathbb Z_p$ and $v_p(e_{N,j})\to+\infty$.
Equations \eqref{V-A6:ordinary-node} and the weight quotient give
\[
\mu_N(pJ)=-p\Gamma_p(1/4)^4\mu_M(J)
\left(1+\sum_{j\ge1}e_{N,j}J^{4j}\right).
\]
Each positive moment at scale $M$ is in $p^{5(r-1)}\mathbb Z_p$
by \eqref{V-A6:moment}. Summing over the finitely many coarse nodes,
then the convergent correction series, proves
\[
\sum_{p\mid H}\mu_N(H)
\equiv-p\Gamma_p(1/4)^4 A(M)\pmod{p^{5r}}.
\]
For $r=1$ this uses the same identity at $M=1$, where every positive
moment equals one; the coefficient bound still supplies $p^5$.
Finally add \eqref{V-A6:units}. We have proved
\begin{equation}\label{V-A6:ordinary-main}
\boxed{\quad A(p^r)\equiv-p\Gamma_p(1/4)^4 A(p^{r-1})
\pmod{p^{5r}}\qquad(p\equiv1\pmod4,\ r\ge1).\quad}
\end{equation}
This is the full ordinary branch of Swisher's A.3, including $p=5$.

\section{Exact interpolation and the ordinary D model}\label{loc-V-4}
\label{V-sec:d-ordinary-preliminaries}
For $N\equiv1\pmod3$, set $n=(N-1)/3$ and
\[
 b_k=\frac{(1/3)_k}{k!},\qquad
 D(N)=\sum_{k=0}^n(6k+1)b_k^6.
\]
We prove below that for every prime $p\equiv1\pmod3$ and every
$r\ge1$,
\[
 D(p^r)\equiv-p\Gamma_p(1/3)^9D(p^{r-1})\pmod{p^{6r}}.
\]
The interpolation formulas first hold over the rational numbers;
the subsequent gamma identities are interpreted $p$-adically.

\subsection{Exact fifth-degree interpolation}
Define
\[
 S_N(z)=\sum_{k=0}^n(6k+1)b_k^6
 \prod_{j=0}^{k-1}\frac{1-z^5/(3j+1)^5}{1+z^5/(3j+3)^5},
 \qquad T(H)=S_H(H)\quad(H\equiv1\pmod3).
\]
In the variable $y=z^5$, multiplication by
$\prod_{j=1}^n(1+y/(3j)^5)$ makes $S_N$ a polynomial of degree at
most $n$. At the $n+1$ distinct nodes $y=H^5$,
$H=1,4,\ldots,N$, the rational function satisfies $S_N(H)=T(H)$;
the cleared polynomial has value $\prod_{j=1}^n(1+H^5/(3j)^5)T(H)$. Consequently,
Lagrange interpolation at zero gives
\begin{equation}\label{V-eq:d-lagrange-exact}
 D(N)=\sum_{\substack{1\le H\le N\\H\equiv1\ (3)}}W_N(H)T(H),
 \qquad
 W_N(H)=\frac{\prod_{\substack{1\le E<N\\E\equiv0\ (3)}}(1+H^5/E^5)}
 {\prod_{\substack{1\le K\le N\\K\equiv1\ (3),\ K\ne H}}(1-H^5/K^5)}.
\end{equation}
These are exact rational identities with no vanishing denominator.

\begin{lemma}\label{V-lem:d-weight-scale}
Suppose $p\equiv1\pmod3$, $N=p^r$, and $H\equiv1\pmod3$ satisfies
$1\le H\le N$, $s=v_p(H)\ge1$. Then
\[
 \frac{W_N(H)}{W_{N/p}(H/p)}\equiv1\pmod{p^{r+5s}}.
\]
\end{lemma}
\begin{proof}
Division by $p$ preserves both index classes modulo $3$, so extraction
of the $p$-divisible indices in \eqref{V-eq:d-lagrange-exact} leaves
$W_{N/p}(H/p)$. Pair the remaining numerator index $E$ with the
denominator index $K=N-E$. This bijects the unit indices, including
when $H=N$, since both omitted nodes $H,N$ are divisible by $p$.
The paired quotient has defect
\[
 \frac{1+H^5/E^5}{1-H^5/K^5}-1
 =\frac{H^5(E^5+K^5)}{E^5(K^5-H^5)}.
\]
Its denominator is a unit, and $E^5+(N-E)^5$ is divisible by $N$.
This gives valuation at least $r+5s$ term by term.
\end{proof}

\subsection{The Dougall model at arbitrary nodes}
Let $\zeta$ be a primitive fifth root of unity in an unramified
extension of $\mathbb Q_p$, and put $\sigma=\zeta+\zeta^2+\zeta^3$.
For every positive $H\equiv1\pmod3$, with $m=(H-1)/3$, Dougall's
terminating identity gives
\begin{equation}\label{V-eq:d-arbitrary-dougall}
 T(H)=H\frac{(1/3)_m
 \prod_{1\le i<j\le3}(2/3+H(\zeta^i+\zeta^j)/3)_m}
 {\prod_{i=1}^3(1+H\zeta^i/3)_m
 (1/3+H\sigma/3)_m}.
\end{equation}
This is Dougall's classical terminating ${}_7F_6$ evaluation, applied
to an arbitrary admissible integer before any prime-power
specialization; it is the transformation used for D in Swisher,
\emph{On the supercongruence conjectures of van Hamme},
Research in the Mathematical Sciences \textbf{2} (2015), article~18,
\url{https://doi.org/10.1186/s40687-015-0037-6}.
For $p\mid H$, write $p=3q+1$ and
$m=pm'+q$, $m'=(H/p-1)/3$. In each Pochhammer symbol
$(\alpha+H\beta)_m$, $\alpha\in\{1/3,2/3,1\}$, the divisible factors
have first indices $q,2q,p-1$, respectively, and exactly $m'$ terms.
After division by $p$, they are precisely
$(\alpha+(H/p)\beta)_{m'}$. The exact unit-factor identity
\[
\prod_{\substack{0\le j<m\\\alpha+j\text{ a unit}}}(\alpha+j)
=(-1)^m\frac{\Gamma_p(\alpha+m)}{\Gamma_p(\alpha)}
\]
and gamma reflection then give
\begin{equation}\label{V-eq:d-arbitrary-model-step}
 T(H)=C_pT(H/p)\mathcal U_p(H),\qquad
 C_p=-p\Gamma_p(1/3)^9,
\end{equation}
where
\[
 \mathcal U_p(X)=
 \prod_{a=0}^4\Gamma_p(X\zeta^a/3)
 \prod_{0\le a<b\le4}
 \frac{\Gamma_p(1/3-X(\zeta^a+\zeta^b)/3)}{\Gamma_p(1/3)}.
\]
For clarity, the reflection signs and orbit symmetrization are as
follows. Since $p=3q+1$ is odd, $q$ is even and the least positive
residue of $1/3$ is $2q+1$. Reflection at $1/3$ has sign $-1$;
the six reciprocal gamma factors centered at $2/3$ give net sign
$+1$. The three factors centered at $1$ give net sign $-1$ by
$\Gamma_p(1+x)=-\Gamma_p(x)$ near zero. The relation
$1+\sigma=-\zeta^4$ and $\Gamma_p(x)\Gamma_p(-x)=1$ turn the
five zero-centered shifts into the complete fifth-root orbit.
The other ten shifts are exactly the negative pair sums
$-(\zeta^a+\zeta^b)/3$. This proves the displayed product and its
constant normalization $\mathcal U_p(0)=1$.
The local analytic extension of $\Gamma_p$, including the step and
reflection formulas at these algebraic arguments, is recorded in
Long--Ramakrishna, \emph{Some supercongruences occurring in truncated
hypergeometric series}, Adv. Math. \textbf{290} (2016), 773--808,
Proposition~13 and equations~(2.2)--(2.3),
\url{https://arxiv.org/abs/1403.5232}.

For $s=v_p(H)\ge1$, this is a unit and
\begin{equation}\label{V-eq:d-arbitrary-model-precision}
 \mathcal U_p(H)\equiv1\pmod{p^{5s}}.
\end{equation}
For completeness, multiplication of the argument by $\zeta$
permutes the displayed gamma factors, so the power series of
$\mathcal U_p(X)-1$ has only positive degrees divisible by $5$.
Its degree-five coefficient is integral, because $5<p$.
Long--Ramakrishna's Proposition~13 gives a lower bound
$-j(1/p+1/(p-1))$ for the valuation of the coefficient of degree $j$
in each normalized gamma series, and hence in their finite product.
For $j\ge10$, $p\ge7$, and $s\ge1$, this gives
$js-j(1/p+1/(p-1))\ge5s$. This proves
\eqref{V-eq:d-arbitrary-model-precision}, including $p=7$.

\section{Every positive fifth moment of the D interpolating measure}\label{loc-V-5}
\label{V-sec:d-positive-moments}
Keep the exact weights $W_N(H)$ and terminating model $T(H)$ from
\eqref{V-eq:d-lagrange-exact}. No prime hypothesis is needed in the
following identity. Let $N=3n+1$, $n\ge0$, and $b_n=(1/3)_n/n!$.
\begin{theorem}\label{V-thm:d-positive-moments}
There exist polynomials
\[
 P_j(N)\in\mathbb Z[1/2][N(N+2)],\qquad
 P_0=1,\quad \deg P_j=4j,\quad P_j(1)=1,
\]
such that, for every $j\ge1$,
\[
 \boxed{\sum_{\substack{1\le H\le N\\H\equiv1\ (3)}}
 W_N(H)T(H)H^{5j}=(Nb_n)^6P_{j-1}(N).}
\]
In particular, if $p\equiv1\pmod3$ and $N=p^r$, every positive fifth
moment belongs to $p^{6r}\mathbb Z_p$, including when
$p\equiv1\pmod5$.
\end{theorem}
\begin{proof}
Put
\[
 \mathcal D_N(y)=\prod_{j=1}^n(1+y/(3j)^5),\quad
 \mathcal P_N(y)=\prod_{h=0}^n(1-y/(3h+1)^5),\quad
 F_N(y)=\frac{\mathcal D_N(y)S_N(y^{1/5})}{\mathcal P_N(y)}.
\]
The notation $S_N(y^{1/5})$ means the rational function in $y$;
no root is chosen. Partial fractions, or the interpolation identity,
give
\[
 F_N(y)=\sum_{\substack{1\le H\le N\\H\equiv1\ (3)}}
 \frac{W_N(H)T(H)}{1-y/H^5}.
\]
For $n\ge1$, separating the last summand of $S_N$ gives the exact
rational-function recurrence
\begin{equation}\label{V-eq:d-moment-rational-recurrence}
 (1-y/N^5)F_N(y)=(1+y/(N-1)^5)F_{N-3}(y)+(2N-1)b_n^6.
\end{equation}
Write the Laurent expansion at infinity as
$F_N(y)=-\sum_{j\ge1}M_j(N)y^{-j}$, where
$M_j(N)=\sum_HW_N(H)T(H)H^{5j}$.
The constant term of \eqref{V-eq:d-moment-rational-recurrence} gives
\[
 M_1(N)=N^5\left((2N-1)b_n^6-
                  \frac{M_1(N-3)}{(N-1)^5}\right).
\]
Since $b_n/b_{n-1}=(N-3)/(N-1)$ and $M_1(1)=1$, induction proves
$M_1(N)=(Nb_n)^6$.
For $j\ge1$, comparison of the coefficient of $y^{-j}$ gives
\[
 M_{j+1}(N)=N^5\bigl(M_j(N)-M_j(N-3)\bigr)
                -\frac{N^5}{(N-1)^5}M_{j+1}(N-3).
\]
After substituting the claimed normalization, its polynomial form is
\begin{equation}\label{V-eq:d-moment-polynomial-recurrence}
 NP_j(N)+(N-1)P_j(N-3)
 =N^6P_{j-1}(N)-(N-1)^6P_{j-1}(N-3).
\end{equation}
We now prove that this determines polynomials over $\mathbb Z[1/2]$,
without dividing by any odd integer.

Let $R=\mathbb Z[1/2]$, $X=N(N+2)$, $Y=N(N-1)$, and define
$\mathcal L(P)=NP(N)+(N-1)P(N-3)$.
The invariant polynomials for the involution $N\mapsto-N-2$ are
exactly $R[X]$. The anti-invariant polynomials for $N\mapsto1-N$
are exactly $(2N-1)R[Y]$. Both statements follow by shifting the
variable to the center of the corresponding reflection, which is
allowed because $2$ is invertible in $R$.
For $P\in R[X]$, substitution shows that $\mathcal L(P)$ is
anti-invariant. Moreover
\[
 \frac{\mathcal L(X^k)}{2N-1}
\]
is a polynomial in $Y$ of degree $k$ with leading coefficient $1$:
the numerator has degree $2k+1$ and leading coefficient $2$.
Thus $\mathcal L:R[X]\to(2N-1)R[Y]$ is an isomorphism of filtered
$R$-modules, with a triangular matrix whose diagonal consists of
ones.
If $P_{j-1}\in R[X]$, the right side of
\eqref{V-eq:d-moment-polynomial-recurrence} is anti-invariant under
$N\mapsto1-N$. Therefore it lies in $(2N-1)R[Y]$ and has a unique
preimage $P_j\in R[X]$. Its degree is $\deg P_{j-1}+4$; the leading
coefficient of its right side before applying $\mathcal L^{-1}$ is
$(6+3\deg P_{j-1})$ times the nonzero leading coefficient of
$P_{j-1}$, so this degree does not drop. Evaluating at $N=1$ gives
$P_j(1)=P_{j-1}(1)=1$.
The moment recurrence with this initial value now proves the exact
identity by induction on $n$ and $j$.
Finally, $b_{(p^r-1)/3}$ is a $p$-adic unit for $p\equiv1\pmod3$: at
each power $p^a\le p^r$, the numerator factors $3j+1$,
$0\le j<(p^r-1)/3$, and the denominator factors $3j$,
$1\le j\le(p^r-1)/3$, have the same number of multiples of $p^a$.
Since $P_j$ has no odd prime in its coefficients' denominators,
the asserted valuation follows.
\end{proof}
\section{The complete ordinary D.3 recursion}\label{loc-V-6}
\label{V-sec:d-ordinary-complete}
\begin{theorem}\label{V-thm:d-ordinary-complete}
For every prime $p\equiv1\pmod3$ and every $r\ge1$,
\[
 \boxed{D(p^r)\equiv-p\Gamma_p(1/3)^9D(p^{r-1})\pmod{p^{6r}}.}
\]
Moreover $v_p(D(p^r))=r$ for all $r\ge0$.
\end{theorem}
\begin{proof}
Write $N=p^r$, $M=p^{r-1}$ and $C_p=-p\Gamma_p(1/3)^9$.
We use the exact interpolation identity \eqref{V-eq:d-lagrange-exact},
the arbitrary-node model identity \eqref{V-eq:d-arbitrary-model-step},
and the positive moments of Theorem~\ref{V-thm:d-positive-moments}.

First remove all unit nodes at once. Put $j_a=(p-1)p^a$.
As $a\to\infty$, $H^{5j_a}\to1$ if $p\nmid H$, and
$H^{5j_a}\to0$ if $p\mid H$. Since the interpolating sum has finitely
many rational terms, its limit is
\[
 \sum_{\substack{1\le H\le N\\H\equiv1\ (3),\ p\nmid H}}
 W_N(H)T(H)
 =\lim_{a\to\infty}\sum_H W_N(H)T(H)H^{5j_a}
 \in p^{6r}\mathbb Z_p.
\]
The final inclusion follows because every positive moment belongs to
the displayed closed ideal. It remains to transport the divisible
nodes, which we write $H=pJ$.

Define the unit-index rational function
\[
 B_N(X)=
 \frac{\prod_{\substack{1\le E<N\\E\equiv0\ (3),\ p\nmid E}}(1+X/E^5)}
 {\prod_{\substack{1\le K<N\\K\equiv1\ (3),\ p\nmid K}}(1-X/K^5)}.
\]
Direct extraction of $p$-divisible factors in the exact weights gives
\[
 W_N(pJ)=W_M(J)B_N(p^5J^5).
\]
Pair $E$ with $K=N-E$. Each paired factor is
\[
 \frac{1+X/E^5}{1-X/K^5}
 =1+X\frac{E^5+K^5}{E^5K^5}\frac1{1-X/K^5}.
\]
Since $E^5+K^5$ is divisible by $N$, this proves the coefficientwise
statement
\begin{equation}\label{V-eq:d-weight-series}
 B_N(X)\in1+p^rX\mathbb Z_p[[X]].
\end{equation}
In particular $B_N(p^5J^5)$ is a convergent power series in $J^5$,
whose nonconstant coefficients have valuation at least $r+5$ and
tend to infinity with their index.

The exact gamma factor $\mathcal U_p$ in
\eqref{V-eq:d-arbitrary-model-step} has a series
\[
 \mathcal U_p(X)=1+\sum_{j\ge1}u_jX^{5j}.
\]
For $p\ge7$, its coefficient $u_1$ is integral, and the general
coefficient bound is
$v_p(u_j)\ge-5j(1/p+1/(p-1))$.
Thus
\[
 \mathcal U_p(pJ)=1+\sum_{j\ge1}p^{5j}u_jJ^{5j},
\]
where all nonconstant coefficients belong to $p^5\mathcal O$ and
tend to zero $p$-adically. At $j=1$ use the integrality of $u_1$;
for $j\ge2$ use
$5j(1-1/p-1/(p-1))\ge10(1-13/42)>5$.
Here $\mathcal O$ is the ring of integers of the unramified extension
containing the fifth roots of unity. Combining this fact with
\eqref{V-eq:d-weight-series} gives a convergent series
\begin{equation}\label{V-eq:d-transport-coefficients}
 B_N(p^5J^5)\mathcal U_p(pJ)-1
 =\sum_{j\ge1}c_jJ^{5j},\qquad
 c_j\in p^5\mathcal O,\quad c_j\longrightarrow0.
\end{equation}
This convergence is uniform for $J\in\mathbb Z_p$.

The difference of the divisible-node contributions from the two
levels is therefore exactly
\begin{align*}
 &\sum_{\substack{1\le H\le N\\H\equiv1\ (3),\ p\mid H}}
 W_N(H)T(H)-C_pD(M)\\
 &\qquad=C_p\sum_{\substack{1\le J\le M\\J\equiv1\ (3)}}
 W_M(J)T(J)\sum_{j\ge1}c_jJ^{5j}\\
 &\qquad=C_p\sum_{j\ge1}c_j
 \sum_{\substack{1\le J\le M\\J\equiv1\ (3)}}W_M(J)T(J)J^{5j}.
\end{align*}
Interchanging the finite sum with the convergent series is valid
even if individual interpolation weights have negative valuations.
Every inner positive moment belongs to $p^{6(r-1)}\mathbb Z_p$.
Since $v_p(C_p)=1$ and $c_j\in p^5\mathcal O$, the last expression
belongs to $p^{1+5+6(r-1)}\mathcal O=p^{6r}\mathcal O$.
It is rational over $\mathbb Q_p$, so it belongs to
$p^{6r}\mathbb Z_p$. Together with the unit-node limit this proves
the asserted congruence. Starting with $D(1)=1$, its valuation
statement follows inductively because $6r>r$.
\end{proof}

\section{The entire ordinary branch of G.3 by positive cubic moments}\label{loc-V-7}
\label{V-sec:G6ordinary}
For $N=4n+1$ put
\[
 b_k=\frac{(1/4)_k}{k!},\qquad
 G(N)=\sum_{k=0}^{n}(8k+1)b_k^4.
\]
We prove, for every prime $p\equiv1\pmod4$ and every $r\ge1$,
\begin{equation}\label{V-G6:main}
 \boxed{\quad G(p^r)\equiv p\lambda_pG(p^{r-1})\pmod{p^{4r}},
 \qquad\lambda_p=-(-1)^{(p-1)/4}
 \Gamma_p(1/2)\Gamma_p(1/4)^2.\quad}
\end{equation}
This includes primes $p\equiv1\pmod3$; no estimate of individual
interpolation weights is required.

\subsection{Exact interpolation and its terminating model}
Define the rational function in $y=z^3$
\[
 S_N(z)=\sum_{k=0}^{n}(8k+1)b_k^4
 \prod_{j=0}^{k-1}\frac{1-z^3/(4j+1)^3}{1+z^3/(4j+4)^3}.
\]
Multiplication by $\mathcal D_N(y)=\prod_{j=1}^{n}(1+y/(4j)^3)$
gives a polynomial of degree at most $n$. At the distinct nodes
$y=H^3$, $H=1,5,\ldots,N$, the sum terminates at $(H-1)/4$.
Writing $T(H)=S_H(H)$, interpolation at zero gives
\begin{equation}\label{V-G6:interpolation}
 G(N)=\sum_{H=1,5,\ldots,N}W_N(H)T(H),\qquad
 W_N(H)=
 \frac{\prod_{\substack{1\le E<N\\E\equiv0\ (4)}}(1+H^3/E^3)}
 {\prod_{\substack{1\le K\le N\\K\equiv1\ (4),\ K\ne H}}
        (1-H^3/K^3)}.
\end{equation}
All these are exact finite identities. Their rational denominators
are nonzero, regardless of their $p$-adic valuations.

Let $\zeta$ be a primitive cube root of unity. For $H=4m+1$,
Dougall's terminating very-well-poised ${}_5F_4$ identity gives
\begin{equation}\label{V-G6:node}
 T(H)=(-1)^mH\frac{(1/4)_m(1/2)_m}
 {(1+\zeta H/4)_m(1+\zeta^2H/4)_m}.
\end{equation}
Indeed its parameters are $a=1/4$, $b=(1-\zeta H)/4$,
$c=(1-\zeta^2H)/4$, and $-m=(1-H)/4$. The terminating evaluation is
\[
 \frac{(1+a)_m(1+a-b-c)_m}
 {(1+a-b)_m(1+a-c)_m}
 =\frac{(5/4)_m(1/2-m)_m}
 {(1+\zeta H/4)_m(1+\zeta^2H/4)_m},
\]
which proves \eqref{V-G6:node}. The denominator product is rational
and nonzero.

\subsection{Every positive cubic moment}
There are polynomials $P_j(X)\in\mathbb Z[1/2][X(X+3)]$,
$P_0=1$, $P_j(1)=1$, with degree at most $2j$, such that
\begin{equation}\label{V-G6:moment}
 M_j(N):=\sum_{H=1,5,\ldots,N}W_N(H)T(H)H^{3j}
 =(Nb_n)^4P_{j-1}(N)\qquad(j\ge1).
\end{equation}
To prove this, put
\[
 \mathcal P_N(y)=\prod_{h=0}^{n}(1-y/(4h+1)^3),\qquad
 F_N(y)=\frac{\mathcal D_N(y)S_N(y^{1/3})}{\mathcal P_N(y)}.
\]
Here $S_N(y^{1/3})$ denotes the rational function in $y$.
The partial-fraction decomposition of the interpolation formula is
\[
 F_N(y)=\sum_{H=1,5,\ldots,N}\frac{W_N(H)T(H)}{1-y/H^3}.
\]
Separating the last summand of $S_N$ yields, for $n\ge1$,
\begin{equation}\label{V-G6:rational-rec}
 (1-y/N^3)F_N(y)
 =(1+y/(N-1)^3)F_{N-4}(y)+(2N-1)b_n^4.
\end{equation}
Expand at infinity: $F_N(y)=-\sum_{j\ge1}M_j(N)y^{-j}$.
The constant term of \eqref{V-G6:rational-rec} gives
\[
 M_1(N)=N^3\left((2N-1)b_n^4-
                    \frac{M_1(N-4)}{(N-1)^3}\right).
\]
Since $b_n/b_{n-1}=(N-4)/(N-1)$ and $M_1(1)=1$, this proves
$M_1(N)=(Nb_n)^4$. The remaining coefficients give
\[
 M_{j+1}(N)=N^3\bigl(M_j(N)-M_j(N-4)\bigr)
             -\frac{N^3}{(N-1)^3}M_{j+1}(N-4).
\]
The claimed normalization therefore reduces to
\begin{equation}\label{V-G6:P-rec}
 XP_j(X)+(X-1)P_j(X-4)
 =X^4P_{j-1}(X)-(X-1)^4P_{j-1}(X-4).
\end{equation}

Here polynomial existence requires no division by an odd prime.
Set $R=\mathbb Z[1/2]$ and
$\mathcal LP=XP(X)+(X-1)P(X-4)$.
The invariant polynomials under $X\mapsto-X-3$ are exactly
$R[X(X+3)]$. The anti-invariant polynomials under $X\mapsto1-X$
are exactly $(2X-1)R[X(X-1)]$.
Substitution shows that $\mathcal L$ maps the former module into
the latter. In the corresponding bases
$[X(X+3)]^d$ and $(2X-1)[X(X-1)]^d$ its matrix is triangular with
diagonal entries one, since both leading terms are $2X^{2d+1}$.
It is consequently an isomorphism on every finite degree filtration.
If $P_{j-1}$ is invariant, the right side of \eqref{V-G6:P-rec} is
anti-invariant and has degree at most $\deg P_{j-1}+3$.
Thus the equation uniquely defines an invariant $P_j\in R[X]$
of degree at most $2j$. At $X=1$ it gives $P_j(1)=P_{j-1}(1)=1$.
The moment recurrence now proves \eqref{V-G6:moment} by induction.

For $N=p^r$ with $p\equiv1\pmod4$, $b_n$ is a $p$-adic unit:
at each power $p^a\le N$, its numerator and denominator have the
same number $(p^{r-a}-1)/4$ of divisible factors. Hence
\begin{equation}\label{V-G6:moment-bound}
 M_j(p^r)\in p^{4r}\mathbb Z_p\qquad(j\ge1).
\end{equation}
This also holds at $r=0$, where every moment equals one.

\subsection{Discarding the unit nodes by their moments}
Write $\mu_N(H)=W_N(H)T(H)$. For $e_v=(p-1)p^v$,
\[
 H^{3e_v}\longrightarrow
 \begin{cases}1,&p\nmid H,\\0,&p\mid H\end{cases}
 \quad\text{in }\mathbb Q_p.
\]
There are finitely many nodes, so \eqref{V-G6:moment-bound} implies
\begin{equation}\label{V-G6:unit-nodes}
 \sum_{\substack{H=1,5,\ldots,N\\p\nmid H}}\mu_N(H)
 \in p^{4r}\mathbb Z_p\qquad(N=p^r).
\end{equation}
This step uses the sum of unit nodes; their individual valuations
need not satisfy the displayed bound.

\subsection{Exact scaling of the model}
Now let $p=4q+1$, $p\mid H$, $J=H/p$, and
$m=(H-1)/4=pm'+q$, $m'=(J-1)/4$.
In $(1/4)_m$, $(1/2)_m$, and $(1+\zeta^iH/4)_m$, the indices of
the first $p$-divisible factors are respectively $q$, $2q$, and
$p-1$. Each list has exactly $m'$ divisible factors; after division
by $p$ they give the identical factors for the model at $J$.
Using the functional equation for the remaining unit factors gives
\[
 \frac{T(H)}{pT(J)}=(-1)^q
 \frac{\Gamma_p(H/4)\Gamma_p(1/4+H/4)
       \Gamma_p(1+\zeta H/4)\Gamma_p(1+\zeta^2H/4)}
 {\Gamma_p(1/4)\Gamma_p(1/2)
       \Gamma_p(3/4-\zeta H/4)\Gamma_p(3/4-\zeta^2H/4)}.
\]
The two functional equations $\Gamma_p(1+u)=-\Gamma_p(u)$
apply because $u\in p\mathcal O_K$.
Reflection at $1/4+\zeta^iH/4$ cancels the two equal signs.
Finally $\Gamma_p(1/2)^2=-1$ for $p\equiv1\pmod4$.
Therefore the exact model scaling is
\begin{equation}\label{V-G6:model-scale}
 T(H)=p\lambda_pT(J)U(H),\qquad
 U(X)=\prod_{i=0}^{2}\Gamma_p(\zeta^iX/4)
       \prod_{i=0}^{2}
       \frac{\Gamma_p(1/4+\zeta^iX/4)}{\Gamma_p(1/4)}.
\end{equation}
These gamma values lie in the unramified extension containing
$\zeta$, and the identities are valid on the disks $p\mathcal O_K$.
No integrality assertion for an individual $T(J)$ is needed.

The scaled correction has a restricted series
\begin{equation}\label{V-G6:restricted}
 U(pX)=1+\sum_{j\ge1}u_jX^{3j},\qquad
 u_j\in p^3\mathbb Z_p,\quad v_p(u_j)\longrightarrow+\infty.
\end{equation}
Indeed multiplication of $X$ by $\zeta$ permutes each triple, so
only degrees divisible by three occur. Galois invariance puts the
coefficients in $\mathbb Q_p$.
For the normalized Taylor coefficients of each gamma factor,
Long--Ramakrishna's bounds give integrality in degrees $k<p$ and
valuation at least $-k(1/p+1/(p-1))$ in all degrees.
The degree-three coefficient of the scaled product is thus in
$p^3\mathbb Z_p$, as $p\ge5$. For total degree $k\ge6$ its valuation
is at least
\[
 k\left(1-\frac1p-\frac1{p-1}\right)
 \ge\frac{11k}{20}>3.
\]
These bounds also prove convergence on $X\in\mathbb Z_p$ and the
coefficient limit in \eqref{V-G6:restricted}.

\subsection{Transport using all positive moments}
Put $N=p^r$, $M=N/p$. Stripping the $p$-divisible indices from
\eqref{V-G6:interpolation}, using $p\equiv1\pmod4$, gives exactly
\[
 R_N(J):=\frac{W_N(pJ)}{W_M(J)}
 =\frac{\prod_{\substack{1\le E<N\\E\equiv0\ (4),\ p\nmid E}}
              (1+p^3J^3/E^3)}
 {\prod_{\substack{1\le K<N\\K\equiv1\ (4),\ p\nmid K}}
              (1-p^3J^3/K^3)}.
\]
Its constant coefficient is one and its coefficient of $J^{3j}$
belongs to $p^{3j}\mathbb Z_p$. Thus
\[
 R_N(J)U(pJ)=1+\sum_{j\ge1}e_{N,j}J^{3j},\qquad
 e_{N,j}\in p^3\mathbb Z_p,\quad v_p(e_{N,j})\to+\infty.
\]
Equations \eqref{V-G6:model-scale} and the weight quotient show
\[
 \sum_{p\mid H}\mu_N(H)
 =p\lambda_p\left(G(M)+\sum_{j\ge1}e_{N,j}M_j(M)\right).
\]
The interchange is legitimate: there are finitely many coarse
nodes, and the restricted series converges at each of them.
Every correction term has valuation at least
$1+3+4(r-1)=4r$ by \eqref{V-G6:moment-bound}.
The ideal $p^{4r}\mathbb Z_p$ is closed, so the entire correction
belongs to it. Adding \eqref{V-G6:unit-nodes} proves
\eqref{V-G6:main}, including $r=1$.

\section{The even supersingular branch of G.3 through an alternate family}\label{loc-V-8}
\label{V-sec:g-even-alternate}
Let $p>3$ be prime with $p\equiv3\pmod4$. For $a\in\{1,3\}$ and
$N\equiv a\pmod4$, put $n=(N-a)/4$ and define
\[
 b_{a,k}=\frac{(a/4)_k}{k!},\qquad
 G_a(N)=\sum_{k=0}^{n}(8k+a)b_{a,k}^4.
\]
In particular $G_1(p^{2s})$ is the original even-level G.3 sum.
We prove
\begin{equation}\label{V-Ge:main}
 \boxed{G_1(p^{2s})\equiv-p^3G_1(p^{2s-2})
 \pmod{p^{8s-2}}\qquad(s\ge1).}
\end{equation}
The alternate family uses the weight $8k+3$, rather than
$(8k+3)/3$; thus $G_3(3)=3$.

\subsection{Interpolation and moments for both node classes}
Define
\[
 S_{a,N}(z)=\sum_{k=0}^{n}(8k+a)b_{a,k}^4
 \prod_{j=0}^{k-1}
 \frac{1-z^3/(4j+a)^3}{1+z^3/(4j+4)^3},\qquad
 T_a(H)=S_{a,H}(H).
\]
Multiplication by $\prod_{j=1}^{n}(1+y/(4j)^3)$ makes
$S_{a,N}(y^{1/3})$ a polynomial of degree at most $n$.
At the nodes $H=a,a+4,\ldots,N$, the undeclared root is unnecessary:
the rational function in $y$ takes the value $T_a(H)$ at $y=H^3$.
Lagrange interpolation therefore gives the exact identity
\begin{equation}\label{V-Ge:interpolation}
G_a(N)=\sum_{\substack{1\le H\le N\\H\equiv a\ (4)}}
W_{a,N}(H)T_a(H),
\end{equation}
where
\[
W_{a,N}(H)=
\frac{\prod_{j=1}^{n}(1+H^3/(4j)^3)}
{\prod_{\substack{0\le j\le n\\4j+a\ne H}}
 (1-H^3/(4j+a)^3)}.
\]
All denominators are nonzero rational numbers.

For every $j\ge1$, the positive moments have the exact form
\begin{equation}\label{V-Ge:moment}
\sum_H W_{a,N}(H)T_a(H)H^{3j}
=(Nb_{a,n})^4P^{(a)}_{j-1}(N),
\end{equation}
where
\[
 P^{(a)}_j\in\mathbb Z[1/2][N(N+4-a)],\quad
 P^{(a)}_0=1,\quad \deg P^{(a)}_j=2j,
 \quad P^{(a)}_j(a)=a^{3j}.
\]
Here is a proof including the alternate normalization.
Let $F_{a,N}(y)$ be the proper rational function obtained by
multiplying $S_{a,N}(y^{1/3})$ by
\[
\frac{\prod_{j=1}^{n}(1+y/(4j)^3)}
 {\prod_{j=0}^{n}(1-y/(4j+a)^3)}.
\]
Its partial fractions are
$F_{a,N}(y)=\sum_H W_{a,N}(H)T_a(H)/(1-y/H^3)$.
For $n\ge1$, separating the last summand gives
\[
(1-y/N^3)F_{a,N}(y)
=(1+y/(N-a)^3)F_{a,N-4}(y)+(2N-a)b_{a,n}^4.
\]
Write $F_{a,N}(y)=-\sum_{j\ge1}M_j(N)y^{-j}$ at infinity.
Since $b_{a,n}/b_{a,n-1}=(N-4)/(N-a)$, comparison of constant
terms proves $M_1(N)=(Nb_{a,n})^4$, starting from
$F_{a,a}(y)=a/(1-y/a^3)$.
The higher coefficients give the recurrence
\begin{equation}\label{V-Ge:moment-polynomial}
 NP^{(a)}_j(N)+(N-a)P^{(a)}_j(N-4)
 =N^4P^{(a)}_{j-1}(N)-(N-a)^4P^{(a)}_{j-1}(N-4).
\end{equation}
Over $R=\mathbb Z[1/2]$, the operator
$\mathcal L_aP=NP(N)+(N-a)P(N-4)$ maps
\[
R[N(N+4-a)]\longrightarrow(2N-a)R[N(N-a)].
\]
The domain is invariant under $N\mapsto a-4-N$, and the codomain
is anti-invariant under $N\mapsto a-N$. In the displayed power bases
this map is triangular with diagonal one, because both leading
terms are $2N^{2d+1}$. It is therefore an isomorphism over $R$ on
every finite degree filtration. The right side of
\eqref{V-Ge:moment-polynomial} is anti-invariant, so it defines a
unique $P^{(a)}_j$ over $R$. If the preceding degree is $d$ with
leading coefficient $c$, the right side has leading term
$4(a+d)cN^{d+3}$, so its degree increases by two after applying
$\mathcal L_a^{-1}$. Evaluating at $N=a$ gives
$P^{(a)}_j(a)=a^3P^{(a)}_{j-1}(a)$.
Thus the initial values and the rational-function recurrence prove
\eqref{V-Ge:moment}.

If $N=p^r\equiv a\pmod4$, the factor $b_{a,n}$ is a $p$-adic unit.
Indeed, for every $p^\ell\mid N$, the first numerator index
$k_\ell$ satisfying $p^\ell\mid4k_\ell+a$ is precisely the
residue of $n$ modulo $p^\ell$. Since the numerator product stops
at $n-1$, its number of divisible factors equals
$\lfloor n/p^\ell\rfloor$, the same as the denominator's count.
Consequently every positive moment in \eqref{V-Ge:moment} belongs to
$p^{4r}\mathbb Z_p$.
For $e_v=(p-1)p^v$, the limit of $H^{3e_v}$ is one at unit nodes
and zero at nodes divisible by $p$. Applying this to the finite
moment sum shows
\begin{equation}\label{V-Ge:unitnodes}
\sum_{p\nmid H}W_{a,p^r}(H)T_a(H)\in p^{4r}\mathbb Z_p.
\end{equation}
This argument does not require individual nodes to be integral.

\subsection{The two terminating models and their exact scaling}
Let $\zeta$ be a primitive cube root in an unramified extension of
$\mathbb Q_p$. Dougall's terminating ${}_5F_4$ identity, with central
parameter $a/4$ and other parameters
$(a-H)/4,(a-\zeta H)/4,(a-\zeta^2H)/4$, gives
\begin{equation}\label{V-Ge:models}
 T_a(H)=(-1)^mH\frac{(a/4)_m(a/2)_m}
 {(1+\zeta H/4)_m(1+\zeta^2H/4)_m},\qquad m=(H-a)/4.
\end{equation}
The factor $a$ converting the very-well-poised weight to $8k+a$
is included here. Indeed,
$a(1+a/4)_m=H(a/4)_m$, and the other numerator is
$(1-a/2-m)_m=(-1)^m(a/2)_m$.

Put $q=(p-3)/4$, and write $\Gamma$ for Morita's $\Gamma_p$.
Set
\[
 C_1=\frac{(-1)^q p^2\Gamma(1/4)^2}{2\Gamma(1/2)},\qquad
 C_3=-\frac{2(-1)^q p\Gamma(3/4)^2}{\Gamma(1/2)},
\]
\[
 U_a(X)=\prod_{j=0}^{2}\Gamma(X\zeta^j/4)
 \prod_{j=0}^{2}\frac{\Gamma(a/4+X\zeta^j/4)}{\Gamma(a/4)}.
\]
For every admissible $H$ divisible by $p$, exact factor extraction gives
\begin{equation}\label{V-Ge:modelsteps}
 T_1(H)=C_1T_3(H/p)U_1(H),\qquad
 T_3(H)=C_3T_1(H/p)U_3(H).
\end{equation}
To check counts and constants, let $H/p=4v+3$ in the first case.
Then $m=pv+3q+2$. The divisible counts in
$(1/4)_m,(1/2)_m$ are $v,v+1$, while each denominator count is $v$.
The extracted numerator pair is
$(3/4)_v(1/2)_{v+1}=(1/2)(3/4)_v(3/2)_v$.
The external $H$ therefore contributes $p^2/2$ relative to $T_3$.
In the second case $H/p=4v+1$ and $m=pv+q$; every count is $v$,
and the numerator parameters become $1/4,1/2$, giving the external
factor $p$. In both cases $(-1)^{m-v}=(-1)^q$.

The remaining gamma quotient in case $a$ is
\[
\frac{\Gamma(H/4)\Gamma(a/4+H/4)
 \Gamma(1+\zeta H/4)\Gamma(1+\zeta^2H/4)}
 {\Gamma(a/4)\Gamma(a/2)
  \Gamma(1-a/4-\zeta H/4)\Gamma(1-a/4-\zeta^2H/4)}.
\]
The step formula at the two arguments centered at one and reflection
at the two arguments centered at $1-a/4$ convert this to
$\Gamma(a/4)^2U_a(H)/\Gamma(a/2)$; the paired signs cancel.
Finally $\Gamma(3/2)=-\Gamma(1/2)/2$, proving
\eqref{V-Ge:modelsteps}. The same identities show
\begin{equation}\label{V-Ge:constantproduct}
 v_p(C_1)=2,\qquad v_p(C_3)=1,\qquad C_1C_3=-p^3.
\end{equation}
For the last equality use
$(\Gamma(1/4)\Gamma(3/4))^2=1$ and
$\Gamma(1/2)^2=1$ for $p\equiv3\pmod4$.

We need the whole correction series, not merely its value at a
single node. The normalized analytic gamma series and their
functional identities are valid on the displayed residue disks.
Long--Ramakrishna's Proposition 13 gives integral coefficients in
degrees less than $p$, and the general coefficient bound
$-k(1/p+1/(p-1))$. Multiplication of $X$ by $\zeta$ permutes all
factors in $U_a$, so its positive degrees are multiples of three.
Since $p\ge7$, the degree-three coefficient is integral. Every
scaled coefficient in $U_a(pX)$ of degree $d\ge6$ has valuation at
least
\[
 d\left(1-\frac1p-\frac1{p-1}\right)
 \ge\frac{29d}{42}>3.
\]
These bounds tend to infinity. Hence
\begin{equation}\label{V-Ge:restricted}
 U_a(pX)=1+\sum_{j\ge1}u_{a,j}X^{3j},\qquad
 u_{a,j}\in p^3\mathbb Z_p,\qquad v_p(u_{a,j})\longrightarrow\infty.
\end{equation}
The coefficients lie in $\mathbb Q_p$ by Galois symmetry of the
complete cube-root products. All estimates can equivalently be
made in the unramified extension and then restricted to $\mathbb Q_p$.

\subsection{Transport between the two families}
Let $N=p^r\equiv a\pmod4$, $r\ge1$, and put $M=N/p$ and
$b=4-a$, so $M\equiv b\pmod4$. Dividing the $p$-divisible indices
in the exact interpolation weight by $p$ exchanges the two residue
classes and gives
\[
\frac{W_{a,N}(pJ)}{W_{b,M}(J)}
=\frac{\prod_{\substack{1\le E<N\\E\equiv0\ (4),\ p\nmid E}}
 (1+p^3J^3/E^3)}
{\prod_{\substack{1\le K\le N\\K\equiv a\ (4),\ p\nmid K}}
 (1-p^3J^3/K^3)}
=:R_{a,N}(J).
\]
The numerator range is precisely $E=4,8,\ldots,N-a$; the displayed
congruence description has that same endpoint for $a\in\{1,3\}$.
This restricted series has constant term one and its coefficient
of $J^{3j}$ lies in $p^{3j}\mathbb Z_p$. By \eqref{V-Ge:restricted},
\[
 R_{a,N}(J)U_a(pJ)=1+\sum_{j\ge1}e_{a,N,j}J^{3j},\qquad
 e_{a,N,j}\in p^3\mathbb Z_p,\quad
 v_p(e_{a,N,j})\longrightarrow\infty.
\]
Writing $\mu_{a,N}(H)=W_{a,N}(H)T_a(H)$, the exact model steps imply
\[
\mu_{a,N}(pJ)=C_a\mu_{b,M}(J)
\left(1+\sum_{j\ge1}e_{a,N,j}J^{3j}\right).
\]
Each positive coarse moment belongs to $p^{4(r-1)}\mathbb Z_p$.
The aggregated correction therefore has valuation at least
$v_p(C_a)+3+4(r-1)\ge4r$. The infinite series may be summed after
the finite node sum because its coefficients tend to zero and
there are only finitely many rational node values. Adding the
unit-node contribution \eqref{V-Ge:unitnodes} proves
\begin{equation}\label{V-Ge:one-step-sums}
 G_a(p^r)\equiv C_aG_{4-a}(p^{r-1})\pmod{p^{4r}}.
\end{equation}
For the smallest odd level $r=1$, $a=3$, the coarse level is
$M=1$, $b=1$; its positive moments equal one, so the same bound
still applies.

Apply \eqref{V-Ge:one-step-sums} first at $r=2s$, $a=1$, and then at
$r=2s-1$, $a=3$. The second error has valuation at least $8s-4$;
after multiplication by $C_1$ it has valuation at least $8s-2$.
The first error belongs to $p^{8s}\mathbb Z_p$. Equation
\eqref{V-Ge:constantproduct} now gives \eqref{V-Ge:main}.
The exclusion of $p=3$ is necessary for this even-level statement;
no assertion for that prime is made here.


\begin{corollary}
For $p>3$, $p\equiv3\pmod4$ and $s\ge0$,
$v_p(G(p^{2s}))=3s$.
\end{corollary}
\begin{proof}
The base case is $G(1)=1$. If the preceding value has valuation
$3s-3$, its product with $-p^3$ has valuation $3s$.
The error in \eqref{V-Ge:main} has valuation at least $8s-2>3s$
for $s\ge1$, which completes the induction.
\end{proof}

\section{What remains and how the results were checked}\label{loc-V-9}
The new results remove all remaining A.3 levels, all ordinary D.3
levels, all ordinary G.3 levels, and every even supersingular G.3
level for $p>3$. Together with the earlier companion results,
A.3, D.3 and G.3 (with the stated prime domain) are complete within
the research series. The remaining target in the active scope is
E.3 at $p\equiv2\pmod3$, odd $r\ge3$, with modulus $p^{3r-1}$.
Explicitly, for
\[
 E_r=\sum_{k=0}^{(p^r-2)/3}(-1)^k(6k+1)
 \frac{(1/3)_k^3}{(k!)^3},
\]
the unresolved assertion is
$E_r\equiv p^2E_{r-2}\pmod{p^{3r-1}}$.
Exact parameter reflections and vanishing deformations are
recorded in the archive, but removal of those deformations at
the target precision has not been proved.
The prime $3$ requires a domain distinction for the even G formula.
Direct calculation gives
\[
 G(9)+27=\frac{29407617}{1048576}
 =\frac{3^5\cdot121019}{2^{20}},
\]
so the literal modulus-$p^6$ even-level assertion fails at
$p=3,r=2$. All even supersingular G claims in this manuscript
therefore explicitly assume $p>3$.
Swisher's Table~2 does not display this exception in its G.3 row;
our completion claim for G.3 is for the corrected prime domain $p>3$.

\paragraph{Reproducibility.}
The companion archive contains this single-file source, the PDF,
the proof fragments, exact-arithmetic programs and their recorded
outputs. The checks include rational interpolation identities,
polynomial moment recurrences and evaluations, node-model formulas,
weight valuations, and original prime-power congruences.
The earlier second-level ordinary G calculation also checked all
38 primes $p\le401$ in its residue class: 1,862 centered blocks
and 496,238 fine terms. It is retained as an independent finite-level
check, although the all-level proof no longer relies on it.
The supersingular A recurrence was checked in 26 cases, including
$p=3$ through $r=8$, $p=7$ through $r=5$, and $p=11$ through $r=4$.
There are also 136 exact A interpolation and factorization checks.
Finite calculations verify the formulas; the algebraic and
$p$-adic proofs establish the universal assertions.

\begin{thebibliography}{9}
\bibitem{V-Swisher}
H. Swisher, \emph{On the supercongruence conjectures of van Hamme},
Research in the Mathematical Sciences \textbf{2} (2015), Article 18.
\url{https://doi.org/10.1186/s40687-015-0037-6}.
\bibitem{V-DLMF}
NIST Digital Library of Mathematical Functions, Section~16.4,
\emph{Argument Unity}, especially equations 16.4.7, 16.4.9 and 16.4.10.
\url{https://dlmf.nist.gov/16.4}.
\bibitem{V-Dwork1982}
B. Dwork, \emph{A note on the $p$-adic gamma function},
Groupe d'etude d'analyse ultrametrique \textbf{9} (1981--1982),
no.~3, Expose J5, J1--J10, especially pp.~J5-01--J5-03.
\url{https://numdam.org/item/GAU_1981-1982__9_3_A6_0/}.
\bibitem{V-LongRamakrishna}
L. Long and R. Ramakrishna, \emph{Some supercongruences occurring in
truncated hypergeometric series}, Advances in Mathematics
\textbf{290} (2016), 773--808.
\url{https://arxiv.org/abs/1403.5232} (version 2).
\bibitem{V-Corpus}
Verifiable Mathematical Reasoning Corpus, entries [131] and [132],
motivating research material.
\url{https://github.com/QEEDDD/Verifiable-Mathematical-Reasoning-Corpus}.
\bibitem{V-Companion}
\emph{Supersingular Advances for Swisher's A.3, D.3 and G.3:
A second-level congruence, a complete D.3 residue branch,
and every odd level of a G.3 residue branch}.
Research manuscript, 10 November 2026, included in the companion
archive as \texttt{Swisher\_ADG\_Supersingular\_Advances.pdf} and
its single-file \TeX\ source. Not externally peer reviewed.
\end{thebibliography}

\clearpage\endgroup

%%%%%%%%%%%%%%%%%%%%%%%%%%%%%%%%%%%%%%%%%%%%%%%%%%%%%%%%%%%%%%%%%%%%%%%%%%
% VI | Swisher_E3_Odd_Branch_Exact_Reductions.tex
%%%%%%%%%%%%%%%%%%%%%%%%%%%%%%%%%%%%%%%%%%%%%%%%%%%%%%%%%%%%%%%%%%%%%%%%%%
\clearpage\begingroup
\SetPaper{VI}{Exact reductions for the odd branch of E.3}
\RestoreCounters
\PaperTitle{Exact Reductions for the Remaining Odd Branch\\of Swisher's E.3 Supercongruence}{Research checkpoint}{11 November 2026}

\begin{abstract}
We study only the unresolved branch with $p>3$, $p\equiv2\pmod3$
and odd level $r\ge3$. We prove an exact symmetric companion
transformation, an interpolation formula on three grids adapted
to the short cutoff, a rational-polynomial formula for every
positive interpolation moment, and an exact two-step residual
identity. The first two moments vanish; the third moment contains
a factor of valuation $(7r-3)/2$. We also exhibit an explicit
failure of a proposed coefficient bound and verify the cancellation
between two contributions at the first lifting level. These are
proved structural reductions and finite checks. The full E.3
congruence remains unproved, and no new infinite range of its
prime-power instances is claimed in this note.
\end{abstract}
\noindent\textbf{Status.} The exact identities were independently
checked internally. The note has not undergone external peer review.
The distinction between identities, numerical evidence, and the
remaining conjectural divisibility is explicit throughout.
\setcounter{tocdepth}{1}

\begin{LocalContents}
\LocalEntry{loc-VI-1}{1}{The single remaining target}
\LocalEntry{loc-VI-2}{2}{An exact symmetric companion for the odd E.3 cutoff}
\LocalEntry{loc-VI-3}{3}{Three exact interpolation grids at the short E cutoff}
\LocalEntry{loc-VI-4}{4}{Exact moments for the three-grid short E deformation}
\LocalEntry{loc-VI-5}{5}{Exact two-step transport and the remaining cancellation}
\LocalEntry{loc-VI-6}{6}{Verification and unresolved obligations}
\end{LocalContents}
\clearpage

\section{The single remaining target}\label{loc-VI-1}
The Pochhammer symbol is $(a)_k=\prod_{j=0}^{k-1}(a+j)$, with
$(a)_0=1$. Write $v_p(p)=1$ and $v_p(0)=+\infty$.
For every integer $L\equiv2\pmod3$, $L\ge2$, define
\[
 E(L)=\sum_{k=0}^{(L-2)/3}(-1)^k(6k+1)
       \frac{(1/3)_k^3}{(k!)^3}.
\]
The target inherited from Swisher's Table~2, label E.3~\cite{VI-Swisher}, is
\begin{equation}\label{VI-eq:target}
 E(p^r)-p^2E(p^{r-2})\in p^{3r-1}\mathbb Z_p,
 \qquad p>3,\ p\equiv2\pmod3,\ r\ge3\text{ odd}.
\end{equation}
All other labels and all previously completed branches are outside
this continuation. In particular, no proof from the earlier A, D,
or G manuscripts is used as a black-box solution of \eqref{VI-eq:target}.

\paragraph{What is proved.}
The symmetric finite expression below removes the alternating weight
by an exact transformation. The three-grid formula supplies explicitly
evaluated nodes at the actual short cutoff. Its moments satisfy
$M_1=M_2=0$ and
\[
 M_3(L)=(-1)^{m+1}L^2(L-1)^3
          \left(\frac{(1/3)_m}{m!}\right)^3,
 \qquad m=(L-2)/3.
\]
For $L=p^r$ in the target branch this gives
$v_p(M_3(L))=(7r-3)/2$. This valuation concerns a moment of the
interpolation weights, not the original E.3 recurrence difference.
The final section gives an exact equation identifying the cancellation
still needed to establish \eqref{VI-eq:target}.

\section{An exact symmetric companion for the odd E.3 cutoff}\label{loc-VI-2}
The identities in this section are exact. They do not, by themselves,
prove the two-step prime-power recurrence in E.3.

For any nonnegative integer $n$, set
\[
 S_\alpha(n)=\sum_{k=0}^n(-1)^k\frac{2k+\alpha}{\alpha}
 \frac{(\alpha)_k^3}{(k!)^3},\qquad
 W_\alpha(n)=\sum_{k=0}^n\frac{(\alpha)_k(-n-\alpha)_k}{(k!)^2}.
\]
The finite transformation is
\begin{equation}\label{VI-E6r:finite-transform}
 S_\alpha(n)=\frac{(1+\alpha)_n}{n!}W_\alpha(n).
\end{equation}
Here is a short self-contained verification. For $n\ge1$, put
\[
 q_n(k)=\frac{(-n-\alpha)_k(\alpha)_k}{(k!)^2},\qquad
 R_n(k)=-\frac{k^2}{(n+\alpha)^2}.
\]
Cancellation of consecutive factors verifies
\[
 q_n(k+1)R_n(k+1)-q_n(k)R_n(k)
 =q_n(k)-\frac{n}{n+\alpha}q_{n-1}(k).
\]
Sum over $0\le k\le n$. The upper boundary and the extra term
$q_{n-1}(n)$ both equal $(-1)^n((\alpha)_n/n!)^2$.
Consequently
\[
 W_\alpha(n)-\frac{n}{n+\alpha}W_\alpha(n-1)
 =(-1)^n\frac{2n+\alpha}{n+\alpha}
 \left(\frac{(\alpha)_n}{n!}\right)^2.
\]
Multiplying by $(1+\alpha)_n/n!$ proves that the right side of
\eqref{VI-E6r:finite-transform} has exactly the defining increment of
$S_\alpha(n)$. Both sides equal one at $n=0$, proving the identity.
These calculations hold first away from zero denominators and then
as rational-function identities in $\alpha$.
There is a second exact transformation
\begin{equation}\label{VI-E6r:Euler}
 W_\alpha(n)=(-1)^n\sum_{k=0}^n
 \frac{(-n)_k(n+\alpha+1)_k(1-\alpha)_k}{(k!)^3}.
\end{equation}
Here is a formal power-series proof, which avoids convergence issues.
The left side is the coefficient of $z^n$ in
\[
 (1-z)^{-1}{}_2F_1(\alpha,-n-\alpha;1;z).
\]
Euler's transformation, valid as a formal identity, rewrites this as
\[
 (1-z)^n{}_2F_1(1-\alpha,n+\alpha+1;1;z).
\]
The coefficient of $z^{n-k}$ in $(1-z)^n$ is
$(-1)^{n-k}\binom nk=(-1)^n(-n)_k/k!$.
Taking the coefficient of $z^n$ gives \eqref{VI-E6r:Euler}.

Now let $L\equiv5\pmod6$, $n=(L-2)/3$ and $t=L/3$.
Then $n$ is odd and $-n=2/3-t$. Put
\[
 P_n=\frac{(4/3)_n}{n!},\qquad
 H_n(t)=\sum_{k=0}^n
 \frac{(2/3-t)_k(2/3)_k(2/3+t)_k}{(k!)^3}.
\]
The original E.3 sum is exactly
\begin{equation}\label{VI-E6r:symmetric}
 \boxed{\quad E(L)=-P_nH_n(L/3).\quad}
\end{equation}
In particular the new finite hypergeometric expression has a genuine
quadratic symmetry in its deformation parameter. With
$c_k=(2/3)_k/k!$, it is
\begin{equation}\label{VI-E6r:quadratic}
 H_n(L/3)=\sum_{k=0}^n c_k^3
 \prod_{j=0}^{k-1}\left(1-\frac{L^2}{(3j+2)^2}\right).
\end{equation}
No denominator deformation occurs in this expression.

\subsection{Exact prefactor valuation and a sharp termwise obstruction}
Let $p>3$, $p\equiv2\pmod3$, $r=2s+1\ge3$, $L=p^r$ and
$n=(p^r-2)/3$. Then
\begin{equation}\label{VI-E6r:prefactor}
 v_p(P_n)=s.
\end{equation}
Indeed $P_n=(3n+1)(1/3)_n/n!=(p^r-1)(1/3)_n/n!$.
For each $\ell\le r$, the denominator contributes
$\lfloor n/p^\ell\rfloor$ factors divisible by $p^\ell$.
If $\ell$ is odd, the first numerator index is
$(2p^\ell-1)/3$, while $n\bmod p^\ell=(p^\ell-2)/3$;
there is no extra numerator factor. If $\ell$ is even, the first
index is $(p^\ell-1)/3$, while
$n\bmod p^\ell=2(p^\ell-1)/3$; there is exactly one extra factor.
There are $s$ positive even indices below $r$. No higher level
contributes, proving \eqref{VI-E6r:prefactor}.

Every $c_k$ is $p$-integral (equivalently,
$(2/3)_k/k!=(-1)^k\binom{-2/3}{k}$ is a $p$-adic binomial
coefficient). The factors $3j+2$ in \eqref{VI-E6r:quadratic} satisfy
$0<3j+2<L$ for $j<n$, so each displayed deformation factor is
in $1+p^2\mathbb Z_p$. Consequently, with
$U_n=\sum_{k=0}^n c_k^3$,
\begin{equation}\label{VI-E6r:weak-freeze}
 H_n(L/3)-U_n\in p^2\mathbb Z_p,
 \qquad E(L)+P_nU_n\in p^{s+2}\mathbb Z_p.
\end{equation}
The first bound cannot be improved by bounding each summand
separately. Take $k=p^{r-1}$, which lies below $n$ since $p\ge5$.
The coefficient $c_k$ is a $p$-adic unit: at every level
$\ell\le r-1$, the numerator and denominator counts both equal
$p^{r-1-\ell}$, and at level $r$ neither product has a divisible
factor. Among the integers $3j+2$, $0\le j<k$, the only one
of valuation $r-1$ is $2p^{r-1}$. All others have valuation at
most $r-2$. Therefore
\[
 \prod_{j=0}^{k-1}\left(1-\frac{p^{2r}}{(3j+2)^2}\right)-1
 \equiv-\frac{p^2}{4}\pmod{p^4},
\]
and this single summand of $H_n(L/3)-U_n$ has valuation exactly
$2$. Thus any improvement needed for E.3 must use cancellation
between different values of $k$.

\subsection{A related exactly vanishing Dixon deformation}
For the same $n$ and $t=n+2/3$, define
\[
 D_n(t)=\sum_{k=0}^n
 \frac{(-n)_k(2/3)_k(2/3+t)_k}
 {k!(1-t)_k(1-2t)_k}.
\]
This sum vanishes exactly. More generally,
\[
 \sum_{k=0}^n\frac{(-n)_k(b)_k(c)_k}
 {k!(1-n-b)_k(1-n-c)_k}=0\qquad(n\text{ odd}),
\]
whenever the finite denominators are nonzero. Reversing $k$ to
$n-k$ multiplies every term by $(-1)^n$, proving the identity.
Here $b=2/3$, $c=2/3+t$ give precisely the displayed denominators.
Their linear shifts are not symmetric about one; consequently
this zero identity does not justify removing the denominator
shifts at the E.3 target precision.

\paragraph{Precise remaining target in this representation.}
Let $m=(p^{r-2}-2)/3$ and $t'=p^{r-2}/3$. The unsolved E.3
recurrence is equivalent, by \eqref{VI-E6r:symmetric}, to
\[
 P_nH_n(p^r/3)-p^2P_mH_m(t')\in p^{3r-1}\mathbb Z_p.
\]
For $r\ge3$, $v_p(P_n)=s$ and $v_p(P_m)=s-1$.
Thus after dividing by $P_n$ one needs the aggregate congruence
\[
 H_n(p^r/3)-p^2\frac{P_m}{P_n}H_m(t')
 \in p^{5s+2}\mathbb Z_p,
\]
where the multiplier $p^2P_m/P_n$ has valuation one.
This is an exact reformulation with a symmetric deformation,
not a proof of the required aggregate congruence.

\subsection{Exact two-step stripping of the prefactor}
There is also an exact symmetric gamma description of the prefactor
ratio. Write $L=p^r$, $M=L/p^2$, $n=(L-2)/3$, $m=(M-2)/3$,
and $n_1=(L/p-1)/3$. Let $\Gamma=\Gamma_p$ and put, for
$X\in p\mathbb Z_p$,
\[
 R_p(X)=\frac1{\Gamma(1/3-X/3)\Gamma(1/3+X/3)}.
\]
Then
\begin{equation}\label{VI-E6r:prefactor-strip}
 P_n=-pP_m\frac{R_p(L)}{R_p(L/p)}.
\end{equation}
To check both the constant and the sign, first extract the factors
divisible by $p$ from $(4/3)_n/n!$. The numerator and denominator
have $n_1$ such factors, and their normalized ratio is
$(2/3)_{n_1}/n_1!$. The remaining unit ratio is
\[
 -\frac{\Gamma(2/3+L/3)}{\Gamma(4/3)\Gamma(1/3+L/3)}
 =\frac{3\Gamma(2/3+L/3)}
 {\Gamma(1/3)\Gamma(1/3+L/3)}.
\]
In the second extraction, the numerator has $m+1$ factors divisible
by $p$ and the denominator has $m$. Their ratio is
\[
 p\frac{(1/3)_{m+1}}{m!}=\frac p3P_m.
\]
The unit ratio in this step is
\[
 -\frac{\Gamma(1/3+L/(3p))}
 {\Gamma(2/3)\Gamma(2/3+L/(3p))}.
\]
For $p\equiv2\pmod3$, the least positive residue of $1/3$ modulo
$p$ is $(p+1)/3$, which is even. Thus reflection gives
$\Gamma(1/3)\Gamma(2/3)=1$ and, for $X\in p\mathbb Z_p$,
\[
 \frac{\Gamma(2/3+X/3)}{\Gamma(1/3+X/3)}=R_p(X).
\]
Combining the two extractions proves
\eqref{VI-E6r:prefactor-strip}. In particular $R_p$ is even, so this
second deformation has no linear term either.

Using \eqref{VI-E6r:prefactor-strip}, the exact remaining aggregate
congruence can be written
\[
 H_n(L/3)+p\frac{R_p(L/p)}{R_p(L)}H_m(M/3)
 \in p^{5s+2}\mathbb Z_p,\qquad r=2s+1.
\]
The identities above do not establish this final congruence.

\section{Three exact interpolation grids at the short E cutoff}\label{loc-VI-3}
\label{VI-sec:E6models}
For an integer $m\ge0$ put $N=3m+2$ and
\[
 \mathcal S_m(z)=\sum_{k=0}^{m}(-1)^k(6k+1)
 \frac{(1/3)_k(1/3+z/6)_k(1/3-z/3)_k(2/3-z/3)_k}
 {k!(1-z/6)_k(1+z/3)_k(2/3+z/3)_k}.
\]
At $z=0$ this is the original E sum with cutoff $m$.
Define
\[
 D_m(z)=\prod_{j=1}^{m}
 (1-z/(6j))(1+z/(3j))(1+z/(3j-1)),
\]
\[
 Q_m(z)=\prod_{h=0}^{m}
 (1-z/(3h+1))(1-z/(3h+2))(1+z/(6h+2)).
\]
The cleared expression $P_m(z)=D_m(z)\mathcal S_m(z)$ is a
polynomial of degree at most $3m$. Its $3m+3$ distinct interpolation
nodes are
\[
 C_h=3h+1,\qquad F_h=3h+2,\qquad B_h=-6h-2
 \quad(0\le h\le m).
\]
At these nodes the exact cleared values are
\begin{equation}\label{VI-E6:models}
 P_m(H)=D_m(H)T(H),\qquad
 \begin{cases}
 T(C_h)=C_h,&h\text{ even},\\
 T(C_h)=0,&h\text{ odd},\\
 T(F_h)=(-1)^hF_h/2,\\
 T(B_h)=-B_h/2.
 \end{cases}
\end{equation}

\subsection{The two positive grids}
For complex $z$ with positive real part and away from denominator
poles, apply the nonterminating Whipple ${}_6F_5(-1)$ transformation~\cite{VI-Bailey}
with
\[
 a=1/3,\quad b=1/3+z/6,\quad c=1/3-z/3,
 \quad e=2/3,\quad f=2/3-z/3.
\]
The equal factors at $e=1+a-e=2/3$ cancel. The transformed Dixon
${}_3F_2(1)$ has upper parameters
$2/3,2/3+z/6,2/3-z/3$ and lower parameters
$1-z/6,1+z/3$.
The required convergence conditions reduce to
$\operatorname{Re}(z/3)>0$ and $\operatorname{Re}(z/6)>0$.
Whipple followed by Dixon and gamma reflection evaluates the series as
\[
 \frac z2\,
 \frac{\sin(\pi(1/3+z/6))}{\sin(\pi z/6)}.
\]
At $z=C_h$ or $F_h$, the original series terminates at $h$;
its lower parameters are never nonpositive integers at these
positive nodes. The displayed expression is respectively
$C_h$ or zero according to the parity of $h$, and
$(-1)^hF_h/2$. This proves the first three lines of
\eqref{VI-E6:models}. It also explains why the same sine expression
must not be substituted on the negative grid: there a lower
parameter becomes a nonpositive integer, and analytic continuation
does not represent the specified finite truncation.

\subsection{A finite WZ certificate for the negative grid}
At $z=B_h$, the finite terminating sum through $h$ is
\[
 \sum_{k=0}^{h} t(h,k),\qquad
 t(h,k)=(-1)^k(6k+1)
 \frac{(1/3)_k(-h)_k(2h+1)_k(2h+4/3)_k}
 {k!(h+4/3)_k(1/3-2h)_k(-2h)_k}.
\]
For $0\le k\le h$ the displayed denominator has no zero;
at $h=0$ only $k=0$ is used. We prove
\begin{equation}\label{VI-E6:Bmodel}
 \sum_{k=0}^{h}t(h,k)=3h+1.
\end{equation}
Put $u(h,k)=t(h,k)/(3h+1)$ and define
\[
 \begin{split}
 A(h,k)={}&864h^4+3024h^3-144h^2k^2+96h^2k+3912h^2\\
 &-264hk^2+176hk+2212h+9k^4-12k^3-119k^2+82k+460,
 \end{split}
\]
\[
 R(h,k)=
 -\frac{k(3k-1)(2h-k+1)(6h-3k+2)A(h,k)}
 {16(h+1)(2h+1)^2(3h+2)(6h+5)(6h+7)(6k+1)(h-k+1)}.
\]
Direct rational simplification gives the certificate
\begin{equation}\label{VI-E6:WZ}
 u(h+1,k)-u(h,k)
 =R(h,k+1)u(h,k+1)-R(h,k)u(h,k).
\end{equation}
To interpret the upper boundary, cancel the apparent pole using
\[
 \frac{(-h)_k}{h-k+1}=
 \frac{(-h-1)_k}{h+1}.
\]
Thus the product $R(h,k)u(h,k)$ has a well-defined value at
$k=h+1$. Direct cancellation gives
\[
 R(h,0)u(h,0)=0,\qquad
 R(h,h+1)u(h,h+1)=-u(h+1,h+1).
\]
Summing \eqref{VI-E6:WZ} for $0\le k\le h$ and adding the last term
$u(h+1,h+1)$ proves that the sum of $u(h,k)$ over $0\le k\le h$
is independent of $h$. Its value at $h=0$ is one, proving
\eqref{VI-E6:Bmodel}.

\subsection{The apparent poles in the cleared interpolation}
At $z=B_h$, the only possible zero in $D_m$ is its factor indexed
by $j=2h+1$ in $1+z/(3j-1)$.
If $2h\ge m$, this index is outside $D_m$; the original finite
sum terminates at $h$, and \eqref{VI-E6:Bmodel} gives
$P_m(B_h)=D_m(B_h)(3h+1)$.
If $2h<m$, then $D_m(B_h)=0$. In each cleared summand with
$k\le h$ the same zero remains in the denominator-clearing tail.
In each summand with $k\ge h+1$ the numerator contains the factor
$(1/3+z/6+h)$, which vanishes at $B_h$.
Every cleared summand is therefore zero, so $P_m(B_h)=0$.
This is again exactly $D_m(B_h)(3h+1)$ and proves
\eqref{VI-E6:models} at every negative node, without evaluating any
$0/0$ summand.

Consequently interpolation of $P_m$ at all the listed nodes gives
the exact identity
\[
 \mathcal S_m(0)=\sum_H W_m(H)T(H),\qquad
 W_m(H)=\frac{D_m(H)}{\prod_{K\ne H}(1-H/K)},
\]
with $K,H$ ranging over the three grids. The weights may be zero;
their rational denominators are never zero.
Equivalently,
\[
 \frac{D_m(z)\mathcal S_m(z)}{Q_m(z)}
 =\sum_H\frac{W_m(H)T(H)}{1-z/H}.
\]
The left side is $O(z^{-3})$ at infinity, proving the exact
vanishing of its first and second positive moments.

\section{Exact moments for the three-grid short E deformation}\label{loc-VI-4}
\label{VI-sec:e-three-grid-moments}
The statements in this section are exact. The uniform $p$-integrality
assertion needed by the proposed moment argument is not assumed or
proved here.
Let $L=3m+2$, $m\ge0$, and $b_k=(1/3)_k/k!$. Define
\[
 S_m(z)=\sum_{k=0}^m(-1)^k(6k+1)b_k^3
 \prod_{h=0}^{k-1}
 \frac{(1-z/(3h+1))(1-z/(3h+2))(1+z/(6h+2))}
 {(1-z/(6h+6))(1+z/(3h+3))(1+z/(3h+2))}.
\]
Put
\[
 q_h(z)=(1-z/(3h+1))(1-z/(3h+2))(1+z/(6h+2)),
\]
\[
 d_h(z)=(1-z/(6h+6))(1+z/(3h+3))(1+z/(3h+2)),
\]
\[
 Q_m(z)=\prod_{h=0}^m q_h(z),\qquad
 D_m(z)=\prod_{h=0}^{m-1}d_h(z),\qquad
 F_m(z)=\frac{D_m(z)S_m(z)}{Q_m(z)}.
\]
The numerator $D_mS_m$ is a polynomial of degree at most $3m$,
whereas $Q_m$ has degree $3m+3$ and distinct roots. Thus $F_m$ is
proper and has a Laurent expansion at infinity
\[
 F_m(z)=-\sum_{j\ge1}M_j(L)z^{-j},\qquad M_1(L)=M_2(L)=0.
\]
Equivalently, if $F_m(z)=\sum_H\mu_m(H)/(1-z/H)$ is its partial
fraction expansion at its three node grids, then
$M_j(L)=\sum_H\mu_m(H)H^j$. This definition allows a cancelled pole
to have residue weight zero; it does not assume that all values of
$S_m$ at lower-parameter collisions are terminating evaluations.

\begin{proposition}\label{VI-prop:e-three-grid-moments}
Set
\[
 C(L)=(-1)^{m+1}L^2(L-1)^3b_m^3.
\]
There are uniquely determined even polynomials $P_j(L)\in\mathbb Q[L^2]$,
$j\ge0$, with $P_0=1$ and $\deg P_j\le2\lfloor2j/3\rfloor$, such that
\[
 M_{j+3}(L)=C(L)P_j(L).
\]
In particular,
\[
 P_1=1,\qquad P_2=2L^2-3,\qquad
 P_3=-\frac{L^4-13L^2+16}{4}.
\]
\end{proposition}
\begin{proof}
At infinity the ratio of the leading terms of $q_h$ and $d_h$ is
$-(3h+3)^2/(3h+1)^2$. Hence
\[
 S_m(\infty)=\sum_{k=0}^m(6k+1)b_k
 =\frac{L(L-1)b_m}{2}.
\]
The last equality telescopes: the difference between
$(3k+1)(3k+2)b_k/2$ and its value at $k-1$ is $(6k+1)b_k$.
The leading coefficient of $D_m/Q_m$ is
$(-1)^m2L(L-1)^2b_m^2z^{-3}$. This proves $M_3=C(L)$.

Separating the last summand of $S_m$ gives
\begin{equation}\label{VI-eq:e-three-grid-rational}
 q_m(z)F_m(z)=d_{m-1}(z)F_{m-1}(z)
                         +(-1)^m(2L-3)b_m^3.
\end{equation}
Write
\[
 A=L(L-1),\quad B=(L-2)(L-3),\quad
 a=(L-1)(3L-2),\quad c=(L-2)(3L-7).
\]
The two cubic factors have expansions
\[
 q_m(z)=\frac{z^3-z^2-az+2L(L-1)^2}{2L(L-1)^2},
\]
\[
 d_{m-1}(z)=
 -\frac{z^3-z^2-cz-2(L-3)(L-2)^2}{2(L-3)(L-2)^2}.
\]
Since $b_m/b_{m-1}=(L-4)/(L-2)$, comparison of coefficients in
\eqref{VI-eq:e-three-grid-rational}, after dividing by $C(L)$, yields
\begin{align*}
 A P_j(L)-B P_j(L-3)
 &=A P_{j-1}(L)-B P_{j-1}(L-3)\\
 &\quad+A a P_{j-2}(L)-B c P_{j-2}(L-3)\\
 &\quad-2A L(L-1)^2P_{j-3}(L)
        -2B(L-3)(L-2)^2P_{j-3}(L-3),
\end{align*}
for $j\ge1$, where $P_{-1}=P_{-2}=P_{-3}=0$ and $P_0=1$.
This also provides an independent recursive definition of the
polynomials. Indeed the operator
\[
 \mathcal L(P)=L(L-1)P(L)-(L-2)(L-3)P(L-3)
\]
maps $\mathbb Q[L^2]$ to
$(2L-3)\mathbb Q[L(L-3)]$. The first space consists of even
polynomials; the second consists of polynomials anti-invariant
under $L\mapsto3-L$. In the displayed power bases,
$\mathcal L(L^{2k})/(2L-3)$ has degree $k$ and leading coefficient
$3k+2$. Consequently this operator is an isomorphism of the filtered
vector spaces over $\mathbb Q$.
For even preceding polynomials, the right side of the recurrence is
anti-invariant under $L\mapsto3-L$; the two cubic terms exchange
with opposite signs under this reflection. Thus the recurrence
uniquely defines an even $P_j$ at every step. It gives the degree bound
\[
 \deg P_j\le\max\{\deg P_{j-1},\ \deg P_{j-2}+2,
                                    \deg P_{j-3}+4\}
 \le2\lfloor2j/3\rfloor.
\]
At $L=2$, evaluation of the polynomial recurrence gives the
coefficients of
\[
 \sum_{j\ge0}P_j(2)t^j=\frac1{1-t-4t^2+4t^3}.
\]
This is exactly the normalized expansion of $F_0=1/q_0$.
The recurrence in $m$ now proves the moment identities for every
integer $m\ge0$.
\end{proof}

\paragraph{Prime-power valuation of the first moments.}
For $p>3$, $p\equiv2\pmod3$, and odd $r$, take $L=p^r$.
Then
\[
 v_p(b_{(p^r-2)/3})=\frac{r-1}{2},\qquad
 v_p(C(p^r))=\frac{7r-3}{2}.
\]
For the first equality count multiples of $p^\ell$ in the numerator
and denominator of $b_m$. If $\ell$ is even, the numerator acquires
one additional multiple, while for odd $\ell$ the counts agree;
there are $(r-1)/2$ positive even levels at most $r$.
Therefore $M_1=M_2=0$, and $M_3,M_4,M_5,M_6$ all have exactly the
valuation $(7r-3)/2$ for these primes. This exceeds or equals the
required $3r-1$.

\paragraph{The remaining arithmetic issue.}
The polynomial operator is not an isomorphism over
$\mathbb Z[1/2]$: its diagonal is $3k+2$, not a unit there.
For example
\[
 P_6(L)=\frac{L^8+290L^6-3325L^4+16666L^2-29900}{28},
\]
and $P_9$ has denominator $280$. Thus the rational-polynomial
existence proof alone does not establish $P_j(p^r)\in\mathbb Z_p$
for all $j$ at primes $p\equiv2\pmod3$. Such an integrality theorem,
or an alternative sufficient uniform bound for the full moments,
remains to be proved before using an Euler-power limit to discard
all unit-node contributions. The recurrence and the first moment
valuations above do not resolve that remaining issue.

\section{Exact two-step transport and the remaining cancellation}\label{loc-VI-5}
Let $p>3$, $p\equiv2\pmod3$, and $r\ge3$ be odd. Set
$L=p^r$, $M=L/p^2$, and $m=(L-2)/3$. Use the node and denominator
root sets $\mathcal H_L$ and $\mathcal D_L$ of the three-grid
construction, and write $\mu_L(H)=W_L(H)T(H)$.
Define the restricted rational power series
\[
 R_{L,p}(X)=
 \frac{\prod_{d\in\mathcal D_L,\,p^2\nmid d}(1-p^2X/d)}
 {\prod_{h\in\mathcal H_L,\,p^2\nmid h}(1-p^2X/h)}
 =1+\sum_{j\ge1}c_jX^j.
\]
Every retained root has valuation zero or one. Hence each linear
coefficient $p^2/d$ or $p^2/h$ belongs to $p\mathbb Z_p$.
Expanding the finitely many inverse factors proves
\[
 c_j\in p^j\mathbb Z_p,\qquad v_p(c_j)\longrightarrow+\infty.
\]
The series converges at every integral coarse node $J$.

\begin{proposition}[Exact residual identity]\label{VI-prop:e-residual}
Put
\[
 U_{L,p}=\sum_{\substack{H\in\mathcal H_L\\p^2\nmid H}}\mu_L(H).
\]
Then
\begin{equation}\label{VI-eq:e-residual}
 \boxed{E(L)-p^2E(M)
 =U_{L,p}+p^2\sum_{j\ge3}c_j M_j(M).}
\end{equation}
No assertion that either term on the right is separately divisible
by the target modulus is made.
\end{proposition}
\begin{proof}
Since $p^2\equiv1\pmod6$, dividing all $p^2$-divisible nodes by
$p^2$ gives exactly $\mathcal H_M$. Dividing the corresponding
denominator roots gives exactly $\mathcal D_M$. For example the
positive grids end at $L-1$ and $L$, and become the positive grids
ending at $M-1$ and $M$; the negative grid is $-2$ times the first
positive grid. The positive denominator roots $6i$ and the negative
roots $-3i$ both have coarse index range
$1\le i\le(M-2)/3$. The third denominator grid ends at $-(L-3)$
and becomes the grid ending at $-(M-3)$.

The exact model values satisfy $T(p^2J)=p^2T(J)$, because
multiplication by $p^2$ preserves the residues modulo six and the
sign of each node. Separating divisible factors in the weights
therefore gives the product identities
\[
 W_L(p^2J)=W_M(J)R_{L,p}(J),\qquad
 \mu_L(p^2J)=p^2\mu_M(J)R_{L,p}(J).
\]
These are identities of products, including at zero weights;
no division by a vanishing weight is needed. Sum over coarse nodes
and add the nodes not divisible by $p^2$. Interchanging the finite
node sum with the convergent series is legitimate. The identities
$M_1(M)=M_2(M)=0$ remove its first two positive terms and give
\eqref{VI-eq:e-residual}.
\end{proof}

For $r=2s+1$, the previously proved moment formula rewrites the
remaining assertion as the single explicit divisibility requirement
\[
 U_{L,p}+p^2C(M)\sum_{j\ge3}c_jP_{j-3}(M)
 \in p^{6s+2}\mathbb Z_p.
\]
This is equivalent to the original conjecture, not a completed
proof of it. The polynomial identities do not yet control the
whole series at this precision. Moreover $U_{L,p}$ contains both
the unit nodes and the nodes divisible by $p$ but not by $p^2$.
Eliminating only the first class does not eliminate the second.

\paragraph{A verified cancellation that must be proved uniformly.}
Exact modular arithmetic at $p=11$, $r=3$ gives
\[
 v_{11}\!\left(\sum_{11\nmid H}\mu_{11^3}(H)\right)=9,
 \qquad
 v_{11}\!\left(\sum_{v_{11}(H)=1}\mu_{11^3}(H)\right)=7,
\]
while the second term on the right of \eqref{VI-eq:e-residual} also
has valuation seven. Dividing the last two contributions by $11^7$
gives residues $8$ and $3$ modulo $11$. They cancel, and the
original recurrence difference has valuation eight. These are
finite checks, not a uniform proof. They show why separately
discarding both contributions would give a false proof.

\subsection*{Audit of the cubic coefficient in the three-grid transport}

Let $p>3$ satisfy $p\equiv2\pmod3$, put $L=p^3$ and
$m=(L-2)/3$.  With multiplicities, let
\begin{align*}
 \mathcal H&=\{3j+1,3j+2,-2(3j+1):0\le j\le m\},\\
 \mathcal D&=\{6j,-3j:1\le j\le m\}
       \cup\{-(3j+2):0\le j<m\}.
\end{align*}
The exact factor left after stripping the $p^2$-divisible nodes and
denominator roots is
\[
 R(X)=\frac{\prod_{d\in\mathcal D,\,p^2\nmid d}(1-p^2X/d)}
 {\prod_{h\in\mathcal H,\,p^2\nmid h}(1-p^2X/h)}.
\]
Write $R(X)=\sum c_jX^j$.  The elementary estimate
$c_j\in p^j\mathbb Z_p$ holds, but the proposed improvement
$c_3\in p^4\mathbb Z_p$ is false in general.

For a direct exact test, define
\[
 S_i=\sum_{\substack{h\in\mathcal H\\v_p(h)=1}}(h/p)^{-i}
 -\sum_{\substack{d\in\mathcal D\\v_p(d)=1}}(d/p)^{-i}
 \quad\hbox{in }\mathbb F_p.
\]
The formal logarithm of $R$ through degree three gives
\[
 \frac{c_3}{p^3}\equiv
 \frac{S_1^3+3S_1S_2+2S_3}{6}\pmod p.
\]
Roots of valuation zero cannot affect this congruence, since their
degree-$i$ logarithmic contribution is divisible by $p^{2i}$.
At $p=11$, the differences of multiplicities at the nonzero residues
$1,\ldots,10$ are
\[
 (-1,0,-1,0,2,0,-1,0,-1,0).
\]
Thus $(S_1,S_2,S_3)=(0,1,8)$ and $c_3/p^3\equiv-1\pmod{11}$.
At $p=17$ the residue is also $-1$, and at $p=23$ it is $3$.
The companion script computes these coefficients directly modulo $p^4$
as well as by the displayed Newton-sum formula.

This is an obstruction to the proposed transport argument, not a
counterexample to the original E.3 congruence.  A unit-node Euler
projector removes the nodes $p\nmid H$, but it does not remove nodes
with $v_p(H)=1$.  A $p^2$-strip argument must retain those nodes or give
an additional estimate for their combined contribution.

There is also a removable-specialization detail at the negative nodes.
For $H=-2(3j+1)$ and $2j<m$, a denominator root equals $H$.
The cleared numerator polynomial vanishes there: terms with $k\le j$
retain that clearing factor, whereas terms with $k>j$ contain a zero
upper factor.  These nodes therefore have interpolation weight zero.
Any stripping statement should be written as the exact product identity
$W_L(p^2J)=W_{L/p^2}(J)R(J)$, rather than as a quotient of weights
which can both vanish.


\section{Verification and unresolved obligations}\label{loc-VI-6}
The archive contains the single-file source and reproducible
exact-arithmetic programs. Independent checks include 123 terminating
node evaluations, 693 cleared-node values (including denominator
collisions), 21 complete interpolation identities, 84 low-moment
identities, and a second implementation of 90 moment identities
from direct Laurent division. The negative-grid WZ identity and
its boundary term were also checked symbolically.
Direct modular evaluations at $p=5,11,17$, $r=3$ cross-check the
three-grid decomposition against the original truncated sums.
Their recurrence-difference valuations are respectively $8,8,9$.
These checks support the formulas at the specified inputs; they
are not substitutes for an all-prime, all-level proof.

The remaining work is arithmetic: a sufficient uniform estimate
for all the moments, together with control of the combined
valuation-one-node contribution and the analytic transport
correction. The fact that the first few moments have large
valuation does not justify an estimate for all of them. The
explicit $p=11$ cubic-coefficient calculation also rules out
discarding the correction terms separately.

Thus \eqref{VI-eq:target} remains open within this research effort.
The note records exact reductions and the precise surviving
divisibility requirement, without labeling it a completed proof.

\begin{thebibliography}{9}
\bibitem{VI-Swisher}
H. Swisher, \emph{On the supercongruence conjectures of van Hamme},
Research in the Mathematical Sciences \textbf{2} (2015), Article 18,
Section~7, Table~2.
\url{https://doi.org/10.1186/s40687-015-0037-6}.
\bibitem{VI-Bailey}
W. N. Bailey, \emph{Generalized Hypergeometric Series}, Cambridge
University Press, 1935. Classical Whipple and Dixon transformations.
\bibitem{VI-DLMF}
NIST Digital Library of Mathematical Functions, Sections 15.8
and 16.4. Euler transformations and hypergeometric summation formulas.
\url{https://dlmf.nist.gov/15.8}; \url{https://dlmf.nist.gov/16.4}.
\end{thebibliography}

\clearpage\endgroup

%%%%%%%%%%%%%%%%%%%%%%%%%%%%%%%%%%%%%%%%%%%%%%%%%%%%%%%%%%%%%%%%%%%%%%%%%%
% VII | Swisher_E3_High_Odd_Levels_Proof.tex
%%%%%%%%%%%%%%%%%%%%%%%%%%%%%%%%%%%%%%%%%%%%%%%%%%%%%%%%%%%%%%%%%%%%%%%%%%
\clearpage\begingroup
\SetPaper{VII}{Higher odd levels of E.3}
\RestoreCounters
\PaperTitle{Finite Polynomial Separation and\\Higher Odd Levels of Swisher's E.3}{Keio research notes}{11 November 2026}

\begin{abstract}
For primes $p>3$ with $p\equiv2\pmod3$, we prove the odd-level
Swisher E.3 recurrence with its original cutoff and original modulus
at every odd level $r\ge5$.
The mechanism combines two exact finite interpolation formulas,
polynomial moment recurrences, a uniform logarithmic bound on
denominator loss, and a finite polynomial separating unit nodes.
An additional finite-field reflection produces a full extra factor
of $p$ in the companion numerator. The leakage of the unit polynomial onto divisible nodes is estimated
through the other moment family, rather than discarded termwise.
All identities and valuation estimates needed for these ranges
are included, with explicit finite certificates for the companion
interpolation. An explicitly exhausted finite set of small-prime
cases is settled by exact integer certificates. The third level
$r=3$ remains unresolved. Historical priority is not asserted.
\end{abstract}
\noindent\textbf{Scope and verification.} This is a proof manuscript
with internal independent checks and exact symbolic certificates;
it has not undergone external peer review. Finite computations
are identified separately from uniform theorems. In particular,
this manuscript does not claim the whole of Swisher's E.3.
\setcounter{tocdepth}{1}

\begin{LocalContents}
\LocalEntry{loc-VII-1}{1}{The unchanged problem and the proved range}
\LocalEntry{loc-VII-2}{2}{Three exact interpolation grids at the short E cutoff}
\LocalEntry{loc-VII-3}{3}{Exact moments for the three-grid short E deformation}
\LocalEntry{loc-VII-4}{4}{An exact negative-parameter companion and its first three moments}
\LocalEntry{loc-VII-5}{5}{A uniform logarithmic valuation bound for the primary E moments}
\LocalEntry{loc-VII-6}{6}{Polynomial structure of every companion moment}
\LocalEntry{loc-VII-7}{7}{A uniform logarithmic loss for all companion moments}
\LocalEntry{loc-VII-8}{8}{A uniform lower bound for every companion interpolation mass}
\LocalEntry{loc-VII-9}{9}{One-step transport of the valuation-one stratum}
\LocalEntry{loc-VII-10}{10}{The original odd E.3 congruence at all sufficiently high levels}
\LocalEntry{loc-VII-11}{11}{A reflection gain for companion interpolation samples}
\LocalEntry{loc-VII-12}{12}{Exact finite certificates for the remaining higher-level cases}
\LocalEntry{loc-VII-13}{13}{Exact finite completion of the level \texorpdfstring{$r=5$}{r=5}}
\LocalEntry{loc-VII-14}{14}{A uniform first-step congruence at the third level}
\LocalEntry{loc-VII-15}{15}{Exact two-step transport and the remaining cancellation}
\LocalEntry{loc-VII-16}{16}{The leading kernel coefficients from three short harmonic sums}
\LocalEntry{loc-VII-17}{17}{A uniform coarse-moment estimate and a rigorous transport-tail bound}
\LocalEntry{loc-VII-18}{18}{The single remaining odd level}
\LocalEntry{loc-VII-A}{A}{Finite certificates on the two positive companion grids}
\LocalEntry{loc-VII-B}{B}{A finite certificate for the negative companion grid}
\LocalEntry{loc-VII-C}{C}{Reproducible integer certificate evaluator}
\LocalEntry{loc-VII-D}{D}{Reproducible third-level unit-mass evaluator}
\end{LocalContents}
\clearpage

\section{The unchanged problem and the proved range}\label{loc-VII-1}
Write $(a)_k=\prod_{j=0}^{k-1}(a+j)$, $(a)_0=1$, and
$v_p(p)=1$, $v_p(0)=+\infty$. For $L\equiv2\pmod3$, $L\ge2$, set
\[
 E(L)=\sum_{k=0}^{(L-2)/3}(-1)^k(6k+1)
       \frac{(1/3)_k^3}{k!^3}.
\]
The odd branch of Swisher's Table 2, E.3~\cite{VII-Swisher}, asks for
\begin{equation}\label{VII-eq:target}
 E(p^r)-p^2E(p^{r-2})\in p^{3r-1}\mathbb Z_p,
 \qquad p>3,\quad p\equiv2\pmod3,\quad r\ge3\text{ odd}.
\end{equation}
Only this previously remaining branch is studied here. No result
from a different Swisher label is assumed as a proof of this one.

\begin{theorem}[Range established in this manuscript]\label{VII-thm:main-range}
Congruence \eqref{VII-eq:target} holds for every prime $p>3$ with
$p\equiv2\pmod3$ and every odd integer $r\ge5$:
\[
 \boxed{E(p^r)\equiv p^2E(p^{r-2})\pmod{p^{3r-1}},\qquad r\ge5\text{ odd}.}
\]
The range is established by uniform arguments plus exact finite
certificates for the explicitly enumerated exceptions.
\end{theorem}

The proof is supplied below after the exact interpolation and
valuation lemmas. Its central estimate is stronger than an
initial-moment calculation: every normalized moment loses at most
$O(\log j)$ powers of $p$ at moment order $j$. A polynomial of
degree $O(pr)$ can therefore separate the unit nodes at the target
precision. An exact companion formula controls its contribution
on the other nodes. This is the step that turns the previous
exact reductions into a proof on an infinite subrange.

\section{Three exact interpolation grids at the short E cutoff}\label{loc-VII-2}
\label{VII-sec:E6models}
For an integer $m\ge0$ put $N=3m+2$ and
\[
 \mathcal S_m(z)=\sum_{k=0}^{m}(-1)^k(6k+1)
 \frac{(1/3)_k(1/3+z/6)_k(1/3-z/3)_k(2/3-z/3)_k}
 {k!(1-z/6)_k(1+z/3)_k(2/3+z/3)_k}.
\]
At $z=0$ this is the original E sum with cutoff $m$.
Define
\[
 D_m(z)=\prod_{j=1}^{m}
 (1-z/(6j))(1+z/(3j))(1+z/(3j-1)),
\]
\[
 Q_m(z)=\prod_{h=0}^{m}
 (1-z/(3h+1))(1-z/(3h+2))(1+z/(6h+2)).
\]
The cleared expression $P_m(z)=D_m(z)\mathcal S_m(z)$ is a
polynomial of degree at most $3m$. Its $3m+3$ distinct interpolation
nodes are
\[
 C_h=3h+1,\qquad F_h=3h+2,\qquad B_h=-6h-2
 \quad(0\le h\le m).
\]
At these nodes the exact cleared values are
\begin{equation}\label{VII-E6:models}
 P_m(H)=D_m(H)T(H),\qquad
 \begin{cases}
 T(C_h)=C_h,&h\text{ even},\\
 T(C_h)=0,&h\text{ odd},\\
 T(F_h)=(-1)^hF_h/2,\\
 T(B_h)=-B_h/2.
 \end{cases}
\end{equation}

\subsection{The two positive grids}
For complex $z$ with positive real part and away from denominator
poles, apply the nonterminating Whipple ${}_6F_5(-1)$ transformation~\cite{VII-Bailey}
with
\[
 a=1/3,\quad b=1/3+z/6,\quad c=1/3-z/3,
 \quad e=2/3,\quad f=2/3-z/3.
\]
The equal factors at $e=1+a-e=2/3$ cancel. The transformed Dixon
${}_3F_2(1)$ has upper parameters
$2/3,2/3+z/6,2/3-z/3$ and lower parameters
$1-z/6,1+z/3$.
The required convergence conditions reduce to
$\operatorname{Re}(z/3)>0$ and $\operatorname{Re}(z/6)>0$.
Whipple followed by Dixon and gamma reflection evaluates the series as
\[
 \frac z2\,
 \frac{\sin(\pi(1/3+z/6))}{\sin(\pi z/6)}.
\]
At $z=C_h$ or $F_h$, the original series terminates at $h$;
its lower parameters are never nonpositive integers at these
positive nodes. The displayed expression is respectively
$C_h$ or zero according to the parity of $h$, and
$(-1)^hF_h/2$. This proves the first three lines of
\eqref{VII-E6:models}. It also explains why the same sine expression
must not be substituted on the negative grid: there a lower
parameter becomes a nonpositive integer, and analytic continuation
does not represent the specified finite truncation.

\subsection{A finite WZ certificate for the negative grid}
At $z=B_h$, the finite terminating sum through $h$ is
\[
 \sum_{k=0}^{h} t(h,k),\qquad
 t(h,k)=(-1)^k(6k+1)
 \frac{(1/3)_k(-h)_k(2h+1)_k(2h+4/3)_k}
 {k!(h+4/3)_k(1/3-2h)_k(-2h)_k}.
\]
For $0\le k\le h$ the displayed denominator has no zero;
at $h=0$ only $k=0$ is used. We prove
\begin{equation}\label{VII-E6:Bmodel}
 \sum_{k=0}^{h}t(h,k)=3h+1.
\end{equation}
Put $u(h,k)=t(h,k)/(3h+1)$ and define
\[
 \begin{split}
 A(h,k)={}&864h^4+3024h^3-144h^2k^2+96h^2k+3912h^2\\
 &-264hk^2+176hk+2212h+9k^4-12k^3-119k^2+82k+460,
 \end{split}
\]
\[
 R(h,k)=
 -\frac{k(3k-1)(2h-k+1)(6h-3k+2)A(h,k)}
 {16(h+1)(2h+1)^2(3h+2)(6h+5)(6h+7)(6k+1)(h-k+1)}.
\]
Direct rational simplification gives the certificate
\begin{equation}\label{VII-E6:WZ}
 u(h+1,k)-u(h,k)
 =R(h,k+1)u(h,k+1)-R(h,k)u(h,k).
\end{equation}
To interpret the upper boundary, cancel the apparent pole using
\[
 \frac{(-h)_k}{h-k+1}=
 \frac{(-h-1)_k}{h+1}.
\]
Thus the product $R(h,k)u(h,k)$ has a well-defined value at
$k=h+1$. Direct cancellation gives
\[
 R(h,0)u(h,0)=0,\qquad
 R(h,h+1)u(h,h+1)=-u(h+1,h+1).
\]
Summing \eqref{VII-E6:WZ} for $0\le k\le h$ and adding the last term
$u(h+1,h+1)$ proves that the sum of $u(h,k)$ over $0\le k\le h$
is independent of $h$. Its value at $h=0$ is one, proving
\eqref{VII-E6:Bmodel}.

\subsection{The apparent poles in the cleared interpolation}
At $z=B_h$, the only possible zero in $D_m$ is its factor indexed
by $j=2h+1$ in $1+z/(3j-1)$.
If $2h\ge m$, this index is outside $D_m$; the original finite
sum terminates at $h$, and \eqref{VII-E6:Bmodel} gives
$P_m(B_h)=D_m(B_h)(3h+1)$.
If $2h<m$, then $D_m(B_h)=0$. In each cleared summand with
$k\le h$ the same zero remains in the denominator-clearing tail.
In each summand with $k\ge h+1$ the numerator contains the factor
$(1/3+z/6+h)$, which vanishes at $B_h$.
Every cleared summand is therefore zero, so $P_m(B_h)=0$.
This is again exactly $D_m(B_h)(3h+1)$ and proves
\eqref{VII-E6:models} at every negative node, without evaluating any
$0/0$ summand.

Consequently interpolation of $P_m$ at all the listed nodes gives
the exact identity
\[
 \mathcal S_m(0)=\sum_H W_m(H)T(H),\qquad
 W_m(H)=\frac{D_m(H)}{\prod_{K\ne H}(1-H/K)},
\]
with $K,H$ ranging over the three grids. The weights may be zero;
their rational denominators are never zero.
Equivalently,
\[
 \frac{D_m(z)\mathcal S_m(z)}{Q_m(z)}
 =\sum_H\frac{W_m(H)T(H)}{1-z/H}.
\]
The left side is $O(z^{-3})$ at infinity, proving the exact
vanishing of its first and second positive moments.

\section{Exact moments for the three-grid short E deformation}\label{loc-VII-3}
\label{VII-sec:e-three-grid-moments}
The statements in this section are exact. The uniform $p$-integrality
assertion needed by the proposed moment argument is not assumed or
proved here.
Let $L=3m+2$, $m\ge0$, and $b_k=(1/3)_k/k!$. Define
\[
 S_m(z)=\sum_{k=0}^m(-1)^k(6k+1)b_k^3
 \prod_{h=0}^{k-1}
 \frac{(1-z/(3h+1))(1-z/(3h+2))(1+z/(6h+2))}
 {(1-z/(6h+6))(1+z/(3h+3))(1+z/(3h+2))}.
\]
Put
\[
 q_h(z)=(1-z/(3h+1))(1-z/(3h+2))(1+z/(6h+2)),
\]
\[
 d_h(z)=(1-z/(6h+6))(1+z/(3h+3))(1+z/(3h+2)),
\]
\[
 Q_m(z)=\prod_{h=0}^m q_h(z),\qquad
 D_m(z)=\prod_{h=0}^{m-1}d_h(z),\qquad
 F_m(z)=\frac{D_m(z)S_m(z)}{Q_m(z)}.
\]
The numerator $D_mS_m$ is a polynomial of degree at most $3m$,
whereas $Q_m$ has degree $3m+3$ and distinct roots. Thus $F_m$ is
proper and has a Laurent expansion at infinity
\[
 F_m(z)=-\sum_{j\ge1}M_j(L)z^{-j},\qquad M_1(L)=M_2(L)=0.
\]
Equivalently, if $F_m(z)=\sum_H\mu_m(H)/(1-z/H)$ is its partial
fraction expansion at its three node grids, then
$M_j(L)=\sum_H\mu_m(H)H^j$. This definition allows a cancelled pole
to have residue weight zero; it does not assume that all values of
$S_m$ at lower-parameter collisions are terminating evaluations.

\begin{proposition}\label{VII-prop:e-three-grid-moments}
Set
\[
 C(L)=(-1)^{m+1}L^2(L-1)^3b_m^3.
\]
There are uniquely determined even polynomials $P_j(L)\in\mathbb Q[L^2]$,
$j\ge0$, with $P_0=1$ and $\deg P_j\le2\lfloor2j/3\rfloor$, such that
\[
 M_{j+3}(L)=C(L)P_j(L).
\]
In particular,
\[
 P_1=1,\qquad P_2=2L^2-3,\qquad
 P_3=-\frac{L^4-13L^2+16}{4}.
\]
\end{proposition}
\begin{proof}
At infinity the ratio of the leading terms of $q_h$ and $d_h$ is
$-(3h+3)^2/(3h+1)^2$. Hence
\[
 S_m(\infty)=\sum_{k=0}^m(6k+1)b_k
 =\frac{L(L-1)b_m}{2}.
\]
The last equality telescopes: the difference between
$(3k+1)(3k+2)b_k/2$ and its value at $k-1$ is $(6k+1)b_k$.
The leading coefficient of $D_m/Q_m$ is
$(-1)^m2L(L-1)^2b_m^2z^{-3}$. This proves $M_3=C(L)$.

Separating the last summand of $S_m$ gives
\begin{equation}\label{VII-eq:e-three-grid-rational}
 q_m(z)F_m(z)=d_{m-1}(z)F_{m-1}(z)
                         +(-1)^m(2L-3)b_m^3.
\end{equation}
Write
\[
 A=L(L-1),\quad B=(L-2)(L-3),\quad
 a=(L-1)(3L-2),\quad c=(L-2)(3L-7).
\]
The two cubic factors have expansions
\[
 q_m(z)=\frac{z^3-z^2-az+2L(L-1)^2}{2L(L-1)^2},
\]
\[
 d_{m-1}(z)=
 -\frac{z^3-z^2-cz-2(L-3)(L-2)^2}{2(L-3)(L-2)^2}.
\]
Since $b_m/b_{m-1}=(L-4)/(L-2)$, comparison of coefficients in
\eqref{VII-eq:e-three-grid-rational}, after dividing by $C(L)$, yields
\begin{align*}
 A P_j(L)-B P_j(L-3)
 &=A P_{j-1}(L)-B P_{j-1}(L-3)\\
 &\quad+A a P_{j-2}(L)-B c P_{j-2}(L-3)\\
 &\quad-2A L(L-1)^2P_{j-3}(L)
        -2B(L-3)(L-2)^2P_{j-3}(L-3),
\end{align*}
for $j\ge1$, where $P_{-1}=P_{-2}=P_{-3}=0$ and $P_0=1$.
This also provides an independent recursive definition of the
polynomials. Indeed the operator
\[
 \mathcal L(P)=L(L-1)P(L)-(L-2)(L-3)P(L-3)
\]
maps $\mathbb Q[L^2]$ to
$(2L-3)\mathbb Q[L(L-3)]$. The first space consists of even
polynomials; the second consists of polynomials anti-invariant
under $L\mapsto3-L$. In the displayed power bases,
$\mathcal L(L^{2k})/(2L-3)$ has degree $k$ and leading coefficient
$3k+2$. Consequently this operator is an isomorphism of the filtered
vector spaces over $\mathbb Q$.
For even preceding polynomials, the right side of the recurrence is
anti-invariant under $L\mapsto3-L$; the two cubic terms exchange
with opposite signs under this reflection. Thus the recurrence
uniquely defines an even $P_j$ at every step. It gives the degree bound
\[
 \deg P_j\le\max\{\deg P_{j-1},\ \deg P_{j-2}+2,
                                    \deg P_{j-3}+4\}
 \le2\lfloor2j/3\rfloor.
\]
At $L=2$, evaluation of the polynomial recurrence gives the
coefficients of
\[
 \sum_{j\ge0}P_j(2)t^j=\frac1{1-t-4t^2+4t^3}.
\]
This is exactly the normalized expansion of $F_0=1/q_0$.
The recurrence in $m$ now proves the moment identities for every
integer $m\ge0$.
\end{proof}

\paragraph{Prime-power valuation of the first moments.}
For $p>3$, $p\equiv2\pmod3$, and odd $r$, take $L=p^r$.
Then
\[
 v_p(b_{(p^r-2)/3})=\frac{r-1}{2},\qquad
 v_p(C(p^r))=\frac{7r-3}{2}.
\]
For the first equality count multiples of $p^\ell$ in the numerator
and denominator of $b_m$. If $\ell$ is even, the numerator acquires
one additional multiple, while for odd $\ell$ the counts agree;
there are $(r-1)/2$ positive even levels at most $r$.
Therefore $M_1=M_2=0$, and $M_3,M_4,M_5,M_6$ all have exactly the
valuation $(7r-3)/2$ for these primes. This exceeds or equals the
required $3r-1$.

\paragraph{Why a uniform denominator estimate is needed.}
The polynomial operator is not an isomorphism over
$\mathbb Z[1/2]$: its diagonal is $3k+2$, not a unit there.
For example
\[
 P_6(L)=\frac{L^8+290L^6-3325L^4+16666L^2-29900}{28},
\]
and $P_9$ has denominator $280$. Thus the rational-polynomial
existence proof alone does not establish $P_j(p^r)\in\mathbb Z_p$
for all $j$ at primes $p\equiv2\pmod3$. The logarithmic bounds proved below supply the uniform estimate
needed in this manuscript. They are used with finite polynomial
separation, so no Euler-power limiting argument is required.

\section{An exact negative-parameter companion and its first three moments}\label{loc-VII-4}
\label{VII-sec:E7companion}
This section proves the exact finite companion identity used in the
higher-level argument. The identity alone is not a supercongruence.

Let $M=3n+1$, $n\ge0$, and put
\[
 \mathcal H_n=\{3h+1:0\le h\le n\}
 \cup\{3h+2:0\le h<n\}\cup\{-2(3h+2):0\le h<n\},
\]
\[
 \mathcal D_n=\{6j,-3j:1\le j\le n\}
 \cup\{-3h-1:0\le h<n\}.
\]
All $3n+1$ nodes in $\mathcal H_n$ are distinct and nonzero. Define
\[
 U(3h+1)=-(-1)^h(3h+1)/2,\qquad
 U(3h+2)=\begin{cases}3h+2,&h\text{ odd},\\0,&h\text{ even},\end{cases}
 \qquad U(-2(3h+2))=3h+2,
\]
\[
 D_n(z)=\prod_{D\in\mathcal D_n}(1-z/D),\qquad
 Q_n(z)=\prod_{H\in\mathcal H_n}(1-z/H),\qquad
 \widetilde W_n(H)=\frac{D_n(H)}{\prod_{K\in\mathcal H_n,K\ne H}(1-H/K)}.
\]
The companion and its positive moments are
\[
 \mathcal A(M)=\sum_H\widetilde W_n(H)U(H),\qquad
 \widetilde M_j(M)=\sum_H\widetilde W_n(H)U(H)H^j.
\]
Set $d_k=(-1/3)_k/k!$ and $c_k=(2/3)_k/k!$.

\begin{theorem}\label{VII-E7:companion-exact}
For every integer $n\ge0$,
\begin{align*}
 \mathcal A(3n+1)
 &=-\frac12\sum_{k=0}^n(-1)^k(6k-1)(9k^2-3k-1)d_k^3\\
 &=\sum_{k=0}^n(-1)^k(6k-1)d_k^3+\frac{(-1)^n}{2}c_n^3.
\end{align*}
Moreover, with $C_n=(-1)^{n+1}M^2c_n^3/2$,
\[
 \widetilde M_1(M)=C_n,\qquad
 \widetilde M_2(M)=C_n(M^2+1)/2,\qquad
 \widetilde M_3(M)=C_n(2M^2-1).
\]
\end{theorem}

\subsection{A rational deformation with exactly the companion node values}
Put
\[
 K_k(z)=\prod_{j=0}^{k-1}
 \frac{(1+z/[2(3j-1)])(1-z/(3j-1))(1-z/(3j+1))}
 {(1-z/[2(3j+3)])(1+z/(3j+3))(1+z/(3j+1))},
\]
\begin{equation}\label{VII-E7:V}
 V_n(z)=-\frac1{2(2-z)}\sum_{k=0}^n
 (-1)^k(6k-1)(z+18k^2-6k-2)d_k^3K_k(z).
\end{equation}
Although this expression displays $2-z$ in its denominator, its
$k=0$ term is the constant $-1/2$. Every $k\ge1$ term has the
factor $1-z/2$ in its numerator, cancelling this additional denominator.
Consequently $D_n(z)V_n(z)$ is a polynomial of degree at most $3n$.
At $z=2$, the terms with $k\ge2$ vanish after cancellation; the $k=1$
term simplifies to $-5(z-1)(z+10)/(6(z-6)(z+3))$ and has value $1/2$.
Thus $V_n(2)=0$ for every $n\ge1$.

The three finite certificate proofs supplied below establish
\begin{equation}\label{VII-E7:cleared-all}
 [D_n(z)V_n(z)]_{z=H}=D_n(H)U(H)
 \qquad(H\in\mathcal H_n).
\end{equation}
On the first positive grid the sum terminates at $h$; on the second
it terminates at $h+1$, with $z=2$ handled above. The lower parameters
are regular on both positive grids. On the negative grid the finite
certificate gives the sum through $h+1$; where a later denominator
would vanish, the cleared identity is proved term by term. Thus
\eqref{VII-E7:cleared-all} never assigns a value to a raw $0/0$ summand.

Interpolation of the degree-$3n$ polynomial $D_nV_n$ at its $3n+1$
distinct nodes now proves
\[
 V_n(0)=\sum_H\widetilde W_n(H)U(H)=\mathcal A(M),\qquad
 F_n(z):=\frac{D_n(z)V_n(z)}{Q_n(z)}
 =\sum_H\frac{\widetilde W_n(H)U(H)}{1-z/H}.
\]
Evaluation of \eqref{VII-E7:V} at zero proves the first equality in
Theorem~\ref{VII-E7:companion-exact}.

\subsection{An elementary telescoping simplification}
Let $u_k=(-1)^kd_k^3$ and $B_k=k^3u_k$. Since
$B_{k+1}=-(k-1/3)^3u_k$, direct expansion gives
\[
 (6k-1)(9k^2-3k-1)u_k
 =-27(B_{k+1}-B_k)-2(6k-1)u_k.
\]
Summing uses $B_0=0$ and $27B_{n+1}=(-1)^nc_n^3$ and proves the
second equality of the theorem, without any asymptotic argument.

\subsection{The adjacent-cutoff recurrence and the three moments}
For $n\ge1$, write
\[
 A_n(z)=(1-z/M)(1-z/(M-2))(1+z/[2(M-2)]),
\]
\[
 B_n(z)=(1-z/[2(M-1)])(1+z/(M-1))(1+z/(M-3)).
\]
These are precisely $Q_n/Q_{n-1}$ and $D_n/D_{n-1}$.
The numerator of $K_n(z)$ is
$(1+z)(1-z/2)Q_{n-1}(z)$. Separating the last summand of
\eqref{VII-E7:V} therefore gives the exact polynomial recurrence
\begin{equation}\label{VII-E7:companion-recurrence}
 A_n(z)F_n(z)-B_n(z)F_{n-1}(z)=G_n(z),
\end{equation}
\[
 G_n(z)=-\frac{(-1)^n(6n-1)d_n^3}{4}
 (z+1)(z+18n^2-6n-2).
\]
This derivation uses the summand formula, not the moment identities
being proved.

At infinity, write
$F_n(z)=-\sum_{j\ge1}\widetilde M_j(M)z^{-j}$.
The coefficients of $z^2,z^1,z^0$ in
\eqref{VII-E7:companion-recurrence} determine its first three moments
successively. For complete verification, write $A_n=\sum a_i z^i$,
$B_n=\sum b_i z^i$, and $G_n=\sum g_i z^i$. Those three equations are
\[
 -a_3\widetilde M_1+b_3\widetilde M'_1=g_2,
\]
\[
 -a_2\widetilde M_1-a_3\widetilde M_2
 +b_2\widetilde M'_1+b_3\widetilde M'_2=g_1,
\]
\[
 -a_1\widetilde M_1-a_2\widetilde M_2-a_3\widetilde M_3
 +b_1\widetilde M'_1+b_2\widetilde M'_2+b_3\widetilde M'_3=g_0,
\]
where primes mean the preceding cutoff $M-3$. Here
\[
 a_3=\frac1{2M(M-2)^2},\quad a_2=-2a_3,\quad
 a_1=-\frac{3M-4}{2M(M-2)},
\]
\[
 b_3=-\frac1{2(M-1)^2(M-3)},\quad b_2=-2b_3,\quad
 b_1=\frac{3M-5}{2(M-1)(M-3)}.
\]
Substitution of the three claimed formulas satisfies these rational
identities, using
\[
 \frac{C_{n-1}}{C_n}
 =-\frac{(M-3)^2(M-1)^3}{M^2(M-2)^3}.
\]
At $n=0$, $F_0(z)=-1/[2(1-z)]$, so every positive moment is $-1/2$,
exactly the claimed initial values. Since $a_3\ne0$ for $n\ge1$,
induction proves all three formulas. This completes the theorem.


\section{A uniform logarithmic valuation bound for the primary E moments}\label{loc-VII-5}
\label{VII-sec:e-uniform-log-bound}
Let $P_j(L)\in\mathbb Q[L^2]$ be the primary moment polynomials defined
by the exact three-grid recurrence, normalized by $P_0=1$. They satisfy
\[
 \deg P_j\le d_j:=2\left\lfloor\frac{2j}{3}\right\rfloor,
 \qquad M_{j+3}(L)=C(L)P_j(L),
\]
where, for $L=3m+2$,
\[
 C(L)=(-1)^{m+1}L^2(L-1)^3b_m^3,\qquad b_m=\frac{(1/3)_m}{m!}.
\]
For a prime $p>3$, $p\equiv2\pmod3$, define
\[
 a_j=\left\lfloor\log_p(3d_j+2)\right\rfloor,
 \qquad e_j=a_j+\left\lfloor\frac{a_j}{2}\right\rfloor.
\]
Here the logarithm only denotes the integer characterized by
$p^{a_j}\le3d_j+2<p^{a_j+1}$.

\begin{theorem}\label{VII-thm:e-log-bound}
For every $j\ge0$ and every $x\in\mathbb Z_p$,
\[
 \boxed{v_p(P_j(x))\ge-e_j.}
\]
Equivalently, $p^{e_j}P_j$ is integer-valued on $\mathbb Z_p$.
Moreover $e_0=0$, and $e_j\le j-1$ for every $j\ge1$.
\end{theorem}
This is an all-degree theorem with a logarithmic loss. It does not
assert the stronger, still unproved statement
$P_j\in\operatorname{Int}(\mathbb Z_p)$ for all $j$.

\subsection{An integral-coefficient formula before normalization}
For $h\ge0$, put
\begin{align*}
 \widehat q_h(t)&=(1-(3h+1)t)(1-(3h+2)t)(1+(6h+2)t),\\
 \widehat d_h(t)&=(1-(6h+6)t)(1+(3h+3)t)(1+(3h+2)t).
\end{align*}
For $0\le k\le m$, define integers $\phi_j(k,m)$ by
\[
 \frac{\prod_{h=k}^{m-1}\widehat d_h(t)}
      {\prod_{h=k}^{m}\widehat q_h(t)}
 =\sum_{j\ge0}\phi_j(k,m)t^j.
\]
All these coefficients are integers because every denominator
polynomial has constant term one. The exact normalized moment formula is
\begin{equation}\label{VII-eq:e-integral-phi}
 P_j(3m+2)=\frac{2}{L(L-1)b_m}
          \sum_{k=0}^m(6k+1)b_k\phi_j(k,m),\qquad L=3m+2.
\end{equation}
To verify it directly, expand the proper rational function
$F_m=D_mS_m/Q_m$ term by term at $z=\infty$, using $t=1/z$.
The $k$th summand is
\[
 (-1)^k(6k+1)b_k^3
 \frac{\prod_{h=k}^{m-1}d_h(z)}{\prod_{h=k}^{m}q_h(z)}.
\]
Its leading coefficient is
$(-1)^m2L(L-1)^2b_m^2(6k+1)b_kz^{-3}$.
Dividing the reversed numerator and denominator polynomials by
their leading coefficients gives exactly the series defining
$\phi_j(k,m)$. Comparison with
$F_m=-C(L)\sum_{j\ge0}P_j(L)z^{-j-3}$ proves
\eqref{VII-eq:e-integral-phi}. In particular this argument holds for
every coefficient, not only for any prescribed finite range.

\subsection{Valuation of the normalizing denominator}
\begin{lemma}\label{VII-lem:e-normalization-bound}
Let $m\ge0$, $L=3m+2$, and $a=\lfloor\log_p L\rfloor$. Then
\[
 b_k\in\mathbb Z_p\quad(k\ge0),\qquad
 v_p\bigl(L(L-1)b_m\bigr)\le a+\lfloor a/2\rfloor.
\]
\end{lemma}
\begin{proof}
At the level $P=p^\ell$, let $u_\ell$ be the least nonnegative
solution of $3u_\ell+1\equiv0\pmod P$. Explicitly,
\[
 u_\ell=\begin{cases}(P-1)/3,&\ell\text{ even},\\
                       (2P-1)/3,&\ell\text{ odd}.
          \end{cases}
\]
The difference between the counts of multiples of $P$ in
$\prod_{h=0}^{m-1}(3h+1)$ and in $m!$ is
\[
 \delta_\ell=\begin{cases}1,&m\bmod P>u_\ell,\\0,&m\bmod P\le u_\ell.
                  \end{cases}
\]
The factors of $3$ in $b_m$ are units. Thus
$v_p(b_m)=\sum_{\ell\ge1}\delta_\ell\ge0$, proving integrality.
Also $\delta_\ell=0$ for $p^\ell>L$: all positive numerator
factors are at most $L-4$ and all denominator indices are at most $m$.

If $t=v_p(L)>0$, then for $\ell\le t$ the congruence
$3m\equiv-2\pmod{p^\ell}$ gives
\[
 m\bmod p^\ell=\begin{cases}
 2(p^\ell-1)/3,&\ell\text{ even},\\
 (p^\ell-2)/3,&\ell\text{ odd}.
 \end{cases}
\]
Consequently $\delta_\ell=1$ precisely at the even levels among
$1,\ldots,t$. Since $L-1$ is a unit,
\[
 v_p(L(L-1)b_m)
 \le t+\lfloor t/2\rfloor+(a-t)
 =a+\lfloor t/2\rfloor\le a+\lfloor a/2\rfloor.
\]
If instead $u=v_p(L-1)>0$, then $m\bmod p^\ell=u_\ell$ for
$\ell\le u$. All those $\delta_\ell$ are zero, and
$v_p(L(L-1)b_m)\le u+(a-u)=a$.
If both $L$ and $L-1$ are units, the same bound follows from
$\sum\delta_\ell\le a$. These cases include $m=0$, where $L=2$
and $a=0$.
\end{proof}

\begin{proof}[Proof of Theorem~\ref{VII-thm:e-log-bound}]
Formula \eqref{VII-eq:e-integral-phi} and the preceding lemma imply, for
every nonnegative integer $m$,
\[
 v_p(P_j(3m+2))\ge
 -\left(\lfloor\log_p(3m+2)\rfloor+
 \left\lfloor\frac{\lfloor\log_p(3m+2)\rfloor}{2}\right\rfloor\right).
\]
Use this for $m=0,1,\ldots,d_j$. Every one of these values lies in
$p^{-e_j}\mathbb Z_p$. Set $R_j(X)=P_j(3X+2)$, whose degree is at
most $d_j$. Its exact Newton expansion is
\[
 R_j(X)=\sum_{k=0}^{d_j}\Delta^kR_j(0)\binom Xk.
\]
Each finite difference is an integer linear combination of the
sampled values, so it too belongs to $p^{-e_j}\mathbb Z_p$.
For every $X\in\mathbb Z_p$, the binomial polynomial $\binom Xk$
is $p$-integral, by continuity from the nonnegative integers.
Taking $X=(x-2)/3\in\mathbb Z_p$ proves the uniform bound.
No division by $k!$ is being applied to the finite-difference
coefficient: that factorial belongs to the integer-valued basis.

Finally, for $j\ge1$ the elementary inequality
\[
 6\lfloor2j/3\rfloor+2<5^{\lfloor(2j-1)/3\rfloor+1}
\]
follows by checking $j=1,2,3$ and then increasing $j$ by three;
the left side increases by $12$, while the right side is multiplied
by $25$. Since $p\ge5$, this implies
$a_j\le\lfloor(2j-1)/3\rfloor$, and hence
$a_j+\lfloor a_j/2\rfloor\le j-1$. Also $d_0=0$, so $e_0=0$.
\end{proof}

\subsection{The double-strip transport is controlled at every higher odd level}
\begin{corollary}\label{VII-cor:e-double-strip-transport-bound}
Let $r\ge3$ be odd, $M=p^{r-2}$, and let $c_j\in\mathbb Q_p$ satisfy
$v_p(c_j)\ge j$ for every $j\ge1$. Then the series
\[
 V_r=p^2\sum_{j\ge1}c_jM_j(M)
\]
converges. For every odd $r\ge5$,
\[
 V_r\in p^{3r-1}\mathbb Z_p.
\]
At the remaining level $r=3$, one has the precise reduction
\[
 V_3\equiv p^2c_3M_3(p)\pmod{p^8}.
\]
\end{corollary}
\begin{proof}
The first two moments vanish exactly. For $j\ge3$, put $n=j-3$.
The exact primary moment normalization and the proven valuation
$v_p(C(p^{r-2}))=(7r-17)/2$ give
\[
 v_p\bigl(p^2c_jM_j(M)\bigr)
 \ge2+j+\frac{7r-17}{2}-e_n
 =\frac{7r-7}{2}+n-e_n.
\]
The right side tends to infinity, because $e_n=O(\log n)$, so
convergence follows. For $n=0$ it is at least $(7r-7)/2$;
for $n\ge1$ it is at least $(7r-5)/2$.
When $r\ge5$, the former bound is at least $3r-1$. When $r=3$,
all terms with $n\ge1$ belong to $p^8\mathbb Z_p$, leaving only
the displayed cubic term.
\end{proof}


\section{Polynomial structure of every companion moment}\label{loc-VII-6}
\label{VII-sec:e-companion-polynomials}
Let $\widetilde M_j(M)$ be the companion moments from the proved
negative-parameter companion identity, and put
\[
 C_n=\frac{(-1)^{n+1}M^2}{2}
       \left(\frac{(2/3)_n}{n!}\right)^3,\qquad M=3n+1.
\]
\begin{proposition}\label{VII-prop:e-companion-polynomials}
There exist even polynomials $R_j(M)\in\mathbb Q[M^2]$, $j\ge0$,
with
\[
 \deg R_j\le2\left\lfloor\frac{2j+1}{3}\right\rfloor,
 \qquad R_j(1)=1,
\]
such that, for every integer $n\ge0$,
\[
 \widetilde M_{j+1}(3n+1)=C_nR_j(3n+1).
\]
Their first values are
$R_0=1$, $R_1=(M^2+1)/2$, and $R_2=2M^2-1$.
\end{proposition}
\begin{proof}
Start with the already proved rational-function recurrence
$A_n(z)F_n(z)-B_n(z)F_{n-1}(z)=G_n(z)$, where $G_n$ has degree two.
Its first three Laurent equations give the three initial moments
just stated. Set
\[
 A=M(M-2),\quad B=(M-3)(M-1),\quad
 a=(M-2)(3M-4),\quad c=(M-1)(3M-5).
\]
The monic forms of its cubic factors are
\[
 A_n(z)=\frac{z^3-2z^2-az+2M(M-2)^2}{2M(M-2)^2},
\]
\[
 B_n(z)=-\frac{z^3-2z^2-cz-2(M-3)(M-1)^2}
                       {2(M-3)(M-1)^2}.
\]
Use the exact ratio
\[
 \frac{C_{n-1}}{C_n}
 =-\frac{(M-3)^2(M-1)^3}{M^2(M-2)^3}.
\]
The homogeneous Laurent equations, starting at the coefficient of
$z^{-1}$, yield, for every $j\ge3$,
\begin{align*}
 A R_j(M)-B R_j(M-3)
 &=2\bigl[A R_{j-1}(M)-B R_{j-1}(M-3)\bigr]\\
 &\quad+A a R_{j-2}(M)-B c R_{j-2}(M-3)\\
 &\quad-2A M(M-2)^2R_{j-3}(M)
       -2B(M-3)(M-1)^2R_{j-3}(M-3).
\end{align*}
We show this recurrence uniquely defines even rational polynomials.
The operator
\[
 \mathcal L'(P)=M(M-2)P(M)-(M-3)(M-1)P(M-3)
\]
maps $\mathbb Q[M^2]$ to $(2M-3)\mathbb Q[M(M-3)]$.
Indeed its image is anti-invariant under $M\mapsto3-M$.
For the basis element $M^{2k}$, division of its image by $2M-3$
gives a polynomial of degree $k$ in $M(M-3)$, with leading
coefficient $3k+1$. Thus $\mathcal L'$ is an isomorphism of the
filtered vector spaces over $\mathbb Q$.
The right side of the displayed recurrence is anti-invariant when
the preceding polynomials are even. It therefore has a unique even
preimage. Its degree is bounded by
\[
 \deg R_j\le\max\{\deg R_{j-1},\deg R_{j-2}+2,
                                    \deg R_{j-3}+4\}.
\]
The given initial degrees imply the asserted bound.
At $M=1$ the recurrence reduces to
\[
 R_j(1)=2R_{j-1}(1)+R_{j-2}(1)-2R_{j-3}(1),
\]
so all these values are one. This agrees with the exact base
$F_0(z)=-1/[2(1-z)]$, whose every positive moment equals $-1/2$.
Finally the Laurent recurrence determines the moments uniquely
successively in the cutoff and the moment index, since its leading
cubic coefficient is nonzero for $n\ge1$. Hence the constructed
polynomials give the actual moments for all $n$ and $j$.
\end{proof}
The diagonal $3k+1$ need not be a $p$-adic unit, so this proposition
asserts rational-polynomial structure and a uniform degree bound;
it does not assert coefficient integrality or integer-valuedness.

\section{A uniform logarithmic loss for all companion moments}\label{loc-VII-7}
\label{VII-sec:E7companionlog}
Use the polynomials $R_j$ of Proposition~\ref{VII-prop:e-companion-polynomials},
so that
\[
 \widetilde M_{j+1}(3n+1)=C_nR_j(3n+1),\qquad
 d_j:=2\left\lfloor\frac{2j+1}{3}\right\rfloor\ge\deg R_j.
\]
For $p>3$, $p\equiv2\pmod3$, define
\[
 a_j=\left\lfloor\log_p(3d_j+1)\right\rfloor,
 \qquad \eta_j=a_j+\left\lceil\frac{a_j}{2}\right\rceil
 =\left\lceil\frac{3a_j}{2}\right\rceil.
\]
\begin{theorem}\label{VII-E7:companion-log}
For every $j\ge0$ and every $x\in\mathbb Z_p$,
\[
 v_p(R_j(x))\ge-\eta_j.
\]
In particular, if $M=p^{2s}$, then for every $j\ge1$,
\[
 v_p(\widetilde M_j(M))\ge7s-e_p(j),\qquad
 e_p(j)=\left\lceil\frac32\left\lfloor\log_p(4j+1)\right\rfloor\right\rceil.
\]
\end{theorem}
\begin{proof}
We first prove a finite-value formula with integral coefficients before
normalization. Write $w=1/z$. For each integer root $a$, reversing the
factor $1-z/a$ gives its leading monomial times $1-aw$.
Thus the reversed products in $D_n/Q_n$ and in each $K_k$ have integer
coefficients as formal series in $w$: their denominator factors all have
constant term one and integral coefficients. The additional factor
$(z+18k^2-6k-2)/(z-2)$ becomes
\[
 \frac{1+(18k^2-6k-2)w}{1-2w}\in\mathbb Z[[w]].
\]
The leading coefficient of $K_k$ is $(-1)^k/d_k^2$, where
$d_k=(-1/3)_k/k!$. Also
\[
 \frac{D_n(z)}{Q_n(z)}
 =-\nu_nw\,(1+O(w)),\qquad
 \nu_n=(-1)^nM c_n^2,
 \quad M=3n+1,\quad c_n=(2/3)_n/n!.
\]
It follows from the proved summand formula for $V_n$ that there exist
integers $\phi_j(k,n)$, for every $j\ge0$, such that
\[
 F_n(z)=-\frac{\nu_n}{2}\sum_{j\ge0}
 \left(\sum_{k=0}^n(6k-1)d_k\phi_j(k,n)\right)w^{j+1}.
\]
Here $\phi_j$ is simply the coefficient of the product of the reversed
factors just described, so integrality holds for every degree, not only
for a finite initial range. Since $C_n=-\nu_n M c_n/2$, comparison of
Laurent coefficients proves the exact formula
\begin{equation}\label{VII-E7:companion-integral-values}
 R_j(3n+1)=-\frac1{M c_n}
 \sum_{k=0}^n(6k-1)d_k\phi_j(k,n).
\end{equation}

Both $d_k$ and $c_k$ belong to $\mathbb Z_p$: for any
$\alpha\in\mathbb Z_p$, $(\alpha)_k/k!=(-1)^k\binom{-\alpha}{k}$
is integral. Put $b=\lfloor\log_p M\rfloor$ and $t=v_p(M)$.
At the level $q=p^\ell$, the excess numerator count in $c_n$
is zero or one; it is zero when $q>M$.
For $\ell\le t$, the residue of $n=(M-1)/3$ is
$(q-1)/3$ when $\ell$ is even, and $(2q-1)/3$ when $\ell$ is odd.
The first numerator index for $3k+2\equiv0\pmod q$ is respectively
$(2q-2)/3$ and $(q-2)/3$. Hence exactly the odd levels among
$1,\ldots,t$ contribute an excess. The remaining $b-t$ levels each
contribute at most one. Consequently
\[
 v_p(Mc_n)\le t+\lceil t/2\rceil+(b-t)
 \le b+\lceil b/2\rceil.
\]
Equation~\eqref{VII-E7:companion-integral-values} therefore bounds every
finite value $R_j(3n+1)$ by the negative of this quantity.

Apply this bound at $n=0,1,\ldots,d_j$. All these values belong to
$p^{-\eta_j}\mathbb Z_p$. The polynomial $T_j(X)=R_j(3X+1)$ has
exact Newton expansion
\[
 T_j(X)=\sum_{k=0}^{d_j}\Delta^kT_j(0)\binom Xk.
\]
Each finite difference is an integer linear combination of the sampled
values. Each $\binom Xk$ is integral on $\mathbb Z_p$, by continuity
from the nonnegative integers. Taking $X=(x-1)/3$ proves the uniform
bound for $R_j(x)$.

Finally, at $M=p^{2s}$ the excess-count argument gives
$v_p(c_n)=s$ exactly, so $v_p(C_n)=7s$.
For the moment with index $j$, its polynomial is $R_{j-1}$, and
\[
 3d_{j-1}+1
 =6\left\lfloor\frac{2j-1}{3}\right\rfloor+1
 \le4j-1\le4j+1.
\]
Thus $\eta_{j-1}\le e_p(j)$, proving the last assertion.
\end{proof}


\section{A uniform lower bound for every companion interpolation mass}\label{loc-VII-8}
\label{VII-sec:E7companionmass}
Let $p>3$, $p\equiv2\pmod3$, $s\ge1$, $M=p^{2s}$ and
$n=(M-1)/3$. Retain the companion grids, denominator roots, and models
from Theorem~\ref{VII-E7:companion-exact}, and put
\[
 \mu_M(H)=\widetilde W_n(H)U(H).
\]
\begin{proposition}\label{VII-E7:companion-mass-bound}
For every companion node $H$, with $t=v_p(H)$,
\[
 \boxed{\quad v_p(\mu_M(H))\ge3s-\lceil t/2\rceil\ge2s.\quad}
\]
Zero masses have valuation $+\infty$. In particular every unit-node
mass belongs to $p^{3s}\mathbb Z_p$.
\end{proposition}
\begin{proof}
We may suppose $U(H)\ne0$ and $H\notin\mathcal D_n$; otherwise the
mass is zero. At an active node, $U(H)/H$ is one of
$1,1/2,-1/2$, hence is a unit. Direct multiplication of the grids gives
\[
 \frac{\prod_{K\in\mathcal H_n}K}{\prod_{D\in\mathcal D_n}D}
 =(-1)^nM\left(\frac{(2/3)_n}{n!}\right)^2.
\]
The Pochhammer ratio has valuation $s$: at each odd level
$p^\ell$, $1\le\ell\le2s$, its numerator has one extra divisible
factor, and at each even level its numerator and denominator counts
agree. Thus the displayed normalization has valuation $4s$.
The exact mass, up to an irrelevant sign, is this normalization times
\[
 \frac{U(H)}H
 \frac{\prod_{D\in\mathcal D_n}(D-H)}
 {\prod_{K\in\mathcal H_n,K\ne H}(K-H)}.
\]
All nonzero differences here have absolute value less than $4M$.
Since $p\ge5$, no power $p^\ell$ with $\ell>2s$ divides any of them.

For $q=p^\ell$ and $x=H\bmod q$, define
\[
 \delta_q(x)=1+\#\{D\in\mathcal D_n:D\equiv x\pmod q\}
 -\#\{K\in\mathcal H_n:K\equiv x\pmod q\}.
\]
The initial one removes the node $K=H$ from the denominator count.
Consequently
\begin{equation}\label{VII-E7:mass-count}
 v_p(\mu_M(H))=4s+\sum_{\ell=1}^{2s}\delta_{p^\ell}(H\bmod p^\ell).
\end{equation}
All three arithmetic progressions have steps prime to $q$.
Their complete blocks of $q$ consecutive indices contribute equally
to the numerator and denominator counts and cancel. If $\ell$ is
even, the remaining cutoff is $b=(q-1)/3$; if $\ell$ is odd,
it is $b=(2q-1)/3$. The residual grids and denominators are exactly
the companion sets with $n=b$, respectively $M'=q$ or $M'=2q$.

Here are the complete residual counts for nonzero residues
$1\le x<q$, listed by $x\bmod6$:
\[
\begin{array}{c|rrrrrr}
x\bmod6&0&1&2&3&4&5\\\hline
\#K\ (q\equiv1\pmod6)&0&1&1&1&2&1\\
\#D\ (q\equiv1\pmod6)&2&1&0&1&1&1\\
\delta_q\ (q\equiv1\pmod6)&3&1&0&1&0&1\\\hline
\#K\ (q\equiv5\pmod6)&2&2&2&1&2&3\\
\#D\ (q\equiv5\pmod6)&2&3&2&2&2&1\\
\delta_q\ (q\equiv5\pmod6)&1&2&1&2&1&-1
\end{array}
\]
These counts follow directly by solving the three progression
congruences in the residual ranges. For example, in the even case
$M'=q$, the negative node grid contributes once for residues
$x\equiv3,4\pmod6$, while the positive denominator grid contributes
once for $x\equiv0,5\pmod6$. In the odd case $M'=2q$, those same
contributions are respectively $(1,1,0,0,1,1)$ and
$(1,1,1,1,0,0)$ across the six residue classes. The other two grids
have step three, yielding the displayed totals.
At the zero residue, the even case has one node and no denominator
root, so $\delta_q(0)=0$; the odd case has three nodes and no
denominator root, so $\delta_q(0)=-2$.

Thus every even level contributes a nonnegative amount. Each odd
level contributes at least $-1$, with a possible extra loss of one
precisely when $p^\ell\mid H$. There are $s$ odd levels altogether
and $\lceil t/2\rceil$ of them divide $H$; here $t\le2s$ because
$0<|H|<2M<p^{2s+1}$. Equation~\eqref{VII-E7:mass-count} proves
\[
 v_p(\mu_M(H))\ge4s-s-\lceil t/2\rceil,
\]
as asserted.
\end{proof}

This is an individual-mass estimate. It is independent of the unresolved
cancellation in E.3 and can be combined with exact moment formulas or
polynomial projectors without assuming any positive-moment integrality.

\subsection{The corresponding primary mass bound}
Let $L=p^{2s+1}$, $m=(L-2)/3$, and let $\mu_L(H)$ be the primary
three-grid mass. Then
\begin{equation}\label{VII-E7:primary-mass-bound}
 \boxed{v_p(\mu_L(H))\ge3s+1-\lfloor v_p(H)/2\rfloor\ge2s+1.}
\end{equation}
Again zero masses are harmless. We give the full analogous count,
so this bound does not rely on numerical interpolation tests.
The primary normalization is
\[
 \frac{\prod K}{\prod D}
 =(-1)^{m+1}2L(L-1)^2\left(\frac{(1/3)_m}{m!}\right)^2.
\]
Its valuation is $4s+1$, since the Pochhammer ratio has valuation $s$.
At an active node $T(H)/H$ is a unit, and all nonzero differences
have absolute value less than $4L<pL$.
The same formula $1+\#D_x-\#K_x$ gives each level's correction.
For odd $\ell$, the residual cutoff is $(q-2)/3$ with $q=p^\ell$;
for even $\ell$, it is $(2q-2)/3$. The corresponding residual primary
endpoint is $q$ or $2q$ respectively. Complete blocks cancel as before.
For $1\le x<q$, the exact tables are
\[
\begin{array}{c|rrrrrr}
x\bmod6&0&1&2&3&4&5\\\hline
\#K\ (\ell\text{ odd})&0&1&2&1&1&1\\
\#D\ (\ell\text{ odd})&2&1&1&1&0&1\\
\delta_q\ (\ell\text{ odd})&3&1&0&1&0&1\\\hline
\#K\ (\ell\text{ even})&2&3&2&1&2&2\\
\#D\ (\ell\text{ even})&2&1&2&2&2&3\\
\delta_q\ (\ell\text{ even})&1&-1&1&2&1&2
\end{array}
\]
For example, at an odd level the negative node grid contributes
in residue classes $2,3\pmod6$, while the positive denominator grid
contributes in classes $0,1\pmod6$. At an even level their respective
six-class contributions are $(1,1,1,0,0,1)$ and $(1,0,0,1,1,1)$.
The remaining step-three grids give the totals shown above.
At $x=0$ the correction is zero for odd $\ell$ and $-2$ for even
$\ell$: there are respectively one or three nodes and no denominator
root at that residue. Thus odd levels cost nothing, while each of the
$s$ even levels costs at most one, plus one more when it divides $H$.
There are $\lfloor v_p(H)/2\rfloor$ such divisible even levels.
Subtracting these losses from $4s+1$ proves
\eqref{VII-E7:primary-mass-bound}. In particular unit-node masses belong
to $p^{3s+1}\mathbb Z_p$.

\subsection{The unit-node contribution at level three}
\begin{corollary}\label{VII-E7:primary-unit-r3}
For every prime $p>3$ with $p\equiv2\pmod3$,
\[
 \sum_{H\in\mathcal H_{(p^3-2)/3},\ p\nmid H}\mu_{p^3}(H)
 \in p^8\mathbb Z_p.
\]
\end{corollary}
\begin{proof}
First assume $p\ge41$. Let $g(Y)=3Y^2-2Y^3$ and $\Pi(X)=g(g(X^{p-1}))$. This polynomial
has integer coefficients, no terms of degree below $4(p-1)$, and
degree $9(p-1)$. For a unit integer $H$, Fermat's congruence and
$g(1+u)=1-3u^2-2u^3$ give $\Pi(H)\equiv1\pmod{p^4}$.
For $p\mid H$, one has $\Pi(H)\in p^{4(p-1)}\mathbb Z_p$.
The preceding proposition, with $s=1$, therefore gives
\[
 \sum_{p\nmid H}\mu_{p^3}(H)
 \equiv\sum_H\mu_{p^3}(H)\Pi(H)\pmod{p^8}.
\]
The right-hand side is an integer linear combination of the moments
$M_j(p^3)$ with $4(p-1)\le j\le9(p-1)$. Their common normalization
$C(p^3)$ has valuation nine. For every such $j$, the logarithmic
moment bound in Theorem~\ref{VII-thm:e-log-bound} has loss at most one:
indeed, putting $n=j-3$,
\[
 3d_n+2=6\lfloor2n/3\rfloor+2\le4j-10
 \le36(p-1)-10<p^2 \qquad(p\ge41).
\]
Consequently $M_j(p^3)\in p^8\mathbb Z_p$.

The remaining primes are exactly $5,11,17,23,29$. They can be checked
by finite integer arithmetic in the exact mass formula above. Here is
an explicit modular certificate, computed modulo $p^{10}$:
\[
\begin{array}{c|rrrrr}
p&5&11&17&23&29\\\hline
p^{-9}\sum_{p\nmid H}\mu_{p^3}(H)\pmod p&1&6&6&18&10.
\end{array}
\]
In particular every one of those sums lies in $p^9\mathbb Z_p$.
The reproducible certificate is \nolinkurl{fast_primary_masses.py}:
it evaluates the displayed product formula modulo $p^{10}$,
removing the exact power of $p$ from every factor before taking a
modular inverse. Prefix products in each residue class modulo three
and six evaluate every arithmetic progression in the formula. A zero
factor is omitted only at the unique denominator occurrence $K=H$;
numerator zeroes are excluded by testing $H\in\mathcal D_m$ first.
Thus all modular inverses are applied to units and the finite checks
are exact. They complete the five exceptional cases.
\end{proof}

The corollary controls only the unit-node contribution. The shell
$v_p(H)=1$ and its cancellation with the cubic double-strip term
remain separate obligations.

\section{One-step transport of the valuation-one stratum}\label{loc-VII-9}
\label{VII-sec:E7one-step}
Let $L=p^{2s+1}$, $M=L/p=p^{2s}$, $p>3$, $p\equiv2\pmod3$.
Write $\mu_L(H)$ for the primary interpolation masses and
$\widetilde\mu_M(J)$ for the proved companion masses.
Let $\mathcal H_L,\mathcal D_L$ denote the primary root sets and set
\[
 R^{(1)}_{L,p}(X)=
 \frac{\prod_{d\in\mathcal D_L,\,p\nmid d}(1-pX/d)}
 {\prod_{h\in\mathcal H_L,\,p\nmid h}(1-pX/h)}
 =\sum_{j\ge0}a_jX^j.
\]
Every retained root is a unit, so $a_0=1$ and $a_j\in p^j\mathbb Z_p$.
The series is restricted and converges at every integral node.

The exact one-step product identity is
\begin{equation}\label{VII-E7:one-step-masses}
 \mu_L(pJ)=p\widetilde\mu_M(J)R^{(1)}_{L,p}(J).
\end{equation}
Indeed dividing the divisible primary nodes by $p$ gives respectively
\[
 \{2,5,\ldots,M-2\},\quad
 \{1,4,\ldots,M\},\quad
 -2\{2,5,\ldots,M-2\},
\]
exactly the companion nodes. Dividing the divisible denominator roots
gives $6j,-3j$ for $1\le j\le(M-1)/3$ and
$-1,-4,\ldots,-(M-3)$, exactly the companion denominator roots.
For the first two denominator grids this follows from
$\lfloor(L-2)/(3p)\rfloor=(M-1)/3$; the last grid has fine endpoint
$-(L-3)$, so its largest divisible quotient is $-(M-3)$.
Since $p\equiv5\pmod6$, the primary model formulas give
$T(pJ)=pU(J)$ on each companion grid, including zero model values.
Separating the divisible and unit factors in the interpolation products
proves \eqref{VII-E7:one-step-masses}; no division by a zero weight is used.

In particular, the omitted part of the two-step transport has the exact
split
\begin{equation}\label{VII-E7:U-one-step}
 U_{L,p}
 =\sum_{p\nmid H}\mu_L(H)
 +p\sum_{p\nmid J}\widetilde\mu_M(J)R^{(1)}_{L,p}(J).
\end{equation}
This identity places the formerly correlated valuation-one stratum
inside a weighted unit sum of the exact companion family.


\section{The original odd E.3 congruence at all sufficiently high levels}\label{loc-VII-10}
\label{VII-sec:e-high-level-closure}
Let $p>3$ be a prime with $p\equiv2\pmod3$, and define
\[
 E(L)=\sum_{k=0}^{(L-2)/3}(-1)^k(6k+1)
          \left(\frac{(1/3)_k}{k!}\right)^3
 \qquad(L\equiv2\pmod3).
\]
All interpolation identities and moment bounds used below are exact
or have been proved in the preceding sections. In particular, no
assertion of integral normalized moments of every degree is required.

\begin{theorem}\label{VII-thm:e-high-level-closure}
Put $r=2s+1$, with $s\ge2$, and suppose
\begin{equation}\label{VII-eq:e-improved-threshold}
 \left\lceil\frac32\left\lfloor
 \log_p\bigl(4(3s+1)(p-1)+1\bigr)\right\rfloor\right\rceil
 \le s-1.
\end{equation}
Then the original E.3 congruence holds:
\[
 \boxed{E(p^{2s+1})\equiv p^2E(p^{2s-1})\pmod{p^{6s+2}}.}
\]
In particular it holds in each of the following ranges:
\begin{enumerate}
 \item every $p>3$, $p\equiv2\pmod3$, and every odd $r\ge13$;
 \item every such $p\ge11$ and every odd $r\ge9$;
 \item every such $p\ge41$ and every odd $r\ge7$.
\end{enumerate}
\end{theorem}

\subsection{The two exact one-step identities}
Set $N=p^{2s+1}$, $M=p^{2s}$, and $K=p^{2s-1}$.
Write $\mu_N(H)$ and $\mu_K(H)$ for the primary three-grid masses,
and $\widetilde\mu_M(J)$ for the companion masses. Their sums are
$E(N)$, $E(K)$, and $\mathcal A(M)$, respectively. Their positive
moments are denoted by $M_j$ and $\widetilde M_j$.

For either grid family, form the quotient of unit denominator-root
factors by unit node factors, evaluated at $pX$:
\[
 R(X)=
 \frac{\prod_{d\in\mathcal D,\ p\nmid d}(1-pX/d)}
      {\prod_{h\in\mathcal H,\ p\nmid h}(1-pX/h)}.
\]
Let $R_0$ use the primary family at $N$, and $R_1$ the companion
family at $M$. Both have expansions
\begin{equation}\label{VII-eq:e-kernel-coefficient-bound}
 R_i(X)=1+\sum_{j\ge1}a_{i,j}X^j,
 \qquad a_{i,j}\in p^j\mathbb Z_p\quad(i=0,1).
\end{equation}
These are restricted power series and converge at every $X\in\mathbb Z_p$.
Separating divisible and unit roots in the interpolation weights gives
\begin{equation}\label{VII-eq:e-two-one-step-identities}
 \mu_N(pJ)=p\widetilde\mu_M(J)R_0(J),\qquad
 \widetilde\mu_M(pH)=p\mu_K(H)R_1(H).
\end{equation}
Here the correspondence between divisible nodes and the lower family
is exact, including endpoints and zero model values.

For completeness, the first correspondence is the one established in
\eqref{VII-E7:one-step-masses}. For the reverse correspondence, divide the
divisible companion nodes by $p$. The positive grids become
$\{2,5,\ldots,K\}$ and $\{1,4,\ldots,K-1\}$, and the negative grid
becomes $-2\{1,4,\ldots,K-1\}$, precisely the primary grids at $K$.
The divided denominator roots are
\[
 \{6j,-3j:1\le j\le(K-2)/3\}
 \ \cup\ \{-2,-5,\ldots,-(K-3)\},
\]
also precisely the primary denominator roots. The explicit models
satisfy $U(pH)=pT(H)$ on these grids because $p\equiv5\pmod6$.
Thus the same product separation proves the second identity in
\eqref{VII-eq:e-two-one-step-identities}, without dividing by a model
value that might vanish.

\subsection{Moment estimates in the forms needed for transport}
The proved primary normalizations and logarithmic bounds give
\begin{align}
 M_1(N)=M_2(N)&=M_1(K)=M_2(K)=0,\\
 v_p(M_j(N))&\ge7s+2-e_{j-3} &&(j\ge3),\\
 v_p(M_j(K))&\ge7s-5-e_{j-3} &&(j\ge3),
 \label{VII-eq:e-coarse-moment-bound}
\end{align}
where $e_0=0$, $e_h\le h-1$ for $h\ge1$, and $e_h=O(\log(h+2))$.
The companion logarithmic bound gives
\[
 v_p(\widetilde M_j(M))\ge7s-\eta_{j-1}.
\]
In particular,
\begin{equation}\label{VII-eq:e-companion-linear-loss}
 v_p(\widetilde M_j(M))\ge7s-(j-1)\qquad(j\ge1).
\end{equation}
For $j=1,2$ this follows from the exact normalized polynomials
$R_0^{\rm mom}=1$ and $R_1^{\rm mom}=(M^2+1)/2$. For the other
indices use $\eta_h\le h$ for $h\ge2$. To verify this last elementary
inequality, the companion definition is
\[
 \eta_h=\left\lceil\frac32
 \left\lfloor\log_p\left(6\left\lfloor\frac{2h+1}{3}\right\rfloor+1\right)
 \right\rfloor\right\rceil.
\]
The inequality
\[
 6\left\lfloor\frac{2h+1}{3}\right\rfloor+1
 <5^{\lfloor2h/3\rfloor+1}
\]
holds at $h=2,3,4$ and is preserved by increasing $h$ by three:
the left side increases by $12$, whereas the right side is multiplied
by $25$. Since $p\ge5$, it implies $\eta_h\le h$.

A common nondecreasing logarithmic bound is
\begin{equation}\label{VII-eq:e-common-loss}
 \epsilon_p(j)=\left\lceil\frac32\left\lfloor
                       \log_p(4j+1)\right\rfloor\right\rceil.
\end{equation}
For $j\ge3$ it bounds $e_{j-3}$, and for $j\ge1$ it bounds
$\eta_{j-1}$. Consequently the positive moments of the two fine
families satisfy
\[
 v_p(M_j(N))\ge7s+2-\epsilon_p(j),\qquad
 v_p(\widetilde M_j(M))\ge7s-\epsilon_p(j).
\]
The first inequality also holds for $j=1,2$, whose moments vanish.
The proved individual unit-node estimates are
\begin{equation}\label{VII-eq:e-unit-mass-bounds}
 v_p(\mu_N(H))\ge3s+1\ (p\nmid H),\qquad
 v_p(\widetilde\mu_M(J))\ge3s\ (p\nmid J).
\end{equation}

\subsection{A low-degree unit projector and its leakage}
Put
\[
 t=3s+1,\qquad D=t(p-1),\qquad
 \chi_t(X)=1-(1-X^{p-1})^t\in\mathbb Z[X].
\]
Its nonzero monomials have degrees $i(p-1)$ for $1\le i\le t$,
so its smallest degree is $p-1\ge4$ and its degree is $D$.
For every integral unit $x$, Fermat's congruence gives
$1-\chi_t(x)\in p^t\mathbb Z_p$.
Let $B=\epsilon_p(D)$; the theorem assumes $B\le s-1$.

First consider the primary unit sum
$U_N=\sum_{p\nmid H}\mu_N(H)$.
Replacing $1$ by $\chi_t(H)$ on its unit nodes changes it by an
element of $p^{6s+2}\mathbb Z_p$, by
\eqref{VII-eq:e-unit-mass-bounds}. The full-node polynomial sum satisfies
\[
 \sum_H\mu_N(H)\chi_t(H)\in p^{7s+2-B}\mathbb Z_p
                              \subseteq p^{6s+3}\mathbb Z_p.
\]
Its contribution from nonunit nodes is, by the first exact strip,
\[
 \sum_{p\mid H}\mu_N(H)\chi_t(H)
   =p\sum_J\widetilde\mu_M(J)R_0(J)\chi_t(pJ).
\]
Write $R_0(X)\chi_t(pX)=\sum_{j\ge p-1}b_jX^j$.
Every $b_j$ belongs to $p^j\mathbb Z_p$: both factors have this
coefficient property, and their convolution preserves it. Hence
\eqref{VII-eq:e-companion-linear-loss} shows that each term of its
moment expansion has valuation at least
\[
 1+j+7s-(j-1)=7s+2.
\]
The expansion converges, since the sharper logarithmic bound makes
these valuations tend to infinity. We have proved
\begin{equation}\label{VII-eq:e-primary-unit-sum}
 U_N\in p^{6s+2}\mathbb Z_p.
\end{equation}

Apply the same projector to the companion unit sum
$\widetilde U_M=\sum_{p\nmid J}\widetilde\mu_M(J)$.
The approximation error on unit nodes belongs to
$p^{3s+t}\mathbb Z_p=p^{6s+1}\mathbb Z_p$.
The full-node polynomial sum has valuation at least
$7s-B\ge6s+1$. Its nonunit leakage is
\[
 p\sum_H\mu_K(H)R_1(H)\chi_t(pH).
\]
Again the product series starts in degree $p-1\ge4$ and its
degree-$j$ coefficient has valuation at least $j$.
For all such $j$, $e_{j-3}\le j-4$. Therefore each leakage term has
valuation at least
\[
 1+j+(7s-5)-(j-4)=7s\ge6s+1.
\]
Convergence follows from the logarithmic loss as before. Thus
\begin{equation}\label{VII-eq:e-companion-unit-sum}
 \widetilde U_M\in p^{6s+1}\mathbb Z_p.
\end{equation}
This calculation explicitly controls the nonunit leakage; the projector
is not being treated as a literal indicator at nonunit nodes.

\subsection{Completion of the two one-step congruences}
Sum the first identity in \eqref{VII-eq:e-two-one-step-identities} and
add the primary unit nodes. It gives exactly
\[
 E(N)=U_N+p\mathcal A(M)
             +p\sum_{j\ge1}a_{0,j}\widetilde M_j(M).
\]
Every term of the last series has valuation at least $7s+2$ by
\eqref{VII-eq:e-kernel-coefficient-bound} and
\eqref{VII-eq:e-companion-linear-loss}. Thus
\begin{equation}\label{VII-eq:e-primary-one-step-congruence}
 E(N)\equiv p\mathcal A(M)\pmod{p^{6s+2}}.
\end{equation}
The reverse strip similarly gives
\[
 \mathcal A(M)=\widetilde U_M+pE(K)
                         +p\sum_{j\ge3}a_{1,j}M_j(K),
\]
where the first two moments have vanished exactly. The $j=3$ term
has valuation at least $7s-1$; every $j\ge4$ term has valuation at
least $7s$. Since $s\ge2$, these are at least $6s+1$.
Together with \eqref{VII-eq:e-companion-unit-sum}, this proves
\begin{equation}\label{VII-eq:e-companion-one-step-congruence}
 \mathcal A(M)\equiv pE(K)\pmod{p^{6s+1}}.
\end{equation}
All series manipulations are legitimate: there are finitely many
nodes, the root-product series converge on $\mathbb Z_p$, and the
logarithmic moment estimates also give convergence after rearranging
by degree. Combining the two displayed congruences proves the
assertion of Theorem~\ref{VII-thm:e-high-level-closure}.

\subsection{Verification of the explicit ranges}
For integer $s\ge2$, condition \eqref{VII-eq:e-improved-threshold} is
equivalent to the strict inequality
\begin{equation}\label{VII-eq:e-threshold-power-form}
 (12s+4)(p-1)+1
 <p^{\lfloor2(s-1)/3\rfloor+1}.
\end{equation}
For fixed $s$ let $A=12s+4$ and $a=\lfloor2(s-1)/3\rfloor+1$.
In each range considered below, $a\ge2$. The function
$(Ap-A+1)/p^a$ is decreasing for real $p\ge5$, because its
derivative has numerator $Ap(1-a)+a(A-1)<0$. It therefore suffices
to check the smallest value of $p$ in each claimed range.

For $p\ge5$, the left side at $p=5$ is $48s+17$.
At $s=6,7,8$ the required comparisons are
\[
 305<625,\qquad353<3125,\qquad401<3125.
\]
Increasing $s$ by three adds $144$ to the left side and multiplies
the right side by $25$, preserving the inequality. Thus it holds
for all $s\ge6$.

For $p\ge11$, the left side at $p=11$ is $120s+41$.
At $s=4,5,6$ the comparisons are
\[
 521<1331,\qquad641<1331,\qquad761<14641.
\]
The step $s\mapsto s+3$ adds $360$ to the left side and multiplies
the right side by $121$. Thus the inequality holds for all $s\ge4$.

For $p\ge41$, the left side at $p=41$ is $480s+161$.
At $s=3,4,5$ the comparisons are
\[
 1601<1681,\qquad2081<68921,\qquad2561<68921.
\]
The step $s\mapsto s+3$ adds $1440$ to the left side and multiplies
the right side by $1681$. This proves the final stated range.


\section{A reflection gain for companion interpolation samples}\label{loc-VII-11}
\label{VII-sec:e-companion-sample-gain}
Retain $d_k=(-1/3)_k/k!$, $c_n=(2/3)_n/n!$, $M=3n+1$,
and the exact normalized companion moments $R_j(M)$.

\begin{lemma}\label{VII-lem:e-companion-numerator-gain}
Let $p>3$, $p\equiv2\pmod3$. For every $n\ge0$ with
$p\mid3n+1$, every coefficient of the integral numerator series
\[
 \mathcal N_n(w)=\sum_{k=0}^n(6k-1)d_k\phi_k(w),
 \qquad \phi_k(w)=\sum_{j\ge0}\phi_j(k,n)w^j,
\]
from \eqref{VII-E7:companion-integral-values} belongs to $p\mathbb Z_p$.
\end{lemma}
\begin{proof}
Use the monic reversed product of the companion summand kernel:
\[
 \widehat K_k(w)=\prod_{j=0}^{k-1} f_j(w),\qquad
 f_j(w)=
 \frac{(1+(6j-2)w)(1-(3j-1)w)(1-(3j+1)w)}
 {(1-(6j+6)w)(1+(3j+3)w)(1+(3j+1)w)}.
\]
The coefficient series $\phi_k$ has the form
\[
 \phi_k(w)=A_n(w)\widehat K_k(w)
                \frac{1+(18k^2-6k-2)w}{1-2w},
 \qquad A_n(w)\in\mathbb Z[[w]],\quad A_n(0)=1.
\]
Here $A_n$ is the common reversed denominator-root/node-root
quotient. Its precise roots are immaterial for this argument;
every root is integral, so $A_n$ is an integral unit series.
It suffices to prove that
\[
 \sum_{k=0}^n(6k-1)d_k\widehat K_k(w)
                       (1+(18k^2-6k-2)w)
\]
vanishes in $\mathbb F_p[[w]]$.

Put $a=(p+1)/3$, which is an even integer. The binomial identity
$d_k=(-1)^k\binom{1/3}{k}$ and the $p$-adic Lucas congruence give
\begin{equation}\label{VII-eq:e-companion-lucas}
 d_{pb+c}\equiv
 \frac{(1/3)_b}{b!}\,(-1)^c\binom ac\pmod p
 \qquad(0\le c<p).
\end{equation}
Indeed $1/3=a+p(-1/3)$ in $\mathbb Z_p$; the ordinary Lucas
identity extends to binomial coefficients with $p$-adic top argument
by continuity from nonnegative integers. This also proves that the
left side vanishes when $c>a$.

For any consecutive $p$ indices, each of the three linear root
progressions in the numerator and denominator of $f_j$ runs through
all residues of $\mathbb F_p$. Thus their products agree, and
\[
 \prod_{j=b p}^{(b+1)p-1}f_j(w)=1
       \quad\hbox{in }\mathbb F_p[[w]],
 \qquad \widehat K_{pb+c}(w)=\widehat K_c(w).
\]
The remaining linear and quadratic factors in $k$ depend only on
$k\bmod p$. Therefore every complete block reduces, up to the scalar
$(1/3)_b/b!$, to
\[
 B(w)=\sum_{c=0}^a(6c-1)(-1)^c\binom ac\widehat K_c(w)
                       (1+(18c^2-6c-2)w).
\]
The same holds for the final partial block: since $3n+1\equiv0$
modulo $p$, its last residue is $(2p-1)/3=2a-1\ge a$ and
therefore includes the entire nonzero Lucas support.

It remains to show $B(w)=0$. The exact root lists defining $f_j$
give, using $3a=1$ in $\mathbb F_p$,
\[
 f_{a-1-j}(w)=f_j(w)^{-1},\qquad
 \widehat K_a(w)=1,\qquad
 \widehat K_{a-c}(w)=\widehat K_c(w).
\]
The first equality pairs the numerator roots
$2-6(a-1-j),3(a-1-j)-1,3(a-1-j)+1$
with the denominator roots $6j+6,-3j-3,-3j-1$, and conversely.
Because $a$ is even, $(-1)^c\binom ac$ is unchanged under
$c\mapsto a-c$. Also
\[
 6(a-c)-1=-(6c-1),\qquad
 18(a-c)^2-6(a-c)-2=18c^2-6c-2
 \quad\hbox{in }\mathbb F_p.
\]
Consequently the summands cancel in reflected pairs. If there is a
middle summand, its factor $6c-1$ is zero. Hence $B(w)=0$, proving
the claimed coefficientwise divisibility.
\end{proof}

\begin{corollary}\label{VII-cor:e-companion-improved-log}
With the notation of Theorem~\ref{VII-E7:companion-log}, the bound improves to
\[
 v_p(R_j(x))\ge-\max\{0,\eta_j-1\}
 \qquad(j\ge0,\ x\in\mathbb Z_p).
\]
\end{corollary}
\begin{proof}
At a sample $M=3n+1$, put $b=\lfloor\log_p M\rfloor$.
If $p\nmid M$, its normalizing denominator $Mc_n$ has valuation
at most $b$. If $p\mid M$, the previous count and the new numerator
gain give a loss of at most $b+\lceil b/2\rceil-1$.
For $b\ge1$ the latter dominates $b$; for $b=0$ there is no loss.
Apply these bounds to the same Newton samples as in
Theorem~\ref{VII-E7:companion-log}.
\end{proof}

\begin{corollary}\label{VII-cor:e-companion-short-sample}
For every integer $n\ge0$ with $1\le M=3n+1<p^2$, and every $j\ge0$,
\[
 v_p(R_j(M))\ge-1.
\]
Consequently, if $3d_j+1<p^2$, where
$d_j=2\lfloor(2j+1)/3\rfloor$, then
\[
 pR_j(x)\in\mathbb Z_p\qquad\hbox{for every }x\in\mathbb Z_p.
\]
\end{corollary}
\begin{proof}
Use the exact formula
$R_j(M)=-(Mc_n)^{-1}[w^j]\mathcal N_n(w)$.
If $p\nmid M$, the excess-count proof gives $v_p(c_n)\le1$
since $M<p^2$. Thus its normalization costs at most one power of $p$.
If $p\mid M$, then $v_p(M)=v_p(c_n)=1$, so the normalizing
denominator costs two powers, but
Lemma~\ref{VII-lem:e-companion-numerator-gain} restores one.
This proves the first assertion, including $M=1$.
For the second, apply it at the Newton interpolation samples
$M=3n+1$, $0\le n\le d_j$. Their maximum is below $p^2$.
All their finite differences belong to $p^{-1}\mathbb Z_p$,
and the binomial basis is integral on $\mathbb Z_p$.
\end{proof}

\begin{theorem}\label{VII-thm:e-level-five-large-primes}
For every prime $p\ge29$ with $p\equiv2\pmod3$,
\[
 \boxed{E(p^5)\equiv p^2E(p^3)\pmod{p^{14}}.}
\]
\end{theorem}
\begin{proof}
Repeat the proof of Theorem~\ref{VII-thm:e-high-level-closure} with
$s=2$, $t=7$, and $D=7(p-1)$. The leakage and one-step transport
estimates in that proof are valid for every $s\ge2$ and require
no logarithmic threshold. Only its two full-node projector sums
need the following improved estimates.

For each companion moment index $1\le j\le D$, its polynomial
is $R_{j-1}$ and its largest interpolation sample satisfies
\[
 3d_{j-1}+1
 =6\left\lfloor\frac{2j-1}{3}\right\rfloor+1
 \le4j-1\le28(p-1)-1<p^2.
\]
The last inequality holds for $p\ge29$. The preceding corollary
therefore gives
$v_p(\widetilde M_j(p^4))\ge14-1=13$.
The full-node companion projector sum belongs to $p^{13}\mathbb Z_p$,
as required.

The primary moments have their first two terms zero. For $3\le j\le D$,
the largest sample for $P_{j-3}$ is at most
$4(j-3)+2\le28(p-1)-10<p^2$.
The primary logarithmic bound is therefore at most one, giving
$v_p(M_j(p^5))\ge16-1=15$. This exceeds the required primary
projector precision $14$. The unchanged leakage calculations now
give $E(p^5)\equiv p\mathcal A(p^4)\pmod{p^{14}}$ and
$\mathcal A(p^4)\equiv pE(p^3)\pmod{p^{13}}$.
Their combination proves the theorem.
\end{proof}

\section{Exact finite certificates for the remaining higher-level cases}\label{loc-VII-12}
\label{VII-sec:E7finiteexceptions}
The analytic theorem leaves only the following pairs $(s,p)$ at
odd levels $r=2s+1\ge7$:
\[
 (3,5),(3,11),(3,17),(3,23),(3,29),(4,5),(5,5).
\]
This is a finite exhaustive list, not a numerical sample of the
remaining higher-level range. We give exact polynomial-interpolation
certificates for the moments required by the proved projector argument.

Put $t=3s+1$. The projector $1-(1-X^{p-1})^t$ uses precisely the
moment indices $j=i(p-1)$, $1\le i\le t$. Its proof only needs
\begin{equation}\label{VII-E7:finite-required}
 v_p\bigl(R_{i(p-1)-1}(p^{2s})\bigr)\ge-(s-1),\qquad
 v_p\bigl(P_{i(p-1)-3}(p^{2s+1})\bigr)\ge-s.
\end{equation}
The first bound controls the full companion projector sum, and the
second controls the full primary projector sum. The bounds on unit-node
approximation, nonunit leakage, and the infinite transport corrections
remain the previously proved uniform estimates; none of those is
replaced by a finite truncation.

\subsection{The integer recurrences used for certification}
For the companion define a formal series
\[
 Y_n(w)=-M c_n\sum_{j\ge0}R_j(M)w^j,\qquad M=3n+1.
\]
The proved rational-function recurrence gives, for $n\ge1$,
\begin{align*}
 &(1-2w-(M-2)(3M-4)w^2+2M(M-2)^2w^3)Y_n(w)\\
 &\quad=(1-2w-(M-1)(3M-5)w^2
            -2(M-3)(M-1)^2w^3)Y_{n-1}(w)\\
 &\qquad\quad +(6n-1)d_n(1+w)(1+(18n^2-6n-2)w),
\end{align*}
with $Y_0(w)=-1/(1-w)$, $d_n=(-1/3)_n/n!$, and $c_n=(2/3)_n/n!$.
All polynomial coefficients are integers and the left constant term
is one. Thus each coefficient is determined using only addition and
multiplication and already known lower coefficients.

For the primary family let $L=3n+2$, $b_n=(1/3)_n/n!$, and
\[
 Z_n(w)=\frac{L(L-1)b_n}{2}\sum_{j\ge0}P_j(L)w^j.
\]
Its integer recurrence is
\begin{align*}
 &(1-w-(L-1)(3L-2)w^2+2L(L-1)^2w^3)Z_n(w)\\
 &\quad=(1-w-(L-2)(3L-7)w^2
            -2(L-2)^2(L-3)w^3)Z_{n-1}(w)+(6n+1)b_n.
\end{align*}
At $n=0$ take the preceding series as zero, giving
$Z_0(w)=1/(1-w-4w^2+4w^3)$.
This recurrence follows equally from the integral reversed-factor
formula for the primary normalized moments.

To compute a Pochhammer ratio modulo a power of $p$, each numerator
and denominator factor is written as a power of $p$ times a unit.
Only the unit denominator is inverted; the accumulated valuation is
tracked as an integer. This computes $b_n,c_n,d_n$ without inverting
$p$ modulo its own power.

\subsection{Degree bounds, precision, and exact Newton evaluation}
For any selected normalized index $j$, the proved degree bounds are
\[
 d_j^R=2\lfloor(2j+1)/3\rfloor,\qquad
 d_j^P=2\lfloor2j/3\rfloor.
\]
Therefore its values at the small cutoffs $n=0,\ldots,d_j$ determine
the desired polynomial exactly. Let $d_{\max}$ be the largest degree
needed for a fixed family and pair $(s,p)$. Choose the common scale
\[
 A_R=\left\lceil\frac32\lfloor\log_p(3d_{\max}+1)\rfloor\right\rceil,
 \qquad
 A_P=\left\lfloor\frac32\lfloor\log_p(3d_{\max}+2)\rfloor\right\rfloor,
\]
and compute modulo $p^{A+2}$ in the corresponding family.
The previously proved normalization bounds ensure
$v_p(Mc_n)\le A_R$ and $v_p(L(L-1)b_n)\le A_P$ throughout those samples.
If a normalization is $p^vu$ with $u$ a unit, the scaled sample is
computed as $p^{A-v}$ times the unnormalized coefficient times $u^{-1}$,
with the indicated factor $-1$ or $2$. No modular division by $p$ is
performed and no precision is lost: the exponent $A-v$ is nonnegative.

For either polynomial, let $f(n)$ denote its scaled value at the
appropriate affine cutoff. Its exact Newton formula is
\[
 f(X)=\sum_{k=0}^{d_j}\Delta^kf(0)\binom Xk.
\]
Here $X=(p^{2s}-1)/3$ for the companion and
$X=(p^{2s+1}-2)/3$ for the primary. Both are nonnegative integers.
The binomial coefficients are computed as exact integers before
reduction modulo $p^{A+2}$. All finite differences and final sums are
then integer operations in that explicitly specified residue ring.
This evaluates the actual polynomial value, not an approximation
obtained by guessing a polynomial from too few samples.

If the scaled residue is nonzero with valuation $v<A+2$, its unscaled
value has exactly valuation $v-A$. If the residue is zero, the
unscaled value has valuation at least two. Thus the residue arithmetic
certifies every inequality in \eqref{VII-E7:finite-required}.

\subsection{Certificate output and exhaustion of the exceptional set}
The complete reproducible program is
\nolinkurl{work_e7_companion/finite_projector_certificates.py}; its
integer outputs are in the adjacent JSON file. It verifies 158 selected
moment indices across the seven pairs and both families. The following
table gives the largest small sample index and the minimum valuation
of the scaled residues in each family; the scale $A$ is displayed
alongside, so the minimum unscaled valuation is their difference.
\[
\begin{array}{c|c|r|c|r|r|c|r}
s&p&d_{\max}^R&A_R&\min v(p^{A_R}R)&d_{\max}^P&A_P&\min v(p^{A_P}P)\\\hline
3&5 &52 &5&5&48 &4&4\\
3&11&132&3&2&128&3&3\\
3&17&212&3&2&208&3&3\\
3&23&292&3&2&288&3&3\\
3&29&372&3&2&368&3&3\\
4&5 &68 &5&5&64 &4&4\\
5&5 &84 &5&5&80 &4&4
\end{array}
\]
The moduli used are respectively $p^{A_R+2}$ and $p^{A_P+2}$.
In particular every companion minimum is at least $-1$ after removing
the scale, and every primary minimum is at least zero. These are
stronger than the requirements in \eqref{VII-E7:finite-required}.

Applying the already proved projector-and-leakage argument with these
158 certified moment bounds proves the unchanged E.3 congruence for
all seven exceptional pairs. Combining them with the analytic ranges
therefore gives the unconditional computationally certified conclusion
\[
 \boxed{E(p^r)\equiv p^2E(p^{r-2})\pmod{p^{3r-1}}
 \quad\text{for every }p>3,\ p\equiv2\pmod3,
 \text{ and every odd }r\ge7.}
\]
The finite computations establish the finite exceptional set only;
the assertion for every other prime and level comes from the uniform
analytic theorem. The fifth level is supplied by the reflection argument and the
additional finite certificates in the next section.

\section{Exact finite completion of the level \texorpdfstring{$r=5$}{r=5}}\label{loc-VII-13}
\label{VII-sec:E7r5finite}
The uniform reflection-gain argument proves the original $r=5$
congruence for every prime $p\ge29$ with $p\equiv2\pmod3$.
The only primes satisfying $3<p<29$ and $p\equiv2\pmod3$ are
$5,11,17,23$. This section certifies that exhaustive finite exception
set using the already proved polynomial-moment method.

Here $s=2$, so the projector $1-(1-X^{p-1})^7$ uses exactly the
seven positive moment indices $j=i(p-1)$, $1\le i\le7$.
The sufficient bounds from the proved projector-and-leakage argument are
\[
 v_p\bigl(R_{i(p-1)-1}(p^4)\bigr)\ge-1,
 \qquad
 v_p\bigl(P_{i(p-1)-3}(p^5)\bigr)\ge-2.
\]
All nonunit leakage and all infinite transport terms continue to use
the uniform estimates; the calculation below only supplies these
fourteen polynomial values per prime.

Use exactly the integer Laurent recurrences, proved degree bounds,
and scaled Newton interpolation described in
Section~\ref{VII-sec:E7finiteexceptions}. In all eight family-prime cases,
the common scaling exponent is $A=3$ and the calculation is modulo
$p^5=p^{A+2}$. The normalization divisors at every sampled small cutoff
have valuation at most three. Thus multiplication by $p^{A-v}$ and
inversion of the unit part supplies each scaled sample without losing
precision. The Newton evaluation uses exact integer binomial
coefficients at $(p^4-1)/3$ or $(p^5-2)/3$.

The complete reproducible invocation, from the archive root, is
\begin{verbatim}
python work_e7_companion/r5_finite_closure_driver.py
\end{verbatim}
It imports the shared, independently audited function
\texttt{certify} from \nolinkurl{finite_projector_certificates.py}
and writes \nolinkurl{r5_finite_closure.json}. It does not overwrite
the earlier 158-index higher-level certificate. Only the Python
standard library is required.

The output for all 56 required values is summarized below. The
minimum scaled valuation is the valuation of the exact residue of
$p^3R$ or $p^3P$ modulo $p^5$; subtracting three gives the certified
unscaled valuation. All minima shown are below the precision cap,
so their listed valuations are exact.
\[
\begin{array}{c|r|r|r|r}
p&d_{\max}^{R}&\min v_p(p^3R)&d_{\max}^{P}&\min v_p(p^3P)\\\hline
5 &36 &3&32 &3\\
11&92 &2&88 &3\\
17&148&2&144&3\\
23&204&2&200&3
\end{array}
\]
Thus every companion value has valuation at least $-1$, and every
primary value has valuation at least zero, which is stronger than the
required primary bound. The finite sufficient conditions hold for
all four primes. The uniform leakage argument therefore gives
\[
 E(p^5)\equiv p\mathcal A(p^4)\pmod{p^{14}},\qquad
 \mathcal A(p^4)\equiv pE(p^3)\pmod{p^{13}}.
\]
Combining them proves the unchanged $r=5$ target for this exhaustive
finite exception set. Together with the separate uniform argument
for $p\ge29$, it yields
\[
 \boxed{E(p^5)\equiv p^2E(p^3)\pmod{p^{14}}
 \qquad\text{for every prime }p>3,\ p\equiv2\pmod3.}
\]
The finite output is not being extrapolated to other primes: those
primes are covered by the uniform reflection-gain proof. This section
makes no claim about the still separate $r=3$ case.

\section{A uniform first-step congruence at the third level}\label{loc-VII-14}
\label{VII-sec:E7firststepr3}
\begin{corollary}\label{VII-E7:first-step-r3}
For every prime $p>3$, $p\equiv2\pmod3$, the exact companion sum
of Theorem~\ref{VII-E7:companion-exact} satisfies
\[
 \boxed{E(p^3)\equiv p\mathcal A(p^2)\pmod{p^8}.}
\]
In fact the positive-moment transport part of this difference belongs
to $p^9\mathbb Z_p$.
\end{corollary}
\begin{proof}
Write $R^{(1)}_{p^3,p}(X)=1+\sum_{j\ge1}a_jX^j$ for the exact
one-step kernel. Its coefficients satisfy $v_p(a_j)\ge j$.
Summing the proved identity
$\mu_{p^3}(pJ)=p\widetilde\mu_{p^2}(J)R^{(1)}_{p^3,p}(J)$
over the companion nodes gives
\[
 E(p^3)-p\mathcal A(p^2)
 =\sum_{p\nmid H}\mu_{p^3}(H)
  +p\sum_{j\ge1}a_j\widetilde M_j(p^2).
\]
The first sum belongs to $p^8\mathbb Z_p$ by
Corollary~\ref{VII-E7:primary-unit-r3}.

The common companion moment normalization has valuation seven.
The first two explicit moments thus have valuation at least seven,
so their transport terms have valuation at least nine. For $j\ge3$,
Theorem~\ref{VII-E7:companion-log} gives
\[
 v_p\bigl(\widetilde M_j(p^2)\bigr)\ge7-\eta_{j-1},\qquad
 \eta_{j-1}=\left\lceil\frac32
 \left\lfloor\log_p(6\lfloor(2j-1)/3\rfloor+1)\right\rfloor
 \right\rceil.
\]
Here $\eta_{j-1}\le j-1$. Indeed
$6\lfloor(2j-1)/3\rfloor+1\le4j-1$, and
\[
 4j-1<5^{\lfloor2(j-1)/3\rfloor+1}\qquad(j\ge3).
\]
The latter inequality is checked for $j=3,4,5$; increasing $j$ by
three adds twelve on the left and multiplies the right by twenty-five.
It follows that each remaining transport term has valuation at least
$1+j+7-(j-1)=9$. The logarithmic growth of $\eta_{j-1}$ proves
convergence of the series. This proves the corollary.
\end{proof}

Consequently the still-unproved third-level E.3 congruence is exactly
equivalent to the shorter companion estimate
\[
 \mathcal A(p^2)-pE(p)\in p^7\mathbb Z_p.
\]
This equivalence includes an independently proved precision bound for
one entire transport step. It does not assert the remaining companion
estimate.

\section{Exact two-step transport and the remaining cancellation}\label{loc-VII-15}
Let $p>3$, $p\equiv2\pmod3$, and $r\ge3$ be odd. Set
$L=p^r$, $M=L/p^2$, and $m=(L-2)/3$. Use the node and denominator
root sets $\mathcal H_L$ and $\mathcal D_L$ of the three-grid
construction, and write $\mu_L(H)=W_L(H)T(H)$.
Define the restricted rational power series
\[
 R_{L,p}(X)=
 \frac{\prod_{d\in\mathcal D_L,\,p^2\nmid d}(1-p^2X/d)}
 {\prod_{h\in\mathcal H_L,\,p^2\nmid h}(1-p^2X/h)}
 =1+\sum_{j\ge1}c_jX^j.
\]
Every retained root has valuation zero or one. Hence each linear
coefficient $p^2/d$ or $p^2/h$ belongs to $p\mathbb Z_p$.
Expanding the finitely many inverse factors proves
\[
 c_j\in p^j\mathbb Z_p,\qquad v_p(c_j)\longrightarrow+\infty.
\]
The series converges at every integral coarse node $J$.

\begin{proposition}[Exact residual identity]\label{VII-prop:e-residual}
Put
\[
 U_{L,p}=\sum_{\substack{H\in\mathcal H_L\\p^2\nmid H}}\mu_L(H).
\]
Then
\begin{equation}\label{VII-eq:e-residual}
 \boxed{E(L)-p^2E(M)
 =U_{L,p}+p^2\sum_{j\ge3}c_j M_j(M).}
\end{equation}
No assertion that either term on the right is separately divisible
by the target modulus is made.
\end{proposition}
\begin{proof}
Since $p^2\equiv1\pmod6$, dividing all $p^2$-divisible nodes by
$p^2$ gives exactly $\mathcal H_M$. Dividing the corresponding
denominator roots gives exactly $\mathcal D_M$. For example the
positive grids end at $L-1$ and $L$, and become the positive grids
ending at $M-1$ and $M$; the negative grid is $-2$ times the first
positive grid. The positive denominator roots $6i$ and the negative
roots $-3i$ both have coarse index range
$1\le i\le(M-2)/3$. The third denominator grid ends at $-(L-3)$
and becomes the grid ending at $-(M-3)$.

The exact model values satisfy $T(p^2J)=p^2T(J)$, because
multiplication by $p^2$ preserves the residues modulo six and the
sign of each node. Separating divisible factors in the weights
therefore gives the product identities
\[
 W_L(p^2J)=W_M(J)R_{L,p}(J),\qquad
 \mu_L(p^2J)=p^2\mu_M(J)R_{L,p}(J).
\]
These are identities of products, including at zero weights;
no division by a vanishing weight is needed. Sum over coarse nodes
and add the nodes not divisible by $p^2$. Interchanging the finite
node sum with the convergent series is legitimate. The identities
$M_1(M)=M_2(M)=0$ remove its first two positive terms and give
\eqref{VII-eq:e-residual}.
\end{proof}


\section{The leading kernel coefficients from three short harmonic sums}\label{loc-VII-16}
\label{VII-sec:E7harmonic}
Let $p>3$ be prime with $p\equiv2\pmod3$, let $r\ge3$ be odd,
and put $L=p^r$, $q=(p-2)/3$. Set
\[
 H_i=\sum_{a=1}^{q}a^{-i}\in\mathbb F_p\quad(1\le i\le3).
\]
Use the exact two-step kernel
\[
 R_{L,p}(X)=
 \frac{\prod_{d\in\mathcal D_L,\ p^2\nmid d}(1-p^2X/d)}
 {\prod_{h\in\mathcal H_L,\ p^2\nmid h}(1-p^2X/h)}
 =1+\sum_{j\ge1}c_jX^j.
\]
All products here count multiplicities. In particular, an overlap
between the node and denominator lists may be canceled but need
not be removed in defining the products.

\begin{proposition}\label{VII-E7:harmonic-kernel}
For every such $p,r$,
\begin{align*}
 c_1/p&\equiv-H_1/2\pmod p,\\
 c_2/p^2&\equiv H_1^2/8+5H_2/24\pmod p,\\
 \kappa_p:=c_3/p^3&\equiv
 -\frac{3H_1^3+15H_1H_2+2H_3}{144}\pmod p.
\end{align*}
Thus these leading coefficients are independent of the odd level $r$.
\end{proposition}
\begin{proof}
For $i\le3$ define the normalized root-power difference
\[
 s_i=p^{-i}\left(
 \sum_{h\in\mathcal H_L,\ p^2\nmid h}(p^2/h)^i
 -\sum_{d\in\mathcal D_L,\ p^2\nmid d}(p^2/d)^i\right).
\]
Every term is integral at $p$. Modulo $p$, the unit roots
contribute zero, leaving exactly the roots of valuation one.
After dividing those roots by $p$, their lists are the unit roots
of the companion grids at $M=p^{r-1}=3n+1$:
\[
 \begin{split}
 \mathcal H'_M={}&\{3j+1:0\le j\le n\}
 \cup\{3j+2,-2(3j+2):0\le j<n\},\\
 \mathcal D'_M={}&\{6j,-3j:1\le j\le n\}
 \cup\{-(3j+1):0\le j<n\}.
 \end{split}
\]
These identifications follow directly by solving the original
linear root congruences modulo $p$; they respect multiplicities.

Put $t=(2p-1)/3=2q+1$. Since $n\equiv t\pmod p$, complete
blocks of $p$ consecutive indices contribute zero modulo $p$ to
each inverse-power sum under consideration. Indeed,
$\sum_{u\in\mathbb F_p^\times}u^{-i}=0$ for $1\le i\le3<p-1$.
Let
\[
 T_i=\sum_{j=1}^{t}j^{-i},\qquad
 A_i=\sum_{j=0}^{t-1}(3j+1)^{-i},\qquad
 B_i=\sum_{\substack{0\le j<t\\j\ne q}}(3j+2)^{-i}.
\]
The extra endpoint $3t+1=2p$ of the first node list is excluded
as a nonunit. The involution $j\mapsto t-1-j$ gives
$A_i=(-3)^{-i}T_i$. In the $B$ list, the residues are
$3a$ for $a=-q,\ldots,-1,1,\ldots,q$. Consequently
\[
 B_1=B_3=0,\qquad B_2=2H_2/9.
\]
Since $t=p-1-q$, the full nonzero-residue sum also gives
$T_i=(-1)^{i+1}H_i$.
Taking the difference of the companion node and denominator
power sums yields
\[
 s_i\equiv
 (1-(-1)^i)A_i+(1+(-2)^{-i})B_i
 -(6^{-i}+(-3)^{-i})T_i\pmod p.
\]
Substitution gives
\[
 s_1=-H_1/2,\qquad s_2=5H_2/12,\qquad s_3=-H_3/24
 \quad\text{in }\mathbb F_p.
\]
Finally the logarithmic product expansion, or equivalently the
first three Newton identities, gives
\[
 \frac{c_1}{p}=s_1,\quad
 \frac{c_2}{p^2}=\frac{s_1^2+s_2}{2},\quad
 \frac{c_3}{p^3}=\frac{s_1^3+3s_1s_2+2s_3}{6}.
\]
Here these identities are exact before the reductions of $s_i$.
The stated congruences follow; all displayed denominators are
units because $p>3$.
\end{proof}

This proposition computes one contribution to the residual and
does not prove its cancellation with the other contribution.
For example $p=11$ gives $(H_1,H_2,H_3)=(0,9,6)$ modulo $11$,
so $\kappa_{11}=-1$. In particular the cubic kernel coefficient
is not always divisible by $p^4$.

\section{A uniform coarse-moment estimate and a rigorous transport-tail bound}\label{loc-VII-17}
\label{VII-sec:E7coarsemoments}
Let $p>3$ be prime, $p\equiv2\pmod3$, and put $q=(p-2)/3$.
Then $q$ is odd. Use the original three-grid interpolation at
$N=p$, and denote its positive moments by $M_j(p)$.

\begin{proposition}\label{VII-E7:all-coarse-moments}
For every positive integer $j$,
\[
 \boxed{M_j(p)\in p^2\mathbb Z_p.}
\]
This bound is uniform in $j$; it is not inferred from finitely many
polynomial moment formulas.
\end{proposition}
\begin{proof}
Write $P_q(z)=D_q(z)\mathcal S_q(z)$ for the cleared polynomial
of the three-grid construction. All its coefficients belong to
$\mathbb Z_p$: the fixed denominators $k!$, $(2/3)_k$, and the
normalizing roots of $D_q$ are units for $0\le k\le q$.
We prove the stronger coefficientwise assertion
\begin{equation}\label{VII-E7:cleared-modp}
 P_q(z)\in p\mathbb Z_p[z].
\end{equation}

Work first in $\mathbb F_p(z)$. Insert the canceled numerator and
denominator parameter $e=2/3$ in $\mathcal S_q$ to exhibit the
very-well-poised ${}_6F_5(-1)$ form.
Modulo $p$ its central parameter satisfies
$a=1/3=-2q-1$, and the inserted parameter satisfies $e=2/3=-q$.
The finite canceled factors $(-q)_k$ are nonzero for $k\le q$,
so this insertion is legitimate.
Apply Whipple's terminating transformation using this $-q$ as the
terminating upper parameter. The other parameters may be written
\[
 b=a+z/6,\quad c=a-z/3,\quad e'=-q-z/3,\quad f=-q.
\]
The transformed terminating ${}_3F_2$ is
\[
 {}_3F_2\left[\begin{matrix}
 -q,\ -q+z/6,\ -q-z/3\\1-z/6,\ 1+z/3
 \end{matrix};1\right]_q.
\]
Its prefactor is $(-2q)_q/(-q+z/3)_q$, a well-defined element of
$\mathbb F_p(z)$. The displayed sum is a terminating Dixon sum
with odd central parameter $-q$: the two lower parameters equal
$1-q-(-q+z/6)$ and $1-q-(-q-z/3)$.
Reversing the summation index $k\mapsto q-k$ negates every term.
The sum is therefore zero. All transformations are finite rational
identities whose denominators are nonzero in $\mathbb F_p(z)$;
no analytic continuation in characteristic $p$ is used.
It follows that the reduction of $\mathcal S_q$ is zero as a
rational function. Since $D_q$ is a nonzero polynomial modulo $p$,
this proves \eqref{VII-E7:cleared-modp}.

Let $\widehat Q_q(z)$ be the monic polynomial whose roots are the
three interpolation grids. These roots are integers. Exactly one
root is divisible by $p$, namely the positive node $p$ itself;
every other root is a $p$-adic unit. Thus
$Q_q(z)=\ell_q\widehat Q_q(z)$ with $v_p(\ell_q)=-1$.
By \eqref{VII-E7:cleared-modp},
\[
 \frac{D_q(z)\mathcal S_q(z)}{Q_q(z)}
 =\frac{A_q(z)}{\widehat Q_q(z)},\qquad
 A_q(z)=\ell_q^{-1}P_q(z)\in p^2\mathbb Z_p[z].
\]
The inverse Laurent series at infinity of a monic polynomial in
$\mathbb Z_p[z]$ has integral coefficients after its leading power
of $z$ is removed. Hence every Laurent coefficient of this proper
rational function belongs to $p^2\mathbb Z_p$.
Those coefficients are $-M_j(p)$, which proves the proposition.
\end{proof}

\begin{corollary}\label{VII-E7:r3tail}
At $L=p^3$, let $R_{L,p}(X)=1+\sum_{j\ge1}c_jX^j$ be the exact
two-step weight kernel. Then
\[
 p^2\sum_{j\ge4}c_jM_j(p)\in p^8\mathbb Z_p.
\]
In particular the entire transport error, including its infinite
tail, satisfies
\[
 p^2\sum_{j\ge3}c_jM_j(p)
 \equiv p^2c_3M_3(p)\pmod{p^8}.
\]
\end{corollary}
\begin{proof}
The established coefficient estimate is $c_j\in p^j\mathbb Z_p$.
Proposition~\ref{VII-E7:all-coarse-moments} bounds every summand for
$j\ge4$ by $p^{2+j+2}$, hence by $p^8$. These bounds tend to
infinity, so the entire convergent tail belongs to that closed ideal.
\end{proof}

For an explicit leading term put
\[
 b_q=\frac{(1/3)_q}{q!},\qquad
 \kappa_p\equiv c_3/p^3\pmod p.
\]
Since $q$ is odd, the exact first nonzero moment is
$M_3(p)=p^2(p-1)^3b_q^3$. The corollary therefore gives
\[
 \boxed{\quad
 p^2\sum_{j\ge3}c_jM_j(p)
 \equiv-p^7\kappa_pb_q^3\pmod{p^8}.
 \quad}
\]
Together with the separately derived formula
\[
 \kappa_p\equiv-
 \frac{3H_1^3+15H_1H_2+2H_3}{144}\pmod p,
 \qquad H_i=\sum_{a=1}^{q}a^{-i},
\]
this determines the full transport contribution modulo the target
modulus by three harmonic sums of length $(p-2)/3$.

The remaining unit and valuation-one node contribution has not been
bounded here. In particular this proposition and corollary do not
claim the uniform $r=3$ recurrence: they prove that all orders beyond
the cubic transport term are already harmless, leaving one explicit
finite-node cancellation to establish.


\section{The single remaining odd level}\label{loc-VII-18}
The preceding uniform proofs and exact finite certificates establish
every odd level $r\ge5$. The sole level not settled uniformly in this
manuscript is $r=3$, for primes $p>3$, $p\equiv2\pmod3$.
The independently proved first-step congruence reduces it exactly to
\[
 \boxed{\mathcal A(p^2)-pE(p)\in p^7\mathbb Z_p.}
\]
This is a smaller finite sum but is still a conjectural uniform
estimate; its statement is not counted as its proof.

At $r=3$, let $q=(p-2)/3$, $b_q=(1/3)_q/q!$, and use the
coefficient $\kappa_p$ evaluated in Section~\ref{VII-sec:E7harmonic}.
The proved exact residual and the uniform tail bound reduce
\eqref{VII-eq:target} to
\[
 U_{p^3,p}\equiv p^7\kappa_p b_q^3\pmod{p^8}.
\]
This congruence remains to be proved uniformly. The cubic coefficient
need not vanish: $\kappa_{11}=-1$ modulo $11$. Thus it would be
incorrect to eliminate the two contributions separately.

\paragraph{Novelty and success criteria.}
The named problem, cutoff, recurrence factor and modulus have been
preserved. The proof is unconditional in its stated ranges and
contains a uniform working estimate rather than merely an equivalent
formulation. It does not meet the requirement of a complete proof of
the whole odd branch. No claim that the technique or subrange result
has historical priority is made; that requires a separate comparison
with the full relevant literature.

\paragraph{Verification.}
The symbolic WZ certificates, boundary identities, polynomial
recurrences, root-count tables, normalized moment formulas and
valuation budgets were checked independently. The accompanying
archive contains the actual derivation and validation scripts.
The higher-level finite reduction certifies 158 specified moments,
and the fifth-level finite reduction certifies 56 further moments.
Both batches use proved degree bounds and exact integer Newton
interpolation, with modular normalization performed without division
by $p$. The complete certificate evaluator is reproduced below.
Finite modular checks of the harmonic kernel coefficients used
$p=5,11,17,23,29,41$ at level three and $p=5$ at level five.
Those checks do not replace the proofs for arbitrary parameters.

\LocalAppendix
\section{Finite certificates on the two positive companion grids}\label{loc-VII-A}

\subsection{A finite certificate on the grid $H=3h+1$}
Retain the rational deformation $V_n(z)$ defined in the companion
construction. For every $h\ge0$ and $n\ge h$,
\begin{equation}\label{VII-E7positiveone:main}
 V_n(3h+1)=-\frac{(-1)^h(3h+1)}2.
\end{equation}
All lower parameters at $z=3h+1$ are regular: in particular
$(5-3h)/6$ is never an integer. The factor $(-h)_k$ makes the
sum terminate at $k=h$, so it suffices to prove the identity for $n=h$.

Set
\[
 P(h,k)=18k^2-6k+3h-1,
\]
\[
 u(h,k)=\frac{(-1)^{k-h}(6k-1)P(h,k)}{(1-3h)(3h+1)}
 \frac{(-1/3)_k((3h-1)/6)_k(-h-2/3)_k(-h)_k}
 {k!((5-3h)/6)_k(h+4/3)_k(h+2/3)_k}.
\]
This is exactly the summand of $V_h(3h+1)$ divided by the
right side of \eqref{VII-E7positiveone:main}. We prove
$\sum_{k=0}^hu(h,k)=1$ by a recurrence with step two in $h$.

Define the polynomial
\begin{align*}
 A(h,k)={}&486h^5+(1458k^2-1944k+3564)h^4\\
 &+(7857k^2-10476k+10233)h^3\\
 &+(-972k^4+2592k^3+13149k^2-19836k+14229)h^2\\
 &+(-2511k^4+6696k^3+7137k^2-15468k+9486)h\\
 &+243k^6-972k^5-189k^4+3384k^3+414k^2-4072k+2392.
\end{align*}
The rational certificate is
\[
 R(h,k)=\frac{-k(3k-2)(3h-6k+1)A(h,k)}{D_+(h,k)},
\]
where
\begin{align*}
 D_+(h,k)={}&(6k-1)(h-k+1)(h-k+2)(3h-3k+5)(3h-3k+8)\\
 &\quad\cdot(3h+3k+2)(3h+3k+4)P(h,k).
\end{align*}
Cancellation of consecutive rising factors and expansion of the
remaining polynomials verifies the rational identity
\begin{equation}\label{VII-E7positiveone:certificate}
 u(h+2,k)-u(h,k)
 =R(h,k+1)u(h,k+1)-R(h,k)u(h,k).
\end{equation}
For an independently checkable form, if
$q(h,k)=u(h,k+1)/u(h,k)$ and $r(h,k)=u(h+2,k)/u(h,k)$,
then the certificate is the polynomial equality obtained by clearing
all denominators in
$q(h,k)R(h,k+1)-R(h,k)=r(h,k)-1$.
The ratio $q$ is explicitly
\[
 q(h,k)=-\frac{(6k+5)P(h,k+1)}{(6k-1)P(h,k)}
 \frac{(k-1/3)(k+(3h-1)/6)(k-h-2/3)(k-h)}
 {(k+1)(k+(5-3h)/6)(k+h+4/3)(k+h+2/3)}.
\]

Here are the upper-boundary details. Put $g(h,k)=R(h,k)u(h,k)$.
The apparent poles at $k=h+1,h+2$ cancel because
\[
 \frac{(-h)_k}{(h-k+1)(h-k+2)}
 =\frac{(-h-2)_k}{(h+1)(h+2)}.
\]
Thus $g$ has a finite value at both boundary indices. Also
$g(h,0)=0$. Direct cancellation gives
\[
 \left.\frac{R(h,k)}{r(h,k)}\right|_{k=h+1}
 =-\frac{4(3h+1)(9h^3+39h^2+46h+15)}
 {(h+2)(6h+5)(18h^2+33h+17)},
\]
\[
 q(h+2,h+1)=
 -\frac{5(3h+2)(9h+11)}
 {(h+2)(6h+5)(18h^2+33h+17)}.
\]
The first displayed quotient plus $1+q(h+2,h+1)$ is identically
zero. Multiplication by $u(h+2,h+1)$ therefore gives
\[
 g(h,h+1)=-u(h+2,h+1)-u(h+2,h+2).
\]
Summing \eqref{VII-E7positiveone:certificate} for $0\le k\le h$
and adding these two final terms proves
\[
 \sum_{k=0}^{h+2}u(h+2,k)=\sum_{k=0}^hu(h,k).
\]
At $h=0$ the sum is $1$; at $h=1$ its terms are $-1/4$ and
$5/4$, again summing to one. Induction separately on the two
parity classes proves \eqref{VII-E7positiveone:main} for every $h$.

This proves one exact node family of the companion interpolation.
It does not, by itself, prove any prime-power lifting congruence.


\subsection{The second positive grid: a complete finite certificate}
For $h\ge1$, let $t_h(k)$ be the summand of $V_{h+1}(3h+2)$
and normalize it by $u_h(k)=t_h(k)/(3h+2)$. Explicitly,
\[
 u_h(k)=\frac{(-1)^k(6k-1)(h+6k^2-2k)}{2h(3h+2)}
 \frac{(-1/3,h/2,-h-1,-h-1/3)_k}
 {k!(2/3-h/2,h+5/3,h+1)_k}.
\]
The sum terminates at $h+1$. The desired exact assertion is
\[
 \sum_{k=0}^{h+1}u_h(k)=
 \begin{cases}1,&h\text{ odd},\\0,&h\text{ even}.\end{cases}
\]
All denominator Pochhammers are nonzero on this range.

Put
\begin{align*}
 J(h,k)={}&54h^5+162h^4k^2-216h^4k+486h^4
 +1089h^3k^2-1452h^3k+1725h^3\\
 &-108h^2k^4+288h^2k^3+2442h^2k^2-3512h^2k+3002h^2\\
 &-351hk^4+936hk^3+2082hk^2-3608hk+2549h\\
 &+27k^6-108k^5-126k^4+656k^3+507k^2-1316k+840,
\end{align*}
\[
 R(h,k)=-\frac{k(3k-2)(3h-6k+2)J(h,k)}{D_2(h,k)},
\]
where
\begin{align*}
 D_2(h,k)={}&(6k-1)(h-k+2)(h-k+3)(h+k+1)(h+6k^2-2k)\\
 &\quad\cdot(3h-3k+4)(3h-3k+7)(3h+3k+5).
\end{align*}
Direct rational cancellation proves the finite WZ certificate
\begin{equation}\label{VII-E7:F-certificate}
 u_{h+2}(k)-u_h(k)
 =R(h,k+1)u_h(k+1)-R(h,k)u_h(k).
\end{equation}
For a compact independently reproducible verification, the consecutive
ratios are
\[
 q(h,k)=\frac{u_h(k+1)}{u_h(k)}
 =\frac{(h+2k)(3k-1)(6k+5)(h-k+1)(3h-3k+1)
 (h+6k^2+10k+4)}
 {(k+1)(6k-1)(h+k+1)(h+6k^2-2k)
 (3h-6k-4)(3h+3k+5)},
\]
\begin{align*}
 r(h,k)&=\frac{u_{h+2}(k)}{u_h(k)}\\
 &=\frac{(h+1)(h+2)(h+3)(h+2k)(3h+4)(3h+5)(3h+7)}
 {(h-k+2)(h-k+3)(h+k+1)(h+k+2)(h+6k^2-2k)}\\
 &\quad\times\frac{(3h-6k+2)(h+6k^2-2k+2)}
 {(3h-3k+4)(3h-3k+7)(3h+3k+5)(3h+3k+8)}.
\end{align*}
Substitution in $q(h,k)R(h,k+1)-R(h,k)=r(h,k)-1$ leaves
identically zero after clearing denominators.

At the lower boundary $R(h,0)u_h(0)=0$. At the upper boundary,
cancel the apparent zeros and poles before evaluating. For instance,
\[
 \frac{(-h-1)_k}{(h-k+2)(h-k+3)}
 =\frac{(-h-3)_k}{(h+2)(h+3)}.
\]
The needed endpoint can also be checked entirely from rational ratios:
\[
 \left.\frac{R(h,k)}{r(h,k)}\right|_{k=h+2}
 =-\frac{4(9h^3+60h^2+127h+87)}{(h+3)(6h+11)^2},
\]
\[
 \frac{u_{h+2}(h+3)}{u_{h+2}(h+2)}
 =-\frac{3(3h+5)}{(h+3)(6h+11)^2}.
\]
The first displayed ratio plus one plus the second ratio is exactly
zero. Hence the upper boundary is
$-u_{h+2}(h+2)-u_{h+2}(h+3)$. Summing
\eqref{VII-E7:F-certificate} over $0\le k\le h+1$ therefore proves that
the normalized sum is unchanged when $h$ increases by two.
The initial normalized terms are
\[
 h=1:\quad -\frac1{10},\ \frac54,\ -\frac3{20},
 \qquad
 h=2:\quad-\frac1{16},\ -\frac{105}{176},\ \frac{11}{16},\ -\frac5{176}.
\]
Their sums are respectively one and zero, proving both parity cases.
At $h=0$ the removable value $V_n(2)=0$ for $n\ge1$ supplies the
remaining case. For an ambient truncation $n\ge h+1$, all subsequent
terms vanish and their denominators remain nonzero; the same node
value therefore holds for every admissible ambient cutoff.


\section{A finite certificate for the negative companion grid}\label{loc-VII-B}
\label{VII-sec:E7negative}
The following establishes the negative grid of the proposed companion
interpolation. It is an exact finite identity and makes no claim about
the remaining prime-power error estimate.

For $h\ge0$, set
\[
 t(h,k)=-\frac{(-1)^k(6k-1)(3k^2-k-h-1)}{2(h+1)}
 \frac{(-1/3)_k(-h-1)_k(2h+1)_k(2h+5/3)_k}
 {k!(h+5/3)_k(-2h-1/3)_k(-2h-1)_k}.
\]
All denominator factors are nonzero for $0\le k\le h+1$.
We prove
\begin{equation}\label{VII-E7neg:sum}
 \sum_{k=0}^{h+1}t(h,k)=3h+2.
\end{equation}
Put $u(h,k)=t(h,k)/(3h+2)$ and define the polynomial
\[
\begin{split}
 A(h,k)={}&1728h^5-2592h^4k^2+3456h^4k+9504h^4\\
 &-12024h^3k^2+16032h^3k+20712h^3\\
 &+432h^2k^4-1152h^2k^3-19704h^2k^2+27296h^2k+22408h^2\\
 &+1026hk^4-2736hk^3-13308hk^2+20176hk+12074h\\
 &-27k^6+108k^5+450k^4-1520k^3-3027k^2+5444k+2604.
\end{split}
\]
The certificate is
\[
 R(h,k)=-\frac{k(3k-2)(2h-k+2)(6h-3k+4)A(h,k)}{D_-(h,k)},
\]
where
\begin{align*}
 D_-(h,k)={}&16(h+1)(2h+1)(2h+3)(3h+4)(6h+5)(6h+7)\\
 &\quad\cdot(6k-1)(h-k+2)(h-3k^2+k+1).
\end{align*}
Direct rational cancellation gives
\begin{equation}\label{VII-E7neg:WZ}
 u(h+1,k)-u(h,k)
 =R(h,k+1)u(h,k+1)-R(h,k)u(h,k).
\end{equation}
For reproducibility, this is the polynomial identity obtained after
substituting the displayed $R$ and the consecutive-factor ratios of
$u$; the exact residual is zero. The quadratic factor
$h-3k^2+k+1$ cancels from the product $Ru$, so roots of that factor
cause no singularity.

We specify the upper boundary without interpreting an undefined
terminating summand. Write $r(h,k)=u(h+1,k)/u(h,k)$ as its rational
consecutive-parameter ratio. Cancellation gives
\[
 \left.\frac{R(h,k)}{r(h,k)}\right|_{k=h+2}=-1.
\]
The canceled quotient has no pole at this substitution for any
integer $h\ge0$. Thus the continued product at the upper boundary is
\[
 [R(h,k)u(h,k)]_{k=h+2}=-u(h+1,h+2).
\]
This can equivalently be obtained by cancelling
$(-h-1)_k/(h-k+2)=(-h-2)_k/(h+2)$ before specializing; the quotient
$R/r$ also covers the coincident zero factors at $h=0$.
At the lower boundary $R(h,0)u(h,0)=0$.
Sum \eqref{VII-E7neg:WZ} over $0\le k\le h+1$ and add the new endpoint
$u(h+1,h+2)$. The sums of $u$ at successive $h$ are equal.
For $h=0$, the unnormalized terms are $-1/2$ and $5/2$, whose sum
is two. Hence the normalized sum equals one, proving
\eqref{VII-E7neg:sum}.

\subsection{Embedding in the cleared interpolation}
At the negative node $H=-2(3h+2)$, the candidate $V_n(z)$ has
exactly the summand $t(h,k)$ above. Its numerator parameter
$-h-1$ terminates the finite sum at $k=h+1$.
For this node to belong to the companion grid one has $n\ge h+1$.
The companion denominator is
\[
 D_n(z)=\prod_{j=1}^n(1-z/(6j))(1+z/(3j))
         \prod_{j=0}^{n-1}(1+z/(3j+1)).
\]
The only possible zero of $D_n(H)$ is the last product's factor
at $j=2h+1$. It occurs precisely when $n\ge2h+2$.
If $h+1\le n\le2h+1$, all original summand denominators are nonzero,
the sum terminates at $h+1$, and \eqref{VII-E7neg:sum} gives
$V_n(H)=3h+2$.

If $n\ge2h+2$, treat $D_n(z)V_n(z)$ term by term after clearing
its finite denominator factors. In a summand $k\le h+1$, the
factor indexed by $j=2h+1$ remains in the denominator-clearing
tail and vanishes at $H$. In a summand $k\ge h+2$, the numerator
contains the factor at $j=h+1$ in
$1+z/[2(3j-1)]$, and that factor vanishes at $H$.
The additional factor $2-z$ in $V_n$ is nonzero at $H$.
Thus every cleared summand is zero, giving
\[
 [D_n(z)V_n(z)]_{z=H}=0=D_n(H)(3h+2).
\]
Combining the two ranges proves the full cleared negative-grid identity
\begin{equation}\label{VII-E7neg:cleared}
 \boxed{[D_n(z)V_n(z)]_{z=-2(3h+2)}
 =D_n(-2(3h+2))(3h+2)\qquad(0\le h<n).}
\end{equation}
No raw $0/0$ hypergeometric summand is used in this argument.

\section{Reproducible integer certificate evaluator}\label{loc-VII-C}
The following complete program is the evaluator used for the 158
higher-level moment certificates. For the 56 fifth-level certificates,
the exhaustive case list is instead
$[(5,2),(11,2),(17,2),(23,2)]$; all other calculations are
identical. These lists supplement the uniform proofs and do not
extrapolate from a finite sample to an infinite family.
\begin{lstlisting}[language=Python]
"""Exact finite certificates for exceptional polynomial-projector parameters.

The proved degree and logarithmic denominator bounds reduce each required
prime-power moment to Newton interpolation from small integer cutoffs.
All calculations use Python integers modulo explicit prime powers; no
floating point, symbolic heuristic, or original enormous truncated sum.
"""
from pathlib import Path
from math import comb
import json


def strip(x,p):
    assert x
    v=0
    while x%p==0:x//=p;v+=1
    return x,v


def ilog(x,p):
    k=0
    while x>=p:x//=p;k+=1
    return k


def vmod(x,p,cap):
    return cap if x==0 else strip(x,p)[1]


def poch_step(unit,val,num,den,p,mod):
    nu,nv=strip(num,p);du,dv=strip(den,p)
    return unit*nu*pow(du,-1,mod)%mod,val+nv-dv


def solve_series(a,b,previous,source,J,mod):
    out=[0]*(J+1)
    for j in range(J+1):
        rhs=sum(b[i]*previous[j-i] for i in range(min(3,j)+1))
        if j<len(source):rhs+=source[j]
        rhs-=sum(a[i]*out[j-i] for i in range(1,min(3,j)+1))
        out[j]=rhs%mod
    return out


def newton_value(values,X,mod):
    delta=values[:];ans=0;binom=1
    for k in range(len(values)):
        if k:binom=binom*(X-k+1)//k
        ans=(ans+delta[0]*(binom%mod))%mod
        delta=[(delta[i+1]-delta[i])%mod for i in range(len(delta)-1)]
    return ans


def certify(p,s,family):
    indices=[i*(p-1)-(1 if family=='companion' else 3) for i in range(1,3*s+2)]
    degrees={j:2*((2*j+1)//3) if family=='companion' else 2*((2*j)//3) for j in indices}
    J=max(indices);nmax=max(degrees.values())
    log=ilog(3*nmax+(1 if family=='companion' else 2),p)
    A=(3*log+1)//2 if family=='companion' else (3*log)//2
    cap=A+2;mod=p**cap
    sample={j:[] for j in indices}
    previous=[0]*(J+1)
    du=cu=bu=1;dv=cv=bv=0
    max_norm_v=0
    for n in range(nmax+1):
        if n:
            du,dv=poch_step(du,dv,3*n-4,3*n,p,mod)
            cu,cv=poch_step(cu,cv,3*n-1,3*n,p,mod)
            bu,bv=poch_step(bu,bv,3*n-2,3*n,p,mod)
        assert min(dv,cv,bv)>=0
        if family=='companion':
            M=3*n+1
            if n==0:
                current=[mod-1]*(J+1)
            else:
                a=[1,-2,-(M-2)*(3*M-4),2*M*(M-2)**2]
                b=[1,-2,-(M-1)*(3*M-5),-2*(M-3)*(M-1)**2]
                t=(6*n-1)*du*pow(p,dv,mod)%mod
                alpha=18*n*n-6*n-2
                current=solve_series(a,b,previous,[t,t*(alpha+1),t*alpha],J,mod)
            mu,mv=strip(M,p);norm_v=mv+cv;norm_unit=mu*cu%mod
            sign=-1
        else:
            L=3*n+2
            a=[1,-1,-(L-1)*(3*L-2),2*L*(L-1)**2]
            b=[1,-1,-(L-2)*(3*L-7),-2*(L-2)**2*(L-3)]
            t=(6*n+1)*bu*pow(p,bv,mod)%mod
            current=solve_series(a,b,previous,[t],J,mod)
            lu,lv=strip(L,p);vu,vv=strip(L-1,p)
            norm_v=lv+vv+bv;norm_unit=lu*vu*bu%mod;sign=2
        assert norm_v<=A,(family,p,s,n,norm_v,A)
        max_norm_v=max(max_norm_v,norm_v)
        scale=sign*pow(p,A-norm_v,mod)*pow(norm_unit,-1,mod)%mod
        for j in indices:
            if n<=degrees[j]:sample[j].append(current[j]*scale%mod)
        previous=current
    X=(p**(2*s)-1)//3 if family=='companion' else (p**(2*s+1)-2)//3
    rows=[]
    for j in indices:
        result=newton_value(sample[j],X,mod)
        valuation=vmod(result,p,cap)-A
        target=-(s-1) if family=='companion' else -s
        assert valuation>=target,(p,s,family,j,valuation,target)
        rows.append({'normalized_index':j,'polynomial_degree_bound':degrees[j],
            'scaled_value_residue':result,'valuation_lower_bound':valuation,
            'required_lower_bound':target,'passed':True})
    return {'p':p,'s':s,'r':2*s+1,'family':family,'precision_exponent':cap,
        'common_scale_exponent':A,'max_sample_cutoff':nmax,
        'max_normalizer_valuation':max_norm_v,'number_of_moments':len(rows),
        'checks':rows}


if __name__=='__main__':
    cases=[(5,3),(11,3),(17,3),(23,3),(29,3),(5,4),(5,5)]
    records=[]
    for p,s in cases:
        for family in ['companion','primary']:
            row=certify(p,s,family);records.append(row)
            print(json.dumps({k:v for k,v in row.items() if k!='checks'}),flush=True)
    result={'status':'All exact finite projector-moment certificates passed.',
        'pair_count':len(cases),'family_pair_count':len(records),
        'moment_count':sum(r['number_of_moments'] for r in records),
        'method':'Proved small-cutoff Laurent recurrence, common p-power scaling, exact Newton interpolation at the prime-power cutoff.',
        'records':records}
    Path(__file__).with_suffix('.json').write_text(json.dumps(result,indent=2)+'\n')
\end{lstlisting}
\section{Reproducible third-level unit-mass evaluator}\label{loc-VII-D}
This evaluator supplies the five finite exceptions in the proof of
the third-level first-step congruence. The final companion step at
that level is not asserted by this computation.
\begin{lstlisting}[language=Python]
"""Exact modular interpolation masses, using interval products in residue classes.
No floating arithmetic; products remove powers of p before inversion.
"""
from pathlib import Path
from collections import Counter
import json,time

class Products:
    def __init__(self,p,L,N):
        self.p=p;self.mod=p**L;self.tables={}
        B=4*N
        for step in (3,6):
            for residue in range(step):
                low=(-B-residue+step-1)//step;high=(B-residue)//step
                vs=[0];us=[1]; factors=[]
                for j in range(low,high+1):
                    x=residue+step*j;v=0
                    if x==0:x=1
                    else:
                        while x%p==0:x//=p;v+=1
                    factors.append(x%self.mod)
                    vs.append(vs[-1]+v);us.append(us[-1]*x%self.mod)
                inv=[0]*len(us);inv[-1]=pow(us[-1],-1,self.mod)
                for j in range(len(factors)-1,-1,-1):
                    inv[j]=inv[j+1]*factors[j]%self.mod
                self.tables[(step,residue)]=(low,vs,us,inv)
    def product(self,lo,hi,step):
        assert (hi-lo)%step==0 and hi>=lo
        residue=lo%step;base,vs,us,inv=self.tables[(step,residue)]
        a=(lo-residue)//step-base;b=(hi-residue)//step-base+1
        return vs[b]-vs[a],us[b]*inv[a]%self.mod

def analyze(p,r,L=10):
    N=p**r;m=(N-2)//3;mod=p**L;pr=Products(p,L,N)
    def prods(H,numer):
        ranges=([(6-H,6*m-H,6),(-3*m-H,-3-H,3),(-3*m+1-H,-2-H,3)] if numer
                else [(1-H,3*m+1-H,3),(2-H,3*m+2-H,3),(-6*m-2-H,-2-H,6)])
        v=0;u=1
        for lo,hi,step in ranges:
            a,b=pr.product(lo,hi,step);v+=a;u=u*b%mod
        return v,u
    vq,uq=prods(0,False);vd,ud=prods(0,True)
    cv=vq-vd;cu=uq*pow(ud,-1,mod)%mod
    denominator_roots=set(z for j in range(1,m+1) for z in (6*j,-3*j,-3*j+1))
    mass=[0]*(r+1);counts=[0]*(r+1);minimum=[None]*(r+1);dist=[Counter() for _ in range(r+1)]
    def vp(x):
        n=0
        if not x:return L
        while x%p==0:n+=1;x//=p
        return n
    half=pow(2,-1,mod)
    for h in range(m+1):
        for H,c in [(3*h+1,1 if h%2==0 else 0),(3*h+2,(-1)**h*half),(-6*h-2,-half)]:
            if not c or H in denominator_roots:continue
            a,b=prods(H,True);d,e=prods(H,False)
            v=cv+a-d;assert v>=0
            u=cu*c*b*pow(e,-1,mod)%mod
            term=u*p**v%mod;t=vp(H)
            mass[t]=(mass[t]+term)%mod;counts[t]+=1;dist[t][v]+=1
            minimum[t]=v if minimum[t] is None else min(minimum[t],v)
    return {'p':p,'r':r,'modulus_exponent':L,'mass_mod':mass,
            'mass_valuations':[vp(z) for z in mass],'counts':counts,
            'minimum_individual_valuations':minimum,
            'individual_valuation_counts':[dict(z) for z in dist]}
if __name__=='__main__':
    rows=[]
    for p in (5,11,17,23,29):
        start=time.time();row=analyze(p,3);rows.append(row)
        print(p,row['mass_valuations'],row['minimum_individual_valuations'],'seconds',round(time.time()-start,2),flush=True)
    Path(__file__).with_suffix('.json').write_text(json.dumps(rows,indent=2))
\end{lstlisting}


\begin{thebibliography}{9}
\bibitem{VII-Swisher}
H. Swisher, \emph{On the supercongruence conjectures of van Hamme},
Research in the Mathematical Sciences \textbf{2} (2015), Article 18,
Section 7, Table 2.
\url{https://doi.org/10.1186/s40687-015-0037-6}.
\bibitem{VII-Bailey}
W. N. Bailey, \emph{Generalized Hypergeometric Series}, Cambridge
University Press, 1935. Classical Whipple and Dixon transformations.
\bibitem{VII-GuoSchlosser}
V. J. W. Guo and M. J. Schlosser,
\emph{New q-Analogues of Van Hamme's (E.2) Supercongruence and of
a Supercongruence by Swisher}, Results in Mathematics (2023).
\url{https://doi.org/10.1007/s00025-023-01885-8}.
\bibitem{VII-Corpus}
QEEDDD, \emph{Verifiable Mathematical Reasoning Corpus},
research context supplied for this investigation.
\url{https://github.com/QEEDDD/Verifiable-Mathematical-Reasoning-Corpus}.
\end{thebibliography}

\clearpage\endgroup

%%%%%%%%%%%%%%%%%%%%%%%%%%%%%%%%%%%%%%%%%%%%%%%%%%%%%%%%%%%%%%%%%%%%%%%%%%
% VIII | Swisher_E3_Odd_Branch_Completion.tex
%%%%%%%%%%%%%%%%%%%%%%%%%%%%%%%%%%%%%%%%%%%%%%%%%%%%%%%%%%%%%%%%%%%%%%%%%%
\clearpage\begingroup
\SetPaper{VIII}{Reflection, rational rigidity and the odd branch of E.3}
\RestoreCounters
\PaperTitle{Reflection, Rational Rigidity, and\\the Odd Branch of Swisher's E.3}{Keio research notes}{11 November 2026}

\begin{abstract}
For every prime $p>3$ with $p\equiv2\pmod3$ and every odd integer
$r\ge3$, we prove Swisher's E.3 recurrence with its original cutoff
and modulus $p^{3r-1}$. The higher levels follow from exact finite
interpolation, polynomial moments, a logarithmic denominator bound,
and an exhaustive finite set of integer certificates. The third
level is settled uniformly by a finite-field rational identity.
A reflected block derivative has a zero of order $p-1$ at the
origin, only simple finite poles, a specified regular point, and
a known leading coefficient at infinity. These properties determine
the entire rational function. Its unit-pole residue cancels exactly
the cubic transport correction, without evaluating either separately.
All proof ingredients and finite certificates are included in this
single source. The third-level argument includes $p=5$ and $p=11$.
Historical priority is not asserted.
\end{abstract}
\noindent\textbf{Scope and verification.} This manuscript supplies a
proof of the complete odd branch $p\equiv2\pmod3$, $r\ge3$.
It incorporates the earlier higher-level proof and the new uniform
third-level completion. It has received internal independent checks,
but has not undergone external peer review. Other branches and other
Swisher labels are not re-proved here. Finite computations are
identified separately from proofs valid for arbitrary parameters.
\setcounter{tocdepth}{1}

\begin{LocalContents}
\LocalEntry{loc-VIII-1}{1}{The unchanged problem and the proved range}
\LocalEntry{loc-VIII-2}{2}{Three exact interpolation grids at the short E cutoff}
\LocalEntry{loc-VIII-3}{3}{Exact moments for the three-grid short E deformation}
\LocalEntry{loc-VIII-4}{4}{An exact negative-parameter companion and its first three moments}
\LocalEntry{loc-VIII-5}{5}{A uniform logarithmic valuation bound for the primary E moments}
\LocalEntry{loc-VIII-6}{6}{Polynomial structure of every companion moment}
\LocalEntry{loc-VIII-7}{7}{A uniform logarithmic loss for all companion moments}
\LocalEntry{loc-VIII-8}{8}{A uniform lower bound for every companion interpolation mass}
\LocalEntry{loc-VIII-9}{9}{One-step transport of the valuation-one stratum}
\LocalEntry{loc-VIII-10}{10}{The original odd E.3 congruence at all sufficiently high levels}
\LocalEntry{loc-VIII-11}{11}{A reflection gain for companion interpolation samples}
\LocalEntry{loc-VIII-12}{12}{Exact finite certificates for the remaining higher-level cases}
\LocalEntry{loc-VIII-13}{13}{Exact finite completion of the level \texorpdfstring{$r=5$}{r=5}}
\LocalEntry{loc-VIII-14}{14}{A uniform first-step congruence at the third level}
\LocalEntry{loc-VIII-15}{15}{Exact two-step transport and the remaining cancellation}
\LocalEntry{loc-VIII-16}{16}{The leading kernel coefficients from three short harmonic sums}
\LocalEntry{loc-VIII-17}{17}{A uniform coarse-moment estimate and a rigorous transport-tail bound}
\LocalEntry{loc-VIII-18}{18}{The first resonance determines the companion pole residues}
\LocalEntry{loc-VIII-19}{19}{Exact reflection and second-order unit-support reduction}
\LocalEntry{loc-VIII-20}{20}{A coefficientwise valuation-two theorem and a full-block derivative}
\LocalEntry{loc-VIII-21}{21}{Simple poles of the full-block derivative}
\LocalEntry{loc-VIII-22}{22}{The leading infinity coefficient of the full-block derivative}
\LocalEntry{loc-VIII-23}{23}{An integral residue functional for the unit poles}
\LocalEntry{loc-VIII-24}{24}{Closure of the third odd level}
\LocalEntry{loc-VIII-25}{25}{The completed odd branch and verification}
\LocalEntry{loc-VIII-A}{A}{Finite certificates on the two positive companion grids}
\LocalEntry{loc-VIII-B}{B}{A finite certificate for the negative companion grid}
\LocalEntry{loc-VIII-C}{C}{Reproducible integer certificate evaluator}
\LocalEntry{loc-VIII-D}{D}{Reproducible third-level unit-mass evaluator}
\LocalEntry{loc-VIII-E}{E}{Direct checks of the original third-level sum}
\end{LocalContents}\medskip

\section{The unchanged problem and the proved range}\label{loc-VIII-1}
Write $(a)_k=\prod_{j=0}^{k-1}(a+j)$, $(a)_0=1$, and
$v_p(p)=1$, $v_p(0)=+\infty$. For $L\equiv2\pmod3$, $L\ge2$, set
\[
 E(L)=\sum_{k=0}^{(L-2)/3}(-1)^k(6k+1)
       \frac{(1/3)_k^3}{k!^3}.
\]
The odd branch of Swisher's Table 2, E.3~\cite{VIII-Swisher}, asks for
\begin{equation}\label{VIII-eq:target}
 E(p^r)-p^2E(p^{r-2})\in p^{3r-1}\mathbb Z_p,
 \qquad p>3,\quad p\equiv2\pmod3,\quad r\ge3\text{ odd}.
\end{equation}
Only this previously remaining branch is studied here. No result
from a different Swisher label is assumed as a proof of this one.

\begin{theorem}[Range established in this manuscript]\label{VIII-thm:main-range}
Congruence \eqref{VIII-eq:target} holds for every prime $p>3$ with
$p\equiv2\pmod3$ and every odd integer $r\ge3$:
\[
 \boxed{E(p^r)\equiv p^2E(p^{r-2})\pmod{p^{3r-1}},\qquad r\ge3\text{ odd}.}
\]
The higher levels use uniform arguments and exact finite certificates
for explicitly enumerated exceptions. The third level is proved uniformly
in Theorem~\ref{VIII-thm:e8-r3-closed}, with no exceptional prime.
\end{theorem}

The proof is supplied below after the exact interpolation and
valuation lemmas. Its central estimate is stronger than an
initial-moment calculation: every normalized moment loses at most
$O(\log j)$ powers of $p$ at moment order $j$. A polynomial of
degree $O(pr)$ can therefore separate the unit nodes at the target
precision. An exact companion formula controls its contribution
on the other nodes. This is the step that turns the previous
exact reductions into a proof on an infinite subrange.

\section{Three exact interpolation grids at the short E cutoff}\label{loc-VIII-2}
\label{VIII-sec:E6models}
For an integer $m\ge0$ put $N=3m+2$ and
\[
 \mathcal S_m(z)=\sum_{k=0}^{m}(-1)^k(6k+1)
 \frac{(1/3)_k(1/3+z/6)_k(1/3-z/3)_k(2/3-z/3)_k}
 {k!(1-z/6)_k(1+z/3)_k(2/3+z/3)_k}.
\]
At $z=0$ this is the original E sum with cutoff $m$.
Define
\[
 D_m(z)=\prod_{j=1}^{m}
 (1-z/(6j))(1+z/(3j))(1+z/(3j-1)),
\]
\[
 Q_m(z)=\prod_{h=0}^{m}
 (1-z/(3h+1))(1-z/(3h+2))(1+z/(6h+2)).
\]
The cleared expression $P_m(z)=D_m(z)\mathcal S_m(z)$ is a
polynomial of degree at most $3m$. Its $3m+3$ distinct interpolation
nodes are
\[
 C_h=3h+1,\qquad F_h=3h+2,\qquad B_h=-6h-2
 \quad(0\le h\le m).
\]
At these nodes the exact cleared values are
\begin{equation}\label{VIII-E6:models}
 P_m(H)=D_m(H)T(H),\qquad
 \begin{cases}
 T(C_h)=C_h,&h\text{ even},\\
 T(C_h)=0,&h\text{ odd},\\
 T(F_h)=(-1)^hF_h/2,\\
 T(B_h)=-B_h/2.
 \end{cases}
\end{equation}

\subsection{The two positive grids}
For complex $z$ with positive real part and away from denominator
poles, apply the nonterminating Whipple ${}_6F_5(-1)$ transformation~\cite{VIII-Bailey}
with
\[
 a=1/3,\quad b=1/3+z/6,\quad c=1/3-z/3,
 \quad e=2/3,\quad f=2/3-z/3.
\]
The equal factors at $e=1+a-e=2/3$ cancel. The transformed Dixon
${}_3F_2(1)$ has upper parameters
$2/3,2/3+z/6,2/3-z/3$ and lower parameters
$1-z/6,1+z/3$.
The required convergence conditions reduce to
$\operatorname{Re}(z/3)>0$ and $\operatorname{Re}(z/6)>0$.
Whipple followed by Dixon and gamma reflection evaluates the series as
\[
 \frac z2\,
 \frac{\sin(\pi(1/3+z/6))}{\sin(\pi z/6)}.
\]
At $z=C_h$ or $F_h$, the original series terminates at $h$;
its lower parameters are never nonpositive integers at these
positive nodes. The displayed expression is respectively
$C_h$ or zero according to the parity of $h$, and
$(-1)^hF_h/2$. This proves the first three lines of
\eqref{VIII-E6:models}. It also explains why the same sine expression
must not be substituted on the negative grid: there a lower
parameter becomes a nonpositive integer, and analytic continuation
does not represent the specified finite truncation.

\subsection{A finite WZ certificate for the negative grid}
At $z=B_h$, the finite terminating sum through $h$ is
\[
 \sum_{k=0}^{h} t(h,k),\qquad
 t(h,k)=(-1)^k(6k+1)
 \frac{(1/3)_k(-h)_k(2h+1)_k(2h+4/3)_k}
 {k!(h+4/3)_k(1/3-2h)_k(-2h)_k}.
\]
For $0\le k\le h$ the displayed denominator has no zero;
at $h=0$ only $k=0$ is used. We prove
\begin{equation}\label{VIII-E6:Bmodel}
 \sum_{k=0}^{h}t(h,k)=3h+1.
\end{equation}
Put $u(h,k)=t(h,k)/(3h+1)$ and define
\[
 \begin{split}
 A(h,k)={}&864h^4+3024h^3-144h^2k^2+96h^2k+3912h^2\\
 &-264hk^2+176hk+2212h+9k^4-12k^3-119k^2+82k+460,
 \end{split}
\]
\[
 R(h,k)=
 -\frac{k(3k-1)(2h-k+1)(6h-3k+2)A(h,k)}
 {16(h+1)(2h+1)^2(3h+2)(6h+5)(6h+7)(6k+1)(h-k+1)}.
\]
Direct rational simplification gives the certificate
\begin{equation}\label{VIII-E6:WZ}
 u(h+1,k)-u(h,k)
 =R(h,k+1)u(h,k+1)-R(h,k)u(h,k).
\end{equation}
To interpret the upper boundary, cancel the apparent pole using
\[
 \frac{(-h)_k}{h-k+1}=
 \frac{(-h-1)_k}{h+1}.
\]
Thus the product $R(h,k)u(h,k)$ has a well-defined value at
$k=h+1$. Direct cancellation gives
\[
 R(h,0)u(h,0)=0,\qquad
 R(h,h+1)u(h,h+1)=-u(h+1,h+1).
\]
Summing \eqref{VIII-E6:WZ} for $0\le k\le h$ and adding the last term
$u(h+1,h+1)$ proves that the sum of $u(h,k)$ over $0\le k\le h$
is independent of $h$. Its value at $h=0$ is one, proving
\eqref{VIII-E6:Bmodel}.

\subsection{The apparent poles in the cleared interpolation}
At $z=B_h$, the only possible zero in $D_m$ is its factor indexed
by $j=2h+1$ in $1+z/(3j-1)$.
If $2h\ge m$, this index is outside $D_m$; the original finite
sum terminates at $h$, and \eqref{VIII-E6:Bmodel} gives
$P_m(B_h)=D_m(B_h)(3h+1)$.
If $2h<m$, then $D_m(B_h)=0$. In each cleared summand with
$k\le h$ the same zero remains in the denominator-clearing tail.
In each summand with $k\ge h+1$ the numerator contains the factor
$(1/3+z/6+h)$, which vanishes at $B_h$.
Every cleared summand is therefore zero, so $P_m(B_h)=0$.
This is again exactly $D_m(B_h)(3h+1)$ and proves
\eqref{VIII-E6:models} at every negative node, without evaluating any
$0/0$ summand.

Consequently interpolation of $P_m$ at all the listed nodes gives
the exact identity
\[
 \mathcal S_m(0)=\sum_H W_m(H)T(H),\qquad
 W_m(H)=\frac{D_m(H)}{\prod_{K\ne H}(1-H/K)},
\]
with $K,H$ ranging over the three grids. The weights may be zero;
their rational denominators are never zero.
Equivalently,
\[
 \frac{D_m(z)\mathcal S_m(z)}{Q_m(z)}
 =\sum_H\frac{W_m(H)T(H)}{1-z/H}.
\]
The left side is $O(z^{-3})$ at infinity, proving the exact
vanishing of its first and second positive moments.

\section{Exact moments for the three-grid short E deformation}\label{loc-VIII-3}
\label{VIII-sec:e-three-grid-moments}
The statements in this section are exact. The uniform $p$-integrality
assertion needed by the proposed moment argument is not assumed or
proved here.
Let $L=3m+2$, $m\ge0$, and $b_k=(1/3)_k/k!$. Define
\[
 S_m(z)=\sum_{k=0}^m(-1)^k(6k+1)b_k^3
 \prod_{h=0}^{k-1}
 \frac{(1-z/(3h+1))(1-z/(3h+2))(1+z/(6h+2))}
 {(1-z/(6h+6))(1+z/(3h+3))(1+z/(3h+2))}.
\]
Put
\[
 q_h(z)=(1-z/(3h+1))(1-z/(3h+2))(1+z/(6h+2)),
\]
\[
 d_h(z)=(1-z/(6h+6))(1+z/(3h+3))(1+z/(3h+2)),
\]
\[
 Q_m(z)=\prod_{h=0}^m q_h(z),\qquad
 D_m(z)=\prod_{h=0}^{m-1}d_h(z),\qquad
 F_m(z)=\frac{D_m(z)S_m(z)}{Q_m(z)}.
\]
The numerator $D_mS_m$ is a polynomial of degree at most $3m$,
whereas $Q_m$ has degree $3m+3$ and distinct roots. Thus $F_m$ is
proper and has a Laurent expansion at infinity
\[
 F_m(z)=-\sum_{j\ge1}M_j(L)z^{-j},\qquad M_1(L)=M_2(L)=0.
\]
Equivalently, if $F_m(z)=\sum_H\mu_m(H)/(1-z/H)$ is its partial
fraction expansion at its three node grids, then
$M_j(L)=\sum_H\mu_m(H)H^j$. This definition allows a cancelled pole
to have residue weight zero; it does not assume that all values of
$S_m$ at lower-parameter collisions are terminating evaluations.

\begin{proposition}\label{VIII-prop:e-three-grid-moments}
Set
\[
 C(L)=(-1)^{m+1}L^2(L-1)^3b_m^3.
\]
There are uniquely determined even polynomials $P_j(L)\in\mathbb Q[L^2]$,
$j\ge0$, with $P_0=1$ and $\deg P_j\le2\lfloor2j/3\rfloor$, such that
\[
 M_{j+3}(L)=C(L)P_j(L).
\]
In particular,
\[
 P_1=1,\qquad P_2=2L^2-3,\qquad
 P_3=-\frac{L^4-13L^2+16}{4}.
\]
\end{proposition}
\begin{proof}
At infinity the ratio of the leading terms of $q_h$ and $d_h$ is
$-(3h+3)^2/(3h+1)^2$. Hence
\[
 S_m(\infty)=\sum_{k=0}^m(6k+1)b_k
 =\frac{L(L-1)b_m}{2}.
\]
The last equality telescopes: the difference between
$(3k+1)(3k+2)b_k/2$ and its value at $k-1$ is $(6k+1)b_k$.
The leading coefficient of $D_m/Q_m$ is
$(-1)^m2L(L-1)^2b_m^2z^{-3}$. This proves $M_3=C(L)$.

Separating the last summand of $S_m$ gives
\begin{equation}\label{VIII-eq:e-three-grid-rational}
 q_m(z)F_m(z)=d_{m-1}(z)F_{m-1}(z)
                         +(-1)^m(2L-3)b_m^3.
\end{equation}
Write
\[
 A=L(L-1),\quad B=(L-2)(L-3),\quad
 a=(L-1)(3L-2),\quad c=(L-2)(3L-7).
\]
The two cubic factors have expansions
\[
 q_m(z)=\frac{z^3-z^2-az+2L(L-1)^2}{2L(L-1)^2},
\]
\[
 d_{m-1}(z)=
 -\frac{z^3-z^2-cz-2(L-3)(L-2)^2}{2(L-3)(L-2)^2}.
\]
Since $b_m/b_{m-1}=(L-4)/(L-2)$, comparison of coefficients in
\eqref{VIII-eq:e-three-grid-rational}, after dividing by $C(L)$, yields
\begin{align*}
 A P_j(L)-B P_j(L-3)
 &=A P_{j-1}(L)-B P_{j-1}(L-3)\\
 &\quad+A a P_{j-2}(L)-B c P_{j-2}(L-3)\\
 &\quad-2A L(L-1)^2P_{j-3}(L)
        -2B(L-3)(L-2)^2P_{j-3}(L-3),
\end{align*}
for $j\ge1$, where $P_{-1}=P_{-2}=P_{-3}=0$ and $P_0=1$.
This also provides an independent recursive definition of the
polynomials. Indeed the operator
\[
 \mathcal L(P)=L(L-1)P(L)-(L-2)(L-3)P(L-3)
\]
maps $\mathbb Q[L^2]$ to
$(2L-3)\mathbb Q[L(L-3)]$. The first space consists of even
polynomials; the second consists of polynomials anti-invariant
under $L\mapsto3-L$. In the displayed power bases,
$\mathcal L(L^{2k})/(2L-3)$ has degree $k$ and leading coefficient
$3k+2$. Consequently this operator is an isomorphism of the filtered
vector spaces over $\mathbb Q$.
For even preceding polynomials, the right side of the recurrence is
anti-invariant under $L\mapsto3-L$; the two cubic terms exchange
with opposite signs under this reflection. Thus the recurrence
uniquely defines an even $P_j$ at every step. It gives the degree bound
\[
 \deg P_j\le\max\{\deg P_{j-1},\ \deg P_{j-2}+2,
                                    \deg P_{j-3}+4\}
 \le2\lfloor2j/3\rfloor.
\]
At $L=2$, evaluation of the polynomial recurrence gives the
coefficients of
\[
 \sum_{j\ge0}P_j(2)t^j=\frac1{1-t-4t^2+4t^3}.
\]
This is exactly the normalized expansion of $F_0=1/q_0$.
The recurrence in $m$ now proves the moment identities for every
integer $m\ge0$.
\end{proof}

\paragraph{Prime-power valuation of the first moments.}
For $p>3$, $p\equiv2\pmod3$, and odd $r$, take $L=p^r$.
Then
\[
 v_p(b_{(p^r-2)/3})=\frac{r-1}{2},\qquad
 v_p(C(p^r))=\frac{7r-3}{2}.
\]
For the first equality count multiples of $p^\ell$ in the numerator
and denominator of $b_m$. If $\ell$ is even, the numerator acquires
one additional multiple, while for odd $\ell$ the counts agree;
there are $(r-1)/2$ positive even levels at most $r$.
Therefore $M_1=M_2=0$, and $M_3,M_4,M_5,M_6$ all have exactly the
valuation $(7r-3)/2$ for these primes. This exceeds or equals the
required $3r-1$.

\paragraph{Why a uniform denominator estimate is needed.}
The polynomial operator is not an isomorphism over
$\mathbb Z[1/2]$: its diagonal is $3k+2$, not a unit there.
For example
\[
 P_6(L)=\frac{L^8+290L^6-3325L^4+16666L^2-29900}{28},
\]
and $P_9$ has denominator $280$. Thus the rational-polynomial
existence proof alone does not establish $P_j(p^r)\in\mathbb Z_p$
for all $j$ at primes $p\equiv2\pmod3$. The logarithmic bounds proved below supply the uniform estimate
needed in this manuscript. They are used with finite polynomial
separation, so no Euler-power limiting argument is required.

\section{An exact negative-parameter companion and its first three moments}\label{loc-VIII-4}
\label{VIII-sec:E7companion}
This section proves the exact finite companion identity used in the
higher-level argument. The identity alone is not a supercongruence.

Let $M=3n+1$, $n\ge0$, and put
\[
 \mathcal H_n=\{3h+1:0\le h\le n\}
 \cup\{3h+2:0\le h<n\}\cup\{-2(3h+2):0\le h<n\},
\]
\[
 \mathcal D_n=\{6j,-3j:1\le j\le n\}
 \cup\{-3h-1:0\le h<n\}.
\]
All $3n+1$ nodes in $\mathcal H_n$ are distinct and nonzero. Define
\[
 U(3h+1)=-(-1)^h(3h+1)/2,\qquad
 U(3h+2)=\begin{cases}3h+2,&h\text{ odd},\\0,&h\text{ even},\end{cases}
 \qquad U(-2(3h+2))=3h+2,
\]
\[
 D_n(z)=\prod_{D\in\mathcal D_n}(1-z/D),\qquad
 Q_n(z)=\prod_{H\in\mathcal H_n}(1-z/H),\qquad
 \widetilde W_n(H)=\frac{D_n(H)}{\prod_{K\in\mathcal H_n,K\ne H}(1-H/K)}.
\]
The companion and its positive moments are
\[
 \mathcal A(M)=\sum_H\widetilde W_n(H)U(H),\qquad
 \widetilde M_j(M)=\sum_H\widetilde W_n(H)U(H)H^j.
\]
Set $d_k=(-1/3)_k/k!$ and $c_k=(2/3)_k/k!$.

\begin{theorem}\label{VIII-E7:companion-exact}
For every integer $n\ge0$,
\begin{align*}
 \mathcal A(3n+1)
 &=-\frac12\sum_{k=0}^n(-1)^k(6k-1)(9k^2-3k-1)d_k^3\\
 &=\sum_{k=0}^n(-1)^k(6k-1)d_k^3+\frac{(-1)^n}{2}c_n^3.
\end{align*}
Moreover, with $C_n=(-1)^{n+1}M^2c_n^3/2$,
\[
 \widetilde M_1(M)=C_n,\qquad
 \widetilde M_2(M)=C_n(M^2+1)/2,\qquad
 \widetilde M_3(M)=C_n(2M^2-1).
\]
\end{theorem}

\subsection{A rational deformation with exactly the companion node values}
Put
\[
 K_k(z)=\prod_{j=0}^{k-1}
 \frac{(1+z/[2(3j-1)])(1-z/(3j-1))(1-z/(3j+1))}
 {(1-z/[2(3j+3)])(1+z/(3j+3))(1+z/(3j+1))},
\]
\begin{equation}\label{VIII-E7:V}
 V_n(z)=-\frac1{2(2-z)}\sum_{k=0}^n
 (-1)^k(6k-1)(z+18k^2-6k-2)d_k^3K_k(z).
\end{equation}
Although this expression displays $2-z$ in its denominator, its
$k=0$ term is the constant $-1/2$. Every $k\ge1$ term has the
factor $1-z/2$ in its numerator, cancelling this additional denominator.
Consequently $D_n(z)V_n(z)$ is a polynomial of degree at most $3n$.
At $z=2$, the terms with $k\ge2$ vanish after cancellation; the $k=1$
term simplifies to $-5(z-1)(z+10)/(6(z-6)(z+3))$ and has value $1/2$.
Thus $V_n(2)=0$ for every $n\ge1$.

The three finite certificate proofs supplied below establish
\begin{equation}\label{VIII-E7:cleared-all}
 [D_n(z)V_n(z)]_{z=H}=D_n(H)U(H)
 \qquad(H\in\mathcal H_n).
\end{equation}
On the first positive grid the sum terminates at $h$; on the second
it terminates at $h+1$, with $z=2$ handled above. The lower parameters
are regular on both positive grids. On the negative grid the finite
certificate gives the sum through $h+1$; where a later denominator
would vanish, the cleared identity is proved term by term. Thus
\eqref{VIII-E7:cleared-all} never assigns a value to a raw $0/0$ summand.

Interpolation of the degree-$3n$ polynomial $D_nV_n$ at its $3n+1$
distinct nodes now proves
\[
 V_n(0)=\sum_H\widetilde W_n(H)U(H)=\mathcal A(M),\qquad
 F_n(z):=\frac{D_n(z)V_n(z)}{Q_n(z)}
 =\sum_H\frac{\widetilde W_n(H)U(H)}{1-z/H}.
\]
Evaluation of \eqref{VIII-E7:V} at zero proves the first equality in
Theorem~\ref{VIII-E7:companion-exact}.

\subsection{An elementary telescoping simplification}
Let $u_k=(-1)^kd_k^3$ and $B_k=k^3u_k$. Since
$B_{k+1}=-(k-1/3)^3u_k$, direct expansion gives
\[
 (6k-1)(9k^2-3k-1)u_k
 =-27(B_{k+1}-B_k)-2(6k-1)u_k.
\]
Summing uses $B_0=0$ and $27B_{n+1}=(-1)^nc_n^3$ and proves the
second equality of the theorem, without any asymptotic argument.

\subsection{The adjacent-cutoff recurrence and the three moments}
For $n\ge1$, write
\[
 A_n(z)=(1-z/M)(1-z/(M-2))(1+z/[2(M-2)]),
\]
\[
 B_n(z)=(1-z/[2(M-1)])(1+z/(M-1))(1+z/(M-3)).
\]
These are precisely $Q_n/Q_{n-1}$ and $D_n/D_{n-1}$.
The numerator of $K_n(z)$ is
$(1+z)(1-z/2)Q_{n-1}(z)$. Separating the last summand of
\eqref{VIII-E7:V} therefore gives the exact polynomial recurrence
\begin{equation}\label{VIII-E7:companion-recurrence}
 A_n(z)F_n(z)-B_n(z)F_{n-1}(z)=G_n(z),
\end{equation}
\[
 G_n(z)=-\frac{(-1)^n(6n-1)d_n^3}{4}
 (z+1)(z+18n^2-6n-2).
\]
This derivation uses the summand formula, not the moment identities
being proved.

At infinity, write
$F_n(z)=-\sum_{j\ge1}\widetilde M_j(M)z^{-j}$.
The coefficients of $z^2,z^1,z^0$ in
\eqref{VIII-E7:companion-recurrence} determine its first three moments
successively. For complete verification, write $A_n=\sum a_i z^i$,
$B_n=\sum b_i z^i$, and $G_n=\sum g_i z^i$. Those three equations are
\[
 -a_3\widetilde M_1+b_3\widetilde M'_1=g_2,
\]
\[
 -a_2\widetilde M_1-a_3\widetilde M_2
 +b_2\widetilde M'_1+b_3\widetilde M'_2=g_1,
\]
\[
 -a_1\widetilde M_1-a_2\widetilde M_2-a_3\widetilde M_3
 +b_1\widetilde M'_1+b_2\widetilde M'_2+b_3\widetilde M'_3=g_0,
\]
where primes mean the preceding cutoff $M-3$. Here
\[
 a_3=\frac1{2M(M-2)^2},\quad a_2=-2a_3,\quad
 a_1=-\frac{3M-4}{2M(M-2)},
\]
\[
 b_3=-\frac1{2(M-1)^2(M-3)},\quad b_2=-2b_3,\quad
 b_1=\frac{3M-5}{2(M-1)(M-3)}.
\]
Substitution of the three claimed formulas satisfies these rational
identities, using
\[
 \frac{C_{n-1}}{C_n}
 =-\frac{(M-3)^2(M-1)^3}{M^2(M-2)^3}.
\]
At $n=0$, $F_0(z)=-1/[2(1-z)]$, so every positive moment is $-1/2$,
exactly the claimed initial values. Since $a_3\ne0$ for $n\ge1$,
induction proves all three formulas. This completes the theorem.


\section{A uniform logarithmic valuation bound for the primary E moments}\label{loc-VIII-5}
\label{VIII-sec:e-uniform-log-bound}
Let $P_j(L)\in\mathbb Q[L^2]$ be the primary moment polynomials defined
by the exact three-grid recurrence, normalized by $P_0=1$. They satisfy
\[
 \deg P_j\le d_j:=2\left\lfloor\frac{2j}{3}\right\rfloor,
 \qquad M_{j+3}(L)=C(L)P_j(L),
\]
where, for $L=3m+2$,
\[
 C(L)=(-1)^{m+1}L^2(L-1)^3b_m^3,\qquad b_m=\frac{(1/3)_m}{m!}.
\]
For a prime $p>3$, $p\equiv2\pmod3$, define
\[
 a_j=\left\lfloor\log_p(3d_j+2)\right\rfloor,
 \qquad e_j=a_j+\left\lfloor\frac{a_j}{2}\right\rfloor.
\]
Here the logarithm only denotes the integer characterized by
$p^{a_j}\le3d_j+2<p^{a_j+1}$.

\begin{theorem}\label{VIII-thm:e-log-bound}
For every $j\ge0$ and every $x\in\mathbb Z_p$,
\[
 \boxed{v_p(P_j(x))\ge-e_j.}
\]
Equivalently, $p^{e_j}P_j$ is integer-valued on $\mathbb Z_p$.
Moreover $e_0=0$, and $e_j\le j-1$ for every $j\ge1$.
\end{theorem}
This is an all-degree theorem with a logarithmic loss. It does not
assert the stronger, still unproved statement
$P_j\in\operatorname{Int}(\mathbb Z_p)$ for all $j$.

\subsection{An integral-coefficient formula before normalization}
For $h\ge0$, put
\begin{align*}
 \widehat q_h(t)&=(1-(3h+1)t)(1-(3h+2)t)(1+(6h+2)t),\\
 \widehat d_h(t)&=(1-(6h+6)t)(1+(3h+3)t)(1+(3h+2)t).
\end{align*}
For $0\le k\le m$, define integers $\phi_j(k,m)$ by
\[
 \frac{\prod_{h=k}^{m-1}\widehat d_h(t)}
      {\prod_{h=k}^{m}\widehat q_h(t)}
 =\sum_{j\ge0}\phi_j(k,m)t^j.
\]
All these coefficients are integers because every denominator
polynomial has constant term one. The exact normalized moment formula is
\begin{equation}\label{VIII-eq:e-integral-phi}
 P_j(3m+2)=\frac{2}{L(L-1)b_m}
          \sum_{k=0}^m(6k+1)b_k\phi_j(k,m),\qquad L=3m+2.
\end{equation}
To verify it directly, expand the proper rational function
$F_m=D_mS_m/Q_m$ term by term at $z=\infty$, using $t=1/z$.
The $k$th summand is
\[
 (-1)^k(6k+1)b_k^3
 \frac{\prod_{h=k}^{m-1}d_h(z)}{\prod_{h=k}^{m}q_h(z)}.
\]
Its leading coefficient is
$(-1)^m2L(L-1)^2b_m^2(6k+1)b_kz^{-3}$.
Dividing the reversed numerator and denominator polynomials by
their leading coefficients gives exactly the series defining
$\phi_j(k,m)$. Comparison with
$F_m=-C(L)\sum_{j\ge0}P_j(L)z^{-j-3}$ proves
\eqref{VIII-eq:e-integral-phi}. In particular this argument holds for
every coefficient, not only for any prescribed finite range.

\subsection{Valuation of the normalizing denominator}
\begin{lemma}\label{VIII-lem:e-normalization-bound}
Let $m\ge0$, $L=3m+2$, and $a=\lfloor\log_p L\rfloor$. Then
\[
 b_k\in\mathbb Z_p\quad(k\ge0),\qquad
 v_p\bigl(L(L-1)b_m\bigr)\le a+\lfloor a/2\rfloor.
\]
\end{lemma}
\begin{proof}
At the level $P=p^\ell$, let $u_\ell$ be the least nonnegative
solution of $3u_\ell+1\equiv0\pmod P$. Explicitly,
\[
 u_\ell=\begin{cases}(P-1)/3,&\ell\text{ even},\\
                       (2P-1)/3,&\ell\text{ odd}.
          \end{cases}
\]
The difference between the counts of multiples of $P$ in
$\prod_{h=0}^{m-1}(3h+1)$ and in $m!$ is
\[
 \delta_\ell=\begin{cases}1,&m\bmod P>u_\ell,\\0,&m\bmod P\le u_\ell.
                  \end{cases}
\]
The factors of $3$ in $b_m$ are units. Thus
$v_p(b_m)=\sum_{\ell\ge1}\delta_\ell\ge0$, proving integrality.
Also $\delta_\ell=0$ for $p^\ell>L$: all positive numerator
factors are at most $L-4$ and all denominator indices are at most $m$.

If $t=v_p(L)>0$, then for $\ell\le t$ the congruence
$3m\equiv-2\pmod{p^\ell}$ gives
\[
 m\bmod p^\ell=\begin{cases}
 2(p^\ell-1)/3,&\ell\text{ even},\\
 (p^\ell-2)/3,&\ell\text{ odd}.
 \end{cases}
\]
Consequently $\delta_\ell=1$ precisely at the even levels among
$1,\ldots,t$. Since $L-1$ is a unit,
\[
 v_p(L(L-1)b_m)
 \le t+\lfloor t/2\rfloor+(a-t)
 =a+\lfloor t/2\rfloor\le a+\lfloor a/2\rfloor.
\]
If instead $u=v_p(L-1)>0$, then $m\bmod p^\ell=u_\ell$ for
$\ell\le u$. All those $\delta_\ell$ are zero, and
$v_p(L(L-1)b_m)\le u+(a-u)=a$.
If both $L$ and $L-1$ are units, the same bound follows from
$\sum\delta_\ell\le a$. These cases include $m=0$, where $L=2$
and $a=0$.
\end{proof}

\begin{proof}[Proof of Theorem~\ref{VIII-thm:e-log-bound}]
Formula \eqref{VIII-eq:e-integral-phi} and the preceding lemma imply, for
every nonnegative integer $m$,
\[
 v_p(P_j(3m+2))\ge
 -\left(\lfloor\log_p(3m+2)\rfloor+
 \left\lfloor\frac{\lfloor\log_p(3m+2)\rfloor}{2}\right\rfloor\right).
\]
Use this for $m=0,1,\ldots,d_j$. Every one of these values lies in
$p^{-e_j}\mathbb Z_p$. Set $R_j(X)=P_j(3X+2)$, whose degree is at
most $d_j$. Its exact Newton expansion is
\[
 R_j(X)=\sum_{k=0}^{d_j}\Delta^kR_j(0)\binom Xk.
\]
Each finite difference is an integer linear combination of the
sampled values, so it too belongs to $p^{-e_j}\mathbb Z_p$.
For every $X\in\mathbb Z_p$, the binomial polynomial $\binom Xk$
is $p$-integral, by continuity from the nonnegative integers.
Taking $X=(x-2)/3\in\mathbb Z_p$ proves the uniform bound.
No division by $k!$ is being applied to the finite-difference
coefficient: that factorial belongs to the integer-valued basis.

Finally, for $j\ge1$ the elementary inequality
\[
 6\lfloor2j/3\rfloor+2<5^{\lfloor(2j-1)/3\rfloor+1}
\]
follows by checking $j=1,2,3$ and then increasing $j$ by three;
the left side increases by $12$, while the right side is multiplied
by $25$. Since $p\ge5$, this implies
$a_j\le\lfloor(2j-1)/3\rfloor$, and hence
$a_j+\lfloor a_j/2\rfloor\le j-1$. Also $d_0=0$, so $e_0=0$.
\end{proof}

\subsection{The double-strip transport is controlled at every higher odd level}
\begin{corollary}\label{VIII-cor:e-double-strip-transport-bound}
Let $r\ge3$ be odd, $M=p^{r-2}$, and let $c_j\in\mathbb Q_p$ satisfy
$v_p(c_j)\ge j$ for every $j\ge1$. Then the series
\[
 V_r=p^2\sum_{j\ge1}c_jM_j(M)
\]
converges. For every odd $r\ge5$,
\[
 V_r\in p^{3r-1}\mathbb Z_p.
\]
At the remaining level $r=3$, one has the precise reduction
\[
 V_3\equiv p^2c_3M_3(p)\pmod{p^8}.
\]
\end{corollary}
\begin{proof}
The first two moments vanish exactly. For $j\ge3$, put $n=j-3$.
The exact primary moment normalization and the proven valuation
$v_p(C(p^{r-2}))=(7r-17)/2$ give
\[
 v_p\bigl(p^2c_jM_j(M)\bigr)
 \ge2+j+\frac{7r-17}{2}-e_n
 =\frac{7r-7}{2}+n-e_n.
\]
The right side tends to infinity, because $e_n=O(\log n)$, so
convergence follows. For $n=0$ it is at least $(7r-7)/2$;
for $n\ge1$ it is at least $(7r-5)/2$.
When $r\ge5$, the former bound is at least $3r-1$. When $r=3$,
all terms with $n\ge1$ belong to $p^8\mathbb Z_p$, leaving only
the displayed cubic term.
\end{proof}


\section{Polynomial structure of every companion moment}\label{loc-VIII-6}
\label{VIII-sec:e-companion-polynomials}
Let $\widetilde M_j(M)$ be the companion moments from the proved
negative-parameter companion identity, and put
\[
 C_n=\frac{(-1)^{n+1}M^2}{2}
       \left(\frac{(2/3)_n}{n!}\right)^3,\qquad M=3n+1.
\]
\begin{proposition}\label{VIII-prop:e-companion-polynomials}
There exist even polynomials $R_j(M)\in\mathbb Q[M^2]$, $j\ge0$,
with
\[
 \deg R_j\le2\left\lfloor\frac{2j+1}{3}\right\rfloor,
 \qquad R_j(1)=1,
\]
such that, for every integer $n\ge0$,
\[
 \widetilde M_{j+1}(3n+1)=C_nR_j(3n+1).
\]
Their first values are
$R_0=1$, $R_1=(M^2+1)/2$, and $R_2=2M^2-1$.
\end{proposition}
\begin{proof}
Start with the already proved rational-function recurrence
$A_n(z)F_n(z)-B_n(z)F_{n-1}(z)=G_n(z)$, where $G_n$ has degree two.
Its first three Laurent equations give the three initial moments
just stated. Set
\[
 A=M(M-2),\quad B=(M-3)(M-1),\quad
 a=(M-2)(3M-4),\quad c=(M-1)(3M-5).
\]
The monic forms of its cubic factors are
\[
 A_n(z)=\frac{z^3-2z^2-az+2M(M-2)^2}{2M(M-2)^2},
\]
\[
 B_n(z)=-\frac{z^3-2z^2-cz-2(M-3)(M-1)^2}
                       {2(M-3)(M-1)^2}.
\]
Use the exact ratio
\[
 \frac{C_{n-1}}{C_n}
 =-\frac{(M-3)^2(M-1)^3}{M^2(M-2)^3}.
\]
The homogeneous Laurent equations, starting at the coefficient of
$z^{-1}$, yield, for every $j\ge3$,
\begin{align*}
 A R_j(M)-B R_j(M-3)
 &=2\bigl[A R_{j-1}(M)-B R_{j-1}(M-3)\bigr]\\
 &\quad+A a R_{j-2}(M)-B c R_{j-2}(M-3)\\
 &\quad-2A M(M-2)^2R_{j-3}(M)
       -2B(M-3)(M-1)^2R_{j-3}(M-3).
\end{align*}
We show this recurrence uniquely defines even rational polynomials.
The operator
\[
 \mathcal L'(P)=M(M-2)P(M)-(M-3)(M-1)P(M-3)
\]
maps $\mathbb Q[M^2]$ to $(2M-3)\mathbb Q[M(M-3)]$.
Indeed its image is anti-invariant under $M\mapsto3-M$.
For the basis element $M^{2k}$, division of its image by $2M-3$
gives a polynomial of degree $k$ in $M(M-3)$, with leading
coefficient $3k+1$. Thus $\mathcal L'$ is an isomorphism of the
filtered vector spaces over $\mathbb Q$.
The right side of the displayed recurrence is anti-invariant when
the preceding polynomials are even. It therefore has a unique even
preimage. Its degree is bounded by
\[
 \deg R_j\le\max\{\deg R_{j-1},\deg R_{j-2}+2,
                                    \deg R_{j-3}+4\}.
\]
The given initial degrees imply the asserted bound.
At $M=1$ the recurrence reduces to
\[
 R_j(1)=2R_{j-1}(1)+R_{j-2}(1)-2R_{j-3}(1),
\]
so all these values are one. This agrees with the exact base
$F_0(z)=-1/[2(1-z)]$, whose every positive moment equals $-1/2$.
Finally the Laurent recurrence determines the moments uniquely
successively in the cutoff and the moment index, since its leading
cubic coefficient is nonzero for $n\ge1$. Hence the constructed
polynomials give the actual moments for all $n$ and $j$.
\end{proof}
The diagonal $3k+1$ need not be a $p$-adic unit, so this proposition
asserts rational-polynomial structure and a uniform degree bound;
it does not assert coefficient integrality or integer-valuedness.

\section{A uniform logarithmic loss for all companion moments}\label{loc-VIII-7}
\label{VIII-sec:E7companionlog}
Use the polynomials $R_j$ of Proposition~\ref{VIII-prop:e-companion-polynomials},
so that
\[
 \widetilde M_{j+1}(3n+1)=C_nR_j(3n+1),\qquad
 d_j:=2\left\lfloor\frac{2j+1}{3}\right\rfloor\ge\deg R_j.
\]
For $p>3$, $p\equiv2\pmod3$, define
\[
 a_j=\left\lfloor\log_p(3d_j+1)\right\rfloor,
 \qquad \eta_j=a_j+\left\lceil\frac{a_j}{2}\right\rceil
 =\left\lceil\frac{3a_j}{2}\right\rceil.
\]
\begin{theorem}\label{VIII-E7:companion-log}
For every $j\ge0$ and every $x\in\mathbb Z_p$,
\[
 v_p(R_j(x))\ge-\eta_j.
\]
In particular, if $M=p^{2s}$, then for every $j\ge1$,
\[
 v_p(\widetilde M_j(M))\ge7s-e_p(j),\qquad
 e_p(j)=\left\lceil\frac32\left\lfloor\log_p(4j+1)\right\rfloor\right\rceil.
\]
\end{theorem}
\begin{proof}
We first prove a finite-value formula with integral coefficients before
normalization. Write $w=1/z$. For each integer root $a$, reversing the
factor $1-z/a$ gives its leading monomial times $1-aw$.
Thus the reversed products in $D_n/Q_n$ and in each $K_k$ have integer
coefficients as formal series in $w$: their denominator factors all have
constant term one and integral coefficients. The additional factor
$(z+18k^2-6k-2)/(z-2)$ becomes
\[
 \frac{1+(18k^2-6k-2)w}{1-2w}\in\mathbb Z[[w]].
\]
The leading coefficient of $K_k$ is $(-1)^k/d_k^2$, where
$d_k=(-1/3)_k/k!$. Also
\[
 \frac{D_n(z)}{Q_n(z)}
 =-\nu_nw\,(1+O(w)),\qquad
 \nu_n=(-1)^nM c_n^2,
 \quad M=3n+1,\quad c_n=(2/3)_n/n!.
\]
It follows from the proved summand formula for $V_n$ that there exist
integers $\phi_j(k,n)$, for every $j\ge0$, such that
\[
 F_n(z)=-\frac{\nu_n}{2}\sum_{j\ge0}
 \left(\sum_{k=0}^n(6k-1)d_k\phi_j(k,n)\right)w^{j+1}.
\]
Here $\phi_j$ is simply the coefficient of the product of the reversed
factors just described, so integrality holds for every degree, not only
for a finite initial range. Since $C_n=-\nu_n M c_n/2$, comparison of
Laurent coefficients proves the exact formula
\begin{equation}\label{VIII-E7:companion-integral-values}
 R_j(3n+1)=-\frac1{M c_n}
 \sum_{k=0}^n(6k-1)d_k\phi_j(k,n).
\end{equation}

Both $d_k$ and $c_k$ belong to $\mathbb Z_p$: for any
$\alpha\in\mathbb Z_p$, $(\alpha)_k/k!=(-1)^k\binom{-\alpha}{k}$
is integral. Put $b=\lfloor\log_p M\rfloor$ and $t=v_p(M)$.
At the level $q=p^\ell$, the excess numerator count in $c_n$
is zero or one; it is zero when $q>M$.
For $\ell\le t$, the residue of $n=(M-1)/3$ is
$(q-1)/3$ when $\ell$ is even, and $(2q-1)/3$ when $\ell$ is odd.
The first numerator index for $3k+2\equiv0\pmod q$ is respectively
$(2q-2)/3$ and $(q-2)/3$. Hence exactly the odd levels among
$1,\ldots,t$ contribute an excess. The remaining $b-t$ levels each
contribute at most one. Consequently
\[
 v_p(Mc_n)\le t+\lceil t/2\rceil+(b-t)
 \le b+\lceil b/2\rceil.
\]
Equation~\eqref{VIII-E7:companion-integral-values} therefore bounds every
finite value $R_j(3n+1)$ by the negative of this quantity.

Apply this bound at $n=0,1,\ldots,d_j$. All these values belong to
$p^{-\eta_j}\mathbb Z_p$. The polynomial $T_j(X)=R_j(3X+1)$ has
exact Newton expansion
\[
 T_j(X)=\sum_{k=0}^{d_j}\Delta^kT_j(0)\binom Xk.
\]
Each finite difference is an integer linear combination of the sampled
values. Each $\binom Xk$ is integral on $\mathbb Z_p$, by continuity
from the nonnegative integers. Taking $X=(x-1)/3$ proves the uniform
bound for $R_j(x)$.

Finally, at $M=p^{2s}$ the excess-count argument gives
$v_p(c_n)=s$ exactly, so $v_p(C_n)=7s$.
For the moment with index $j$, its polynomial is $R_{j-1}$, and
\[
 3d_{j-1}+1
 =6\left\lfloor\frac{2j-1}{3}\right\rfloor+1
 \le4j-1\le4j+1.
\]
Thus $\eta_{j-1}\le e_p(j)$, proving the last assertion.
\end{proof}


\section{A uniform lower bound for every companion interpolation mass}\label{loc-VIII-8}
\label{VIII-sec:E7companionmass}
Let $p>3$, $p\equiv2\pmod3$, $s\ge1$, $M=p^{2s}$ and
$n=(M-1)/3$. Retain the companion grids, denominator roots, and models
from Theorem~\ref{VIII-E7:companion-exact}, and put
\[
 \mu_M(H)=\widetilde W_n(H)U(H).
\]
\begin{proposition}\label{VIII-E7:companion-mass-bound}
For every companion node $H$, with $t=v_p(H)$,
\[
 \boxed{\quad v_p(\mu_M(H))\ge3s-\lceil t/2\rceil\ge2s.\quad}
\]
Zero masses have valuation $+\infty$. In particular every unit-node
mass belongs to $p^{3s}\mathbb Z_p$.
\end{proposition}
\begin{proof}
We may suppose $U(H)\ne0$ and $H\notin\mathcal D_n$; otherwise the
mass is zero. At an active node, $U(H)/H$ is one of
$1,1/2,-1/2$, hence is a unit. Direct multiplication of the grids gives
\[
 \frac{\prod_{K\in\mathcal H_n}K}{\prod_{D\in\mathcal D_n}D}
 =(-1)^nM\left(\frac{(2/3)_n}{n!}\right)^2.
\]
The Pochhammer ratio has valuation $s$: at each odd level
$p^\ell$, $1\le\ell\le2s$, its numerator has one extra divisible
factor, and at each even level its numerator and denominator counts
agree. Thus the displayed normalization has valuation $4s$.
The exact mass, up to an irrelevant sign, is this normalization times
\[
 \frac{U(H)}H
 \frac{\prod_{D\in\mathcal D_n}(D-H)}
 {\prod_{K\in\mathcal H_n,K\ne H}(K-H)}.
\]
All nonzero differences here have absolute value less than $4M$.
Since $p\ge5$, no power $p^\ell$ with $\ell>2s$ divides any of them.

For $q=p^\ell$ and $x=H\bmod q$, define
\[
 \delta_q(x)=1+\#\{D\in\mathcal D_n:D\equiv x\pmod q\}
 -\#\{K\in\mathcal H_n:K\equiv x\pmod q\}.
\]
The initial one removes the node $K=H$ from the denominator count.
Consequently
\begin{equation}\label{VIII-E7:mass-count}
 v_p(\mu_M(H))=4s+\sum_{\ell=1}^{2s}\delta_{p^\ell}(H\bmod p^\ell).
\end{equation}
All three arithmetic progressions have steps prime to $q$.
Their complete blocks of $q$ consecutive indices contribute equally
to the numerator and denominator counts and cancel. If $\ell$ is
even, the remaining cutoff is $b=(q-1)/3$; if $\ell$ is odd,
it is $b=(2q-1)/3$. The residual grids and denominators are exactly
the companion sets with $n=b$, respectively $M'=q$ or $M'=2q$.

Here are the complete residual counts for nonzero residues
$1\le x<q$, listed by $x\bmod6$:
\[
\begin{array}{c|rrrrrr}
x\bmod6&0&1&2&3&4&5\\\hline
\#K\ (q\equiv1\pmod6)&0&1&1&1&2&1\\
\#D\ (q\equiv1\pmod6)&2&1&0&1&1&1\\
\delta_q\ (q\equiv1\pmod6)&3&1&0&1&0&1\\\hline
\#K\ (q\equiv5\pmod6)&2&2&2&1&2&3\\
\#D\ (q\equiv5\pmod6)&2&3&2&2&2&1\\
\delta_q\ (q\equiv5\pmod6)&1&2&1&2&1&-1
\end{array}
\]
These counts follow directly by solving the three progression
congruences in the residual ranges. For example, in the even case
$M'=q$, the negative node grid contributes once for residues
$x\equiv3,4\pmod6$, while the positive denominator grid contributes
once for $x\equiv0,5\pmod6$. In the odd case $M'=2q$, those same
contributions are respectively $(1,1,0,0,1,1)$ and
$(1,1,1,1,0,0)$ across the six residue classes. The other two grids
have step three, yielding the displayed totals.
At the zero residue, the even case has one node and no denominator
root, so $\delta_q(0)=0$; the odd case has three nodes and no
denominator root, so $\delta_q(0)=-2$.

Thus every even level contributes a nonnegative amount. Each odd
level contributes at least $-1$, with a possible extra loss of one
precisely when $p^\ell\mid H$. There are $s$ odd levels altogether
and $\lceil t/2\rceil$ of them divide $H$; here $t\le2s$ because
$0<|H|<2M<p^{2s+1}$. Equation~\eqref{VIII-E7:mass-count} proves
\[
 v_p(\mu_M(H))\ge4s-s-\lceil t/2\rceil,
\]
as asserted.
\end{proof}

This is an individual-mass estimate. It is independent of the later
cancellation argument for E.3 and can be combined with exact moment formulas or
polynomial projectors without assuming any positive-moment integrality.

\subsection{The corresponding primary mass bound}
Let $L=p^{2s+1}$, $m=(L-2)/3$, and let $\mu_L(H)$ be the primary
three-grid mass. Then
\begin{equation}\label{VIII-E7:primary-mass-bound}
 \boxed{v_p(\mu_L(H))\ge3s+1-\lfloor v_p(H)/2\rfloor\ge2s+1.}
\end{equation}
Again zero masses are harmless. We give the full analogous count,
so this bound does not rely on numerical interpolation tests.
The primary normalization is
\[
 \frac{\prod K}{\prod D}
 =(-1)^{m+1}2L(L-1)^2\left(\frac{(1/3)_m}{m!}\right)^2.
\]
Its valuation is $4s+1$, since the Pochhammer ratio has valuation $s$.
At an active node $T(H)/H$ is a unit, and all nonzero differences
have absolute value less than $4L<pL$.
The same formula $1+\#D_x-\#K_x$ gives each level's correction.
For odd $\ell$, the residual cutoff is $(q-2)/3$ with $q=p^\ell$;
for even $\ell$, it is $(2q-2)/3$. The corresponding residual primary
endpoint is $q$ or $2q$ respectively. Complete blocks cancel as before.
For $1\le x<q$, the exact tables are
\[
\begin{array}{c|rrrrrr}
x\bmod6&0&1&2&3&4&5\\\hline
\#K\ (\ell\text{ odd})&0&1&2&1&1&1\\
\#D\ (\ell\text{ odd})&2&1&1&1&0&1\\
\delta_q\ (\ell\text{ odd})&3&1&0&1&0&1\\\hline
\#K\ (\ell\text{ even})&2&3&2&1&2&2\\
\#D\ (\ell\text{ even})&2&1&2&2&2&3\\
\delta_q\ (\ell\text{ even})&1&-1&1&2&1&2
\end{array}
\]
For example, at an odd level the negative node grid contributes
in residue classes $2,3\pmod6$, while the positive denominator grid
contributes in classes $0,1\pmod6$. At an even level their respective
six-class contributions are $(1,1,1,0,0,1)$ and $(1,0,0,1,1,1)$.
The remaining step-three grids give the totals shown above.
At $x=0$ the correction is zero for odd $\ell$ and $-2$ for even
$\ell$: there are respectively one or three nodes and no denominator
root at that residue. Thus odd levels cost nothing, while each of the
$s$ even levels costs at most one, plus one more when it divides $H$.
There are $\lfloor v_p(H)/2\rfloor$ such divisible even levels.
Subtracting these losses from $4s+1$ proves
\eqref{VIII-E7:primary-mass-bound}. In particular unit-node masses belong
to $p^{3s+1}\mathbb Z_p$.

\subsection{The unit-node contribution at level three}
\begin{corollary}\label{VIII-E7:primary-unit-r3}
For every prime $p>3$ with $p\equiv2\pmod3$,
\[
 \sum_{H\in\mathcal H_{(p^3-2)/3},\ p\nmid H}\mu_{p^3}(H)
 \in p^8\mathbb Z_p.
\]
\end{corollary}
\begin{proof}
First assume $p\ge41$. Let $g(Y)=3Y^2-2Y^3$ and $\Pi(X)=g(g(X^{p-1}))$. This polynomial
has integer coefficients, no terms of degree below $4(p-1)$, and
degree $9(p-1)$. For a unit integer $H$, Fermat's congruence and
$g(1+u)=1-3u^2-2u^3$ give $\Pi(H)\equiv1\pmod{p^4}$.
For $p\mid H$, one has $\Pi(H)\in p^{4(p-1)}\mathbb Z_p$.
The preceding proposition, with $s=1$, therefore gives
\[
 \sum_{p\nmid H}\mu_{p^3}(H)
 \equiv\sum_H\mu_{p^3}(H)\Pi(H)\pmod{p^8}.
\]
The right-hand side is an integer linear combination of the moments
$M_j(p^3)$ with $4(p-1)\le j\le9(p-1)$. Their common normalization
$C(p^3)$ has valuation nine. For every such $j$, the logarithmic
moment bound in Theorem~\ref{VIII-thm:e-log-bound} has loss at most one:
indeed, putting $n=j-3$,
\[
 3d_n+2=6\lfloor2n/3\rfloor+2\le4j-10
 \le36(p-1)-10<p^2 \qquad(p\ge41).
\]
Consequently $M_j(p^3)\in p^8\mathbb Z_p$.

The remaining primes are exactly $5,11,17,23,29$. They can be checked
by finite integer arithmetic in the exact mass formula above. Here is
an explicit modular certificate, computed modulo $p^{10}$:
\[
\begin{array}{c|rrrrr}
p&5&11&17&23&29\\\hline
p^{-9}\sum_{p\nmid H}\mu_{p^3}(H)\pmod p&1&6&6&18&10.
\end{array}
\]
In particular every one of those sums lies in $p^9\mathbb Z_p$.
The reproducible certificate is \nolinkurl{fast_primary_masses.py}:
it evaluates the displayed product formula modulo $p^{10}$,
removing the exact power of $p$ from every factor before taking a
modular inverse. Prefix products in each residue class modulo three
and six evaluate every arithmetic progression in the formula. A zero
factor is omitted only at the unique denominator occurrence $K=H$;
numerator zeroes are excluded by testing $H\in\mathcal D_m$ first.
Thus all modular inverses are applied to units and the finite checks
are exact. They complete the five exceptional cases.
\end{proof}

The corollary controls only the unit-node contribution. The shell
$v_p(H)=1$ and its cancellation with the cubic double-strip term
remain separate obligations.

\section{One-step transport of the valuation-one stratum}\label{loc-VIII-9}
\label{VIII-sec:E7one-step}
Let $L=p^{2s+1}$, $M=L/p=p^{2s}$, $p>3$, $p\equiv2\pmod3$.
Write $\mu_L(H)$ for the primary interpolation masses and
$\widetilde\mu_M(J)$ for the proved companion masses.
Let $\mathcal H_L,\mathcal D_L$ denote the primary root sets and set
\[
 R^{(1)}_{L,p}(X)=
 \frac{\prod_{d\in\mathcal D_L,\,p\nmid d}(1-pX/d)}
 {\prod_{h\in\mathcal H_L,\,p\nmid h}(1-pX/h)}
 =\sum_{j\ge0}a_jX^j.
\]
Every retained root is a unit, so $a_0=1$ and $a_j\in p^j\mathbb Z_p$.
The series is restricted and converges at every integral node.

The exact one-step product identity is
\begin{equation}\label{VIII-E7:one-step-masses}
 \mu_L(pJ)=p\widetilde\mu_M(J)R^{(1)}_{L,p}(J).
\end{equation}
Indeed dividing the divisible primary nodes by $p$ gives respectively
\[
 \{2,5,\ldots,M-2\},\quad
 \{1,4,\ldots,M\},\quad
 -2\{2,5,\ldots,M-2\},
\]
exactly the companion nodes. Dividing the divisible denominator roots
gives $6j,-3j$ for $1\le j\le(M-1)/3$ and
$-1,-4,\ldots,-(M-3)$, exactly the companion denominator roots.
For the first two denominator grids this follows from
$\lfloor(L-2)/(3p)\rfloor=(M-1)/3$; the last grid has fine endpoint
$-(L-3)$, so its largest divisible quotient is $-(M-3)$.
Since $p\equiv5\pmod6$, the primary model formulas give
$T(pJ)=pU(J)$ on each companion grid, including zero model values.
Separating the divisible and unit factors in the interpolation products
proves \eqref{VIII-E7:one-step-masses}; no division by a zero weight is used.

In particular, the omitted part of the two-step transport has the exact
split
\begin{equation}\label{VIII-E7:U-one-step}
 U_{L,p}
 =\sum_{p\nmid H}\mu_L(H)
 +p\sum_{p\nmid J}\widetilde\mu_M(J)R^{(1)}_{L,p}(J).
\end{equation}
This identity places the formerly correlated valuation-one stratum
inside a weighted unit sum of the exact companion family.


\section{The original odd E.3 congruence at all sufficiently high levels}\label{loc-VIII-10}
\label{VIII-sec:e-high-level-closure}
Let $p>3$ be a prime with $p\equiv2\pmod3$, and define
\[
 E(L)=\sum_{k=0}^{(L-2)/3}(-1)^k(6k+1)
          \left(\frac{(1/3)_k}{k!}\right)^3
 \qquad(L\equiv2\pmod3).
\]
All interpolation identities and moment bounds used below are exact
or have been proved in the preceding sections. In particular, no
assertion of integral normalized moments of every degree is required.

\begin{theorem}\label{VIII-thm:e-high-level-closure}
Put $r=2s+1$, with $s\ge2$, and suppose
\begin{equation}\label{VIII-eq:e-improved-threshold}
 \left\lceil\frac32\left\lfloor
 \log_p\bigl(4(3s+1)(p-1)+1\bigr)\right\rfloor\right\rceil
 \le s-1.
\end{equation}
Then the original E.3 congruence holds:
\[
 \boxed{E(p^{2s+1})\equiv p^2E(p^{2s-1})\pmod{p^{6s+2}}.}
\]
In particular it holds in each of the following ranges:
\begin{enumerate}
 \item every $p>3$, $p\equiv2\pmod3$, and every odd $r\ge13$;
 \item every such $p\ge11$ and every odd $r\ge9$;
 \item every such $p\ge41$ and every odd $r\ge7$.
\end{enumerate}
\end{theorem}

\subsection{The two exact one-step identities}
Set $N=p^{2s+1}$, $M=p^{2s}$, and $K=p^{2s-1}$.
Write $\mu_N(H)$ and $\mu_K(H)$ for the primary three-grid masses,
and $\widetilde\mu_M(J)$ for the companion masses. Their sums are
$E(N)$, $E(K)$, and $\mathcal A(M)$, respectively. Their positive
moments are denoted by $M_j$ and $\widetilde M_j$.

For either grid family, form the quotient of unit denominator-root
factors by unit node factors, evaluated at $pX$:
\[
 R(X)=
 \frac{\prod_{d\in\mathcal D,\ p\nmid d}(1-pX/d)}
      {\prod_{h\in\mathcal H,\ p\nmid h}(1-pX/h)}.
\]
Let $R_0$ use the primary family at $N$, and $R_1$ the companion
family at $M$. Both have expansions
\begin{equation}\label{VIII-eq:e-kernel-coefficient-bound}
 R_i(X)=1+\sum_{j\ge1}a_{i,j}X^j,
 \qquad a_{i,j}\in p^j\mathbb Z_p\quad(i=0,1).
\end{equation}
These are restricted power series and converge at every $X\in\mathbb Z_p$.
Separating divisible and unit roots in the interpolation weights gives
\begin{equation}\label{VIII-eq:e-two-one-step-identities}
 \mu_N(pJ)=p\widetilde\mu_M(J)R_0(J),\qquad
 \widetilde\mu_M(pH)=p\mu_K(H)R_1(H).
\end{equation}
Here the correspondence between divisible nodes and the lower family
is exact, including endpoints and zero model values.

For completeness, the first correspondence is the one established in
\eqref{VIII-E7:one-step-masses}. For the reverse correspondence, divide the
divisible companion nodes by $p$. The positive grids become
$\{2,5,\ldots,K\}$ and $\{1,4,\ldots,K-1\}$, and the negative grid
becomes $-2\{1,4,\ldots,K-1\}$, precisely the primary grids at $K$.
The divided denominator roots are
\[
 \{6j,-3j:1\le j\le(K-2)/3\}
 \ \cup\ \{-2,-5,\ldots,-(K-3)\},
\]
also precisely the primary denominator roots. The explicit models
satisfy $U(pH)=pT(H)$ on these grids because $p\equiv5\pmod6$.
Thus the same product separation proves the second identity in
\eqref{VIII-eq:e-two-one-step-identities}, without dividing by a model
value that might vanish.

\subsection{Moment estimates in the forms needed for transport}
The proved primary normalizations and logarithmic bounds give
\begin{align}
 M_1(N)=M_2(N)&=M_1(K)=M_2(K)=0,\\
 v_p(M_j(N))&\ge7s+2-e_{j-3} &&(j\ge3),\\
 v_p(M_j(K))&\ge7s-5-e_{j-3} &&(j\ge3),
 \label{VIII-eq:e-coarse-moment-bound}
\end{align}
where $e_0=0$, $e_h\le h-1$ for $h\ge1$, and $e_h=O(\log(h+2))$.
The companion logarithmic bound gives
\[
 v_p(\widetilde M_j(M))\ge7s-\eta_{j-1}.
\]
In particular,
\begin{equation}\label{VIII-eq:e-companion-linear-loss}
 v_p(\widetilde M_j(M))\ge7s-(j-1)\qquad(j\ge1).
\end{equation}
For $j=1,2$ this follows from the exact normalized polynomials
$R_0^{\rm mom}=1$ and $R_1^{\rm mom}=(M^2+1)/2$. For the other
indices use $\eta_h\le h$ for $h\ge2$. To verify this last elementary
inequality, the companion definition is
\[
 \eta_h=\left\lceil\frac32
 \left\lfloor\log_p\left(6\left\lfloor\frac{2h+1}{3}\right\rfloor+1\right)
 \right\rfloor\right\rceil.
\]
The inequality
\[
 6\left\lfloor\frac{2h+1}{3}\right\rfloor+1
 <5^{\lfloor2h/3\rfloor+1}
\]
holds at $h=2,3,4$ and is preserved by increasing $h$ by three:
the left side increases by $12$, whereas the right side is multiplied
by $25$. Since $p\ge5$, it implies $\eta_h\le h$.

A common nondecreasing logarithmic bound is
\begin{equation}\label{VIII-eq:e-common-loss}
 \epsilon_p(j)=\left\lceil\frac32\left\lfloor
                       \log_p(4j+1)\right\rfloor\right\rceil.
\end{equation}
For $j\ge3$ it bounds $e_{j-3}$, and for $j\ge1$ it bounds
$\eta_{j-1}$. Consequently the positive moments of the two fine
families satisfy
\[
 v_p(M_j(N))\ge7s+2-\epsilon_p(j),\qquad
 v_p(\widetilde M_j(M))\ge7s-\epsilon_p(j).
\]
The first inequality also holds for $j=1,2$, whose moments vanish.
The proved individual unit-node estimates are
\begin{equation}\label{VIII-eq:e-unit-mass-bounds}
 v_p(\mu_N(H))\ge3s+1\ (p\nmid H),\qquad
 v_p(\widetilde\mu_M(J))\ge3s\ (p\nmid J).
\end{equation}

\subsection{A low-degree unit projector and its leakage}
Put
\[
 t=3s+1,\qquad D=t(p-1),\qquad
 \chi_t(X)=1-(1-X^{p-1})^t\in\mathbb Z[X].
\]
Its nonzero monomials have degrees $i(p-1)$ for $1\le i\le t$,
so its smallest degree is $p-1\ge4$ and its degree is $D$.
For every integral unit $x$, Fermat's congruence gives
$1-\chi_t(x)\in p^t\mathbb Z_p$.
Let $B=\epsilon_p(D)$; the theorem assumes $B\le s-1$.

First consider the primary unit sum
$U_N=\sum_{p\nmid H}\mu_N(H)$.
Replacing $1$ by $\chi_t(H)$ on its unit nodes changes it by an
element of $p^{6s+2}\mathbb Z_p$, by
\eqref{VIII-eq:e-unit-mass-bounds}. The full-node polynomial sum satisfies
\[
 \sum_H\mu_N(H)\chi_t(H)\in p^{7s+2-B}\mathbb Z_p
                              \subseteq p^{6s+3}\mathbb Z_p.
\]
Its contribution from nonunit nodes is, by the first exact strip,
\[
 \sum_{p\mid H}\mu_N(H)\chi_t(H)
   =p\sum_J\widetilde\mu_M(J)R_0(J)\chi_t(pJ).
\]
Write $R_0(X)\chi_t(pX)=\sum_{j\ge p-1}b_jX^j$.
Every $b_j$ belongs to $p^j\mathbb Z_p$: both factors have this
coefficient property, and their convolution preserves it. Hence
\eqref{VIII-eq:e-companion-linear-loss} shows that each term of its
moment expansion has valuation at least
\[
 1+j+7s-(j-1)=7s+2.
\]
The expansion converges, since the sharper logarithmic bound makes
these valuations tend to infinity. We have proved
\begin{equation}\label{VIII-eq:e-primary-unit-sum}
 U_N\in p^{6s+2}\mathbb Z_p.
\end{equation}

Apply the same projector to the companion unit sum
$\widetilde U_M=\sum_{p\nmid J}\widetilde\mu_M(J)$.
The approximation error on unit nodes belongs to
$p^{3s+t}\mathbb Z_p=p^{6s+1}\mathbb Z_p$.
The full-node polynomial sum has valuation at least
$7s-B\ge6s+1$. Its nonunit leakage is
\[
 p\sum_H\mu_K(H)R_1(H)\chi_t(pH).
\]
Again the product series starts in degree $p-1\ge4$ and its
degree-$j$ coefficient has valuation at least $j$.
For all such $j$, $e_{j-3}\le j-4$. Therefore each leakage term has
valuation at least
\[
 1+j+(7s-5)-(j-4)=7s\ge6s+1.
\]
Convergence follows from the logarithmic loss as before. Thus
\begin{equation}\label{VIII-eq:e-companion-unit-sum}
 \widetilde U_M\in p^{6s+1}\mathbb Z_p.
\end{equation}
This calculation explicitly controls the nonunit leakage; the projector
is not being treated as a literal indicator at nonunit nodes.

\subsection{Completion of the two one-step congruences}
Sum the first identity in \eqref{VIII-eq:e-two-one-step-identities} and
add the primary unit nodes. It gives exactly
\[
 E(N)=U_N+p\mathcal A(M)
             +p\sum_{j\ge1}a_{0,j}\widetilde M_j(M).
\]
Every term of the last series has valuation at least $7s+2$ by
\eqref{VIII-eq:e-kernel-coefficient-bound} and
\eqref{VIII-eq:e-companion-linear-loss}. Thus
\begin{equation}\label{VIII-eq:e-primary-one-step-congruence}
 E(N)\equiv p\mathcal A(M)\pmod{p^{6s+2}}.
\end{equation}
The reverse strip similarly gives
\[
 \mathcal A(M)=\widetilde U_M+pE(K)
                         +p\sum_{j\ge3}a_{1,j}M_j(K),
\]
where the first two moments have vanished exactly. The $j=3$ term
has valuation at least $7s-1$; every $j\ge4$ term has valuation at
least $7s$. Since $s\ge2$, these are at least $6s+1$.
Together with \eqref{VIII-eq:e-companion-unit-sum}, this proves
\begin{equation}\label{VIII-eq:e-companion-one-step-congruence}
 \mathcal A(M)\equiv pE(K)\pmod{p^{6s+1}}.
\end{equation}
All series manipulations are legitimate: there are finitely many
nodes, the root-product series converge on $\mathbb Z_p$, and the
logarithmic moment estimates also give convergence after rearranging
by degree. Combining the two displayed congruences proves the
assertion of Theorem~\ref{VIII-thm:e-high-level-closure}.

\subsection{Verification of the explicit ranges}
For integer $s\ge2$, condition \eqref{VIII-eq:e-improved-threshold} is
equivalent to the strict inequality
\begin{equation}\label{VIII-eq:e-threshold-power-form}
 (12s+4)(p-1)+1
 <p^{\lfloor2(s-1)/3\rfloor+1}.
\end{equation}
For fixed $s$ let $A=12s+4$ and $a=\lfloor2(s-1)/3\rfloor+1$.
In each range considered below, $a\ge2$. The function
$(Ap-A+1)/p^a$ is decreasing for real $p\ge5$, because its
derivative has numerator $Ap(1-a)+a(A-1)<0$. It therefore suffices
to check the smallest value of $p$ in each claimed range.

For $p\ge5$, the left side at $p=5$ is $48s+17$.
At $s=6,7,8$ the required comparisons are
\[
 305<625,\qquad353<3125,\qquad401<3125.
\]
Increasing $s$ by three adds $144$ to the left side and multiplies
the right side by $25$, preserving the inequality. Thus it holds
for all $s\ge6$.

For $p\ge11$, the left side at $p=11$ is $120s+41$.
At $s=4,5,6$ the comparisons are
\[
 521<1331,\qquad641<1331,\qquad761<14641.
\]
The step $s\mapsto s+3$ adds $360$ to the left side and multiplies
the right side by $121$. Thus the inequality holds for all $s\ge4$.

For $p\ge41$, the left side at $p=41$ is $480s+161$.
At $s=3,4,5$ the comparisons are
\[
 1601<1681,\qquad2081<68921,\qquad2561<68921.
\]
The step $s\mapsto s+3$ adds $1440$ to the left side and multiplies
the right side by $1681$. This proves the final stated range.


\section{A reflection gain for companion interpolation samples}\label{loc-VIII-11}
\label{VIII-sec:e-companion-sample-gain}
Retain $d_k=(-1/3)_k/k!$, $c_n=(2/3)_n/n!$, $M=3n+1$,
and the exact normalized companion moments $R_j(M)$.

\begin{lemma}\label{VIII-lem:e-companion-numerator-gain}
Let $p>3$, $p\equiv2\pmod3$. For every $n\ge0$ with
$p\mid3n+1$, every coefficient of the integral numerator series
\[
 \mathcal N_n(w)=\sum_{k=0}^n(6k-1)d_k\phi_k(w),
 \qquad \phi_k(w)=\sum_{j\ge0}\phi_j(k,n)w^j,
\]
from \eqref{VIII-E7:companion-integral-values} belongs to $p\mathbb Z_p$.
\end{lemma}
\begin{proof}
Use the monic reversed product of the companion summand kernel:
\[
 \widehat K_k(w)=\prod_{j=0}^{k-1} f_j(w),\qquad
 f_j(w)=
 \frac{(1+(6j-2)w)(1-(3j-1)w)(1-(3j+1)w)}
 {(1-(6j+6)w)(1+(3j+3)w)(1+(3j+1)w)}.
\]
The coefficient series $\phi_k$ has the form
\[
 \phi_k(w)=A_n(w)\widehat K_k(w)
                \frac{1+(18k^2-6k-2)w}{1-2w},
 \qquad A_n(w)\in\mathbb Z[[w]],\quad A_n(0)=1.
\]
Here $A_n$ is the common reversed denominator-root/node-root
quotient. Its precise roots are immaterial for this argument;
every root is integral, so $A_n$ is an integral unit series.
It suffices to prove that
\[
 \sum_{k=0}^n(6k-1)d_k\widehat K_k(w)
                       (1+(18k^2-6k-2)w)
\]
vanishes in $\mathbb F_p[[w]]$.

Put $a=(p+1)/3$, which is an even integer. The binomial identity
$d_k=(-1)^k\binom{1/3}{k}$ and the $p$-adic Lucas congruence give
\begin{equation}\label{VIII-eq:e-companion-lucas}
 d_{pb+c}\equiv
 \frac{(1/3)_b}{b!}\,(-1)^c\binom ac\pmod p
 \qquad(0\le c<p).
\end{equation}
Indeed $1/3=a+p(-1/3)$ in $\mathbb Z_p$; the ordinary Lucas
identity extends to binomial coefficients with $p$-adic top argument
by continuity from nonnegative integers. This also proves that the
left side vanishes when $c>a$.

For any consecutive $p$ indices, each of the three linear root
progressions in the numerator and denominator of $f_j$ runs through
all residues of $\mathbb F_p$. Thus their products agree, and
\[
 \prod_{j=b p}^{(b+1)p-1}f_j(w)=1
       \quad\hbox{in }\mathbb F_p[[w]],
 \qquad \widehat K_{pb+c}(w)=\widehat K_c(w).
\]
The remaining linear and quadratic factors in $k$ depend only on
$k\bmod p$. Therefore every complete block reduces, up to the scalar
$(1/3)_b/b!$, to
\[
 B(w)=\sum_{c=0}^a(6c-1)(-1)^c\binom ac\widehat K_c(w)
                       (1+(18c^2-6c-2)w).
\]
The same holds for the final partial block: since $3n+1\equiv0$
modulo $p$, its last residue is $(2p-1)/3=2a-1\ge a$ and
therefore includes the entire nonzero Lucas support.

It remains to show $B(w)=0$. The exact root lists defining $f_j$
give, using $3a=1$ in $\mathbb F_p$,
\[
 f_{a-1-j}(w)=f_j(w)^{-1},\qquad
 \widehat K_a(w)=1,\qquad
 \widehat K_{a-c}(w)=\widehat K_c(w).
\]
The first equality pairs the numerator roots
$2-6(a-1-j),3(a-1-j)-1,3(a-1-j)+1$
with the denominator roots $6j+6,-3j-3,-3j-1$, and conversely.
Because $a$ is even, $(-1)^c\binom ac$ is unchanged under
$c\mapsto a-c$. Also
\[
 6(a-c)-1=-(6c-1),\qquad
 18(a-c)^2-6(a-c)-2=18c^2-6c-2
 \quad\hbox{in }\mathbb F_p.
\]
Consequently the summands cancel in reflected pairs. If there is a
middle summand, its factor $6c-1$ is zero. Hence $B(w)=0$, proving
the claimed coefficientwise divisibility.
\end{proof}

\begin{corollary}\label{VIII-cor:e-companion-improved-log}
With the notation of Theorem~\ref{VIII-E7:companion-log}, the bound improves to
\[
 v_p(R_j(x))\ge-\max\{0,\eta_j-1\}
 \qquad(j\ge0,\ x\in\mathbb Z_p).
\]
\end{corollary}
\begin{proof}
At a sample $M=3n+1$, put $b=\lfloor\log_p M\rfloor$.
If $p\nmid M$, its normalizing denominator $Mc_n$ has valuation
at most $b$. If $p\mid M$, the previous count and the new numerator
gain give a loss of at most $b+\lceil b/2\rceil-1$.
For $b\ge1$ the latter dominates $b$; for $b=0$ there is no loss.
Apply these bounds to the same Newton samples as in
Theorem~\ref{VIII-E7:companion-log}.
\end{proof}

\begin{corollary}\label{VIII-cor:e-companion-short-sample}
For every integer $n\ge0$ with $1\le M=3n+1<p^2$, and every $j\ge0$,
\[
 v_p(R_j(M))\ge-1.
\]
Consequently, if $3d_j+1<p^2$, where
$d_j=2\lfloor(2j+1)/3\rfloor$, then
\[
 pR_j(x)\in\mathbb Z_p\qquad\hbox{for every }x\in\mathbb Z_p.
\]
\end{corollary}
\begin{proof}
Use the exact formula
$R_j(M)=-(Mc_n)^{-1}[w^j]\mathcal N_n(w)$.
If $p\nmid M$, the excess-count proof gives $v_p(c_n)\le1$
since $M<p^2$. Thus its normalization costs at most one power of $p$.
If $p\mid M$, then $v_p(M)=v_p(c_n)=1$, so the normalizing
denominator costs two powers, but
Lemma~\ref{VIII-lem:e-companion-numerator-gain} restores one.
This proves the first assertion, including $M=1$.
For the second, apply it at the Newton interpolation samples
$M=3n+1$, $0\le n\le d_j$. Their maximum is below $p^2$.
All their finite differences belong to $p^{-1}\mathbb Z_p$,
and the binomial basis is integral on $\mathbb Z_p$.
\end{proof}

\begin{theorem}\label{VIII-thm:e-level-five-large-primes}
For every prime $p\ge29$ with $p\equiv2\pmod3$,
\[
 \boxed{E(p^5)\equiv p^2E(p^3)\pmod{p^{14}}.}
\]
\end{theorem}
\begin{proof}
Repeat the proof of Theorem~\ref{VIII-thm:e-high-level-closure} with
$s=2$, $t=7$, and $D=7(p-1)$. The leakage and one-step transport
estimates in that proof are valid for every $s\ge2$ and require
no logarithmic threshold. Only its two full-node projector sums
need the following improved estimates.

For each companion moment index $1\le j\le D$, its polynomial
is $R_{j-1}$ and its largest interpolation sample satisfies
\[
 3d_{j-1}+1
 =6\left\lfloor\frac{2j-1}{3}\right\rfloor+1
 \le4j-1\le28(p-1)-1<p^2.
\]
The last inequality holds for $p\ge29$. The preceding corollary
therefore gives
$v_p(\widetilde M_j(p^4))\ge14-1=13$.
The full-node companion projector sum belongs to $p^{13}\mathbb Z_p$,
as required.

The primary moments have their first two terms zero. For $3\le j\le D$,
the largest sample for $P_{j-3}$ is at most
$4(j-3)+2\le28(p-1)-10<p^2$.
The primary logarithmic bound is therefore at most one, giving
$v_p(M_j(p^5))\ge16-1=15$. This exceeds the required primary
projector precision $14$. The unchanged leakage calculations now
give $E(p^5)\equiv p\mathcal A(p^4)\pmod{p^{14}}$ and
$\mathcal A(p^4)\equiv pE(p^3)\pmod{p^{13}}$.
Their combination proves the theorem.
\end{proof}

\section{Exact finite certificates for the remaining higher-level cases}\label{loc-VIII-12}
\label{VIII-sec:E7finiteexceptions}
The analytic theorem leaves only the following pairs $(s,p)$ at
odd levels $r=2s+1\ge7$:
\[
 (3,5),(3,11),(3,17),(3,23),(3,29),(4,5),(5,5).
\]
This is a finite exhaustive list, not a numerical sample of the
remaining higher-level range. We give exact polynomial-interpolation
certificates for the moments required by the proved projector argument.

Put $t=3s+1$. The projector $1-(1-X^{p-1})^t$ uses precisely the
moment indices $j=i(p-1)$, $1\le i\le t$. Its proof only needs
\begin{equation}\label{VIII-E7:finite-required}
 v_p\bigl(R_{i(p-1)-1}(p^{2s})\bigr)\ge-(s-1),\qquad
 v_p\bigl(P_{i(p-1)-3}(p^{2s+1})\bigr)\ge-s.
\end{equation}
The first bound controls the full companion projector sum, and the
second controls the full primary projector sum. The bounds on unit-node
approximation, nonunit leakage, and the infinite transport corrections
remain the previously proved uniform estimates; none of those is
replaced by a finite truncation.

\subsection{The integer recurrences used for certification}
For the companion define a formal series
\[
 Y_n(w)=-M c_n\sum_{j\ge0}R_j(M)w^j,\qquad M=3n+1.
\]
The proved rational-function recurrence gives, for $n\ge1$,
\begin{align*}
 &(1-2w-(M-2)(3M-4)w^2+2M(M-2)^2w^3)Y_n(w)\\
 &\quad=(1-2w-(M-1)(3M-5)w^2
            -2(M-3)(M-1)^2w^3)Y_{n-1}(w)\\
 &\qquad\quad +(6n-1)d_n(1+w)(1+(18n^2-6n-2)w),
\end{align*}
with $Y_0(w)=-1/(1-w)$, $d_n=(-1/3)_n/n!$, and $c_n=(2/3)_n/n!$.
All polynomial coefficients are integers and the left constant term
is one. Thus each coefficient is determined using only addition and
multiplication and already known lower coefficients.

For the primary family let $L=3n+2$, $b_n=(1/3)_n/n!$, and
\[
 Z_n(w)=\frac{L(L-1)b_n}{2}\sum_{j\ge0}P_j(L)w^j.
\]
Its integer recurrence is
\begin{align*}
 &(1-w-(L-1)(3L-2)w^2+2L(L-1)^2w^3)Z_n(w)\\
 &\quad=(1-w-(L-2)(3L-7)w^2
            -2(L-2)^2(L-3)w^3)Z_{n-1}(w)+(6n+1)b_n.
\end{align*}
At $n=0$ take the preceding series as zero, giving
$Z_0(w)=1/(1-w-4w^2+4w^3)$.
This recurrence follows equally from the integral reversed-factor
formula for the primary normalized moments.

To compute a Pochhammer ratio modulo a power of $p$, each numerator
and denominator factor is written as a power of $p$ times a unit.
Only the unit denominator is inverted; the accumulated valuation is
tracked as an integer. This computes $b_n,c_n,d_n$ without inverting
$p$ modulo its own power.

\subsection{Degree bounds, precision, and exact Newton evaluation}
For any selected normalized index $j$, the proved degree bounds are
\[
 d_j^R=2\lfloor(2j+1)/3\rfloor,\qquad
 d_j^P=2\lfloor2j/3\rfloor.
\]
Therefore its values at the small cutoffs $n=0,\ldots,d_j$ determine
the desired polynomial exactly. Let $d_{\max}$ be the largest degree
needed for a fixed family and pair $(s,p)$. Choose the common scale
\[
 A_R=\left\lceil\frac32\lfloor\log_p(3d_{\max}+1)\rfloor\right\rceil,
 \qquad
 A_P=\left\lfloor\frac32\lfloor\log_p(3d_{\max}+2)\rfloor\right\rfloor,
\]
and compute modulo $p^{A+2}$ in the corresponding family.
The previously proved normalization bounds ensure
$v_p(Mc_n)\le A_R$ and $v_p(L(L-1)b_n)\le A_P$ throughout those samples.
If a normalization is $p^vu$ with $u$ a unit, the scaled sample is
computed as $p^{A-v}$ times the unnormalized coefficient times $u^{-1}$,
with the indicated factor $-1$ or $2$. No modular division by $p$ is
performed and no precision is lost: the exponent $A-v$ is nonnegative.

For either polynomial, let $f(n)$ denote its scaled value at the
appropriate affine cutoff. Its exact Newton formula is
\[
 f(X)=\sum_{k=0}^{d_j}\Delta^kf(0)\binom Xk.
\]
Here $X=(p^{2s}-1)/3$ for the companion and
$X=(p^{2s+1}-2)/3$ for the primary. Both are nonnegative integers.
The binomial coefficients are computed as exact integers before
reduction modulo $p^{A+2}$. All finite differences and final sums are
then integer operations in that explicitly specified residue ring.
This evaluates the actual polynomial value, not an approximation
obtained by guessing a polynomial from too few samples.

If the scaled residue is nonzero with valuation $v<A+2$, its unscaled
value has exactly valuation $v-A$. If the residue is zero, the
unscaled value has valuation at least two. Thus the residue arithmetic
certifies every inequality in \eqref{VIII-E7:finite-required}.

\subsection{Certificate output and exhaustion of the exceptional set}
The complete reproducible program is
\nolinkurl{work_e7_companion/finite_projector_certificates.py}; its
integer outputs are in the adjacent JSON file. It verifies 158 selected
moment indices across the seven pairs and both families. The following
table gives the largest small sample index and the minimum valuation
of the scaled residues in each family; the scale $A$ is displayed
alongside, so the minimum unscaled valuation is their difference.
\[
\begin{array}{c|c|r|c|r|r|c|r}
s&p&d_{\max}^R&A_R&\min v(p^{A_R}R)&d_{\max}^P&A_P&\min v(p^{A_P}P)\\\hline
3&5 &52 &5&5&48 &4&4\\
3&11&132&3&2&128&3&3\\
3&17&212&3&2&208&3&3\\
3&23&292&3&2&288&3&3\\
3&29&372&3&2&368&3&3\\
4&5 &68 &5&5&64 &4&4\\
5&5 &84 &5&5&80 &4&4
\end{array}
\]
The moduli used are respectively $p^{A_R+2}$ and $p^{A_P+2}$.
In particular every companion minimum is at least $-1$ after removing
the scale, and every primary minimum is at least zero. These are
stronger than the requirements in \eqref{VIII-E7:finite-required}.

Applying the already proved projector-and-leakage argument with these
158 certified moment bounds proves the unchanged E.3 congruence for
all seven exceptional pairs. Combining them with the analytic ranges
therefore gives the unconditional computationally certified conclusion
\[
 \boxed{E(p^r)\equiv p^2E(p^{r-2})\pmod{p^{3r-1}}
 \quad\text{for every }p>3,\ p\equiv2\pmod3,
 \text{ and every odd }r\ge7.}
\]
The finite computations establish the finite exceptional set only;
the assertion for every other prime and level comes from the uniform
analytic theorem. The fifth level is supplied by the reflection argument and the
additional finite certificates in the next section.

\section{Exact finite completion of the level \texorpdfstring{$r=5$}{r=5}}\label{loc-VIII-13}
\label{VIII-sec:E7r5finite}
The uniform reflection-gain argument proves the original $r=5$
congruence for every prime $p\ge29$ with $p\equiv2\pmod3$.
The only primes satisfying $3<p<29$ and $p\equiv2\pmod3$ are
$5,11,17,23$. This section certifies that exhaustive finite exception
set using the already proved polynomial-moment method.

Here $s=2$, so the projector $1-(1-X^{p-1})^7$ uses exactly the
seven positive moment indices $j=i(p-1)$, $1\le i\le7$.
The sufficient bounds from the proved projector-and-leakage argument are
\[
 v_p\bigl(R_{i(p-1)-1}(p^4)\bigr)\ge-1,
 \qquad
 v_p\bigl(P_{i(p-1)-3}(p^5)\bigr)\ge-2.
\]
All nonunit leakage and all infinite transport terms continue to use
the uniform estimates; the calculation below only supplies these
fourteen polynomial values per prime.

Use exactly the integer Laurent recurrences, proved degree bounds,
and scaled Newton interpolation described in
Section~\ref{VIII-sec:E7finiteexceptions}. In all eight family-prime cases,
the common scaling exponent is $A=3$ and the calculation is modulo
$p^5=p^{A+2}$. The normalization divisors at every sampled small cutoff
have valuation at most three. Thus multiplication by $p^{A-v}$ and
inversion of the unit part supplies each scaled sample without losing
precision. The Newton evaluation uses exact integer binomial
coefficients at $(p^4-1)/3$ or $(p^5-2)/3$.

The complete reproducible invocation, from the archive root, is
\begin{verbatim}
python work_e7_companion/r5_finite_closure_driver.py
\end{verbatim}
It imports the shared, independently audited function
\texttt{certify} from \nolinkurl{finite_projector_certificates.py}
and writes \nolinkurl{r5_finite_closure.json}. It does not overwrite
the earlier 158-index higher-level certificate. Only the Python
standard library is required.

The output for all 56 required values is summarized below. The
minimum scaled valuation is the valuation of the exact residue of
$p^3R$ or $p^3P$ modulo $p^5$; subtracting three gives the certified
unscaled valuation. All minima shown are below the precision cap,
so their listed valuations are exact.
\[
\begin{array}{c|r|r|r|r}
p&d_{\max}^{R}&\min v_p(p^3R)&d_{\max}^{P}&\min v_p(p^3P)\\\hline
5 &36 &3&32 &3\\
11&92 &2&88 &3\\
17&148&2&144&3\\
23&204&2&200&3
\end{array}
\]
Thus every companion value has valuation at least $-1$, and every
primary value has valuation at least zero, which is stronger than the
required primary bound. The finite sufficient conditions hold for
all four primes. The uniform leakage argument therefore gives
\[
 E(p^5)\equiv p\mathcal A(p^4)\pmod{p^{14}},\qquad
 \mathcal A(p^4)\equiv pE(p^3)\pmod{p^{13}}.
\]
Combining them proves the unchanged $r=5$ target for this exhaustive
finite exception set. Together with the separate uniform argument
for $p\ge29$, it yields
\[
 \boxed{E(p^5)\equiv p^2E(p^3)\pmod{p^{14}}
 \qquad\text{for every prime }p>3,\ p\equiv2\pmod3.}
\]
The finite output is not being extrapolated to other primes: those
primes are covered by the uniform reflection-gain proof. This section
settles the fifth level. The third level is completed below.

\section{A uniform first-step congruence at the third level}\label{loc-VIII-14}
\label{VIII-sec:E7firststepr3}
\begin{corollary}\label{VIII-E7:first-step-r3}
For every prime $p>3$, $p\equiv2\pmod3$, the exact companion sum
of Theorem~\ref{VIII-E7:companion-exact} satisfies
\[
 \boxed{E(p^3)\equiv p\mathcal A(p^2)\pmod{p^8}.}
\]
In fact the positive-moment transport part of this difference belongs
to $p^9\mathbb Z_p$.
\end{corollary}
\begin{proof}
Write $R^{(1)}_{p^3,p}(X)=1+\sum_{j\ge1}a_jX^j$ for the exact
one-step kernel. Its coefficients satisfy $v_p(a_j)\ge j$.
Summing the proved identity
$\mu_{p^3}(pJ)=p\widetilde\mu_{p^2}(J)R^{(1)}_{p^3,p}(J)$
over the companion nodes gives
\[
 E(p^3)-p\mathcal A(p^2)
 =\sum_{p\nmid H}\mu_{p^3}(H)
  +p\sum_{j\ge1}a_j\widetilde M_j(p^2).
\]
The first sum belongs to $p^8\mathbb Z_p$ by
Corollary~\ref{VIII-E7:primary-unit-r3}.

The common companion moment normalization has valuation seven.
The first two explicit moments thus have valuation at least seven,
so their transport terms have valuation at least nine. For $j\ge3$,
Theorem~\ref{VIII-E7:companion-log} gives
\[
 v_p\bigl(\widetilde M_j(p^2)\bigr)\ge7-\eta_{j-1},\qquad
 \eta_{j-1}=\left\lceil\frac32
 \left\lfloor\log_p(6\lfloor(2j-1)/3\rfloor+1)\right\rfloor
 \right\rceil.
\]
Here $\eta_{j-1}\le j-1$. Indeed
$6\lfloor(2j-1)/3\rfloor+1\le4j-1$, and
\[
 4j-1<5^{\lfloor2(j-1)/3\rfloor+1}\qquad(j\ge3).
\]
The latter inequality is checked for $j=3,4,5$; increasing $j$ by
three adds twelve on the left and multiplies the right by twenty-five.
It follows that each remaining transport term has valuation at least
$1+j+7-(j-1)=9$. The logarithmic growth of $\eta_{j-1}$ proves
convergence of the series. This proves the corollary.
\end{proof}

Consequently the third-level E.3 congruence is exactly
equivalent to the shorter companion estimate
\[
 \mathcal A(p^2)-pE(p)\in p^7\mathbb Z_p.
\]
This equivalence includes an independently proved precision bound for
one entire transport step. The companion estimate is proved below
in Theorem~\ref{VIII-thm:e8-r3-closed}.

\section{Exact two-step transport and the remaining cancellation}\label{loc-VIII-15}
Let $p>3$, $p\equiv2\pmod3$, and $r\ge3$ be odd. Set
$L=p^r$, $M=L/p^2$, and $m=(L-2)/3$. Use the node and denominator
root sets $\mathcal H_L$ and $\mathcal D_L$ of the three-grid
construction, and write $\mu_L(H)=W_L(H)T(H)$.
Define the restricted rational power series
\[
 R_{L,p}(X)=
 \frac{\prod_{d\in\mathcal D_L,\,p^2\nmid d}(1-p^2X/d)}
 {\prod_{h\in\mathcal H_L,\,p^2\nmid h}(1-p^2X/h)}
 =1+\sum_{j\ge1}c_jX^j.
\]
Every retained root has valuation zero or one. Hence each linear
coefficient $p^2/d$ or $p^2/h$ belongs to $p\mathbb Z_p$.
Expanding the finitely many inverse factors proves
\[
 c_j\in p^j\mathbb Z_p,\qquad v_p(c_j)\longrightarrow+\infty.
\]
The series converges at every integral coarse node $J$.

\begin{proposition}[Exact residual identity]\label{VIII-prop:e-residual}
Put
\[
 U_{L,p}=\sum_{\substack{H\in\mathcal H_L\\p^2\nmid H}}\mu_L(H).
\]
Then
\begin{equation}\label{VIII-eq:e-residual}
 \boxed{E(L)-p^2E(M)
 =U_{L,p}+p^2\sum_{j\ge3}c_j M_j(M).}
\end{equation}
No assertion that either term on the right is separately divisible
by the target modulus is made.
\end{proposition}
\begin{proof}
Since $p^2\equiv1\pmod6$, dividing all $p^2$-divisible nodes by
$p^2$ gives exactly $\mathcal H_M$. Dividing the corresponding
denominator roots gives exactly $\mathcal D_M$. For example the
positive grids end at $L-1$ and $L$, and become the positive grids
ending at $M-1$ and $M$; the negative grid is $-2$ times the first
positive grid. The positive denominator roots $6i$ and the negative
roots $-3i$ both have coarse index range
$1\le i\le(M-2)/3$. The third denominator grid ends at $-(L-3)$
and becomes the grid ending at $-(M-3)$.

The exact model values satisfy $T(p^2J)=p^2T(J)$, because
multiplication by $p^2$ preserves the residues modulo six and the
sign of each node. Separating divisible factors in the weights
therefore gives the product identities
\[
 W_L(p^2J)=W_M(J)R_{L,p}(J),\qquad
 \mu_L(p^2J)=p^2\mu_M(J)R_{L,p}(J).
\]
These are identities of products, including at zero weights;
no division by a vanishing weight is needed. Sum over coarse nodes
and add the nodes not divisible by $p^2$. Interchanging the finite
node sum with the convergent series is legitimate. The identities
$M_1(M)=M_2(M)=0$ remove its first two positive terms and give
\eqref{VIII-eq:e-residual}.
\end{proof}


\section{The leading kernel coefficients from three short harmonic sums}\label{loc-VIII-16}
\label{VIII-sec:E7harmonic}
Let $p>3$ be prime with $p\equiv2\pmod3$, let $r\ge3$ be odd,
and put $L=p^r$, $q=(p-2)/3$. Set
\[
 H_i=\sum_{a=1}^{q}a^{-i}\in\mathbb F_p\quad(1\le i\le3).
\]
Use the exact two-step kernel
\[
 R_{L,p}(X)=
 \frac{\prod_{d\in\mathcal D_L,\ p^2\nmid d}(1-p^2X/d)}
 {\prod_{h\in\mathcal H_L,\ p^2\nmid h}(1-p^2X/h)}
 =1+\sum_{j\ge1}c_jX^j.
\]
All products here count multiplicities. In particular, an overlap
between the node and denominator lists may be canceled but need
not be removed in defining the products.

\begin{proposition}\label{VIII-E7:harmonic-kernel}
For every such $p,r$,
\begin{align*}
 c_1/p&\equiv-H_1/2\pmod p,\\
 c_2/p^2&\equiv H_1^2/8+5H_2/24\pmod p,\\
 \kappa_p:=c_3/p^3&\equiv
 -\frac{3H_1^3+15H_1H_2+2H_3}{144}\pmod p.
\end{align*}
Thus these leading coefficients are independent of the odd level $r$.
\end{proposition}
\begin{proof}
For $i\le3$ define the normalized root-power difference
\[
 s_i=p^{-i}\left(
 \sum_{h\in\mathcal H_L,\ p^2\nmid h}(p^2/h)^i
 -\sum_{d\in\mathcal D_L,\ p^2\nmid d}(p^2/d)^i\right).
\]
Every term is integral at $p$. Modulo $p$, the unit roots
contribute zero, leaving exactly the roots of valuation one.
After dividing those roots by $p$, their lists are the unit roots
of the companion grids at $M=p^{r-1}=3n+1$:
\[
 \begin{split}
 \mathcal H'_M={}&\{3j+1:0\le j\le n\}
 \cup\{3j+2,-2(3j+2):0\le j<n\},\\
 \mathcal D'_M={}&\{6j,-3j:1\le j\le n\}
 \cup\{-(3j+1):0\le j<n\}.
 \end{split}
\]
These identifications follow directly by solving the original
linear root congruences modulo $p$; they respect multiplicities.

Put $t=(2p-1)/3=2q+1$. Since $n\equiv t\pmod p$, complete
blocks of $p$ consecutive indices contribute zero modulo $p$ to
each inverse-power sum under consideration. Indeed,
$\sum_{u\in\mathbb F_p^\times}u^{-i}=0$ for $1\le i\le3<p-1$.
Let
\[
 T_i=\sum_{j=1}^{t}j^{-i},\qquad
 A_i=\sum_{j=0}^{t-1}(3j+1)^{-i},\qquad
 B_i=\sum_{\substack{0\le j<t\\j\ne q}}(3j+2)^{-i}.
\]
The extra endpoint $3t+1=2p$ of the first node list is excluded
as a nonunit. The involution $j\mapsto t-1-j$ gives
$A_i=(-3)^{-i}T_i$. In the $B$ list, the residues are
$3a$ for $a=-q,\ldots,-1,1,\ldots,q$. Consequently
\[
 B_1=B_3=0,\qquad B_2=2H_2/9.
\]
Since $t=p-1-q$, the full nonzero-residue sum also gives
$T_i=(-1)^{i+1}H_i$.
Taking the difference of the companion node and denominator
power sums yields
\[
 s_i\equiv
 (1-(-1)^i)A_i+(1+(-2)^{-i})B_i
 -(6^{-i}+(-3)^{-i})T_i\pmod p.
\]
Substitution gives
\[
 s_1=-H_1/2,\qquad s_2=5H_2/12,\qquad s_3=-H_3/24
 \quad\text{in }\mathbb F_p.
\]
Finally the logarithmic product expansion, or equivalently the
first three Newton identities, gives
\[
 \frac{c_1}{p}=s_1,\quad
 \frac{c_2}{p^2}=\frac{s_1^2+s_2}{2},\quad
 \frac{c_3}{p^3}=\frac{s_1^3+3s_1s_2+2s_3}{6}.
\]
Here these identities are exact before the reductions of $s_i$.
The stated congruences follow; all displayed denominators are
units because $p>3$.
\end{proof}

This proposition computes one contribution to the residual and
does not prove its cancellation with the other contribution.
For example $p=11$ gives $(H_1,H_2,H_3)=(0,9,6)$ modulo $11$,
so $\kappa_{11}=-1$. In particular the cubic kernel coefficient
is not always divisible by $p^4$.

\section{A uniform coarse-moment estimate and a rigorous transport-tail bound}\label{loc-VIII-17}
\label{VIII-sec:E7coarsemoments}
Let $p>3$ be prime, $p\equiv2\pmod3$, and put $q=(p-2)/3$.
Then $q$ is odd. Use the original three-grid interpolation at
$N=p$, and denote its positive moments by $M_j(p)$.

\begin{proposition}\label{VIII-E7:all-coarse-moments}
For every positive integer $j$,
\[
 \boxed{M_j(p)\in p^2\mathbb Z_p.}
\]
This bound is uniform in $j$; it is not inferred from finitely many
polynomial moment formulas.
\end{proposition}
\begin{proof}
Write $P_q(z)=D_q(z)\mathcal S_q(z)$ for the cleared polynomial
of the three-grid construction. All its coefficients belong to
$\mathbb Z_p$: the fixed denominators $k!$, $(2/3)_k$, and the
normalizing roots of $D_q$ are units for $0\le k\le q$.
We prove the stronger coefficientwise assertion
\begin{equation}\label{VIII-E7:cleared-modp}
 P_q(z)\in p\mathbb Z_p[z].
\end{equation}

Work first in $\mathbb F_p(z)$. Insert the canceled numerator and
denominator parameter $e=2/3$ in $\mathcal S_q$ to exhibit the
very-well-poised ${}_6F_5(-1)$ form.
Modulo $p$ its central parameter satisfies
$a=1/3=-2q-1$, and the inserted parameter satisfies $e=2/3=-q$.
The finite canceled factors $(-q)_k$ are nonzero for $k\le q$,
so this insertion is legitimate.
Apply Whipple's terminating transformation using this $-q$ as the
terminating upper parameter. The other parameters may be written
\[
 b=a+z/6,\quad c=a-z/3,\quad e'=-q-z/3,\quad f=-q.
\]
The transformed terminating ${}_3F_2$ is
\[
 {}_3F_2\left[\begin{matrix}
 -q,\ -q+z/6,\ -q-z/3\\1-z/6,\ 1+z/3
 \end{matrix};1\right]_q.
\]
Its prefactor is $(-2q)_q/(-q+z/3)_q$, a well-defined element of
$\mathbb F_p(z)$. The displayed sum is a terminating Dixon sum
with odd central parameter $-q$: the two lower parameters equal
$1-q-(-q+z/6)$ and $1-q-(-q-z/3)$.
Reversing the summation index $k\mapsto q-k$ negates every term.
The sum is therefore zero. All transformations are finite rational
identities whose denominators are nonzero in $\mathbb F_p(z)$;
no analytic continuation in characteristic $p$ is used.
It follows that the reduction of $\mathcal S_q$ is zero as a
rational function. Since $D_q$ is a nonzero polynomial modulo $p$,
this proves \eqref{VIII-E7:cleared-modp}.

Let $\widehat Q_q(z)$ be the monic polynomial whose roots are the
three interpolation grids. These roots are integers. Exactly one
root is divisible by $p$, namely the positive node $p$ itself;
every other root is a $p$-adic unit. Thus
$Q_q(z)=\ell_q\widehat Q_q(z)$ with $v_p(\ell_q)=-1$.
By \eqref{VIII-E7:cleared-modp},
\[
 \frac{D_q(z)\mathcal S_q(z)}{Q_q(z)}
 =\frac{A_q(z)}{\widehat Q_q(z)},\qquad
 A_q(z)=\ell_q^{-1}P_q(z)\in p^2\mathbb Z_p[z].
\]
The inverse Laurent series at infinity of a monic polynomial in
$\mathbb Z_p[z]$ has integral coefficients after its leading power
of $z$ is removed. Hence every Laurent coefficient of this proper
rational function belongs to $p^2\mathbb Z_p$.
Those coefficients are $-M_j(p)$, which proves the proposition.
\end{proof}

\begin{corollary}\label{VIII-E7:r3tail}
At $L=p^3$, let $R_{L,p}(X)=1+\sum_{j\ge1}c_jX^j$ be the exact
two-step weight kernel. Then
\[
 p^2\sum_{j\ge4}c_jM_j(p)\in p^8\mathbb Z_p.
\]
In particular the entire transport error, including its infinite
tail, satisfies
\[
 p^2\sum_{j\ge3}c_jM_j(p)
 \equiv p^2c_3M_3(p)\pmod{p^8}.
\]
\end{corollary}
\begin{proof}
The established coefficient estimate is $c_j\in p^j\mathbb Z_p$.
Proposition~\ref{VIII-E7:all-coarse-moments} bounds every summand for
$j\ge4$ by $p^{2+j+2}$, hence by $p^8$. These bounds tend to
infinity, so the entire convergent tail belongs to that closed ideal.
\end{proof}

For an explicit leading term put
\[
 b_q=\frac{(1/3)_q}{q!},\qquad
 \kappa_p\equiv c_3/p^3\pmod p.
\]
Since $q$ is odd, the exact first nonzero moment is
$M_3(p)=p^2(p-1)^3b_q^3$. The corollary therefore gives
\[
 \boxed{\quad
 p^2\sum_{j\ge3}c_jM_j(p)
 \equiv-p^7\kappa_pb_q^3\pmod{p^8}.
 \quad}
\]
Together with the separately derived formula
\[
 \kappa_p\equiv-
 \frac{3H_1^3+15H_1H_2+2H_3}{144}\pmod p,
 \qquad H_i=\sum_{a=1}^{q}a^{-i},
\]
this determines the full transport contribution modulo the target
modulus by three harmonic sums of length $(p-2)/3$.

These estimates show that all orders beyond the cubic transport term
are harmless. The following sections evaluate the unit-node contribution
and prove its exact cancellation with the cubic term.



\section{The first resonance determines the companion pole residues}\label{loc-VIII-18}
\label{VIII-sec:e-first-resonance}
Let $R_j(X)$ be the normalized companion moment polynomials, and set
$\mathcal R(X,t)=\sum_{j\ge0}R_j(X)t^j$.
The following statements are unconditional finite algebraic results.
Their low-order integrality consequence will force a high-order
zero of the block derivative at the origin.

\subsection{An exactly diagonal factorial basis}
For $k\ge0$, define monic polynomials
\[
 P_k(X)=\prod_{h=0}^{k-1}\bigl(X^2-(3h+1)^2\bigr),\qquad
 Q_k(Y)=\prod_{h=0}^{k-1}\bigl(Y-(3h+1)(3h-2)\bigr).
\]
These are auxiliary basis polynomials, not the primary moment
polynomials denoted by $P_j$ elsewhere. Put $Y=X(X-3)$ and
\[
 \mathcal L'f=X(X-2)f(X)-(X-3)(X-1)f(X-3).
\]
Then
\begin{equation}\label{VIII-eq:e-exact-diagonal-basis}
 \mathcal L'P_k=(3k+1)(2X-3)Q_k(Y).
\end{equation}
Indeed the two products on the left have the common factor $Q_k(Y)$,
and their remaining factors are respectively
$X(X+3k-2)$ and $(X-3)(X-3k-1)$. Their difference is
$(3k+1)(2X-3)$. This proves the identity over $\mathbb Z[X]$.

Define
\[
 g_X(t)=(1-Xt)(1-(X-2)t)(1+2(X-2)t).
\]
The exact companion recurrence, including its first three initial
moments, is equivalent to
\begin{equation}\label{VIII-eq:e-full-generating-operator}
 \begin{split}
 &X(X-2)g_X(t)\mathcal R(X,t)
 -(X-3)(X-1)g_{3-X}(t)\mathcal R(X-3,t)\\
 &\hspace{12mm}=(2X-3)\{1+(2Y+3)t+(2Y+2)t^2\}.
 \end{split}
\end{equation}
For example, the $t^0,t^1,t^2$ equations use
$R_0=1$, $R_1=(X^2+1)/2$, and $R_2=2X^2-1$; all later
equations are precisely the proved homogeneous companion recurrence.

Expand in the factorial basis:
\[
 \mathcal R(X,t)=\sum_{k\ge0}C_k(t)P_k(X).
\]
Each coefficient of $t^j$ is a finite sum, since
$\deg R_j\le2\lfloor(2j+1)/3\rfloor$.
Direct multiplication using
$YQ_k=Q_{k+1}+(3k+1)(3k-2)Q_k$ gives
\begin{align*}
 &\frac{X(X-2)g_XP_k-(X-3)(X-1)g_{3-X}P_k(X-3)}{2X-3}\\
 &\quad=(3k+1)g_{3k+1}(t)Q_k(Y)\\
 &\qquad+t^2\{-9k-2+2(18k^2+18k+5)t\}Q_{k+1}(Y)
                  +2t^3Q_{k+2}(Y).
\end{align*}
Consequently the coefficient functions are the following explicit
rational functions:
\begin{equation}\label{VIII-eq:e-factorial-coefficient-recurrence}
 \begin{split}
 C_0(t)&=\frac1{1-t},\\
 C_1(t)&=\frac{t(1+3t)}{2(1-t)(1-16t^2)},\\
 C_k(t)&=
 \frac{t^2\{9k-7-2(18k^2-18k+5)t\}C_{k-1}(t)-2t^3C_{k-2}(t)}
 {(3k+1)g_{3k+1}(t)}\qquad(k\ge2).
 \end{split}
\end{equation}
Here all $g_{3k+1}$ have constant term one. The recurrence also gives
$\operatorname{ord}_t C_k\ge\lfloor3k/2\rfloor$, so it independently
justifies the finite coefficient sums.

\subsection{A rigorous first-resonance and local-constancy theorem}
\begin{theorem}\label{VIII-thm:e-first-resonance}
Let $p>3$, $p\equiv2\pmod3$, and put $k_0=(2p-1)/3$.
If $0\le j\le(p^2-3)/2$, then for every $x\in p\mathbb Z_p$,
$pR_j(x)\in\mathbb Z_p$ and
\begin{equation}\label{VIII-eq:e-first-resonance-residue}
 pR_j(x)\equiv
 P_{k_0}(0)\,[t^j]\bigl(pC_{k_0}(t)\bigr)\pmod p.
\end{equation}
In particular $pR_j(x)$ modulo $p$ is constant on $p\mathbb Z_p$.
For every $p\ge11$ in this residue class, this applies to all
$j\le4p-5$, and hence
\[
 pR_j(0)\equiv pR_j(p^2)\equiv pR_j(2p)\pmod p
 \qquad(0\le j\le4p-5).
\]
Also $pR_j(x)\equiv0\pmod p$ for $j<p-1$.
\end{theorem}
\begin{proof}
For a fixed $k$, let
\[
 m_k=\sum_{i=1}^k v_p(3i+1).
\]
Because the nonconstant denominators in
\eqref{VIII-eq:e-factorial-coefficient-recurrence} are unit series,
induction gives the coefficientwise bound
\[
 C_k(t)\in p^{-m_k}\mathbb Z_p[[t]].
\]
For a moment index in the stated range, only
$k\le\lfloor(2j+1)/3\rfloor\le(p^2-4)/3$ can occur.
Therefore $3k+1<p^2$; every positive valuation appearing in
$m_k$ is exactly one. Thus $m_k$ simply counts the indices
$i\le k$ for which $p\mid3i+1$.

If $k<k_0$, then $m_k=0$, and the corresponding term in
$pR_j(x)$ is divisible by $p$. At $k=k_0$, there is exactly
one such divisor, namely $3k_0+1=2p$. Hence $pC_{k_0}$ is
coefficientwise integral and $P_{k_0}(x)\equiv P_{k_0}(0)$
modulo $p$.

Suppose $k>k_0$, and set
$\delta_k=1$ if $p\mid3k+1$ and $\delta_k=0$ otherwise.
Among the factors of $P_k(x)$ there are exactly
$m_k-\delta_k$ indices $h<k$ with $p\mid3h+1$.
Each of those factors is divisible by $p^2$ because $x\in p\mathbb Z_p$.
Consequently every coefficient contribution has valuation at least
\[
 1-m_k+2(m_k-\delta_k)=1+m_k-2\delta_k.
\]
If $\delta_k=0$, this is at least two. If $\delta_k=1$, then
$k>k_0$ implies $m_k\ge2$, so the same expression is at least one.
All these terms vanish modulo $p$, proving
\eqref{VIII-eq:e-first-resonance-residue}.
The argument also proves the asserted integrality.

For $p\ge11$, $4p-5\le(p^2-3)/2$.
Finally $\operatorname{ord}_t C_{k_0}\ge
\lfloor3k_0/2\rfloor=p-1$, proving the last assertion.
\end{proof}


\section{Exact reflection and second-order unit-support reduction}\label{loc-VIII-19}
\label{VIII-sec:E8-fine-reflection}
This section proves an exact reflection identity and a valuation-two
reduction of the unit-support part of the level-three numerator.
The next sections combine this identity with the carry-support
contribution and rational-function rigidity.

Let $p>3$ be prime with $p\equiv2\pmod3$, and put
\[
 a=(p+1)/3,\qquad q=a-1,\qquad B_b=(1/3)_b/b!.
\]
Then $a$ is even. Work in $\mathcal R=\mathbb Z_p[[w]]$. Define
\[
 f_x(w)=\frac{(1+(6x-2)w)(1-(3x-1)w)(1-(3x+1)w)}
 {(1-(6x+6)w)(1+(3x+3)w)(1+(3x+1)w)},
 \qquad
 v_j(t,w)=\frac{j-1/3+t}{j+1+t}f_{j+t}(w),
\]
and, for $0\le c\le a$,
\[
 T_c(t,w)=\prod_{j=0}^{c-1}v_j(t,w),\qquad
 F(x,w)=(6x-1)(1+(18x^2-6x-2)w).
\]
\[
 \mathcal B(t,w)=\sum_{c=0}^aT_c(t,w)F(c+t,w).
\]
Every numerical factor in $v_j(t,w)$ is a $p$-adic unit for
$0\le j<a$ and $t\in p\mathbb Z_p$. In fact these rational functions
have integral Taylor coefficients in $t$ on this residue disk,
coefficientwise in $w$. All formal denominators have constant term
one in $w$.

\begin{lemma}[Exact reflection]\label{VIII-E8:exact-reflection}
As rational functions of $t,w$, with $p=3a-1$,
\[
 v_{a-1-j}(t,w)=v_j(-p/3-t,w)^{-1},\qquad 0\le j<a.
\]
Consequently, if $u=-p/3-t$,
\[
 T_{a-c}(t,w)=T_a(t,w)T_c(u,w),\qquad
 T_a(t,w)T_a(u,w)=1,
\]
and
\[
 \boxed{\mathcal B(t,w)=-T_a(t,w)\mathcal B(u,w).}
\]
At $t_0=-p/6$ one has the exact identities
\[
 T_a(t_0,w)=1,\qquad \boxed{\mathcal B(t_0,w)=0}.
\]
\end{lemma}
\begin{proof}
First,
\[
 \frac{a-1-j-1/3+t}{a-j+t}
 =\frac{j+1-p/3-t}{j-1/3-p/3-t}.
\]
For the kernel, the three numerator factors of $f_{a-1-j+t}$ are
\[
 1-(6j+6)w+2pw+6tw,\quad
 1+(3j+3)w-pw-3tw,\quad
 1+(3j+1)w-pw-3tw.
\]
They are precisely the denominator factors of $f_{j+u}$; its three
numerator factors similarly equal the reflected denominator factors.
This proves the single-factor identity. Reflecting the final $c$
factors of $T_a(t)$ gives the identity for $T_{a-c}(t)$, and taking
$c=a$ proves self-reciprocity.

A direct polynomial calculation gives
$F(1/3-x,w)=-F(x,w)$. Since
$a-c+t=1/3-(c+u)$, replacing $c$ by $a-c$ in the finite sum proves
the displayed reflection of $\mathcal B$.
Finally $u=t=t_0$ and the $a$ factors of $T_a(t_0)$ pair into reciprocal
pairs. There is no unpaired factor because $a$ is even. Thus
$T_a(t_0)=1$ exactly, and $\mathcal B(t_0)=-\mathcal B(t_0)$ gives zero.
\end{proof}

Write
\[
 g=\left.\partial_t\log T_a(t,w)\right|_{t=0},\qquad
 B_i=\left.\partial_t^i\mathcal B(t,w)\right|_{t=0}
 \quad (i=0,1,2).
\]
Logarithmic derivatives mean rational derivatives and require no
choice of a logarithm of the constant factor. The quantities $g,B_i$
are in $\mathcal R$.

\begin{corollary}[Reflection jets]\label{VIII-E8:reflection-jets}
One has
\[
 B_2\equiv gB_1\pmod{p\mathcal R},\qquad
 B_0\equiv\frac p6 B_1-\frac{p^2}{72}B_2\pmod{p^3\mathcal R}.
\]
For every integer $b$,
\begin{equation}\label{VIII-E8:B-shift}
 \mathcal B(pb,w)\equiv
 p\left(b+\frac16\right)B_1
 +p^2\left(\frac{b^2}{2}-\frac1{72}\right)gB_1
 \pmod{p^3\mathcal R}.
\end{equation}
\end{corollary}
\begin{proof}
Differentiate the exact reflection twice, using $u'=-1$.
At $t=t_0$, its value is zero and $T_a(t_0)=1$, so
\[
 \mathcal B''(t_0,w)
 =\bigl(\partial_t\log T_a\bigr)(t_0,w)\mathcal B'(t_0,w).
\]
Reducing modulo $p$ yields the first congruence, because $t_0\in p\mathbb Z_p$.
Taylor expansion of the exact equality $\mathcal B(-p/6,w)=0$ gives
\[
 0\equiv B_0-\frac p6B_1+\frac{p^2}{72}B_2\pmod{p^3}.
\]
Taylor expansion at $pb$ and substitution prove \eqref{VIII-E8:B-shift}.
The Taylor remainders are integral multiples of $p^3$, since all
rational factors have unit numerical denominators on the residue disk.
\end{proof}

For comparison with the original length-$p$ block, let
\[
 \widehat K_k(w)=\prod_{j=0}^{k-1}f_j(w),\qquad d_k=(-1/3)_k/k!.
\]
Let $Z(t,w)$ be the complete block after the numerical factors divisible
by $p$ have been stripped:
\[
 Z(t,w)=
 \frac{\prod_{0\le j<p,\ j\ne a}(j-1/3+t)}
 {\prod_{1\le j\le p,\ j\ne p}(j+t)}
 \prod_{j=0}^{p-1}f_{j+t}(w).
\]
The complete-block residue calculation proves
\[
 Z(0,w)\equiv1\pmod p,\qquad
 \partial_t\log Z(0,w)\in p\mathcal R.
\]
For completeness, the numerical numerator and denominator residues
both cover $\mathbb F_p^\times$ once. Their reciprocal first-power
sums vanish. Each affine root progression in the kernel covers
$\mathbb F_p$ once, so its complete product is the same for all three
numerator and denominator progressions. In the logarithmic derivative,
the numerator slope sum is $6-3-3=0$ and the denominator slope sum is
$-6+3+3=0$. Thus both assertions hold coefficientwise.
Put
\[
 \beta=\frac{Z(0,w)-1}{p}\in\mathcal R.
\]

\begin{proposition}[The whole unit-support block has valuation two]
\label{VIII-E8:unit-support}
Define
\[
 \mathcal U_p(w)=\sum_{b=0}^q\sum_{c=0}^a
 d_{pb+c}\widehat K_{pb+c}(w)F(pb+c,w).
\]
Then $\mathcal U_p\in p^2\mathcal R$, with the explicit residue
\[
 \boxed{\frac{\mathcal U_p(w)}{p^2}
 \equiv B_q B_1(w)
 \left(-\frac1{12}+\frac{\beta(w)}{28}+\frac{g(w)}{56}\right)
 \pmod{p\mathcal R}.}
\]
\end{proposition}
\begin{proof}
For $c\le a$, extraction of the numerical factors divisible by $p$
gives the exact factorization
\[
 d_{pb+c}\widehat K_{pb+c}
 =B_b\left(\prod_{\ell=0}^{b-1}Z(p\ell,w)\right)T_c(pb,w).
\]
No partial-block numerator factor divisible by $p$ has yet appeared;
the first such index is $c=a+1$.
Since $Z(p\ell,w)\equiv Z(0,w)\pmod{p^2}$,
\[
 \prod_{\ell=0}^{b-1}Z(p\ell,w)\equiv1+pb\beta\pmod{p^2}.
\]
Equation \eqref{VIII-E8:B-shift} puts $\mathcal B(pb,w)$ in $p\mathcal R$.
Therefore multiplying the last approximation by the entire fine-support
sum loses only a multiple of $p^3$. We obtain
\begin{align*}
 \mathcal U_p\equiv {}&
 pB_1\sum_{b=0}^q B_b(b+1/6)\\
 &+p^2B_1\left\{
 \beta\sum_{b=0}^qB_b b(b+1/6)
 +g\sum_{b=0}^qB_b(b^2/2-1/72)\right\}\pmod{p^3}.
\end{align*}
The generalized hockey-stick identity gives
\[
 \sum_{b=0}^q\binom bjB_b
 =\frac{(p-1)\binom qj}{3j+1}B_q\qquad (j=0,1,2).
\]
Using $q=(p-2)/3$ yields the exact equalities
\begin{align*}
 \sum_{b=0}^qB_b(b+1/6)
 &=\frac{p(p-1)}{12}B_q,\\
 \sum_{b=0}^qB_b b(b+1/6)
 &=\frac{(p-2)(p-1)(8p+9)}{504}B_q,\\
 \sum_{b=0}^qB_b(b^2/2-1/72)
 &=\frac{(p-1)(4p^2-7p-9)}{504}B_q.
\end{align*}
The last two coefficients reduce to $1/28$ and $1/56$, respectively.
All their denominators are units for the primes at issue: $p>3$ and
$p\equiv2\pmod3$ exclude $p=7$. This proves the proposition.
\end{proof}

The fine first derivative in this result is explicit:
\[
 B_1(w)=\sum_{c=0}^aT_c(0,w)
 \left(F'(c,w)+F(c,w)
 \left.\partial_t\log T_c(t,w)\right|_{t=0}\right).
\]
The proposition removes the coarse-index sums from the unit-support
contribution. The following identity simplifies its correction terms.

\begin{lemma}[Cancellation of the complete-block correction]
\label{VIII-E8:beta-cancellation-lemma}
The complete-block correction satisfies
\[
 \beta\equiv-g/2\pmod{p\mathcal R}.
\]
Consequently the unit-support formula simplifies to
\[
 \boxed{\mathcal U_p(w)/p^2\equiv-B_qB_1(w)/12\pmod p.}
\]
\end{lemma}
\begin{proof}
Set $K=p+a=4a-1$, and extend $T_c(t,w)$ by the same original product
when $c>a$. Define its unitized value
\[
 V(t,w)=T_K(t,w)\frac{p+t}{p/3+t}.
\]
The product up to $K$ contains just one numerical numerator factor
divisible by $p$, at index $a$, and just one denominator factor
divisible by $p$, at index $p$. Thus $V$ is analytic and integral
on $t\in p\mathbb Z_p$ with unit value. Splitting the product at $p$
gives the exact equality
\[
 V(t,w)=Z(t,w)T_a(t+p,w).
\]
The single-factor reflection for length $K$ is
\[
 v_{K-1-j}(t,w)=v_j(-4p/3-t,w)^{-1}.
\]
At $t_*=-2p/3$, all factors pair into reciprocals except the central
one, because $K$ is odd. Its argument is
$(K-1)/2+t_*=-1/3$, for which $f_{-1/3}=1$ and the numerical ratio
is $-1$. Hence $T_K(t_*,w)=-1$. Since
$(p+t_*)/(p/3+t_*)=-1$, it follows that $V(t_*,w)=1$.

The logarithmic derivative of $V$ at zero reduces to $g$ modulo $p$:
this follows from $V=ZT_a(t+p)$ and the complete-block derivative
$\partial_t\log Z(0)\in p\mathcal R$. Therefore
\[
 V(0,w)\equiv1+\frac{2p}{3}g\pmod{p^2}.
\]
On the other hand $T_a(-p/6,w)=1$, so
\[
 T_a(p,w)\equiv1+\frac{7p}{6}g\pmod{p^2},\qquad
 V(0,w)=Z(0,w)T_a(p,w)
 \equiv1+p\beta+\frac{7p}{6}g\pmod{p^2}.
\]
Comparison gives $\beta\equiv-g/2\pmod p$. Substitution into
Proposition~\ref{VIII-E8:unit-support} cancels its last two terms.
\end{proof}

\section{A coefficientwise valuation-two theorem and a full-block derivative}\label{loc-VIII-20}
\label{VIII-sec:E8-full-block-derivative}
This coefficientwise identity is the bridge from the companion
numerator to a rational function over the finite field.

Let $p>3$, $p\equiv2\pmod3$, $a=(p+1)/3$, $q=a-1$, and
$n=(p^2-1)/3$. Use the functions $f_x,F,v_j$ from
Section~\ref{VIII-sec:E8-fine-reflection}, and write
\[
 P_c(t,w)=\prod_{j=0}^{c-1}v_j(t,w),\qquad
 Z_p(t,w)=\sum_{c=0}^{p-1}P_c(t,w)F(c+t,w).
\]
Here $P_c$ has integral Taylor coefficients at $t=0$ whenever $c<p$;
its numerator may contain a nonunit factor, but its denominators do not.
Put $B_b=(1/3)_b/b!$, $d_k=(-1/3)_k/k!$, and
\[
 \mathcal S_p(w)=\sum_{k=0}^n d_k\widehat K_k(w)F(k,w),
 \qquad \widehat K_k(w)=\prod_{j=0}^{k-1}f_j(w).
\]

\begin{theorem}\label{VIII-E8:full-block-derivative}
Coefficientwise in $\mathbb Z_p[[w]]$,
\[
 \mathcal S_p(w)\in p^2\mathbb Z_p[[w]],\qquad
 \boxed{\frac{\mathcal S_p(w)}{p^2}
 \equiv-\frac{B_q}{12}\left.\partial_t Z_p(t,w)\right|_{t=0}
 \pmod p.}
\]
Consequently the integral companion numerator satisfies
\[
 \boxed{\frac{\mathcal N_n(w)}{p^2}
 \equiv-\frac{B_q A_n(w)}{12(1-2w)}
 \left.\partial_t Z_p(t,w)\right|_{t=0}\pmod p,}
\]
where $\mathcal N_n=A_n\mathcal S_p/(1-2w)$ is its exact
reversed-root representation.
\end{theorem}

\begin{proof}
All congruences below are coefficientwise. Let $Z(t,w)$ denote the
complete length-$p$ product after deleting the numerical numerator
factor at $j=a$ and the numerical denominator factor at $j=p-1$.
Thus
\[
 P_p(t,w)=\frac{p/3+t}{p+t}Z(t,w),\qquad
 \beta=(Z(0,w)-1)/p.
\]
The complete-block calculation gives $Z(0,w)\equiv1\pmod p$ and
$\partial_t\log Z(0,w)\equiv0\pmod p$.
Write $T_a=P_a$, $g=\partial_t\log T_a(0,w)$, and
\[
 \mathcal B(t,w)=\sum_{c=0}^aP_c(t,w)F(c+t,w),\qquad
 B_1(w)=\partial_t\mathcal B(0,w).
\]
The exact short reflection proved in the accompanying lemma gives
$T_a(-p/6,w)=1$ and
\begin{equation}\label{VIII-E8:U-intermediate}
 \frac{\mathcal U_p}{p^2}
 \equiv B_qB_1\left(-\frac1{12}+\frac\beta{28}+\frac g{56}\right)
 \pmod p,
\end{equation}
where $\mathcal U_p$ is the part with fine index $0\le c\le a$.
In particular $\mathcal U_p\in p^2\mathbb Z_p[[w]]$.

We now treat every remaining term. Set $K=p+a=4a-1$ and define
\[
 U_c(t,w)=\frac{P_c(t,w)}{p/3+t}\quad(a<c<p),\qquad
 U_K(t,w)=\frac{p+t}{p/3+t}P_K(t,w).
\]
\[
 W_c(t,w)=U_c(t,w)F(c+t,w).
\]
These unitized rational functions have integral Taylor coefficients
on $p\mathbb Z_p$. The general single-factor reflection is
\[
 v_{K-1-j}(t,w)=v_j(-4p/3-t,w)^{-1}.
\]
Writing $t'=-4p/3-t$, canceling the single deleted numerator factor
on each side, and using $F(1/3-x,w)=-F(x,w)$ yields the exact identity
\begin{equation}\label{VIII-E8:carry-reflection}
 W_{K-c}(t,w)=U_K(t,w)W_c(t',w)
 \qquad(a<c<p).
\end{equation}
At $t=-2p/3$, the $K$ factors of $P_K$ pair reciprocally except
the central factor, whose value is $-1$. Hence $P_K=-1$ there;
the removed ratio $(p/3+t)/(p+t)$ is also $-1$. Therefore
$U_K(-2p/3,w)=1$ exactly.

Put $g_K=\partial_t\log U_K(0,w)$. Taylor expansion of the last
identity and \eqref{VIII-E8:carry-reflection} gives, with $W_c,W_c'$
evaluated at zero,
\begin{align*}
 W_{K-c}(0,w)&\equiv W_c+\frac{2p}{3}(g_KW_c-2W_c')\pmod{p^2},\\
 W_{K-c}'(0,w)&\equiv g_KW_c-W_c'\pmod p.
\end{align*}
The exact product factorization
\[
 U_K(t,w)=Z(t,w)T_a(t+p,w)
\]
first gives $g_K\equiv g\pmod p$. It also gives a useful cancellation.
Indeed $T_a(-p/6)=1$ implies
$T_a(p)\equiv1+7pg/6\pmod{p^2}$, whereas
$U_K(-2p/3)=1$ implies
$U_K(0)\equiv1+2pg/3\pmod{p^2}$. Comparing the product at zero yields
\begin{equation}\label{VIII-E8:beta-cancellation}
 \boxed{\beta\equiv-g/2\pmod p.}
\end{equation}

Pair each $c=a+1,\ldots,2a-1$ with $K-c=2a,\ldots,p-1$.
The first member has coarse range $0\le b\le q$, and the second
has range $0\le b\le q-1$. Exact extraction of the nonunit
Pochhammer factors makes their common outside factor
$pB_b(b+1/3)$; their full-block product is
$1+pb\beta\pmod{p^2}$. Define
\[
 H_i^+=\sum_{b=0}^q b^i(b+1/3)B_b,\qquad
 H_i^-=\sum_{b=0}^{q-1} b^i(b+1/3)B_b\quad(i=0,1).
\]
The elementary finite binomial sums give exactly
\begin{align*}
 H_0^+/B_q&=(p-1)(p+2)/12,&
 H_0^-/B_q&=(p-2)(p-1)/12,\\
 H_1^+/B_q&=(p-2)(p-1)(p+2)/63,&
 H_1^-/B_q&=(p-5)(p-2)(p-1)/63.
\end{align*}
In particular $H_0^++H_0^-=p(p-1)B_q/6$.
The paired contribution, divided by $p$, is congruent modulo $p^2$ to
\[
 H_0^+W_c+H_0^-W_{K-c}
 +pH_1^+(W_c'+\beta W_c)
 +pH_1^-(W_{K-c}'+\beta W_{K-c}).
\]
Its constant term is divisible by $p$. On substitution, the
coefficient of $W_c'$ after division by the second $p$ is
$B_q(-2/9+14/63)=0$. The coefficient of $W_c$ is
\[
 B_q\left(-\frac16-\frac{g_K}{21}-\frac{2\beta}{21}\right)
 \equiv-\frac{B_q}{6}\pmod p
\]
by \eqref{VIII-E8:beta-cancellation}. Thus the entire carry-support
contribution $\mathcal V_p$ satisfies
\[
 \mathcal V_p\in p^2\mathbb Z_p[[w]],\qquad
 \frac{\mathcal V_p}{p^2}
 \equiv-\frac{B_q}{6}\sum_{c=a+1}^{2a-1}W_c(0,w)\pmod p.
\]
Equation \eqref{VIII-E8:beta-cancellation} also reduces
\eqref{VIII-E8:U-intermediate} to $\mathcal U_p/p^2\equiv-B_qB_1/12$.

Finally, for $c>a$ one has
$P_c(t)=(p/3+t)U_c(t)$, so
$\partial_t(P_c(t)F(c+t))|_0\equiv W_c(0)\pmod p$.
The reflection gives $W_{K-c}(0)\equiv W_c(0)\pmod p$.
Consequently
\[
 \partial_t Z_p(0,w)
 \equiv B_1(w)+2\sum_{c=a+1}^{2a-1}W_c(0,w)\pmod p.
\]
Adding the unit and carry contributions proves the first formula.
Multiplication by the exact integral factor $A_n/(1-2w)$ proves
the companion-numerator formula. All numerical denominators used
above are units: the hypotheses exclude $2,3,7$.
\end{proof}

The next sections determine this block derivative exactly and use its
unit-pole residue to complete the third-level congruence.

\section{Simple poles of the full-block derivative}\label{loc-VIII-21}
\label{VIII-sec:E8-simple-poles}
Let $p=3a-1>3$ be prime, so $a$ is even. Work over $\mathbb F_p$.
Set
\[
 y_c(t)=\frac{(t-1/3)_c}{(1+t)_c}.
\]
\[
 K_c(t,w)=\prod_{j=0}^{c-1}
 \frac{(1-(2-6j-6t)w)(1-(3j+3t-1)w)(1-(3j+3t+1)w)}
 {(1-(6j+6t+6)w)(1-(-3j-3t-3)w)(1-(-3j-3t-1)w)},
\]
\[
 F(x,w)=(6x-1)(1+(18x^2-6x-2)w),\qquad
 Z_p(t,w)=\sum_{c=0}^{p-1}y_c(t)K_c(t,w)F(c+t,w).
\]
Every derivative in this section is taken first as a rational function
in $t,w$ and then specialized to $t=0$.

\begin{theorem}\label{VIII-E8:simple-poles}
The rational function $\partial_tZ_p(0,w)$ has at most simple poles
at every finite point. Its possible poles are $w=1/x$ with
$x\in\mathbb F_p^\times$. Moreover it is regular at $w=1/2$.
\end{theorem}
\begin{proof}
All denominators in $y_c(t)$ are nonzero at $t=0$, since $c<p$.
For $0\le c\le a$, $y_c(0)$ is nonzero, while for $a<c<p$ it has a
simple zero. Thus the latter range contributes only
$y'_c(0)K_c(0,w)F(c,w)$ to the derivative.

Fix a nonzero root value $x\in\mathbb F_p$. In the complete index
interval $0\le j<p$, the three denominator root progressions of $K$
hit $x$ at unique positions $D_1,D_2,D_3$, and its numerator
progressions hit it at positions $N_1,N_2,N_3$. Positions are chosen
in $\{0,\ldots,p-1\}$. Direct substitution gives
\begin{equation}\label{VIII-E8:root-reflection}
 N_i\equiv a-1-D_i\pmod p \qquad(i=1,2,3).
\end{equation}
The root slopes with respect to $t$ are opposite in each paired
numerator/denominator progression.
The pole order of $K_c(0,w)$ at $1/x$ is exactly
\[
 q_x(c)=\#\{i:D_i<c\}-\#\{i:N_i<c\}.
\]
A negative value means a zero.

\paragraph{The range $0\le c\le a$.}
Only positions in $0\le j<a$ can contribute here. Reflection
\eqref{VIII-E8:root-reflection} preserves this interval, so its numerator
positions are the reflections $a-1-D_i$ of its denominator positions.
There are at most two denominator occurrences. If two occur, their
positions have sum at least $a-1$. Here are the direct checks:
\[
\begin{array}{c|c|c}
\text{coincident progressions}&\text{index equation}&\text{consequence}\\\hline
D_1,D_2&2j_1+j_2=3a-4&j_1+j_2\ge2a-3\ge a-1\\
D_1,D_3&2j_1+j_3=2a-3&j_1\le a-2,\ j_1+j_3\ge a-1\\
D_2,D_3&j_2=a-1,\ j_3=0&j_2+j_3=a-1.
\end{array}
\]
The last pair has root $x=-1$, which does not occur in $D_1$ on this
interval, excluding a triple coincidence. A two-position multiset
whose sum is at least $a-1$ never has more elements in an initial
interval than its reflected multiset. Consequently $q_x(c)\le0$
in that case. If there is exactly one denominator occurrence,
$q_x(c)\le1$.

It follows that a positive unit-range pole is simple and comes from
a unique denominator factor in the entire interval $j<a$.
Whenever $q_x(c)=1$, the reflected index $c'=a-c$ also has
$q_x(c')=1$. The exact finite-field reflection gives
\[
 y_{a-c}(0)K_{a-c}(0,w)=y_c(0)K_c(0,w),\qquad
 F(a-c,w)=-F(c,w).
\]
The first identity follows from $v_{a-1-j}(0,w)=v_j(0,w)^{-1}$
for $v_j=y_{j+1}K_{j+1}/(y_jK_j)$, and the resulting full product
$v_0\cdots v_{a-1}=1$; the evenness of $a$ ensures its constant
value is one. The second uses $a=1/3$ in the field.
The two summands therefore have opposite residues at their common
simple pole. Their unique denominator factor is the same factor,
so under the $t$-deformation their pole location has the same
velocity. Their derivatives' double-pole coefficients consequently
cancel. If $c=a-c$, the residue is already zero because
$6c-1=0$. All summands with $q_x(c)\le0$ acquire at most a simple
pole upon differentiation, since their logarithmic derivative has
at most simple poles. This proves the assertion for the unit range.

\paragraph{The range $a<c<p$.}
Positions $D_i<a$ cancel their reflected numerator positions in
every such prefix. The remaining positions are reflected by
$D\mapsto4a-2-D$ inside $[a,p-1]$. Thus
$q_x(c)=q_x(p+a-c)$, and it suffices to consider $a<c\le2a-1$.
Write $u=D_3$. The other two denominator positions are explicitly
\[
 D_2=\begin{cases}u+a-1,&u\le2a-1,\\u-2a,&u\ge2a,\end{cases}
\]
\[
 D_1=\begin{cases}
 (5a-4-u)/2,&u\text{ even},\\
 (2a-3-u)/2,&u\text{ odd},\ u\le2a-3,\\
 (8a-5-u)/2,&u\text{ odd},\ u\ge2a-1.
 \end{cases}
\]
If $u$ is odd, $D_1$ lies either below $a$ or above $2a-1$;
and $D_2,D_3$ cannot both lie in $[a,2a-2]$. Thus $q_x(c)\le1$.
For even $u<a$, or even $u\ge2a$, at most one denominator position
lies in that interval, with the same conclusion.
It remains to consider even $a\le u\le2a-2$. The two possible
left-hand positions are $u$ and $(5a-4-u)/2$. The third position
is $D_2=u+a-1$, whose reflected numerator position is
$v=3a-1-u$. For both denominator positions to precede $v$, one must
have
\[
 \frac{5a-4-u}{2}<3a-1-u,\quad\hbox{or equivalently }u<a+2.
\]
Thus only $u=a$ can give a double pole. In this case $x=-2$, the
positions are $a,2a-2,2a-1$, and $q_x(c)=2$ only at
$c=2a-1$ and its reflection $c=2a$.
At these two exceptional terms,
\[
 18c^2-6c-2=2\quad\text{in }\mathbb F_p,
\]
so $F(c,w)$ has the factor $1+2w$, which vanishes at $w=1/x=-1/2$.
It removes one pole order. Since $y_c(0)=0$, no derivative of $K_c$
occurs in this range. All carry contributions therefore have at most
simple poles as well.

\paragraph{Regularity at $w=1/2$.}
For $x=2$, the numerator positions are $0,1,a$, and the denominator
positions are $a-2,a-1,p-1$. If $a\ge4$, all unit-range prefixes have
negative pole order except $c=0,a$. The index $c=0$ is regular.
At $c=a$ differentiation can introduce only a simple pole, but
$F(a,1/2)=0$ removes it. Every carry prefix has a numerator excess
and is regular. For $a=2$ ($p=5$), the only additional possible
unit-range pole is at $c=1$, where $6c-1=0$ in $\mathbb F_5$;
the same cancellation proves regularity. This completes the proof.
\end{proof}

\section{The leading infinity coefficient of the full-block derivative}\label{loc-VIII-22}
\label{VIII-sec:E8-leading-infinity}
Let $Z_p(t,w)$ be the full-block function defined above. We prove
its leading derivative coefficient by an exact finite summation
with a rational telescoping certificate.

\begin{lemma}[A finite very-well-poised identity]\label{VIII-E8:finite-vwp}
For $N\ge0$, as a rational-function identity in $A,B$,
\[
 \sum_{k=0}^N(-1)^k\frac{A+2k}{A}
 \frac{(A)_k(B)_k(-N)_k}
 {k!(1+A-B)_k(1+A+N)_k}
 =\frac{(1+A)_N}{(1+A-B)_N}.
\]
\end{lemma}
\begin{proof}
Write the summand as $t_{N,k}$. For $N\ge1$ put
\[
 r_N=\frac{A+N}{A-B+N},\qquad
 R_{N,k}=-\frac{k(A-B+k)(A+N+k)}
 {(A+2k)N(A-B+N)}.
\]
The successive-term ratios verify the exact rational identity
\[
 R_{N,k+1}t_{N,k+1}-R_{N,k}t_{N,k}
 =t_{N,k}-r_Nt_{N-1,k}.
\]
Indeed the numerator simplification required is just
\[
 -(A+k)(B+k)(k-N)-k(A-B+k)(A+N+k)
 =-(A+2k)(Ak-BN+k^2).
\]
At $k=0$ the boundary vanishes by the factor $k$, and at $k=N+1$
it vanishes by termination. Summing yields $S_N=r_NS_{N-1}$,
with $S_0=1$, proving the displayed product. The argument is first
made at generic parameters; the resulting rational identity extends
to every specialization at which the original two sides are defined.
\end{proof}

\begin{proposition}\label{VIII-E8:leading-infinity-jet}
For every prime $p>3$ with $p\equiv2\pmod3$,
\[
 \boxed{[w^1]_\infty\,\partial_tZ_p(0,w)=-24\quad\text{in }\mathbb F_p.}
\]
Here the coefficient refers to the Laurent expansion at infinity.
\end{proposition}
\begin{proof}
Put $a=(p+1)/3$, so $a$ is even and $a\equiv1/3\pmod p$.
For $0\le c\le a$, all leading kernel roots are units modulo $p$,
and the leading coefficient of the product is
\[
 P_c(t,w)=u_c(t)+O(w^{-1}),\qquad
 u_c(t)=(-1)^c\left(\frac{(t-1/3)_c}{(1+t)_c}\right)^3,
\]
to first order in $t$ at zero. This follows directly from
$f_x(w)=-((x-1/3)/(x+1))^2+O(w^{-1})$ in this range.
For $c>a$, the scalar factor has a simple zero at $t=0$ and the
kernel at zero has degree $-2$ at infinity. Consequently the
derivative of its weighted summand has degree at most $-1$;
these terms contribute nothing to the coefficient of $w^1$.

Write $x=c+t$ and $G_c(t)=x^3u_c(t)$. Directly from the ratio
$u_{c+1}/u_c=-((x-1/3)/(x+1))^3$ one obtains
\[
 (6x-1)(18x^2-6x-2)u_c
 =-54\bigl(G_{c+1}-G_c\bigr)-4(6x-1)u_c.
\]
After reduction modulo $p$, replace $1/3$ by $a$. The derivative
of $G_0=t^3$ at zero is zero. Also $u_{a+1}$ contains the factor
$t^3$, so the derivative of $G_{a+1}$ is zero. Therefore the desired
coefficient is $-4L'(0)$, where
\[
 L(t)=\sum_{c=0}^a(6c+6t-3a)(-1)^c
                 \frac{(t-a)_c^3}{(1+t)_c^3}.
\]

To evaluate this derivative, use the terminating deformation
\[
 V(t)=\sum_{c=0}^a3(A+2c)(-1)^c
       \frac{(A)_c(B)_c(-a)_c}
       {c!(1+A-B)_c(1+A+a)_c},\qquad
 A=-a+2t,\quad B=-a+t.
\]
At $t=0$, each summand equals its counterpart in $L(0)$.
Its derivative also equals that counterpart's derivative: the
upper-parameter velocities are $2,1,0$ and the lower velocities
are $0,1,2$, yielding the same logarithmic derivative as the
three identical upper and three identical lower parameters of
$L$. The weight has derivative $6$ in both sums. None of the
numerical factors in these terms vanishes for $c\le a$ except
possibly the explicit linear weight; equality of derivatives
there follows by the product rule without division by that weight.
Thus $L'(0)=V'(0)$.

Lemma~\ref{VIII-E8:finite-vwp} gives the exact rational identity
\[
 V(t)=3(-a+2t)\frac{(1-a+2t)_a}{(1+t)_a}.
\]
There is exactly one zero numerator factor at $t=0$. As $a$ is even,
\[
 V'(0)=3(-a)\frac{2(-1)^{a-1}(a-1)!}{a!}=6.
\]
All denominators used at zero are units modulo $p$. It follows
that the leading coefficient is $-4\cdot6=-24$, as asserted.
\end{proof}


\section{An integral residue functional for the unit poles}\label{loc-VIII-23}
The following lemma justifies passage from a coefficientwise
congruence to the total residue at the unit poles, even when their
reductions collide.

\begin{lemma}[Integral unit-pole residue functional]
Let $B,C\in\mathbb Z_p[w]$ have constant term one. Suppose the leading
coefficient of $B$ is a unit and all its roots are $p$-adic units, while
$C(w)=\prod_i(1-u_iw)$ with $u_i\in p\mathbb Z_p$. For
$P\in\mathbb Z_p[w]$, the sum of the residues of $P/(BC)$ at the
unit poles belongs to $\mathbb Z_p$. This functional preserves
coefficientwise congruences modulo $p^a$ after retaining the same
denominator factors.

For $a=3$, replace every factor involving a root divisible by $p$,
including inverse factors, by its expansion modulo $p^3$. The resulting
expression has only unit poles. The sum of its finite residues, equivalently
its coefficient of $w^{-1}$ in the Laurent expansion at infinity, agrees
modulo $p^3$ with the original unit-pole residue sum.
\end{lemma}
\begin{proof}
Since $C\equiv1\pmod p$ and $B$ has unit leading coefficient, their
resultant is a unit. Bezout decomposition separates the parts with
denominators $B$ and $C$ over $\mathbb Z_p[w]$. The part with denominator
$B$ has an integral Laurent expansion at infinity, proving integrality
of its total finite residue. The other part has no unit pole.

The coefficient of $w^j$ in $C^{-1}$ is divisible by $p^j$. Thus its
truncation through degree two differs from $C^{-1}$ by $p^3H/C$ for
an integral polynomial $H$. Numerator factors with roots in $p\mathbb Z_p$
have the same divisibility property. Apply the first assertion to the
resulting errors. Repeated poles and cancelled factors cause no change
to this argument.
\end{proof}


\section{Closure of the third odd level}\label{loc-VIII-24}
\label{VIII-sec:e8-r3-closure}
For a prime \(p>3\) with \(p\equiv2\pmod3\), put
\[
 q=(p-2)/3,\qquad a=q+1,\qquad b_q=(1/3)_q/q!,
\]
and retain the exact primary and companion interpolation families.
The result of this section is unconditional; its rational-function
ingredients are the full-block derivative theorem, the simple-pole
theorem, and the first Laurent-coefficient theorem proved above.
No finite unit projector or large-prime restriction is used here.

\begin{theorem}\label{VIII-thm:e8-r3-closed}
For every prime \(p>3\), \(p\equiv2\pmod3\),
\[
 \boxed{\displaystyle
 \sum_{k=0}^{(p^3-2)/3}(-1)^k(6k+1)
       \left(\frac{(1/3)_k}{k!}\right)^3
 \equiv
 p^2\sum_{k=0}^{(p-2)/3}(-1)^k(6k+1)
       \left(\frac{(1/3)_k}{k!}\right)^3
 \pmod{p^8}.}
\]
\end{theorem}

\subsection{The exact block derivative}
Let \(J_p(w)=\partial_tZ_p(0,w)\) be the rational block derivative.
The simple-pole theorem says that its only finite poles are simple
poles at \(w=1/x\), \(x\in\mathbb F_p^\times\), and that it is regular
at \(w=1/2\).  The first Laurent-coefficient theorem gives
\(\deg_\infty J_p\le1\) and \([w^1]_\infty J_p=-24\).
We record explicitly why its required zero at the origin holds also
for \(p=5\) and \(p=11\).

Set \(M=p^2\), \(n=(p^2-1)/3\), and \(c_n=(2/3)_n/n!\).
We first record its normalizing unit explicitly.  Since
\(n=pq+2q+1\), its numerator has \(q+1\) factors divisible by \(p\),
with quotients \(1,4,\ldots,3q+1\), whereas its denominator has
\(q\), with quotients \(1,\ldots,q\); none is divisible by \(p^2\).
Thus \(v_p(c_n)=1\).  The complete unit blocks cancel modulo \(p\).
The remaining partial unit quotient is
\[
 \frac{3^q q!\prod_{j=0}^{q-1}(3j+2)}{(2q+1)!}.
 \]
Combining it with the stripped factors gives
\[
 \frac{c_n}{p}
 \equiv 3^{-n}\frac{(p-1)!}{q!(2q+1)!}
 \equiv-1\pmod p.
\]
Indeed \(n=(p-1)(p+1)/3\) is a multiple of \(p-1\);
Wilson's theorem and odd \(q\) give \(q!(2q+1)!\equiv1\pmod p\).
The exact moment and numerator normalizations give
\[
 \mathcal N_n(w)=-M c_n\sum_{j\ge0}R_j(M)w^j.
\]
For \(j<p-1\), the degree of \(R_j\) in \(M^2\) is at most
\(\lfloor(2j+1)/3\rfloor\le2q\).  The exact triangular moment
operator has diagonal entries \(3k+1\); all entries with
\(0\le k\le2q\) are \(p\)-adic units.  Its integral recurrence and
integral initial moments therefore give \(R_j\in\mathbb Z_p[M^2]\)
in this range.  Since \(v_p(c_n)=1\),
\(\mathcal N_n/p^2\) has zero coefficients modulo \(p\) in degrees
less than \(p-1\).

Write
\[
 A_n(w)=
 \frac{\prod_{D\in\mathcal D_n}(1-Dw)}
      {\prod_{H\in\mathcal H_n}(1-Hw)}.
\]
This notation denotes the common reversed-root quotient, not the
adjacent-cutoff polynomial also sometimes denoted by \(A_n\).
The full-block derivative theorem gives
\[
 \frac{\mathcal N_n(w)}{p^2}
   \equiv-\frac{b_qA_n(w)}{12(1-2w)}J_p(w)\pmod p.
\]
All factors outside \(J_p\) are invertible at zero, so \(J_p\) has a
zero of order at least \(p-1\).  Consequently
\((w^{p-1}-1)J_p(w)\) is a polynomial of degree at most \(p\),
divisible by \(w^{p-1}\).  The coefficient at infinity and regularity
at \(1/2\) determine its two remaining coefficients:
\begin{equation}\label{VIII-eq:e8-full-J}
 \boxed{J_p(w)=12(1-2w)\frac{w^{p-1}}{w^{p-1}-1}
       \quad\text{in }\mathbb F_p(w).}
\end{equation}
In particular,
\begin{equation}\label{VIII-eq:e8-N-reduced}
 \overline{\mathcal N_n/p^2}
    =-b_q\overline A_n(w)\frac{w^{p-1}}{w^{p-1}-1}.
\end{equation}

\subsection{The unit residue and its normalization}
Let \(\widetilde\mu_M(H)\) denote the companion masses and
\(\nu=(-1)^nMc_n^2\).  The exact partial fractions give
\[
 \sum_{p\nmid H}\widetilde\mu_M(H)
   =-\frac{\nu}{2}
       \sum_{\rho\ {\rm unit\ pole}}
                  \operatorname{Res}_{w=\rho}\mathcal N_n(w).
\]
The full-block theorem proves coefficientwise divisibility of
\(\mathcal N_n\) by \(p^2\).  This is also an integral rational
divisibility statement with its original denominator: multiply the
series by that denominator, which is integral with constant term one.
Every coefficient of the resulting polynomial is divisible by \(p^2\).
The integral unit-pole residue lemma may therefore be applied directly
to \(\mathcal N_n/p^2\) modulo \(p\).  Its reduction has only unit
poles, and its total residue is its coefficient of \(w^{-1}\) at
infinity.  Since \(n\) is even and \(c_n/p\equiv-1\pmod p\),
\(\nu/p^4\equiv1\pmod p\).  Formula \eqref{VIII-eq:e8-N-reduced} yields
\begin{equation}\label{VIII-eq:e8-unit-residue-final}
 p^{-6}\sum_{p\nmid H}\widetilde\mu_{p^2}(H)
 \equiv\frac{b_q}{2}[w^{-1}]_\infty
       \left(\overline A_n(w)\frac{w^{p-1}}{w^{p-1}-1}\right)
       \pmod p.
\end{equation}

For completeness the leading normalization of \(\overline A_n\)
can be found without any harmonic evaluation.  A complete block of
\(p\) indices has the same three residue progressions in the numerator
and denominator of \(A_n\); their products cancel modulo \(p\).
As \(n=pq+(2p-1)/3\), the remaining quotient is the companion quotient
at cutoff \(M_*=2p\), with \(n_*=(2p-1)/3=2q+1\).
Its only nonunit nodes are \(p,2p,-2p\), and all its denominator
roots are units.  Put \(c_*=(2/3)_{n_*}/n_*!\).
Direct multiplication of the root products gives
\[
 \frac{\prod_{H\in\mathcal H_{n_*}}H}
      {\prod_{D\in\mathcal D_{n_*}}D}
       =(-1)^{n_*}(2p)c_*^2=-2pc_*^2.
\]
Removing the three nonunit nodes, whose product is \(-4p^3\), shows
that the leading coefficient of its degree-two unit quotient is
\[
 \lambda=\frac{2}{(c_*/p)^2}\pmod p.
\]
The elementary Wilson normalization
\[
 (c_*/p)b_q\equiv1\pmod p
\]
then gives \(\lambda=2b_q^2\).
Indeed, after deleting the unique factor \(p\) in \(c_*\), the remaining
numerator residues from \(c_*/p\) and \(b_q\) together are exactly
\(1,\ldots,p-2\).  Their denominator is
\(3^{p-1}(2q+1)!q!\), and both \((p-2)!\) and
\((2q+1)!q!\) equal \(1\) modulo \(p\), by Wilson's theorem and
the oddness of \(q\).

Define the reciprocal unit-root quotient by
\[
 Q(X)=
 \frac{\prod_{D\in\mathcal D_n,\ p\nmid D}(1-X/D)}
      {\prod_{H\in\mathcal H_n,\ p\nmid H}(1-X/H)}
       =1+\sum_{j\ge1}e_jX^j\in\mathbb Z_p[[X]].
\]
After reduction, the preceding normalization says
\[
 \overline A_n(w)=2b_q^2w^2\overline Q(w^{-1}).
\]
Since \(p-1\ge4\), the factor
\(w^{p-1}/(w^{p-1}-1)=1+O(w^{-(p-1)})\) cannot affect the coefficient
of \(w^{-1}\) in a function of degree two.  Equation
\eqref{VIII-eq:e8-unit-residue-final} therefore becomes
\begin{equation}\label{VIII-eq:e8-unit-mass-cubic}
 \boxed{\displaystyle
 \sum_{p\nmid H}\widetilde\mu_{p^2}(H)
       \equiv p^6 b_q^3e_3\pmod{p^7}.}
\end{equation}
This computation applies uniformly to \(p=5\), \(p=11\), and all
larger primes in the specified residue class.

\subsection{Cancellation with the exact transport correction}
The companion-to-primary weight identity has kernel \(Q(pX)\):
\[
 \widetilde\mu_{p^2}(pH)=p\mu_p(H)Q(pH).
\]
Summing over the nodes and expanding gives the exact convergent identity
\[
 \mathcal A(p^2)-pE(p)
 =\sum_{p\nmid H}\widetilde\mu_{p^2}(H)
      +p\sum_{j\ge1}p^je_jM_j(p).
\]
The first two primary moments vanish.  The proved uniform estimate
\(M_j(p)\in p^2\mathbb Z_p\) for every \(j\ge1\) puts all terms
with \(j\ge4\) in \(p^7\mathbb Z_p\), and the valuations tend to
infinity.  The exact cubic moment is
\[
 M_3(p)=p^2(p-1)^3b_q^3\equiv-p^2b_q^3\pmod{p^3}.
\]
Thus the cubic transport contribution is
\(-p^6e_3b_q^3\pmod{p^7}\), cancelling
\eqref{VIII-eq:e8-unit-mass-cubic}.  Hence
\[
 \mathcal A(p^2)\equiv pE(p)\pmod{p^7}.
\]
Finally the independently proved first-step congruence
\(E(p^3)\equiv p\mathcal A(p^2)\pmod{p^8}\), valid for every prime
in the present domain, proves Theorem~\ref{VIII-thm:e8-r3-closed}.


\section{The completed odd branch and verification}\label{loc-VIII-25}
The higher-level arguments and their exhaustive finite certificates
prove the target for every odd $r\ge5$. Theorem~\ref{VIII-thm:e8-r3-closed}
proves it for $r=3$ for every permitted prime. Together they establish
Theorem~\ref{VIII-thm:main-range} without changing the sum, cutoff,
recurrence multiplier, or target modulus.

The working mechanism in the final step is specific: reflection
cancels the potential double poles of a parameter derivative;
a first-resonance argument supplies a high-order zero; rational
uniqueness then determines the whole function. This proves the
needed unit residue, which equals the negative of the cubic transport
term. It is not an assumption that either term vanishes separately.
The elementary rational uniqueness argument is classical; no claim
that its ingredients or this application have historical priority
is made. Priority and comparison with every relevant preceding proof
require a separate literature assessment.

The new root classification was also checked by exact arithmetic
for all $48$ primes $p\le500$ in the stated residue class. The
original, unchanged third-level sums were independently evaluated
modulo $p^{10}$, keeping valuations separate from unit inversions.
The following are supplementary checks, not the uniform proof:
\[
\begin{array}{c|rrrrrrrrr}
p&5&11&17&23&29&41&47&53&59\\\hline
v_p(E(p^3)-p^2E(p))&8&8&9&8&8&8&8&8&8
\end{array}
\]
The higher-level proof retains its $158$ and $56$ exact moment
certificates. The accompanying archive contains the derivation
fragments, exact verification programs, their recorded outputs,
and a checksum manifest. This single source embeds the principal
finite-certificate evaluators and the new direct sum evaluator.

\LocalAppendix
\section{Finite certificates on the two positive companion grids}\label{loc-VIII-A}

\subsection{A finite certificate on the grid $H=3h+1$}
Retain the rational deformation $V_n(z)$ defined in the companion
construction. For every $h\ge0$ and $n\ge h$,
\begin{equation}\label{VIII-E7positiveone:main}
 V_n(3h+1)=-\frac{(-1)^h(3h+1)}2.
\end{equation}
All lower parameters at $z=3h+1$ are regular: in particular
$(5-3h)/6$ is never an integer. The factor $(-h)_k$ makes the
sum terminate at $k=h$, so it suffices to prove the identity for $n=h$.

Set
\[
 P(h,k)=18k^2-6k+3h-1,
\]
\[
 u(h,k)=\frac{(-1)^{k-h}(6k-1)P(h,k)}{(1-3h)(3h+1)}
 \frac{(-1/3)_k((3h-1)/6)_k(-h-2/3)_k(-h)_k}
 {k!((5-3h)/6)_k(h+4/3)_k(h+2/3)_k}.
\]
This is exactly the summand of $V_h(3h+1)$ divided by the
right side of \eqref{VIII-E7positiveone:main}. We prove
$\sum_{k=0}^hu(h,k)=1$ by a recurrence with step two in $h$.

Define the polynomial
\begin{align*}
 A(h,k)={}&486h^5+(1458k^2-1944k+3564)h^4\\
 &+(7857k^2-10476k+10233)h^3\\
 &+(-972k^4+2592k^3+13149k^2-19836k+14229)h^2\\
 &+(-2511k^4+6696k^3+7137k^2-15468k+9486)h\\
 &+243k^6-972k^5-189k^4+3384k^3+414k^2-4072k+2392.
\end{align*}
The rational certificate is
\[
 R(h,k)=\frac{-k(3k-2)(3h-6k+1)A(h,k)}{D_+(h,k)},
\]
where
\begin{align*}
 D_+(h,k)={}&(6k-1)(h-k+1)(h-k+2)(3h-3k+5)(3h-3k+8)\\
 &\quad\cdot(3h+3k+2)(3h+3k+4)P(h,k).
\end{align*}
Cancellation of consecutive rising factors and expansion of the
remaining polynomials verifies the rational identity
\begin{equation}\label{VIII-E7positiveone:certificate}
 u(h+2,k)-u(h,k)
 =R(h,k+1)u(h,k+1)-R(h,k)u(h,k).
\end{equation}
For an independently checkable form, if
$q(h,k)=u(h,k+1)/u(h,k)$ and $r(h,k)=u(h+2,k)/u(h,k)$,
then the certificate is the polynomial equality obtained by clearing
all denominators in
$q(h,k)R(h,k+1)-R(h,k)=r(h,k)-1$.
The ratio $q$ is explicitly
\[
 q(h,k)=-\frac{(6k+5)P(h,k+1)}{(6k-1)P(h,k)}
 \frac{(k-1/3)(k+(3h-1)/6)(k-h-2/3)(k-h)}
 {(k+1)(k+(5-3h)/6)(k+h+4/3)(k+h+2/3)}.
\]

Here are the upper-boundary details. Put $g(h,k)=R(h,k)u(h,k)$.
The apparent poles at $k=h+1,h+2$ cancel because
\[
 \frac{(-h)_k}{(h-k+1)(h-k+2)}
 =\frac{(-h-2)_k}{(h+1)(h+2)}.
\]
Thus $g$ has a finite value at both boundary indices. Also
$g(h,0)=0$. Direct cancellation gives
\[
 \left.\frac{R(h,k)}{r(h,k)}\right|_{k=h+1}
 =-\frac{4(3h+1)(9h^3+39h^2+46h+15)}
 {(h+2)(6h+5)(18h^2+33h+17)},
\]
\[
 q(h+2,h+1)=
 -\frac{5(3h+2)(9h+11)}
 {(h+2)(6h+5)(18h^2+33h+17)}.
\]
The first displayed quotient plus $1+q(h+2,h+1)$ is identically
zero. Multiplication by $u(h+2,h+1)$ therefore gives
\[
 g(h,h+1)=-u(h+2,h+1)-u(h+2,h+2).
\]
Summing \eqref{VIII-E7positiveone:certificate} for $0\le k\le h$
and adding these two final terms proves
\[
 \sum_{k=0}^{h+2}u(h+2,k)=\sum_{k=0}^hu(h,k).
\]
At $h=0$ the sum is $1$; at $h=1$ its terms are $-1/4$ and
$5/4$, again summing to one. Induction separately on the two
parity classes proves \eqref{VIII-E7positiveone:main} for every $h$.

This proves one exact node family of the companion interpolation.
It does not, by itself, prove any prime-power lifting congruence.


\subsection{The second positive grid: a complete finite certificate}
For $h\ge1$, let $t_h(k)$ be the summand of $V_{h+1}(3h+2)$
and normalize it by $u_h(k)=t_h(k)/(3h+2)$. Explicitly,
\[
 u_h(k)=\frac{(-1)^k(6k-1)(h+6k^2-2k)}{2h(3h+2)}
 \frac{(-1/3,h/2,-h-1,-h-1/3)_k}
 {k!(2/3-h/2,h+5/3,h+1)_k}.
\]
The sum terminates at $h+1$. The desired exact assertion is
\[
 \sum_{k=0}^{h+1}u_h(k)=
 \begin{cases}1,&h\text{ odd},\\0,&h\text{ even}.\end{cases}
\]
All denominator Pochhammers are nonzero on this range.

Put
\begin{align*}
 J(h,k)={}&54h^5+162h^4k^2-216h^4k+486h^4
 +1089h^3k^2-1452h^3k+1725h^3\\
 &-108h^2k^4+288h^2k^3+2442h^2k^2-3512h^2k+3002h^2\\
 &-351hk^4+936hk^3+2082hk^2-3608hk+2549h\\
 &+27k^6-108k^5-126k^4+656k^3+507k^2-1316k+840,
\end{align*}
\[
 R(h,k)=-\frac{k(3k-2)(3h-6k+2)J(h,k)}{D_2(h,k)},
\]
where
\begin{align*}
 D_2(h,k)={}&(6k-1)(h-k+2)(h-k+3)(h+k+1)(h+6k^2-2k)\\
 &\quad\cdot(3h-3k+4)(3h-3k+7)(3h+3k+5).
\end{align*}
Direct rational cancellation proves the finite WZ certificate
\begin{equation}\label{VIII-E7:F-certificate}
 u_{h+2}(k)-u_h(k)
 =R(h,k+1)u_h(k+1)-R(h,k)u_h(k).
\end{equation}
For a compact independently reproducible verification, the consecutive
ratios are
\[
 q(h,k)=\frac{u_h(k+1)}{u_h(k)}
 =\frac{(h+2k)(3k-1)(6k+5)(h-k+1)(3h-3k+1)
 (h+6k^2+10k+4)}
 {(k+1)(6k-1)(h+k+1)(h+6k^2-2k)
 (3h-6k-4)(3h+3k+5)},
\]
\begin{align*}
 r(h,k)&=\frac{u_{h+2}(k)}{u_h(k)}\\
 &=\frac{(h+1)(h+2)(h+3)(h+2k)(3h+4)(3h+5)(3h+7)}
 {(h-k+2)(h-k+3)(h+k+1)(h+k+2)(h+6k^2-2k)}\\
 &\quad\times\frac{(3h-6k+2)(h+6k^2-2k+2)}
 {(3h-3k+4)(3h-3k+7)(3h+3k+5)(3h+3k+8)}.
\end{align*}
Substitution in $q(h,k)R(h,k+1)-R(h,k)=r(h,k)-1$ leaves
identically zero after clearing denominators.

At the lower boundary $R(h,0)u_h(0)=0$. At the upper boundary,
cancel the apparent zeros and poles before evaluating. For instance,
\[
 \frac{(-h-1)_k}{(h-k+2)(h-k+3)}
 =\frac{(-h-3)_k}{(h+2)(h+3)}.
\]
The needed endpoint can also be checked entirely from rational ratios:
\[
 \left.\frac{R(h,k)}{r(h,k)}\right|_{k=h+2}
 =-\frac{4(9h^3+60h^2+127h+87)}{(h+3)(6h+11)^2},
\]
\[
 \frac{u_{h+2}(h+3)}{u_{h+2}(h+2)}
 =-\frac{3(3h+5)}{(h+3)(6h+11)^2}.
\]
The first displayed ratio plus one plus the second ratio is exactly
zero. Hence the upper boundary is
$-u_{h+2}(h+2)-u_{h+2}(h+3)$. Summing
\eqref{VIII-E7:F-certificate} over $0\le k\le h+1$ therefore proves that
the normalized sum is unchanged when $h$ increases by two.
The initial normalized terms are
\[
 h=1:\quad -\frac1{10},\ \frac54,\ -\frac3{20},
 \qquad
 h=2:\quad-\frac1{16},\ -\frac{105}{176},\ \frac{11}{16},\ -\frac5{176}.
\]
Their sums are respectively one and zero, proving both parity cases.
At $h=0$ the removable value $V_n(2)=0$ for $n\ge1$ supplies the
remaining case. For an ambient truncation $n\ge h+1$, all subsequent
terms vanish and their denominators remain nonzero; the same node
value therefore holds for every admissible ambient cutoff.


\section{A finite certificate for the negative companion grid}\label{loc-VIII-B}
\label{VIII-sec:E7negative}
The following establishes the negative grid of the proposed companion
interpolation. It is an exact finite identity and makes no claim about
the remaining prime-power error estimate.

For $h\ge0$, set
\[
 t(h,k)=-\frac{(-1)^k(6k-1)(3k^2-k-h-1)}{2(h+1)}
 \frac{(-1/3)_k(-h-1)_k(2h+1)_k(2h+5/3)_k}
 {k!(h+5/3)_k(-2h-1/3)_k(-2h-1)_k}.
\]
All denominator factors are nonzero for $0\le k\le h+1$.
We prove
\begin{equation}\label{VIII-E7neg:sum}
 \sum_{k=0}^{h+1}t(h,k)=3h+2.
\end{equation}
Put $u(h,k)=t(h,k)/(3h+2)$ and define the polynomial
\[
\begin{split}
 A(h,k)={}&1728h^5-2592h^4k^2+3456h^4k+9504h^4\\
 &-12024h^3k^2+16032h^3k+20712h^3\\
 &+432h^2k^4-1152h^2k^3-19704h^2k^2+27296h^2k+22408h^2\\
 &+1026hk^4-2736hk^3-13308hk^2+20176hk+12074h\\
 &-27k^6+108k^5+450k^4-1520k^3-3027k^2+5444k+2604.
\end{split}
\]
The certificate is
\[
 R(h,k)=-\frac{k(3k-2)(2h-k+2)(6h-3k+4)A(h,k)}{D_-(h,k)},
\]
where
\begin{align*}
 D_-(h,k)={}&16(h+1)(2h+1)(2h+3)(3h+4)(6h+5)(6h+7)\\
 &\quad\cdot(6k-1)(h-k+2)(h-3k^2+k+1).
\end{align*}
Direct rational cancellation gives
\begin{equation}\label{VIII-E7neg:WZ}
 u(h+1,k)-u(h,k)
 =R(h,k+1)u(h,k+1)-R(h,k)u(h,k).
\end{equation}
For reproducibility, this is the polynomial identity obtained after
substituting the displayed $R$ and the consecutive-factor ratios of
$u$; the exact residual is zero. The quadratic factor
$h-3k^2+k+1$ cancels from the product $Ru$, so roots of that factor
cause no singularity.

We specify the upper boundary without interpreting an undefined
terminating summand. Write $r(h,k)=u(h+1,k)/u(h,k)$ as its rational
consecutive-parameter ratio. Cancellation gives
\[
 \left.\frac{R(h,k)}{r(h,k)}\right|_{k=h+2}=-1.
\]
The canceled quotient has no pole at this substitution for any
integer $h\ge0$. Thus the continued product at the upper boundary is
\[
 [R(h,k)u(h,k)]_{k=h+2}=-u(h+1,h+2).
\]
This can equivalently be obtained by cancelling
$(-h-1)_k/(h-k+2)=(-h-2)_k/(h+2)$ before specializing; the quotient
$R/r$ also covers the coincident zero factors at $h=0$.
At the lower boundary $R(h,0)u(h,0)=0$.
Sum \eqref{VIII-E7neg:WZ} over $0\le k\le h+1$ and add the new endpoint
$u(h+1,h+2)$. The sums of $u$ at successive $h$ are equal.
For $h=0$, the unnormalized terms are $-1/2$ and $5/2$, whose sum
is two. Hence the normalized sum equals one, proving
\eqref{VIII-E7neg:sum}.

\subsection{Embedding in the cleared interpolation}
At the negative node $H=-2(3h+2)$, the candidate $V_n(z)$ has
exactly the summand $t(h,k)$ above. Its numerator parameter
$-h-1$ terminates the finite sum at $k=h+1$.
For this node to belong to the companion grid one has $n\ge h+1$.
The companion denominator is
\[
 D_n(z)=\prod_{j=1}^n(1-z/(6j))(1+z/(3j))
         \prod_{j=0}^{n-1}(1+z/(3j+1)).
\]
The only possible zero of $D_n(H)$ is the last product's factor
at $j=2h+1$. It occurs precisely when $n\ge2h+2$.
If $h+1\le n\le2h+1$, all original summand denominators are nonzero,
the sum terminates at $h+1$, and \eqref{VIII-E7neg:sum} gives
$V_n(H)=3h+2$.

If $n\ge2h+2$, treat $D_n(z)V_n(z)$ term by term after clearing
its finite denominator factors. In a summand $k\le h+1$, the
factor indexed by $j=2h+1$ remains in the denominator-clearing
tail and vanishes at $H$. In a summand $k\ge h+2$, the numerator
contains the factor at $j=h+1$ in
$1+z/[2(3j-1)]$, and that factor vanishes at $H$.
The additional factor $2-z$ in $V_n$ is nonzero at $H$.
Thus every cleared summand is zero, giving
\[
 [D_n(z)V_n(z)]_{z=H}=0=D_n(H)(3h+2).
\]
Combining the two ranges proves the full cleared negative-grid identity
\begin{equation}\label{VIII-E7neg:cleared}
 \boxed{[D_n(z)V_n(z)]_{z=-2(3h+2)}
 =D_n(-2(3h+2))(3h+2)\qquad(0\le h<n).}
\end{equation}
No raw $0/0$ hypergeometric summand is used in this argument.

\section{Reproducible integer certificate evaluator}\label{loc-VIII-C}
The following complete program is the evaluator used for the 158
higher-level moment certificates. For the 56 fifth-level certificates,
the exhaustive case list is instead
$[(5,2),(11,2),(17,2),(23,2)]$; all other calculations are
identical. These lists supplement the uniform proofs and do not
extrapolate from a finite sample to an infinite family.
\begin{lstlisting}[language=Python]
"""Exact finite certificates for exceptional polynomial-projector parameters.

The proved degree and logarithmic denominator bounds reduce each required
prime-power moment to Newton interpolation from small integer cutoffs.
All calculations use Python integers modulo explicit prime powers; no
floating point, symbolic heuristic, or original enormous truncated sum.
"""
from pathlib import Path
from math import comb
import json


def strip(x,p):
    assert x
    v=0
    while x%p==0:x//=p;v+=1
    return x,v


def ilog(x,p):
    k=0
    while x>=p:x//=p;k+=1
    return k


def vmod(x,p,cap):
    return cap if x==0 else strip(x,p)[1]


def poch_step(unit,val,num,den,p,mod):
    nu,nv=strip(num,p);du,dv=strip(den,p)
    return unit*nu*pow(du,-1,mod)%mod,val+nv-dv


def solve_series(a,b,previous,source,J,mod):
    out=[0]*(J+1)
    for j in range(J+1):
        rhs=sum(b[i]*previous[j-i] for i in range(min(3,j)+1))
        if j<len(source):rhs+=source[j]
        rhs-=sum(a[i]*out[j-i] for i in range(1,min(3,j)+1))
        out[j]=rhs%mod
    return out


def newton_value(values,X,mod):
    delta=values[:];ans=0;binom=1
    for k in range(len(values)):
        if k:binom=binom*(X-k+1)//k
        ans=(ans+delta[0]*(binom%mod))%mod
        delta=[(delta[i+1]-delta[i])%mod for i in range(len(delta)-1)]
    return ans


def certify(p,s,family):
    indices=[i*(p-1)-(1 if family=='companion' else 3) for i in range(1,3*s+2)]
    degrees={j:2*((2*j+1)//3) if family=='companion' else 2*((2*j)//3) for j in indices}
    J=max(indices);nmax=max(degrees.values())
    log=ilog(3*nmax+(1 if family=='companion' else 2),p)
    A=(3*log+1)//2 if family=='companion' else (3*log)//2
    cap=A+2;mod=p**cap
    sample={j:[] for j in indices}
    previous=[0]*(J+1)
    du=cu=bu=1;dv=cv=bv=0
    max_norm_v=0
    for n in range(nmax+1):
        if n:
            du,dv=poch_step(du,dv,3*n-4,3*n,p,mod)
            cu,cv=poch_step(cu,cv,3*n-1,3*n,p,mod)
            bu,bv=poch_step(bu,bv,3*n-2,3*n,p,mod)
        assert min(dv,cv,bv)>=0
        if family=='companion':
            M=3*n+1
            if n==0:
                current=[mod-1]*(J+1)
            else:
                a=[1,-2,-(M-2)*(3*M-4),2*M*(M-2)**2]
                b=[1,-2,-(M-1)*(3*M-5),-2*(M-3)*(M-1)**2]
                t=(6*n-1)*du*pow(p,dv,mod)%mod
                alpha=18*n*n-6*n-2
                current=solve_series(a,b,previous,[t,t*(alpha+1),t*alpha],J,mod)
            mu,mv=strip(M,p);norm_v=mv+cv;norm_unit=mu*cu%mod
            sign=-1
        else:
            L=3*n+2
            a=[1,-1,-(L-1)*(3*L-2),2*L*(L-1)**2]
            b=[1,-1,-(L-2)*(3*L-7),-2*(L-2)**2*(L-3)]
            t=(6*n+1)*bu*pow(p,bv,mod)%mod
            current=solve_series(a,b,previous,[t],J,mod)
            lu,lv=strip(L,p);vu,vv=strip(L-1,p)
            norm_v=lv+vv+bv;norm_unit=lu*vu*bu%mod;sign=2
        assert norm_v<=A,(family,p,s,n,norm_v,A)
        max_norm_v=max(max_norm_v,norm_v)
        scale=sign*pow(p,A-norm_v,mod)*pow(norm_unit,-1,mod)%mod
        for j in indices:
            if n<=degrees[j]:sample[j].append(current[j]*scale%mod)
        previous=current
    X=(p**(2*s)-1)//3 if family=='companion' else (p**(2*s+1)-2)//3
    rows=[]
    for j in indices:
        result=newton_value(sample[j],X,mod)
        valuation=vmod(result,p,cap)-A
        target=-(s-1) if family=='companion' else -s
        assert valuation>=target,(p,s,family,j,valuation,target)
        rows.append({'normalized_index':j,'polynomial_degree_bound':degrees[j],
            'scaled_value_residue':result,'valuation_lower_bound':valuation,
            'required_lower_bound':target,'passed':True})
    return {'p':p,'s':s,'r':2*s+1,'family':family,'precision_exponent':cap,
        'common_scale_exponent':A,'max_sample_cutoff':nmax,
        'max_normalizer_valuation':max_norm_v,'number_of_moments':len(rows),
        'checks':rows}


if __name__=='__main__':
    cases=[(5,3),(11,3),(17,3),(23,3),(29,3),(5,4),(5,5)]
    records=[]
    for p,s in cases:
        for family in ['companion','primary']:
            row=certify(p,s,family);records.append(row)
            print(json.dumps({k:v for k,v in row.items() if k!='checks'}),flush=True)
    result={'status':'All exact finite projector-moment certificates passed.',
        'pair_count':len(cases),'family_pair_count':len(records),
        'moment_count':sum(r['number_of_moments'] for r in records),
        'method':'Proved small-cutoff Laurent recurrence, common p-power scaling, exact Newton interpolation at the prime-power cutoff.',
        'records':records}
    Path(__file__).with_suffix('.json').write_text(json.dumps(result,indent=2)+'\n')
\end{lstlisting}
\section{Reproducible third-level unit-mass evaluator}\label{loc-VIII-D}
This evaluator supplies the five finite exceptions in the proof of
the third-level first-step congruence. The final companion step at
that level is not asserted by this computation.
\begin{lstlisting}[language=Python]
"""Exact modular interpolation masses, using interval products in residue classes.
No floating arithmetic; products remove powers of p before inversion.
"""
from pathlib import Path
from collections import Counter
import json,time

class Products:
    def __init__(self,p,L,N):
        self.p=p;self.mod=p**L;self.tables={}
        B=4*N
        for step in (3,6):
            for residue in range(step):
                low=(-B-residue+step-1)//step;high=(B-residue)//step
                vs=[0];us=[1]; factors=[]
                for j in range(low,high+1):
                    x=residue+step*j;v=0
                    if x==0:x=1
                    else:
                        while x%p==0:x//=p;v+=1
                    factors.append(x%self.mod)
                    vs.append(vs[-1]+v);us.append(us[-1]*x%self.mod)
                inv=[0]*len(us);inv[-1]=pow(us[-1],-1,self.mod)
                for j in range(len(factors)-1,-1,-1):
                    inv[j]=inv[j+1]*factors[j]%self.mod
                self.tables[(step,residue)]=(low,vs,us,inv)
    def product(self,lo,hi,step):
        assert (hi-lo)%step==0 and hi>=lo
        residue=lo%step;base,vs,us,inv=self.tables[(step,residue)]
        a=(lo-residue)//step-base;b=(hi-residue)//step-base+1
        return vs[b]-vs[a],us[b]*inv[a]%self.mod

def analyze(p,r,L=10):
    N=p**r;m=(N-2)//3;mod=p**L;pr=Products(p,L,N)
    def prods(H,numer):
        ranges=([(6-H,6*m-H,6),(-3*m-H,-3-H,3),(-3*m+1-H,-2-H,3)] if numer
                else [(1-H,3*m+1-H,3),(2-H,3*m+2-H,3),(-6*m-2-H,-2-H,6)])
        v=0;u=1
        for lo,hi,step in ranges:
            a,b=pr.product(lo,hi,step);v+=a;u=u*b%mod
        return v,u
    vq,uq=prods(0,False);vd,ud=prods(0,True)
    cv=vq-vd;cu=uq*pow(ud,-1,mod)%mod
    denominator_roots=set(z for j in range(1,m+1) for z in (6*j,-3*j,-3*j+1))
    mass=[0]*(r+1);counts=[0]*(r+1);minimum=[None]*(r+1);dist=[Counter() for _ in range(r+1)]
    def vp(x):
        n=0
        if not x:return L
        while x%p==0:n+=1;x//=p
        return n
    half=pow(2,-1,mod)
    for h in range(m+1):
        for H,c in [(3*h+1,1 if h%2==0 else 0),(3*h+2,(-1)**h*half),(-6*h-2,-half)]:
            if not c or H in denominator_roots:continue
            a,b=prods(H,True);d,e=prods(H,False)
            v=cv+a-d;assert v>=0
            u=cu*c*b*pow(e,-1,mod)%mod
            term=u*p**v%mod;t=vp(H)
            mass[t]=(mass[t]+term)%mod;counts[t]+=1;dist[t][v]+=1
            minimum[t]=v if minimum[t] is None else min(minimum[t],v)
    return {'p':p,'r':r,'modulus_exponent':L,'mass_mod':mass,
            'mass_valuations':[vp(z) for z in mass],'counts':counts,
            'minimum_individual_valuations':minimum,
            'individual_valuation_counts':[dict(z) for z in dist]}
if __name__=='__main__':
    rows=[]
    for p in (5,11,17,23,29):
        start=time.time();row=analyze(p,3);rows.append(row)
        print(p,row['mass_valuations'],row['minimum_individual_valuations'],'seconds',round(time.time()-start,2),flush=True)
    Path(__file__).with_suffix('.json').write_text(json.dumps(rows,indent=2))
\end{lstlisting}



\section{Direct checks of the original third-level sum}\label{loc-VIII-E}
The following standalone standard-library program verifies the
original cutoff and target modulus. It does not use the block
identity in its computation.
\begin{lstlisting}[language=Python]
"""Supplementary exact modular checks of the unchanged Swisher E.3 sum.

This program uses no block identity or pole formula from the proof.
It keeps p-valuations separately, so it never inverts a multiple of p.
"""
import json
from pathlib import Path


def strip(x, p):
    v = 0
    while x % p == 0:
        x //= p
        v += 1
    return x, v


def e_sum(p, r, precision):
    modulus = p ** precision
    unit, valuation, total = 1, 0, 0
    cutoff = (p ** r - 2) // 3
    for k in range(cutoff + 1):
        assert valuation >= 0
        if 3 * valuation < precision:
            term = (6 * k + 1) * pow(unit, 3, modulus)
            term *= p ** (3 * valuation)
            total = (total + (-term if k % 2 else term)) % modulus
        if k < cutoff:
            numerator, nv = strip(3 * k + 1, p)
            denominator, dv = strip(3 * (k + 1), p)
            unit = unit * numerator * pow(denominator, -1, modulus) % modulus
            valuation += nv - dv
    return total


def run():
    rows = []
    for p in (5, 11, 17, 23, 29, 41, 47, 53, 59):
        precision = 10
        defect = (e_sum(p, 3, precision) - p * p * e_sum(p, 1, precision)) % p ** precision
        valuation = precision if defect == 0 else strip(defect, p)[1]
        assert valuation >= 8, (p, defect, valuation)
        rows.append({"p": p, "r": 3, "cutoff": (p ** 3 - 2) // 3,
                     "target_exponent": 8, "precision": precision,
                     "defect_residue": defect, "defect_valuation_at_least": valuation,
                     "passed": True})
    return {"method": "direct original sum; exact modular arithmetic", "cases": rows}


if __name__ == "__main__":
    result = run()
    Path(__file__).with_suffix(".json").write_text(json.dumps(result, indent=2) + "\n")
    print(json.dumps(result, indent=2))

\end{lstlisting}
\begin{thebibliography}{9}
\bibitem{VIII-Swisher}
H. Swisher, \emph{On the supercongruence conjectures of van Hamme},
Research in the Mathematical Sciences \textbf{2} (2015), Article 18,
Section 7, Table 2.
\url{https://doi.org/10.1186/s40687-015-0037-6}.
\bibitem{VIII-Bailey}
W. N. Bailey, \emph{Generalized Hypergeometric Series}, Cambridge
University Press, 1935. Classical Whipple and Dixon transformations.
\bibitem{VIII-GuoSchlosser}
V. J. W. Guo and M. J. Schlosser,
\emph{New q-Analogues of Van Hamme's (E.2) Supercongruence and of
a Supercongruence by Swisher}, Results in Mathematics (2023).
\url{https://doi.org/10.1007/s00025-023-01885-8}.
\bibitem{VIII-Corpus}
QEEDDD, \emph{Verifiable Mathematical Reasoning Corpus},
research context supplied for this investigation.
\url{https://github.com/QEEDDD/Verifiable-Mathematical-Reasoning-Corpus}.
\end{thebibliography}

\clearpage\endgroup

%%%%%%%%%%%%%%%%%%%%%%%%%%%%%%%%%%%%%%%%%%%%%%%%%%%%%%%%%%%%%%%%%%%%%%%%%%
% S | work_f3/f_odd_q_proof.tex in the J.3/F.3 research bundle.
%%%%%%%%%%%%%%%%%%%%%%%%%%%%%%%%%%%%%%%%%%%%%%%%%%%%%%%%%%%%%%%%%%%%%%%%%%
\clearpage\begingroup
\SetPaper{S}{Working proof of the odd branch of F.3}
\RestoreCounters
\counterwithout*{equation}{section}\renewcommand{\theequation}{\arabic{equation}}
\counterwithout*{theorem}{section}\renewcommand{\thetheorem}{\arabic{theorem}}
\let\lemma\supplemma\let\endlemma\endsupplemma
\PaperTitle{A short-cutoff creative-microscope proof of the odd branch of F.3}{Working proof for independent audit}{10 November 2026}

Let $(x;q)_k=\prod_{i=0}^{k-1}(1-xq^i)$ and
$[m]_q=(1-q^m)/(1-q)$. Products with several first arguments mean the
product of the respective shifted factorials. Define
\[
 T_L(a,q)=\sum_{k=0}^L c_k(a,q),\qquad
 c_k(a,q)=(-1)^k[8k+1]_q q^k
 \frac{(aq,q/a,q,-q^2;q^4)_k}{(aq^4,q^4/a,q^4,-q^3;q^4)_k}.
\]
The parameter $a$ is an indeterminate throughout every root-of-unity
calculation. In particular, parameter denominators do not vanish there.

\begin{theorem}\label{S-main}
Let $p\equiv3\pmod4$ be prime and $r=2s+1\ge3$. Put
$N=(p^r-3)/4$ and $M=(p^{r-2}-3)/4$. Then
\[
 \sum_{k=0}^{N}(-1)^k(8k+1)\frac{(1/4)_k^3}{k!^3}
 \equiv p^2\sum_{k=0}^{M}(-1)^k(8k+1)\frac{(1/4)_k^3}{k!^3}
 \pmod {p^{3s}}.
\]
Since $3s\ge2s+1=r$, this implies the odd branch of Swisher F.3.
\end{theorem}

\paragraph{The terminating evaluation.}
We use the following terminating very-well-poised basic hypergeometric
evaluation, with $v=4t+1$:
\begin{equation}\label{S-jackson}
 T_t(q^{v},q)=T_t(q^{-v},q)=(-1)^t q^{-2t}[v]_q.
\end{equation}
Here is an exact general identity from which it follows:
\begin{equation}\label{S-jackson-general}
 \sum_{k=0}^{t}[2mk+1]_q
 \frac{(q^{1-v},q^{1+v},q/c,q;q^m)_k}
 {(q^{m-v},q^{m+v},cq^m,q^m;q^m)_k}
 c^k q^{(m-2)k}
 = (c/q)^t\frac{(q^2/c;q^m)_t}{(cq^m;q^m)_t}[v]_q,
 \quad v=mt+1.
\end{equation}
To specify the input identity completely, the terminating
very-well-poised ${}_6\phi_5$ summation is
\[
 \sum_{k=0}^t\frac{1-AQ^{2k}}{1-A}
 \frac{(A,B,C,Q^{-t};Q)_k}
 {(Q,AQ/B,AQ/C,AQ^{t+1};Q)_k}
 \left(\frac{AQ^{t+1}}{BC}\right)^k
 =\frac{(AQ,AQ/(BC);Q)_t}{(AQ/B,AQ/C;Q)_t}.
\]
Take $Q=q^m$, $A=q$, $B=q^{1+v}$, $C=q/c$, and $v=mt+1$.
The argument becomes $cq^{m-2}$. On the right, reversal of the
two finite products with negative powers gives
\[
 \frac{(q^{m+1};q^m)_t}{(q^{m-v};q^m)_t}
 =(-1)^tq^{mt(t-1)/2+t}[v]_q,
\]
\[
 (cq^{m-v-1};q^m)_t
 =(-c)^tq^{-mt(t-1)/2-2t}(q^2/c;q^m)_t,
\]
which proves \eqref{S-jackson-general} from that standard summation.
This is also the specialization used in Guo--Zudilin,
\emph{Dwork-type supercongruences
through a creative $q$-microscope}, arXiv:2001.02311, formula (3.23)
and the proof of Theorem 3.6. Its substitution $m=4,c=-q^{-1}$ gives
\eqref{S-jackson}. Equation \eqref{S-jackson-general} is an exact identity
at the two parameter specializations, rather than a congruence.

\begin{lemma}[Root blocks]\label{S-blocks}
Suppose $h>1$, $h\equiv1\pmod4$, and $\zeta$ is a primitive $h$th root
of unity. Write $t=(h-1)/4$. For $b\ge0$ and $0\le j<h$,
\begin{equation}\label{S-block-formula}
 \lim_{q\to\zeta}c_{bh+j}(a,q)=
 \begin{cases}
 (-1)^b\dfrac{(1/4)_b}{b!}c_j(a,\zeta),&0\le j\le t,\\
 0,&t<j<h.
 \end{cases}
\end{equation}
Furthermore $c_{t-j}(a,\zeta)=-c_j(a,\zeta)$ for $0\le j\le t$.
Consequently, if $L\bmod h\ge t$, then $T_L(a,q)$ vanishes modulo
$\Phi_h(q)$, in the local ring with coefficient field $\mathbb Q(a)$.
\end{lemma}
\begin{proof}
Set
\[
\begin{aligned}
 U_k&=(aq,q/a;q^4)_k/(aq^4,q^4/a;q^4)_k,\\
 Q_k&=(q;q^4)_k/(q^4;q^4)_k,\qquad
 V_k=(-q^2;q^4)_k/(-q^3;q^4)_k.
\end{aligned}
\]
In a complete block of length $h$, the numerator and denominator of
$U$ each run over the same root residues, so their ratio is one.
The analogous full-block ratio for $V$ is one because both products
equal $2$. In $Q$, the numerator zeros occur at indices $t+uh$ and
the denominator zeros at indices $(u+1)h-1$. Their limiting ratio in
the $u$th block is $(4u+1)/(4u+4)$; the products of nonzero factors
cancel by permutation of the nontrivial roots. Thus $b$ full blocks
contribute $(1/4)_b/b!$. The external $(-1)^kq^k$ contributes
$(-1)^b$. A final partial block with $j>t$ has one extra numerator
zero. This proves \eqref{S-block-formula}.

For completeness, reversing the products in the first short block
gives, with all the following expressions evaluated at $q=\zeta$,
\[
 \frac{U_{t-j}}{U_j}=q^{6j-4t(t+1)},\qquad
 \frac{Q_{t-j}}{Q_j}=(-1)^tq^{3j-2t(t+1)},\qquad
 \frac{V_{t-j}}{V_j}=q^{j-2t^2-t}.
\]
These follow by replacing the reversed factors $q^{h-u}$ by $q^{-u}$;
all denominators here are nonzero for generic $a$.
Also
$[8(t-j)+1]_q=-q^{-8j-1}[8j+1]_q$.
Combining these identities and the factor $(-1)^{t-2j}q^{t-2j}$
from the external sign and power yields
\[
 c_{t-j}(a,q)=-q^{-(2t+1)h}c_j(a,q)=-c_j(a,q).
\]
The short sum is zero, and \eqref{S-block-formula} proves the assertion
for every cutoff whose last block includes this short block.
\end{proof}

\paragraph{Parameter interpolation and the common denominator.}
Set $n=p^2$, so $n\equiv1\pmod8$, and write
\[
 N=nM+\frac{3(n-1)}4,\qquad
 \Delta(a,q)=T_N(a,q)-q^{(1-n)/2}[n]_qT_M(a,q^n).
\]
For $0\le j\le M$, let $A_j=(4j+1)n$.
At $a=q^{\pm A_j}$, the left series terminates at
$nj+(n-1)/4\le N$, while the right series terminates at $j$.
By \eqref{S-jackson}, the left value is
$(-1)^{nj+(n-1)/4}q^{(1-A_j)/2}[A_j]_q$.
The right value is
$q^{(1-n)/2}[n]_q(-1)^j q^{-2nj}[4j+1]_{q^n}$.
They agree since $(n-1)/4$ is even and
$[n]_q[4j+1]_{q^n}=[A_j]_q$. Therefore
\[
 P(a,q)=\prod_{j=0}^{M}(1-aq^{A_j})(a-q^{A_j})
\]
divides the parameter numerator of $\Delta$.
Explicitly a common parameter denominator is
\[
 D(a,q)=(aq^4,q^4/a;q^4)_N.
\]
It clears the right denominator too, because each exponent
$4n,8n,\ldots,4nM$ occurs among $4,8,\ldots,4N$.
Hence $a^ND\Delta$ is a polynomial in $a$ with coefficients in
$\mathbb Q(q)$ and is divisible by $P$ in $\mathbb Q(q)[a]$.
The leading coefficient of $P$ is a signed power of $q$.

\paragraph{An extra cyclotomic factor survives interpolation.}
Fix an even $\ell\in\{2,4,\ldots,2s\}$, and put $h=p^\ell$.
Let $R_h=\mathbb Q[q,q^{-1}]_{(\Phi_h)}$.
For any odd integer $d$, the ratio $Q_k=(q;q^4)_k/(q^4;q^4)_k$
has nonnegative $\Phi_d$ valuation: in each period $d$ its numerator
zero occurs at some index $0\le j_0<d$, whereas its denominator zero
occurs at index $d-1$. The factors $1+q^{4i+3}$ have no zero at an
odd-order root. The same facts apply to $q^n$ at a primitive
$p^\ell$ root (at $\ell=2$ the limit is $q^n=1$ and $Q_k$ has the
finite value $(1/4)_k/k!$). Thus
$a^ND\Delta\in R_h[a]$.

At a primitive $h$th root, the left sum vanishes by Lemma
\ref{S-blocks}, since
\[
 N\bmod h=\frac{3(h-1)}4\ge\frac{h-1}4.
\]
For $\ell=2$, the right side vanishes due to $[n]_q$ and the remaining
sum is finite for generic $a$. For $\ell>2$, its base $q^n$ has order
$h/n$, and
$M\bmod(h/n)=3(h/n-1)/4$, so Lemma \ref{S-blocks} again applies.
Consequently $a^ND\Delta\in\Phi_hR_h[a]$.

Since $P$ has unit leading coefficient in $R_h$, polynomial division
shows that $H=a^ND\Delta/P$ belongs to $R_h[a]$. Reducing the equality
$PH=a^ND\Delta$ modulo $\Phi_h$, the polynomial $P$ is nonzero
(indeed its leading coefficient remains a unit). The residue ring
is a polynomial ring over a field and has no zero divisors; hence
$H\in\Phi_hR_h[a]$. Specializing $a=1$ gives
\begin{equation}\label{S-multiplicity}
 \vp_{\Phi_h}\Delta(1,q)\ge1+
 \vp_{\Phi_h}P(1,q)-\vp_{\Phi_h}D(1,q).
\end{equation}

\paragraph{Exact root counts.}
Put $u=h/n=p^{\ell-2}$. Then $u\equiv1\pmod4$ (including $u=1$),
and
\[
 \#\{0\le j\le M:u\mid4j+1\}
 =\lfloor M/u\rfloor+1=\lfloor N/h\rfloor+1.
\]
Indeed the first solution in a period is $(u-1)/4$, while
$M\bmod u=3(u-1)/4$; for $u=1$ both residues are zero.
The last floor equality follows from
$N=nM+3(n-1)/4$ and
$n\cdot3(u-1)/4+3(n-1)/4=3(h-1)/4<h$.
Therefore
\[
 \vp_{\Phi_h}P(1,q)=2(\lfloor N/h\rfloor+1),\qquad
 \vp_{\Phi_h}D(1,q)=2\lfloor N/h\rfloor.
\]
Equation \eqref{S-multiplicity} gives at least three factors for every
even $\ell$, i.e.
\begin{equation}\label{S-qresult}
 \vp_{\Phi_{p^{2j}}}\Delta(1,q)\ge3\quad(1\le j\le s).
\end{equation}

\paragraph{The specialization $q=1$ has a $p$-unit denominator.}
This point is needed in addition to the local multiplicities.
At $a=1$ the Pochhammer portion of each summand is
$Q_k^3V_k$. The preceding zero-count argument shows that after
cancellation $Q_k$ has no odd-order cyclotomic denominator.
Its remaining denominator factors come from $1-q^{4i}$, and so are
even-order cyclotomic polynomials; their values at $1$ are $1$ or $2$.
Likewise every factor of $1+q^{4i+3}$ has even order and value at $1$
equal to $1$ or $2$. Thus there is a Laurent polynomial denominator
$B(q)\in\mathbb Z[q,q^{-1}]$ such that $B(1)$ is a power of $2$,
$A(q)=B(q)\Delta(1,q)\in\mathbb Z[q,q^{-1}]$, and $B$ has no
odd-order cyclotomic factor. This remains true for the right side
after substituting $q\mapsto q^n$ because $n$ is odd: an odd power
cannot turn an even-order root into an odd-order root, and the value
at $q=1$ of the substituted denominator is unchanged.
By \eqref{S-qresult}, $A$ is divisible over $\mathbb Q[q,q^{-1}]$ by
$\prod_{j=1}^s\Phi_{p^{2j}}(q)^3$. Since these factors are monic,
Gauss's lemma gives the same divisibility over
$\mathbb Z[q,q^{-1}]$. Their values at $q=1$ are all $p$.
Finally,
\[
 T_L(1,1)=\sum_{k=0}^L(-1)^k(8k+1)\frac{(1/4)_k^3}{k!^3},
 \qquad q^{(1-n)/2}[n]_q\big|_{q=1}=p^2.
\]
Since $B(1)$ is a $p$-adic unit, \eqref{S-qresult} proves
Theorem \ref{S-main}.

\clearpage\endgroup
\end{document}
